\documentclass[a4,11pt]{aleph-examen}

% -- Paquetes adicionales
\usepackage{enumitem}
\usepackage{aleph-comandos}

% -- Datos
\institucion{Facultad de Ciencias Exactas, Naturales y Ambientales}
\carrera{Catálogo STEM}
\asignatura{Matemática III}
\tema{Ejercicios 01: Matrices}
\autor{Andrés Merino}
\fecha{Periodo 2026-1}

\logouno[0.16\textwidth]{Logos/logoPUCE_04_ac}
\definecolor{colortext}{HTML}{1957d1}
\definecolor{colordef}{HTML}{1957d1}
\fuente{montserrat}

% -- Comandos adicionales
\setlist[enumerate]{label=\roman*.}

\begin{document}

\encabezado

\begin{preguntas}

%%%%%%%%%%%%%%%%%%%%%%%%%%%%%%%%%%%%%%%%%%%
%%%%%%%%%%%%%%%%%%%%%%%%%%%%%%%%%%%%%%%%%%%
%%%%%%%%%%%%%%%%%%%%%%%%%%%%%%%%%%%%%%%%%%%
% Tema 5: Limites y diferenciabilidad
\item
    Dado el campo escalar
    \[
       \funcion{f}{\R^n}{\R}{x}
       {f(x)=
           \begin{cases}
               \|x\|^2 &\text{si }\|x\|<1,\\
               0 &\text{si }\|x\|\geq 1.
           \end{cases}
       }
    \]
    Describir los conjuntos de nivel de $f$.

\begin{respuesta}
    Notemos que
    \[
        0\leq f(x) < 1,
    \]
    para todo $x\in\R^n$, por lo tanto,
    \[
        L_c(f)=\varnothing,
    \]
    si $c<0$ o si $c\geq 1$. Por otro lado, tenemos que
    \begin{align*}
        L_0(f) 
            &= \{ x\in\R^n: f(x)=0 \}\\
            &= \{ x\in\R^n: (\|x\|^2=0 \land \|x\|<1) \lor \|x\|\geq 1 \}\\
            &= \{ x\in\R^n: x=0 \lor \|x\|\geq 1 \}\\
            &= \{ x\in\R^n: \|x\|\geq 1 \} \cup \{0\}\\
            &= B(0,1)^c\cup\{0\}.
    \end{align*}
    Además, para todo $c\in\lop0,1\rop$, se tiene que
    \begin{align*}
        L_c(f)
            &= \{ x\in\R^n: f(x)=c \}\\
            &= \{ x\in\R^n: ||x||^2=c \}\\
            &= S(0,\sqrt{c}).
    \end{align*}
    Así, se tiene que
    \[
        L_c(f)=
        \begin{cases}
            B(0,1)^c\cup\{0\} &\text{si }c=0,\\
            S(0,\sqrt{c}) &\text{si }c\in\lop0,1\rop,\\
            \varnothing &\text{caso contrario.}
        \end{cases}\qedhere
    \]
\end{respuesta}

%%%%%%%%%%%%%%%%%%%%%%%%%%%%%%%%%%%%%%%%%%%
%%%%%%%%%%%%%%%%%%%%%%%%%%%%%%%%%%%%%%%%%%%
%%%%%%%%%%%%%%%%%%%%%%%%%%%%%%%%%%%%%%%%%%%
% Tema 5: Limites y diferenciabilidad
\item
    Dado el campo escalar
    \[
        \funcion{f}{\R^2}{\R}{(x,y)}{
         \begin{cases}
            1&\text{ si }x^2+y^2<1,\\
            0&\text{ si }x^2+y^2\geq 1,
         \end{cases}}
    \]
    describa todos los conjuntos de nivel de $f$ y bosqueje sus gráficas.
  
\begin{respuesta}
    Notemos que
    \[
        f(x)\in \{0,1\},
    \]
    por lo tanto
    \[
        L_c(f)=\varnothing,
    \]
    si $c \not\in \{0,1\}$. Por otro lado, tenemos que
    \begin{align*}
        L_0(f)
            & = \{(x,y)\in\R^2 : f(x,y) = 0\} \\
            & = \{(x,y)\in\R^2 : x^2 + y^2 \geq 1\} \\
            & = \R^2 \setminus B(0,1)
    \end{align*}
    y
    \begin{align*}
        L_1(f)
            & = \{(x,y)\in\R^2 : f(x,y) = 1\} \\
            & = \{(x,y)\in\R^2 : x^2 + y^2 < 1\} \\
            & = B(0,1).
    \end{align*}
    Así, se tiene que
    \[
        L_c(f)=
        \begin{cases}
            \R^2 \setminus B(0,1) & \text{ si } c=0 \\
            B(0,1) & \text{ si } c=1 \\
            \varnothing & \text{caso contrario.}
        \end{cases}
    \]
\end{respuesta}

%%%%%%%%%%%%%%%%%%%%%%%%%%%%%%%%%%%%%%%%%%%
%%%%%%%%%%%%%%%%%%%%%%%%%%%%%%%%%%%%%%%%%%%
%%%%%%%%%%%%%%%%%%%%%%%%%%%%%%%%%%%%%%%%%%%
% Tema 5: Limites y diferenciabilidad
\item
    Encontrar y describir los conjuntos de nivel en $-1$, $0$, $1$, $2$ para los campos escalares dados a continuación. Con esta información, bosquejar las funciones.
    \begin{enumerate}
    \item
        El campo escalar 
        \[
            \funcion{f}{\R^2}{\R}{(x,y)}{x+2y.}
        \]
            
    \item
        El campo escalar
        \[
            \funcion{f}{\R^2}{\R}{(x,y)}{xy.}
        \]
    
    \item
        El campo escalar
        \[
            \funcion{f}{\R^2}{\R}{(x,y)}{x^2-y.}
        \]
   \end{enumerate}

\begin{respuesta}
    \begin{enumerate}
        \item
            Notemos que 
            \begin{align*}
                L_{-1}(f)&=\l\{(x,y)\in\R^2:y=-\frac{x+1}{2}\r\}\\
                L_{0}(f)&=\l\{(x,y)\in\R^2:y=-\frac{x}{2}\r\}\\
                L_{1}(f)&=\l\{(x,y)\in\R^2:y=\frac{1-x}{2}\r\}\\
                L_{2}(f)&=\l\{(x,y)\in\R^2:y=\frac{2-x}{2}\r\}.
            \end{align*}
            Nos podemos percatar que todas estas son rectas, como se muestra en la siguiente figura:
            \begin{center}
                \includegraphics[width=0.8\linewidth]{Graficos/fig02:08a_c.png}
            \end{center}
            Así, la gráfica de la función sería:
            \begin{center}
                \includegraphics[width=0.8\linewidth]{Graficos/fig02:08a_s.png}
            \end{center}
        \item
            Notemos que 
            \begin{align*}
                L_{-1}(f)&=\{(x,y)\in\R^2:xy=-1\}\\
                L_{0}(f)&=\{(x,y)\in\R^2:xy=0\}=\{(x,y)\in\R^2:x=0 \lor y=0\}\\
                L_{1}(f)&=\{(x,y)\in\R^2:xy=1\}\\
                L_{2}(f)&=\{(x,y)\in\R^2:xy=2\}.
            \end{align*}
            Nos podemos percatar que todas, excepto $L_{0}(f)$, son hipérbolas con centro en el origen y asíntotas sobre los ejes, como se muestra en la siguiente figura:
            \begin{center}
                \includegraphics[width=0.8\linewidth]{Graficos/fig02:08b_c.png}
            \end{center}
            Así, la gráfica de la función sería:
            \begin{center}
                \includegraphics[width=0.8\linewidth]{Graficos/fig02:08b_s.png}
            \end{center}
            
        \item
            Notemos que 
            \begin{align*}
                L_{-1}(f)&=\{(x,y)\in\R^2:y=x^2+1\}\\
                L_{0}(f)&=\{(x,y)\in\R^2:y=x^2\}\\
                L_{1}(f)&=\{(x,y)\in\R^2:y=x^2-1\}\\
                L_{2}(f)&=\{(x,y)\in\R^2:y=x^2-2\}.
            \end{align*}
            Nos podemos percatar que todas estas son parábolas cuyo vértice está sobre el eje $x$, como se muestra en la siguiente figura:
            \begin{center}
                \includegraphics[width=0.5\linewidth]{Graficos/fig02:08c_c.pdf}
            \end{center}
            Así, la gráfica de la función sería:
            \begin{center}
                \includegraphics[width=0.6\linewidth]{Graficos/fig02:08c_s.pdf}
            \end{center}
    \end{enumerate}
\end{respuesta}

%%%%%%%%%%%%%%%%%%%%%%%%%%%%%%%%%%%%%%%%%%%
%%%%%%%%%%%%%%%%%%%%%%%%%%%%%%%%%%%%%%%%%%%
%%%%%%%%%%%%%%%%%%%%%%%%%%%%%%%%%%%%%%%%%%%
% Tema 5: Limites y diferenciabilidad
\item 
    Dado el campo escalar
    \[
        \funcion{f}{\R^3}{\R}{(x,y,z)}{x^2+y^2-z,}
    \]
    estudie y describa sus conjuntos de nivel.
      
% \begin{respuesta}
%     Sea $c\in\R$, entonces
%     \begin{align*}
%       L_c(f)
%       &=\{(x,y,z)\in\R^3: f(x,y,z)=c\}\\
%       &=\{(x,y,z)\in\R^3: x^2+y^2-z=c\}\\
%       &=\{(x,y,z)\in\R^3: x^2+y^2-c=z\}.
%     \end{align*}
%     La gráfica de este conjunto es un paraboloide elíptico de vértice en el punto $(0,0,-c)$.
% \end{respuesta}

%%%%%%%%%%%%%%%%%%%%%%%%%%%%%%%%%%%%%%%%%%%
%%%%%%%%%%%%%%%%%%%%%%%%%%%%%%%%%%%%%%%%%%%
%%%%%%%%%%%%%%%%%%%%%%%%%%%%%%%%%%%%%%%%%%%
% Tema 5: Limites y diferenciabilidad
\item 
    Sean $f,g:\R^n\to\R$. Demuestre o refute: 
    
    si $L_c(f)=L_c(g)$ para todo $c\in\R$, entonces $f=g$. 

\begin{respuesta}
    Es verdadero. Supongamos que 
    \[
        L_c(f)=L_c(g) 
    \]
    para todo $c\in\R$. Dado $x\in\R^n$ vamos a demostrar que $f(x)=g(x)$. Puesto que $f(x)\in\R$, por la hipótesis,
    \[
        L_{f(x)}(f)=L_{f(x)}(g).
    \]
    Como $x\in L_{f(x)}(f)$, se tiene que $x\in L_{f(x)}(g)$, de donde
    \[
        g(x)=f(x).
    \]
    Puesto que $x$ es arbitrario, se tiene que
    \[
        f=g.\qedhere
    \]
\end{respuesta}

%%%%%%%%%%%%%%%%%%%%%%%%%%%%%%%%%%%%%%%%%%%
%%%%%%%%%%%%%%%%%%%%%%%%%%%%%%%%%%%%%%%%%%%
%%%%%%%%%%%%%%%%%%%%%%%%%%%%%%%%%%%%%%%%%%%
% Tema 5: Limites y diferenciabilidad
\item
    Encontrar una función lineal $\func{f}{\R^3}{\R}$ tal que su conjunto de nivel en $1$ sea
    \[
        \{(x,y,z)\in\R^3:z+1=2x+3y\}.
    \]

\begin{respuesta}
    Estamos buscando una función $f$ de tal forma que
    \[
        L_1(f)=\{(x,y,z)\in\R^3:z+1=2x+3y\},
    \]
    es decir, buscamos que
    \begin{align*}
        f(x,y,z)=1
            & \Longleftrightarrow z+1=2x+3y \\
            & \Longleftrightarrow 2x+3y-z=1,
    \end{align*}
    para todo $(x,y,z)\in\R^3$. Además, dado que $f$ es lineal, existen $\alpha, \beta, \gamma\in\R$ tal que
    \[
        f(x,y,z)=\alpha x + \beta y +\gamma z
    \]
    para todo $(x,y,z)\in\R^3$. Por lo tanto, buscamos una función tal que
    \[
        \alpha x + \beta y +\gamma z=1 \Longleftrightarrow 2x+3y-z=1,
    \]
    para todo $(x,y,z)\in\R^3$. Así, se tiene que si
    \[
        \alpha=2,
        \qquad
        \beta=3,
        \texty
        \gamma=-1,
    \]
    se tiene la función buscada. Es decir, $f$ es la función definida por
    \[
        f(x,y,z)=2x+3y-z,
    \]
    para todo $(x,y,z)\in\R^3$.
\end{respuesta}

%%%%%%%%%%%%%%%%%%%%%%%%%%%%%%%%%%%%%%%%%%%
%%%%%%%%%%%%%%%%%%%%%%%%%%%%%%%%%%%%%%%%%%%
%%%%%%%%%%%%%%%%%%%%%%%%%%%%%%%%%%%%%%%%%%%
% Tema 1: Funciones vectoriales y curvas
\item
    Calcule el límite de la trayectoria $\alpha$ en $5$, donde
    \[
        \funcion{\alpha}{\R\setminus\{0\}}{\R^3}{y}
        {\left(y^3+5,y^2,\frac{y+1}{y}\right).}
    \]
    
% \begin{respuesta}
%     Primero, definamos las funciones componentes de $\alpha$, tenemos $\func{\alpha_1,\alpha_2,\alpha_3}{\R}{\R}$ definidas por
%     \[
%         \alpha_1(y)=y^3+5,
%         \qquad
%         \alpha_2(y)=y^2
%         \texty
%         \alpha_3(y)=\frac{y+1}{y}.
%     \]
%     para $x\in\R$, así, se tiene que para todo $x\in \R$ 
%     \[
%         \alpha(x)=(\alpha_1(x), \alpha_2(x), \alpha_3(x)).
%     \]
%     Luego, gracias a que
%     \[
%         \lim_{y\to 5}\alpha_1(x)=130,
%         \qquad
%         \lim_{y\to 5}\alpha_2(x)=25
%         \texty
%         \lim_{y\to 5}\alpha_3(x)=\frac{6}{5},
%     \]
%     se sigue que $\dlim_{y\to 5} \alpha(x)$ existe y
%     \begin{align*}
%         \dlim_{y\to 5} \alpha(x)
%             & = \l(\lim_{y\to 5}\alpha_1(x), \lim_{y\to 5}\alpha_2(x), \lim_{y\to 5}\alpha_3(x) \r) \\
%             & = \l(130, 25,\frac{6}{5} \r). \qedhere
%     \end{align*}
% \end{respuesta}

%%%%%%%%%%%%%%%%%%%%%%%%%%%%%%%%%%%%%%%%%%%
%%%%%%%%%%%%%%%%%%%%%%%%%%%%%%%%%%%%%%%%%%%
%%%%%%%%%%%%%%%%%%%%%%%%%%%%%%%%%%%%%%%%%%%
% Tema 1: Funciones vectoriales y curvas
\item
    Dada la trayectoria
    \[
        \funcion{\alpha}{\R}{\R^3}{t}
        {\begin{cases}
            \left(\dfrac{\sen(t)}{t},2t^2,\dfrac{e^t-1}{t}\right) &\text{si }t\neq 0,\\[4mm]
            x&\text{si }t= 0,
        \end{cases}}
    \]
    con $x\in \R^3$. ¿Cuál es el valor de $x$ para que $\alpha$ sea continua en $0$?
\begin{respuesta}
    Para que la trayectoria sea continua en este punto se debe cumplir que
    \[
        \lim_{t \to 0}\alpha(t)=\alpha(0)=x.
    \]
    Por lo tanto, procedemos a calcular el límite:
    \begin{align*}
        \lim_{t \rightarrow 0}\alpha(t)
            &=\left(\lim_{t \rightarrow 0}\frac{\sen(t)}{t},\lim_{t \rightarrow 0}2t^2,\lim_{t \rightarrow 0}\frac{e^t-1}{t}\right) \\
            & =(1,0,1).
    \end{align*}
    Entonces el valor de $x$ para que la trayectoria sea continua es $(1,0,1)$.
\end{respuesta}

%%%%%%%%%%%%%%%%%%%%%%%%%%%%%%%%%%%%%%%%%%% JO
%%%%%%%%%%%%%%%%%%%%%%%%%%%%%%%%%%%%%%%%%%%
%%%%%%%%%%%%%%%%%%%%%%%%%%%%%%%%%%%%%%%%%%%
% Tema 1: Funciones vectoriales y curvas
\item
    Dada la función
    \[
        \funcion{\alpha}{[-1,1]\setminus\{0\}}{\R^2}{t}{\l(\dfrac{\sen(\pi t)}t,e^{3t}-t^2\r).}
    \]
    Estudie la continuidad de la función $\alpha$ en $0$, ¿existe $\Lim_{t\to 0} \alpha(t)$? Si $\Lim_{t\to 0} \alpha(t)$ existe, halle su valor.

% \begin{respuesta}
%     Dado que $0\not\in\dom(\alpha)$, se tiene que $\alpha(0)$ no existe, luego $\alpha$ no es continua en $0$. Por otra parte, para que
%     \[
%       \Lim_{t\to 0} \alpha(t) 
%     \]
%     exista, deben existir
%     \[
%         \Lim_{t\to 0} \dfrac{\sen(\pi t)}t
%         \yds
%         \Lim_{t\to 0} e^{3t}-t^2.
%     \]
%     Ahora bien, como $\Lim_{t\to 0}\sen(\pi t)=0$ y $\Lim_{t\to 0} t=0$, por la regla de L'Hopital, se tiene que
%     \[
%         \Lim_{t\to 0} \dfrac{\sen(\pi t)}t
%         =\Lim_{t\to 0} \dfrac{\pi\cos(\pi t)}1
%         =\pi\cos(0)
%         =\pi.
%     \]
%     Por otra parte, como
%     \[
%     t\mapsto e^{3t}-t^2
%     \]
%     es continua en 0, se tiene que
%     \[
%          \Lim_{t\to 0} e^{3t}-t^2=e^0-0^2=1.
%     \]
%     Por lo tanto, $\Lim_{t\to 0} \alpha(t)$ existe y 
%     \[
%         \Lim_{t\to 0} \alpha(t)=(\pi,1).\qedhere
%     \]
% \end{respuesta}

%%%%%%%%%%%%%%%%%%%%%%%%%%%%%%%%%%%%%%%%%%%
%%%%%%%%%%%%%%%%%%%%%%%%%%%%%%%%%%%%%%%%%%%
%%%%%%%%%%%%%%%%%%%%%%%%%%%%%%%%%%%%%%%%%%%
% Tema 1: Funciones vectoriales y curvas
\item
    Dada la trayectoria
	\[
        \funcion{\alpha}{\R}{\R^3}{x}{(x+2,x^2,x^4+x^2),} 
    \]
    calcule la derivada de $\alpha$ y evalúela en $x=10$.
    
% \begin{respuesta}
%     Primero, definamos las funciones componentes de $\alpha$, tenemos $\func{\alpha_1,\alpha_2,\alpha_3}{\R}{\R}$ definidas por
%     \[
%         \alpha_1(x)=x+2,
%         \qquad
%         \alpha_2(x)=x^2
%         \texty
%         \alpha_3(x)=x^4+x^2.
%     \]
%     para $x\in\R$, así, se tiene que para todo $x\in \R$ 
%     \[
%         \alpha(x)=(\alpha_1(x), \alpha_2(x), \alpha_3(x)).
%     \]
%     Luego, gracias a que $\alpha_1,\alpha_2$ y $\alpha_3$ son derivables en $\R$ (puesto que son funciones polinomiales), se sigue que $\alpha$ es derivable en $\R$ y
%     \begin{align*}
%         \alpha'(x) 
%             & = (\alpha'_{1}(x),\alpha'_{2}(x),\alpha'_{3}(x))\\
%             & = (1,2x,4x^3+2x)
%     \end{align*}
%     para todo $x\in\R$. Finalmente, evaluando en $x=10$ tenemos que
%     \[
%         \alpha'(10)=(1,20,4020).\qedhere
%     \]
% \end{respuesta}

%%%%%%%%%%%%%%%%%%%%%%%%%%%%%%%%%%%%%%%%%%%
%%%%%%%%%%%%%%%%%%%%%%%%%%%%%%%%%%%%%%%%%%%
%%%%%%%%%%%%%%%%%%%%%%%%%%%%%%%%%%%%%%%%%%%
% Tema 1: Funciones vectoriales y curvas
\item 
    Dada la trayectoria 
    \[
        \funcion{\alpha}{\R}{\R^2}{x}{\l(\dfrac{x^2}{x^2+1},x|x|\r).}
    \]
    Demuestre que $\alpha$ es derivable en todo su dominio y halle su derivada.

\begin{respuesta}
    Vamos a probar que cada componente es derivable. Tenemos que sus componentes son
    \[
        \funcion{\alpha_1}{\R}{\R}{x}{\dfrac{x^2}{x^2+1}}
        \qquad \text y \qquad
        \funcion{\alpha_2}{\R}{\R}{x}{x|x|.}
    \]
    La función $\alpha_1$ es derivable en todo $\R$ pues es el cociente de dos polinomios y su denominador nunca se anula. Además,
        \begin{align*}
            \alpha_1'(x)
                &=\l(\dfrac{x^2}{x^2+1}\r)'\\
                &=\dfrac{2x(x^2+1)-2x(x^2)}{(x^2+1)^2}\\
                &=\dfrac{2x}{(x^2+1)^2}
        \end{align*}
        para cada $x\in\R$. 
        
    Por la definición del valor absoluto,
        \[
            \alpha_2(x)=
                \begin{cases}
                    -x^2&\text{ si }x<0,\\
                    x^2&\text{ si }x\geq0.
                \end{cases}
        \]
        Es claro que $\alpha_2$ es derivable para todo $x\neq 0$ y se tiene que
        \[
            \alpha'_2(x)=
                \begin{cases}
                    -2x&\text{ si }x<0,\\
                    2x&\text{ si }x>0.
                \end{cases}
        \]
        Ahora analicemos la derivada de $\alpha_2$ en 0. Para todo $h\neq0$, se tiene que
        \begin{align*}
            \dfrac{\alpha_2(h)-\alpha_2(0)}{h}
                &=\dfrac{h|h|}{h}\\
                &=|h|.
        \end{align*}
        Así,
        \[
            \lim_{h\to 0}\dfrac{\alpha_2(h)-\alpha_2(0)}{h}
            =\lim_{h\to 0}|h|=0
        \]
        pues el valor absoluto es una función continua, y se sigue que 
        \[
            \alpha_2'(0)=0.
        \]
        Entonces,
        \[
            \alpha'_2(x)=
            \begin{cases}
                -2x&\text{ si }x<0,\\
                2x&\text{ si }x\geq 0,
            \end{cases}
        \]
        de donde $\alpha'_2(x)=2|x|$.
    
    Por lo tanto, $\alpha$ es derivable en todo su dominio y
    \[
        \alpha'(x)=\l(\dfrac{2x}{(x^2+1)^2},2|x|\r)
    \]
    para todo $x\in\R$.
\end{respuesta}


%%%%%%%%%%%%%%%%%%%%%%%%%%%%%%%%%%%%%%%%%%%
%%%%%%%%%%%%%%%%%%%%%%%%%%%%%%%%%%%%%%%%%%%
%%%%%%%%%%%%%%%%%%%%%%%%%%%%%%%%%%%%%%%%%%%
% Tema 1: Funciones vectoriales y curvas
\item
    Dada la trayectoria
    \[
        \funcion{\alpha}{\R}{\R^3}{t}{\alpha(t)=
            \frac{2t}{1+t^2}\ \veci+\frac{1-t^2}{1+t^2}\ \vecj+\veck,
        }
    \]
    comprobar que el ángulo que forma $\alpha(t)$ con $\alpha'(t)$ es constante.

\begin{respuesta}
    Notemos que
    \[
        \alpha(t)=\left(\frac{2t}{1+t^2},\frac{1-t^2}{1+t^2},1\right)
    \]
    para todo $t\in\R$, además, podemos apreciar que cada función componente es derivable, por lo tanto
    \[
        \alpha'(t)=\left(-\frac{2 \left(t^2-1\right)}{\left(t^2+1\right)^2},-\frac{4 t}{\left(t^2+1\right)^2},0\right)
    \]
    para todo $t\in\R$. Ahora, para calcular el ángulo entre los dos vectores, tenemos que, para todo $t\in\R$
    \[
        \alpha(t)\cdot \alpha'(t) = 0,
    \]
    además,
    \[
        \alpha(t)\neq 0
        \texty
        \alpha'(t)\neq 0,
    \]
    por lo tanto, podemos calcular el ángulo formado por los vectores y este es
    \[
        \arccos\left( \frac{\alpha(t)\cdot \alpha'(t)}{\|\alpha(t)\|\|\alpha'(t)\|}\right)
        = \arccos(0) = \dfrac{\pi}{2}.
    \]
    para todo $t\in\R$.
\end{respuesta}

%%%%%%%%%%%%%%%%%%%%%%%%%%%%%%%%%%%%%%%%%%%
%%%%%%%%%%%%%%%%%%%%%%%%%%%%%%%%%%%%%%%%%%%
%%%%%%%%%%%%%%%%%%%%%%%%%%%%%%%%%%%%%%%%%%%
% Tema 1: Funciones vectoriales y curvas
\item
    Sea $\func{\alpha}{I}{\R^m}$, con $m>1$, una trayectoria derivable. Si existe una constante $k\in\R\setminus\{0\}$ tal que $\| \alpha(t) \|=k$ para todo $t\in I$, demostrar que $\alpha(t)$ y $\alpha'(t)$ son ortogonales para todo $t\in I$.
    
\begin{respuesta}
    Sea $t\in I$, tenemos que
    \begin{align*}
        (\alpha(t) \cdot \alpha(t))'
            & = \alpha'(t) \cdot \alpha(t) + \alpha(t) \cdot \alpha'(t)\\
            & = 2(\alpha(t) \cdot \alpha'(t)),
    \end{align*}
    o, escrito de manera equivalente,
    \[
        (\|\alpha(t)\|^2)' = 2(\alpha(t) \cdot \alpha'(t)),
    \]
    de donde, gracias a que $\| \alpha(t) \|=k$, se sigue que
    \[
        0=\alpha(t) \cdot \alpha'(t),
    \]
    y por lo tanto, $\alpha(t)$ y $\alpha'(t)$ son ortogonales para cada $t\in I$.
\end{respuesta}

%%%%%%%%%%%%%%%%%%%%%%%%%%%%%%%%%%%%%%%%%%%
%%%%%%%%%%%%%%%%%%%%%%%%%%%%%%%%%%%%%%%%%%%
%%%%%%%%%%%%%%%%%%%%%%%%%%%%%%%%%%%%%%%%%%%
% Tema 1: Funciones vectoriales y curvas
\item
    Dada una trayectoria $\func{\alpha}{\R}{\R^3}$, que tiene tres derivadas, se define la función
    \[
        \funcion{u}{\R}{\R}{t}{u(t)=
            \alpha(t)\cdot(\alpha'(t)\times \alpha''(t)).
        }
    \]
    Comprobar que $u$ es derivable y que
    \[
        u'=\alpha\cdot(\alpha'\times \alpha''').
    \]

\begin{respuesta}
    Dado que $f$ tiene tres derivadas, entonces $f'$ y $f''$ son derivables, por lo tanto, $f'\times f''$ es derivable y por lo tanto $f\cdot(f'\times f'')$ es derivable, es decir, $u$ es derivable. Además, se tiene que:
    \begin{align*}
        u' 
            & = (\alpha\cdot(\alpha'\times \alpha''))'\\
            & = \alpha'\cdot(\alpha'\times \alpha'')+\alpha\cdot(\alpha'\times \alpha'')'\\
            & = 0 + \alpha\cdot(\alpha''\times \alpha'' + \alpha'\times \alpha''')\\
            & = \alpha\cdot(0 + \alpha'\times \alpha''')\\
            & = \alpha\cdot(\alpha'\times \alpha''').\qedhere
    \end{align*}
\end{respuesta}

%%%%%%%%%%%%%%%%%%%%%%%%%%%%%%%%%%%%%%%%%%%
%%%%%%%%%%%%%%%%%%%%%%%%%%%%%%%%%%%%%%%%%%%
%%%%%%%%%%%%%%%%%%%%%%%%%%%%%%%%%%%%%%%%%%%
% Tema 1: Funciones vectoriales y curvas
\item
    Dada una trayectoria $\func{\alpha}{I}{\R^3}$, se define
    \[
        \funcion{\beta}{I}{R^3}{t}{\alpha(t)\times \alpha'(t).}
    \]
    Comprobar que $\beta$ es derivable y hallar su derivada.
% \begin{respuesta}
%     Se tiene que $\alpha$ y $\alpha'$ son derivables, por lo tanto, $\beta$ es derivable. Además, se tiene que
%     \[
%         \beta'(t) = (\alpha(t) \times \alpha'(t))' = \alpha'(t) \times \alpha'(t) + \alpha(t) \times \alpha''(t) = \alpha(t) \times \alpha'(t),
%     \]
%     pues $\alpha'(t) \times \alpha'(t) = 0$.
% \end{respuesta}    

%%%%%%%%%%%%%%%%%%%%%%%%%%%%%%%%%%%%%%%%%%%
%%%%%%%%%%%%%%%%%%%%%%%%%%%%%%%%%%%%%%%%%%%
%%%%%%%%%%%%%%%%%%%%%%%%%%%%%%%%%%%%%%%%%%%
% Tema 1: Funciones vectoriales y curvas
\item 
    Sean $\func{\alpha,\beta,\gamma}{I}{\R^3}$ trayectorias derivables. Comprobar que $\alpha \cdot (\beta \times \gamma)$ es derivable y que
    \[
        (\alpha \cdot (\beta \times \gamma))' = \alpha \cdot (\beta \times \gamma') + \alpha \cdot (\beta' \times \gamma) + \alpha' \cdot (\beta \times \gamma).
    \]
    
\begin{respuesta}
    Dado que $\beta$ y $\gamma$ son derivables, $\beta\times \gamma$ también lo es. Además, dado que $\alpha$ es derivable, se tiene que $\alpha \cdot (\beta \times \gamma)$ es derivable y se tiene:
    \begin{align*}
        (\alpha \cdot (\beta \times \gamma))' 
            &= \alpha' \cdot( \beta \times \gamma) + \alpha \cdot (\beta \times \gamma)'\\
            &= \alpha' \cdot (\beta \times \gamma) + \alpha \cdot (\beta' \times \gamma + \beta \times \gamma')\\
            &= \alpha \cdot (\beta \times \gamma') + \alpha \cdot (\beta' \times \gamma) + \alpha' \cdot (\beta \times \gamma).\qedhere
    \end{align*}
\end{respuesta}

%%%%%%%%%%%%%%%%%%%%%%%%%%%%%%%%%%%%%%%%%%% FC
%%%%%%%%%%%%%%%%%%%%%%%%%%%%%%%%%%%%%%%%%%%
%%%%%%%%%%%%%%%%%%%%%%%%%%%%%%%%%%%%%%%%%%%
% Tema 1: Funciones vectoriales y curvas
\item
    Dada la función
    \[
        \funcion{\alpha}{\R}{\R^3}{t}{(\cos t , 1, -2t)}
    \]
    Si existe, calcule la integral.
    \[
        \int_0^{\pi} \alpha(t) dt
    \]
\begin{respuesta}
    Para que $\alpha(t)=\l(\alpha_1(t),\alpha_2(t).\alpha_3(t)\r)$ con $t \in \R$ sea integrable, sus campos escalares $\alpha_1(t)$, $\alpha_2(t)$ y $\alpha_3(t)$ deben ser integrables. Por ser funciones continuas. son integrables, entonces:
    \[
    \begin{aligned}
        \int_0^{\pi} \alpha_1(t) dt &= \int_0^{\pi} \cos t dt\\
        &= \sen t |_0^{\pi}\\
        &= 0-0\\
        &=0,\\
    \end{aligned}
    \]
    \[
    \begin{aligned}
        \int_0^{\pi} \alpha_2(t) dt &= \int_0^{\pi} dt\\
        &= t |_0^{\pi}\\
        &= \pi-0\\
        &=\pi,\\
    \end{aligned}
    \]
    \[
    \begin{aligned}
        \int_0^{\pi} \alpha_3(t) dt &= \int_0^{\pi} 2t dt\\
        &= t^2 |_0^{\pi}\\
        &= \pi^2-0^2\\
        &=\pi^2.\\
    \end{aligned}
    \]
    
    Por lo tanto, $\int_0^{\pi} \alpha(t) dt$ existe y 
    \[
        \int_0^{\pi} \alpha(t) dt=\l(0,,\pi,\pi^2\r).
    \]
\end{respuesta}

%%%%%%%%%%%%%%%%%%%%%%%%%%%%%%%%%%%%%%%%%%% JR
%%%%%%%%%%%%%%%%%%%%%%%%%%%%%%%%%%%%%%%%%%%
%%%%%%%%%%%%%%%%%%%%%%%%%%%%%%%%%%%%%%%%%%%
% Tema 1: Funciones vectoriales y curvas
\item Dada la trayectoria
\[
\funcion{\alpha}{\R}{\R^3}
{t}{\alpha(t)=\l(t-\cos(t),1-\cos(t),t\r)}
\]
y la función
\[
\funcion{u}{[-1,1]}{\R}
{z}{u(z)=\arcsen(z).}
\]
Utilizando la regla de la cadena para funciones vectoriales determine $(\alpha \circ u)'(z)$ para cada $z\in \l] -1,1 \r[$.

\begin{respuesta}
Notemos que $u$ es derivable en $(-1,1)$ con derivada
\[
    u'(z) = \dfrac{1}{\sqrt{1-z^2}}
\]
para cada $z\in(-1,1)$, además, $\mbox{rec}\,u= \l[-\dfrac{\pi}{2},\dfrac{\pi}{2}\r]$.

Por otro lado, las componentes de $\alpha$ son las funciones $\func{\alpha_1,\alpha_2,\alpha_3}{\R}{\R}$ definidas por
\[
\begin{array}{ccccc}
    \alpha_1(t) = t-\cos(t), & &  \alpha_2(t)=1-\cos(t), & & \alpha_3(t)=t\\
\end{array}
\]
para cada $t\in\R$. Las funciones $\alpha_1,\alpha_2,\alpha_3$ son derivables con derivada
\[
\begin{array}{ccccc}
    \alpha'_1(t) = 1+\sen(t), & &  \alpha'_2(t)=\sen(t), & & \alpha_3(t)=1\\
\end{array}
\]
para cada $t\in\R$, luego, $\alpha$ es derivable con derivada 
\[
\alpha'(t) = \l(1+\sen(t),\sen(t),1\r)
\]
para cada $t\in\R$, en particular, $\alpha$ es derivable en $u(z)$ para cada $z\in(-1,1)$.

Gracias a lo anterior, aplicando la regla de la cadena para funciones vectoriales se tiene que para cada $z\in (-1,1)$
\begin{align*}
    (\alpha \circ u)'(z) &= u'(z)\alpha'(u(z))\\
                         &= \dfrac{1}{\sqrt{1-z^2}}\l(1+\sen(\arcsen(z)),\sen(\arcsen(z)),1\r)\\
                         &= \l(\dfrac{1+z}{\sqrt{1-z^2}},\dfrac{z}{\sqrt{1-z^2}},\dfrac{1}{\sqrt{1-z^2}}\r).
\end{align*}

Por lo tanto, la derivada de $(\alpha \circ u)$ está dada por
\[
    (\alpha \circ u)'(z) \l(\dfrac{1+z}{\sqrt{1-z^2}},\dfrac{z}{\sqrt{1-z^2}},\dfrac{1}{\sqrt{1-z^2}}\r)
\]
para cada $z\in(-1,1)$.
\end{respuesta}


%%%%%%%%%%%%%%%%%%%%%%%%%%%%%%%%%%%%%%%%%%%
%%%%%%%%%%%%%%%%%%%%%%%%%%%%%%%%%%%%%%%%%%%
%%%%%%%%%%%%%%%%%%%%%%%%%%%%%%%%%%%%%%%%%%%
% Tema 1: Funciones vectoriales y curvas
\item
    Dada la trayectoria
    \[
        \funcion {\alpha}{\R}{\R^n}{x}{\alpha(x)=e^{2x}a+e^{-2x}b}
    \]
    con $a,b \in \R^n$ vectores fijos no nulos. Comprobar que $\alpha''(x)$ tiene la misma dirección que $\alpha(x)$, para todo $x\in\R$.
    
\begin{respuesta}
    Se tiene que dos vectores tiene la misma dirección si uno es múltiplo del otro, por lo tanto, para comprobar que $\alpha''(x)$ tiene la misma dirección que $\alpha(x)$, para todo $x\in\R$, procedemos a derivar $\alpha$ dos veces, así
    \begin{align*}
        \alpha'(x)&=2e^{2x}a-2e^{-2x}b,\\
        \alpha''(x)&=4e^{2x}a+4e^{-2x}b,
    \end{align*}
    para todo $x\in \R$. Dado que
    \[
        \alpha''(x) = 4\alpha(x)
    \]
    para todo $x\in \R$, se tiene que $\alpha(x)$ y $\alpha''(x)$ tienen la misma dirección.
\end{respuesta} 

%%%%%%%%%%%%%%%%%%%%%%%%%%%%%%%%%%%%%%%%%%%
%%%%%%%%%%%%%%%%%%%%%%%%%%%%%%%%%%%%%%%%%%%
%%%%%%%%%%%%%%%%%%%%%%%%%%%%%%%%%%%%%%%%%%%

% Tema 1: Funciones vectoriales y curvas
\item
    Determine el vector tangente y las ecuaciones de la recta tangente a la curva $C$ descrita por la trayectoria $\alpha$, en el punto $(1,-1,1)$, donde
    \[
        \funcion{\alpha}{\R}{\R^3}{t}{(e^t(\sen(t)+\cos(t)),e^t(\sen(t)-\cos(t)),e^t)}
    \]
    
\begin{respuesta}
    Para determinar el vector tangente es necesario conocer el valor de $t$ tal que $\alpha(t)=(1,-1,1)$, así, se tiene las siguientes ecuaciones
    \begin{align*}
        1&= e^t(\sen(t)+\cos(t)),\\
        -1&= e^t(\sen(t)-\cos(t)),\\
        1&= e^t.
    \end{align*}
    De la última ecuación se tiene que $t=0$ y se tiene consistencia con las otras ecuaciones, por lo tanto
    \[
        \alpha(0)=(1,-1,1).
    \]
    Ahora, gracias a que las componentes la función $\alpha$ son derivables en $\R$, $\alpha$ es derivable en $\R$, y por lo tanto
    \begin{align*}
        \alpha'(t)
            & = ((e^t(\sen(t)+\cos(t)))', (e^t(\sen(t)-\cos(t)))', (e^t)') \\
            & = (2e^t\cos(t), 2e^t\sin(t), e^t).
    \end{align*}   
    Evaluando en $t=0$, se obtiene que el vector tangente a $C$ en el punto dado es $\alpha'(0)=(2,0,1)$.

    Finalmente, la recta tangente es
    \[
        L((1,-1,1);(2,0,1))=\{(1,-1,1)+t(2,0,1):t\in\R\}.\qedhere
    \]
\end{respuesta}

%%%%%%%%%%%%%%%%%%%%%%%%%%%%%%%%%%%%%%%%%%%
%%%%%%%%%%%%%%%%%%%%%%%%%%%%%%%%%%%%%%%%%%%
%%%%%%%%%%%%%%%%%%%%%%%%%%%%%%%%%%%%%%%%%%%
% Tema 1: Funciones vectoriales y curvas
\item
	Sean
	\[
	    I = [0, 4\pi]
	    \texty
	    J = \l[-\frac{7\pi}{2}, \frac{\pi}{2}\r].
	\]
	Se definen las trayectorias $\alpha$ y $\beta$ por
    \[
        \funcion{\alpha}{I}{\R^3}
        {t}{\alpha(t)=(\sen(t), \cos(t), t)}
	\texty
        \funcion{\beta}{J}{\R^3}
        {t}{\beta(t) = \l(\cos(t), \sen(t), \frac{\pi}{2} - t\r)}.
	\]
    Muestre que $\alpha$ y $\beta$ son dos trayectorias equivalentes.
    
\begin{respuesta}
    Definamos la función $u$ por
    \[
        \funcion{u}{I}{J}{t}{u(t) = \frac{\pi}{2} - t}
	\]
    nótese que $\img(u)=J$, es decir, $u$ es sobreyectiva, además, $u$ es derivable y su derivada es no nula para todo $t\in I$. Luego, para cada $t\in I$,
    \begin{align*}
        \beta(u(t)) 
            & = \l(\cos(u(t)), \sen(u(t)), u(t)\r)\\
            & =\l(\cos{\l(\frac{\pi}{2} - t\r)}, \sen{\l(\frac{\pi}{2} - t\r)}, \frac{\pi}{2} - \l(\frac{\pi}{2} - t\r)\r)\\
            & = (\sen{t}, \cos{t}, t) = \alpha(t).
    \end{align*}
    Así, se tiene que ambas trayectorias son equivalentes.
\end{respuesta}

%%%%%%%%%%%%%%%%%%%%%%%%%%%%%%%%%%%%%%%%%%% PE
%%%%%%%%%%%%%%%%%%%%%%%%%%%%%%%%%%%%%%%%%%%
%%%%%%%%%%%%%%%%%%%%%%%%%%%%%%%%%%%%%%%%%%%
% Tema 1: Funciones vectoriales y curvas
\item 
    Encontrar una función $\func{\alpha}{\l]0,\infty\r[}{\R^n}$ continua, tal que
    \[
        \alpha(x)=xe^xA+\dfrac{1}{x}\int_1^x \alpha(t)dt
    \]
    para todo $x>0$, con $A\in\R^n$.

\begin{respuesta}
    Es equivalente escribir
    \begin{equation}
        \label{eq1}
        x\alpha(x)-x^2e^xA=\int_1^x \alpha(t)dt
    \end{equation}
    para todo $x>0$, con $A\in\R^n$.
    Por el primer teorema fundamental del cálculo se tiene que
    \[
        \alpha(x)+x\alpha'(x)-x^2e^xA-2xe^xA=\alpha(x)
    \]
    para todo $x>0$, simplificando la expresión anterior se obtiene
    \begin{align*}
        \alpha'(x)&=xe^xA+2e^xA\\
        &=(xe^x+e^x)A+e^xA\\
        &=(xe^x)'A+e^xA
    \end{align*}
    para $x>0$. De donde
    \[
        \alpha(x)=xe^xA+e^xA+k
    \]
    para $x>0$ y $k\in\R^n$.
    
    Reemplazando $\alpha(x)$ en (\ref{eq1}) y por el segundo teorema fundamental del cálculo, se tiene
    \begin{align*}
        x^2e^xA+xe^xA+kx-x^2e^xA&=\int_1^x \l(te^tA+e^tA+k\r) dt\\
        xe^xA+kx&=\l.\l(te^tA+kt\r)\r|_1^x\\
        xe^xA+kx&=xe^xA+kx-eA-k
    \end{align*}
    de aquí $k=-eA$.
    
    Finalmente se tiene que
    \[
        \alpha(x)=xe^xA+e^xA-eA
    \]
    para todo $x>0$ y $A\in\R^n$.
\end{respuesta}

%%%%%%%%%%%%%%%%%%%%%%%%%%%%%%%%%%%%%%%%%%%
%%%%%%%%%%%%%%%%%%%%%%%%%%%%%%%%%%%%%%%%%%%
%%%%%%%%%%%%%%%%%%%%%%%%%%%%%%%%%%%%%%%%%%%
% Tema 2: Movimiento en el espacio
\item 
    En cada uno de los literales siguientes, $\alpha(t)$ representa el vector posición en el instante $t$ correspondiente a una partícula que se mueve sobre una curva $C$. En cada caso, determinar los vectores velocidad $v(t)$ y aceleración $a(t)$, para cada $t\in\R$.	
    \begin{enumerate}
    \item 
        \[
            \funcion{\alpha}{\R}{\R^3}{t}{(\cos(t),\sen(t),e^t)}
        \]
    \item
        \[
            \funcion{\alpha}{\R}{\R^3}{t}{t\veci+\sen(t)\vecj+(1-\cos(t))\veck}
        \]
    \item
    	\[
        	\funcion{\alpha}{\R}{\R^3}{t}{(t-\sen(t))\veci+(1-\cos(t))\vecj+4\sen\l(\frac{t}{2}\r)\veck}
       	\]
    \end{enumerate}

\begin{respuesta}
    \begin{enumerate}
    \item
        Tenemos que
        \[
            v(t)=\alpha'(t)=(-\sen(t),\cos(t),e^t)
            \texty
            a(t)=\alpha''(t)=(-\cos(t),-\sen(t),e^t).
        \]
    \item
        Tenemos que
        \[
            v(t)=\alpha'(t)=\veci +\cos(t)\vecj +\sen(t)\veck
            \texty
            a(t)=\alpha''(t)= -\sin(t)\vecj +\cos(t)\veck.
        \]
    \item
    	Tenemos que
        \[
            v(t)=\alpha'(t)=(1-\cos(t))\veci +\sen(t)\vecj +2\cos\l(\frac{t}{2}\r)\veck
        \]
        y
        \[
            a(t)=\alpha''(t)= \sen(t)\veci +\cos(t)\vecj -\sin\l(\frac{t}{2}\r)\veck. \qedhere
        \]
    \end{enumerate}
\end{respuesta}

%%%%%%%%%%%%%%%%%%%%%%%%%%%%%%%%%%%%%%%%%%%
%%%%%%%%%%%%%%%%%%%%%%%%%%%%%%%%%%%%%%%%%%%
%%%%%%%%%%%%%%%%%%%%%%%%%%%%%%%%%%%%%%%%%%%
% Tema 2: Movimiento en el espacio
\item
	Dada la trayectoria
    \[
    	\funcion{\alpha}{\R}{\R^3}{t}{3t\cos(t) \ \veci+3t\sen(t)\   \vecj\ + 4t\ \veck,}
    \]
    la cual representa el vector posición de un movimiento, determinar los vectores velocidad $v(t)$ y aceleración $a(t)$ y la rapidez $\rho(t)$, para cada $t\in\R$.
    
% \begin{respuesta}    
%     El correspondiente vector velocidad está dado por la derivada de la función posición
%     \[
%         v(t)=\alpha'(t)=(3\cos(t)-3t\sen(t))\ \veci+(3\sen(t)+3t\cos(t))\ \vecj+4\ \veck,
%      \]
%     para cada $t\in\R$. Así, la rapidez para cada $t\in\R$ es
%     \begin{align*}
%     	\rho(t)
%     	    & =\| v(t) \| \\
%             & =\sqrt{(3\cos(t)-3t\sen(t))^2+(3\sen(t)+3t\cos(t))^2+4^2}\\
%             & =\sqrt{9+9t^2+16}\\
%             & =\sqrt{25+9t^2}.
% 	\end{align*} 
%     Por otro lado, el vector aceleración es la derivada del vector velocidad, así
% 	\[
%         a(t)=\alpha''(t)=(-6 \sen(t)-3t \cos(t))\ \veci+(6 \cos(t) -3t \sen(t))\ \vecj +0\veck
%     \]
%     para cada $t\in\R$.
% \end{respuesta}


%%%%%%%%%%%%%%%%%%%%%%%%%%%%%%%%%%%%%%%%%%%
%%%%%%%%%%%%%%%%%%%%%%%%%%%%%%%%%%%%%%%%%%%
%%%%%%%%%%%%%%%%%%%%%%%%%%%%%%%%%%%%%%%%%%%
% Tema 2: Movimiento en el espacio
\item 
    Considere las trayectorias $\func{\alpha,\beta}{\l[0,+\infty\r[}{\R^3}$, definidas por $\alpha(t)=(\cos(t),\sen(t),t)$ y $\beta(t)=(1,0,t)$ para todo $t\in\l[0,+\infty\r[$. Estas trayectorias representan los vectores posición en un instante $t$ de dos partículas que se mueven sobre dos curvas descritas por $\alpha$ y $\beta$ respectivamente. ¿En qué puntos, si los hay, estas partículas se encuentran? 

\begin{respuesta}
    Buscamos los elementos $t\in \l[0,+\infty\r[$ tal que
    \[
        \alpha(t)=\beta(t),
    \]
    es decir,
    \[
       (\cos(t),\sen(t),t)=(1,0,t).
    \]
    Así, se tienen las ecuaciones
    \[
        \cos(t)=1,
        \qquad
        \sen(t)=0
        \texty
        t=t.
    \]
    Por lo tanto, se tiene que las partículas se encuentran cuando
    \[
        t\in\{2k\pi:k\in\N\},
    \]
    y lo hacen en los puntos 
    \[
        \{(1,0,2k\pi):k\in\N\}.\qedhere
    \]
\end{respuesta}

%%%%%%%%%%%%%%%%%%%%%%%%%%%%%%%%%%%%%%%%%%% FC
%%%%%%%%%%%%%%%%%%%%%%%%%%%%%%%%%%%%%%%%%%%
%%%%%%%%%%%%%%%%%%%%%%%%%%%%%%%%%%%%%%%%%%%

% Tema 2: Movimiento en el espacio
\item
Suponga que no conocemos la trayectoria de un planeador, solo su vector de
aceleración, que está dada por $a(t)= \l(-3 \cos t, -3 \sen t, 2 \r)$, para $t\in \l[0,+\infty\r[$. También sabemos que inicialmente
(en el instante $t=0$), el planeador partió del punto $(3, 0, 0)$ con una velocidad inicial $v(0)=\l(0,3,0\r)$. Obtenga la posición del planeador como una función de $t$.
    
\begin{respuesta}
    Sea $\func{r}{\l[0,+\infty\r[}{\R^3}$ la posición de la partícula. Conocemos que la velocidad $v$ es la derivada de la posición $r$ y la aceleración $a$, la derivada de $v$. Entonces, teniendo la función aceleración, la integramos dos veces para obtener la posición, entonces
    \[
        a(t)=\dfrac{dv}{dt}(t)=\dfrac{d}{dt}\dfrac{dr}{dt}(t)= \l(-3 \cos t, -3 \sen t, 2 \r)
    \]
    para $t\in\l[0,+\infty\r[$, con condiciones iniciales:
    \[
        v(0)= \l(0,3,0\r)
        \texty
        r(0)= \l(3,0,0\r).
    \]
    Al integrar ambos lados de esta expresión tenemos
    \[
        v(t)= \l(-3 \sen t+C_1, 3 \cos t+C_2, 2t +C_3 \r),
    \]
    utilizando la condición inicial $v(0)= \l(0,3,0\r)$ se determina que $C_1=0$, $C_3=0$ y $C_3=0$.

    Luego, la velocidad del planeador es
    \[
        v(t)= \l(-3 \sen t, 3 \cos t, 2t \r)
    \]
    para todo $t\in\l[0,+\infty\r[$. Al integrar la velocidad $v$ se obtiene la posición del planeador
    \[
        r(t)= \l(3 \cos t + C_4, 3 \sen t + C_5, t^2 + C_6 \r),
    \]
    utilizando la condición inicial $r(0)= \l(3,0,0\r)$ de donde, $C_4=0$, $C_5=0$ y $C_6=0$.

    Por lo tanto la posición del planeador es
    \[
        r(t)= \l(3 \cos t, 3 \sen t, t^2\r)
    \]
    para todo $t\in\l[0,+\infty\r[$.
\end{respuesta}

%%%%%%%%%%%%%%%%%%%%%%%%%%%%%%%%%%%%%%%%%%% JR
%%%%%%%%%%%%%%%%%%%%%%%%%%%%%%%%%%%%%%%%%%%
%%%%%%%%%%%%%%%%%%%%%%%%%%%%%%%%%%%%%%%%%%%
% Tema 2: Movimiento en el espacio
\item Un cohete descompuesto se mueve de acuerdo con la trayectoria
\[
\funcion{\alpha}{[0,+\infty[\setminus \lbrace 0 \rbrace}{\R^3}
{t}{\alpha(t)=\l(e^{2t},3t^3-2t,t-\dfrac{1}{t}\r)}
\]
con la esperanza de llegar a una estación de reparación ubicada en el punto $(7e^4,35,5)$. El tiempo se mide en minutos y las coordenadas espaciales se miden en millas. Los motores del cohete se apagan de súbito a los dos minutos, si sobre el cohete no actúa ninguna fuerza externa, ¿llegará el cohete con su solo impulso a la estación de reparación?

\begin{respuesta}
Notemos que $\alpha$ es derivable con derivada
\[
\alpha'(t) = \l(2e^{2t},9t^2-2,1+\dfrac{1}{t^2}\r)
\]
para cada $t\in \l] 0,\infty \r[ $. Entonces, a los dos minutos, el cohete se encuentra en la posición $\alpha(2)=\l(e^4,20,\dfrac{3}{2}\r)$ con velocidad $\alpha'(2)=\l(2e^4,34,\dfrac{5}{4}\r)$. 

Ya que el cohete deja de funcionar a los dos minutos, este deberá continuar por la trayectoria de la recta tangente a $\alpha$ en 2, es decir, la recta
\begin{align*}
    L(\alpha(2),\alpha'(2)) &= \l\{\l(e^4,20,\dfrac{3}{2}\r) + (t-2)\l(2e^4,34,\dfrac{5}{4}\r)\,:\,t \geq 2\r\}\\
    &= \l\{\l((2t-3)e^4, 34t-48,\dfrac{5}{4}t-1\r)\,:\,t \geq 2\r\}\\
\end{align*}
Si el cohete logra alcanzar la estación de reparación debe tenerse que $(7e^4,35,5)\in L(\alpha(2),\alpha'(2))$, es decir, se tiene
\[
(7e^4,35,5) = \l((2t-3)e^4, 34t-48,\dfrac{5}{4}t-1\r)
\]
para algún $t \geq 2$. En particular tenemos
\[
7e^4 = (2t-3)e^4
\]
lo que se verifica cuando $t=5$. Notemos ahora que si $t=5$ el cohete se encontrará en la posición $\l(7e^4,122,\dfrac{21}{4}\r)\neq(7e^4,35,5)$.

Por lo tanto, el cohete no logrará llegar hasta la estación de reparación.
\end{respuesta}

%%%%%%%%%%%%%%%%%%%%%%%%%%%%%%%%%%%%%%%%%%% JO
%%%%%%%%%%%%%%%%%%%%%%%%%%%%%%%%%%%%%%%%%%%
%%%%%%%%%%%%%%%%%%%%%%%%%%%%%%%%%%%%%%%%%%%

% Tema 2: Movimiento en el espacio
\item 
    Sea la trayectoria 
    \[
        \funcion{\alpha}{[-1,3]}{\R^2}{t}{(t,\sen(\pi t)).}
    \]
    Si $\alpha$ describe la trayectoria de una mosca, halle la recta tangente a su trayectoria en $t=1$. Interprete este resultado. Encuentre el desplazamiento total de la mosca.

%     \begin{respuesta}
%         Sea $L$ la recta tangente a la curva de $\alpha$ en $t=1$. Así, $\alpha(1)=(1,0)\in L$; además, la dirección de $L$ coincide con $\alpha'(1)$. Dado que cada componente de $\alpha$ es derivable, $\alpha$ es derivable y 
%         \[
%             \alpha'(t)=(1,\pi \cos(\pi t))
%         \]
%         para cada $r\in \lop1,3\rop$. En particular,
%         \[
%             \alpha'(1)=(1,-\pi).
%         \]
%         Por lo tanto, $L=L((1,0);(1,-\pi))$; es decir,
%         \[
%             L=\{(1,0)+\lambda (1,-\pi): \lambda\in \R \}.
%         \]
%         Este resultado indica que la velocidad de la mosca tiene dirección $(1,-\pi)$ en el punto $(1,0)$. Finalmente el desplazamiento total de la mosca está dado por la norma de la posición final menos la posición inicial:
%         \[
%             \|\alpha(3)-\alpha(-1)\|
%             =\|(3,0)-(-1,0)\|
%             =\|(4,0)\|
%             =4.\qedhere
%         \]
% \end{respuesta}

%%%%%%%%%%%%%%%%%%%%%%%%%%%%%%%%%%%%%%%%%%%
%%%%%%%%%%%%%%%%%%%%%%%%%%%%%%%%%%%%%%%%%%%
%%%%%%%%%%%%%%%%%%%%%%%%%%%%%%%%%%%%%%%%%%%
% Tema 2: Movimiento en el espacio
\item 
    Dado un  vector fijo no nulo $b\in\R^3$, se tiene un movimiento con vector posición definido por $\func{\alpha}{\R}{\R^3}$, que satisface
    \[
        \alpha'(t) = b\times \alpha(t)
    \]
    para todo $t\in\R$. Demostrar que el vector aceleración es siempre ortogonal a $b$.

\begin{respuesta}
    Calculemos el vector aceleración para $t\in\R$,
    \[
        a(t)=\alpha''(t)=(b\times \alpha(t))'=0+b\times \alpha'(t)=b\times \alpha'(t).
    \]
    Ahora, para comprobar que es ortogonal a $b$, notemos que
    \[
        b\cdot a(t) = b\cdot (b\times \alpha'(t)) = 0
    \]
    para todo $t\in\R$, así, el vector aceleración es siempre ortogonal a $b$.
\end{respuesta}





%%%%%%%%%%%%%%%%%%%%%%%%%%%%%%%%%%%%%%%%%%%
%%%%%%%%%%%%%%%%%%%%%%%%%%%%%%%%%%%%%%%%%%%
%%%%%%%%%%%%%%%%%%%%%%%%%%%%%%%%%%%%%%%%%%%
% Tema 2: Movimiento en el espacio
\item
    Dada la trayectoria
	\[
        \funcion{\alpha}{\R}{\R^3}{(t)}{b\cos(\omega t)\veci+b\sen(\omega t)\vecj +(c\omega t)\veck}
    \]
    donde $\omega$ es una constante positiva y $b,c\in \R$. Compruebe que los vectores velocidad $v(t)$ y aceleración $a(t)$ tienen longitud constante para todo $t\in \R$ y que
    \[
        \frac{\|v(t)\times a(t)\|}{\|v(t)\|^3}=\frac{b}{b^2+c^2},
    \]
    para todo $t\in \R$.
    
% \begin{respuesta}
%     Tenemos que
%     \[
%         v(t)=\alpha'(t)=-b\omega \sen(\omega t)\veci+b\omega\cos(\omega t)\vecj+ c\omega\veck
%     \]
%     y
%     \[
%         a(t)=\alpha''(t)=-b\omega^2 \cos(\omega t)\veci-b\omega^2\sen(\omega t)\vecj+0\veck,
%     \]
%     para todo $t\in\R$, entonces
%     \begin{align*}
%         \|v(t)\|
%             & =\sqrt{(-b\omega \sen(\omega t))^2+(b\omega\cos(\omega t))^2+ (c\omega)^2}\\
%             & =\omega \sqrt{ b^2+c^2},
%     \end{align*}
% 	y
% 	\begin{align*}
% 	    \|a(t)\|
% 	        & = \sqrt{(-b\omega^2 \cos(\omega t))^2+(-b\omega^2\sen(\omega t))^2} \\
% 	        & = b\omega^2,
% 	\end{align*}
% 	para todo $t\in\R$, y por lo tanto son constantes.
	
% 	Por otra parte, se tiene que
% 	\[
%         v(t)\times a(t)= bc\omega^3\sen(\omega t)\veci -bc\omega^3\cos(\omega t)\vecj+b^2\omega^3\veck
%     \]
%     por lo tanto
%     \begin{align*}
%         \|v(t)\times a(t)\|
%             & = \sqrt{(bc\omega^3\sen(\omega t))^2+(bc\omega^3\cos(\omega t))^2+(b^2\omega^3)^2}\\
%             & =b\omega^3\sqrt{b^2+c^2}.
%     \end{align*}
%     Así, tenemos que
%     \[
%         \frac{\|v(t)\times a(t)\|}{\|v(t)\|^3} =\frac{b\omega^3\sqrt{b^2+c^2}}{\omega^3(\sqrt[]{b^2+ c^2})^3}=\frac{b}{b^2+ c^2}
%     \]
%     para todo $t\in\R$.\qedhere
% \end{respuesta}

%%%%%%%%%%%%%%%%%%%%%%%%%%%%%%%%%%%%%%%%%%%
%%%%%%%%%%%%%%%%%%%%%%%%%%%%%%%%%%%%%%%%%%%
%%%%%%%%%%%%%%%%%%%%%%%%%%%%%%%%%%%%%%%%%%%
% Tema 2: Movimiento en el espacio
\item
    Compruebe que, para cualquier movimiento, el producto escalar de los vectores velocidad y aceleración es igual a la mitad de la derivada del cuadrado de la rapidez, es decir, dado un movimiento definido por la trayectoria $\func{\alpha}{I}{\R^n}$, se tiene que
	\[
        v(t)\cdot a(t)=\dfrac{1}{2} \dfrac{d}{dt} (\rho^2(t))
    \]
    para todo $t\in\R$.
    
\begin{respuesta}
   	Se conoce que $\|v(t)\|=\rho(t)$, para todo $t\in I$, por lo tanto
    \begin{align*}
        \rho^2(t) = v(t)\cdot v(t).
	\end{align*}
    Derivando esto respecto al tiempo, tenemos
    \begin{align*}
		\dfrac{d}{dt} (\rho^2(t))
		    & =\dfrac{d}{dt} (v(t) \cdot v(t)) \\
            & = v'(t) \cdot v(t) + v(t) \cdot v(t)'\\
            & = 2 v(t) \cdot v(t)',
	\end{align*}
	para todo $t\in I$. Por lo tanto, como la aceleración es la derivada del vector velocidad, 
	\[
        \dfrac{d}{dt} (\rho^2(t))=2v(t)\cdot a(t),
   	\]
	para todo $t\in I$; reordenamos la expresión y obtenemos
	\[
   		\dfrac{1}{2} \dfrac{d}{dt} (\rho^2(t))=v(t)\cdot a(t),
    \]
    para todo $t\in I$, como queríamos.
\end{respuesta}


%%%%%%%%%%%%%%%%%%%%%%%%%%%%%%%%%%%%%%%%%%%
%%%%%%%%%%%%%%%%%%%%%%%%%%%%%%%%%%%%%%%%%%%
%%%%%%%%%%%%%%%%%%%%%%%%%%%%%%%%%%%%%%%%%%%
% Tema 2: Movimiento en el espacio
\item 
    Dada la trayectoria 
	\[	
        \funcion{\alpha}{\l[0,+\infty\r[}{\R^n}{t}{\alpha(t)=c+\dfrac{1-e^{-kt}}{k}d,}
    \]
	donde $k>0$ y $c,d\in\R^n$.	Muestre que $\alpha$ verifica
    \[
        \alpha''=-k\alpha'.
    \]
	Hallar la relación de $c$ y $d$ con la posición inicial y calcular
    \[
    	\lim_{t\to+\infty} (\alpha(t) - \alpha(0))
    \]
    en términos de la velocidad inicial (a este límite se lo conoce como la velocidad total de desplazamiento de la partícula).

\begin{respuesta}
	Es fácil notar que $\alpha$ es infinitamente derivable y se tiene que
	\[
		\alpha'(t)=e^{-kt}d 
		\texty
		\alpha''(t)=-ke^{-kt}d
	\]
	para cada $t\in\R$. Así,
	\[
		\alpha''=-k\alpha'.
	\]
	Por otro lado, puesto que
	\[
		\alpha(0)=c\quad \text y \quad 
		\alpha'(0)=d,
	\]
	es decir, $c$ es la posición inicial y $d$ es la velocidad inicial. Además, se tiene que
	\[
		\lim_{t\to+\infty}( \alpha(t) - \alpha(0) )
		=\lim_{t\to+\infty} \l(c+\dfrac{1-e^{-kt}}{k}d\r) - c
		=\lim_{t\to+\infty} \dfrac{1-e^{-kt}}{k}d
		=\dfrac{1}{k}d,
	\]
	dado que $k>0$. Finalmente, la velocidad total de desplazamiento de la partícula es
	\[
		\lim_{t\to+\infty}( \alpha(t) - \alpha(0) )
		=\dfrac{1}{k}\alpha'(0). \qedhere
	\]
\end{respuesta}

%%%%%%%%%%%%%%%%%%%%%%%%%%%%%%%%%%%%%%%%%%%
%%%%%%%%%%%%%%%%%%%%%%%%%%%%%%%%%%%%%%%%%%%
%%%%%%%%%%%%%%%%%%%%%%%%%%%%%%%%%%%%%%%%%%%
% Tema 2: Movimiento en el espacio
\item
	Compruebe que, en un movimiento representado por la función $\func{\alpha}{I}{\R^3}$, si el vector aceleración $a(t)$ es cero para todo $t\in I$, el movimiento es rectilíneo.
	
\begin{respuesta}
    Tenemos que
    \[
    	\funcion{\alpha}{I}{\R^3}{t}{(\alpha_1(t),\alpha_2(t),\alpha_3(t)),}
    \]
    y dado que $a(t)=0$ para todo $t\in I$, se tiene que
    \[
        \alpha_1''(t)=0,
        \qquad
        \alpha_2''(t)=0
        \texty
        \alpha_3''(t)=0,
    \]
    y por lo tanto
    \[
        \alpha_1(t)=a_1+tb_1,
        \qquad
        \alpha_2(t)=a_2+tb_2
        \texty
        \alpha_3(t)=a_3+tb_3,
    \]
    donde $a_i,b_i\in\R$ para todo $i\in\{1,2,3\}$.
    Así, el vector posición es
    \[
        \alpha(t)=(a_1+tb_1,a_2+tb_2,a_3+tb_3)=a+tb
    \]
    para todo $t\in I$, donde $a=(a_1,a_2,a_3)$ y $b=(b_1,b_2,b_3)$, y por lo tanto el movimiento descrito por $r$ es rectilíneo.
\end{respuesta}

%%%%%%%%%%%%%%%%%%%%%%%%%%%%%%%%%%%%%%%%%%%
%%%%%%%%%%%%%%%%%%%%%%%%%%%%%%%%%%%%%%%%%%%
%%%%%%%%%%%%%%%%%%%%%%%%%%%%%%%%%%%%%%%%%%% PE
% Tema 3: Geometria de curvas
\item
    Verificar que la trayectoria 
    \[
        \funcion{\alpha}{\R}{\R^3}{t}{(\sen t^2,\cos t^2,e)}
    \]
    describe un movimiento planar.
    
\begin{respuesta}
    Para esto se debe verificar que la dirección del vector binormal es constante, entonces
    \begin{align*}
        \alpha'(t)&=(2t\cos t^2,-2t\sen t^2,0), \\
        \rho(t)&=2t
    \end{align*}
    de donde 
    \[
        T(t)=(\cos t^2,-\sen t^2,0)
    \]
    y
    \begin{align*}
        T'(t)&=(-2t\sen t^2,-2t\cos t^2,0)\\
        \|T'(t)\|&=2t\\
        N(t)&=(-\sen t^2,-\cos t^2,0).
    \end{align*}
    Finalmente
    \begin{align*}
        B(t)&=T(t) \times N(t)\\
        &=(\cos t^2,-\sen t^2,0) \times (-\sen t^2,-\cos t^2,0)\\
        &=(0,0,-1)
    \end{align*}
    con lo que se verifica que el movimiento descrito es planar. 
\end{respuesta}

%%%%%%%%%%%%%%%%%%%%%%%%%%%%%%%%%%%%%%%%%%%
%%%%%%%%%%%%%%%%%%%%%%%%%%%%%%%%%%%%%%%%%%%
%%%%%%%%%%%%%%%%%%%%%%%%%%%%%%%%%%%%%%%%%%%
% Tema 3: Geometria de curvas
\item
    Dada una curva $C$ en $\R^3$, con dos parametrizaciones equivalente
    \[
        \func{\alpha}{I}{\R^3}
        \texty
        \func{\beta}{J}{\R^3}
    \]
    con cambio de parámetro $\func{u}{I}{J}$ creciente, demostrar que los vectores principales no se ven afectados por el cambio de parámetro.
    
\begin{respuesta}
    Dado que $u$ es el cambio de parámetro, se tiene que
    \[
        \beta(u(t))=\alpha(t)
    \]
    para todo $t\in I$. Además, dado que $u$ es creciente, se tiene que $u'(t)\geq 0$ para todo $t\in\R$.
    
    Fijemos $t\in \R$ y llamemos $\tau=u(t)$. Calculemos los vectores principales en en el punto $\beta(\tau)$. Con la parametrización $\beta$, tenemos que
    \[
        T_\beta(\tau)= \frac{\beta'(\tau)}{\|\beta'(\tau)\|},
        \qquad
        N_\beta(\tau)= \frac{T_\beta'(\tau)}{\|T_\beta'(\tau)\|}
        \texty
        B_\beta(\tau)= T_\beta(\tau)\times N_\beta(\tau).
    \]
    Ahora, hallemos los vectores en el mismo punto pero con la parametrización $\alpha$, es decir en el punto $\beta(\tau)=\beta(u(t))=\alpha(t)$ y tenemos
    \[
        T_\alpha(t)= \frac{\alpha'(t)}{\|\alpha'(t)\|},
        \qquad
        N_\alpha(t)= \frac{T_\alpha'(\tau)}{\|T_\alpha'(t)\|}
        \texty
        B_\alpha(t)= T_\alpha(t)\times N_\alpha(t).
    \]
    Pero notemos que
    \begin{align*}
        T_\alpha(t)
            & = \frac{\alpha'(t)}{\|\alpha'(t)\|} = \frac{(\beta(u(t)))'}{\|(\beta(u(t)))'\|}\\[1mm]
            & = \frac{u'(t)\beta'(u(t))}{\|u'(t)\beta'(u(t))\|} = \frac{u'(t)\beta'(u(t))}{|u'(t)|\|\beta'(u(t))\|}\\[1mm]
            & = \frac{\beta'(u(t))}{\|\beta'(u(t))\|} = \frac{\beta'(\tau)}{\|\beta'(\tau)\|}\\[1mm]
            & = T_\beta(\tau),
    \end{align*}
    es decir, $T\alpha(t)=T_\beta(u(t))$. Por otro lado, 
    \begin{align*}
        N_\alpha(t)
            & = \frac{T_\alpha'(t)}{\|T_\alpha'(t)\|} = \frac{(T_\beta(u(t)))'}{\|(T_\beta(u(t)))'\|}\\[1mm]
            & = \frac{u'(t)T_\beta'(u(t))}{\|u'(t)T_\beta'(u(t))\|} = \frac{u'(t)T_\beta'(u(t))}{|\beta'(t)|\|T_\beta'(u(t))\|}\\[1mm]
            & = \frac{T_\beta'(u(t))}{\|T_\beta'(u(t))\|} = \frac{T_\beta'(\tau)}{\|T_\beta'(\tau)\|}\\[1mm]
            &= N_\beta(\tau).
    \end{align*}
    Finalmente
    \[
        B_\alpha(t)= T_\alpha(t)\times N_\alpha(t) 
        = T_\beta(\tau)\times N_\beta(\tau)=B_\beta(\tau),
    \]
    es decir,
    \[
        T_\alpha(t)= T_\beta(\tau),
        \qquad
        N_\alpha(t)= N_\beta(\tau)
        \texty
        B_\alpha(t)= B_\beta(\tau).\qedhere
    \]
\end{respuesta}

%%%%%%%%%%%%%%%%%%%%%%%%%%%%%%%%%%%%%%%%%%%
%%%%%%%%%%%%%%%%%%%%%%%%%%%%%%%%%%%%%%%%%%%
%%%%%%%%%%%%%%%%%%%%%%%%%%%%%%%%%%%%%%%%%%%
% Tema 3: Geometria de curvas
\item 
    Demostrar que el módulo de la componente normal del vector aceleración en $t$ es $\dfrac{\|v(t) \times a(t)\|}{\|v(t)\|}$ para cualquier curva.

\begin{respuesta}
    Sabemos que $a(t)$ se puede escribir como una combinación lineal de los vectores tangente y normal, es decir,
    \[
        a(t)=\alpha T(t)+\beta N(t),
    \]
    donde $\alpha,\beta\in\R$.
    Ahora, tenemos que
    \[
        v(t)\times a(t)=\alpha v(t)\times T(t)+\beta v(t)\times N(t).
    \]
    Notemos que, gracias a que $v(t)$ y $T(t)$ son paralelos, $v(t)\times T(t)=0$, por lo tanto
    \[
        v(t)\times a(t)=\beta v(t)\times N(t).
    \]
    Así, se tiene que
    \[
         \| v(t)\times a(t) \| = |\beta| \| v(t)\times N(t) \|,
    \]
    y por lo tanto
    \begin{align*}
        |\beta| 
            & = \frac{\| v(t)\times a(t) \|}{\| v(t)\times N(t) \|} \\
            & = \frac{\| v(t)\times a(t) \|}{\| v(t) \| \| N(t) \| \sen\l(\frac{\pi}{2}\r)} \\
            & = \frac{\| v(t)\times a(t) \|}{\| v(t) \| \| N(t) \|}\\
            & = \frac{\| v(t)\times a(t) \|}{\| v(t) \|}.
    \end{align*}
    Luego, la componente normal de la aceleración es
    \[
        \frac{\| v(t)\times a(t) \|}{\| v(t) \|} N(t)
    \]
\end{respuesta}

%%%%%%%%%%%%%%%%%%%%%%%%%%%%%%%%%%%%%%%%%%%
%%%%%%%%%%%%%%%%%%%%%%%%%%%%%%%%%%%%%%%%%%%
%%%%%%%%%%%%%%%%%%%%%%%%%%%%%%%%%%%%%%%%%%%
% Tema 3: Geometria de curvas
\item 
    Dado el movimiento representado por 
	\[
		\funcion{\alpha}{\R}{\R^3}{t}{(\sen(t),\cos(t),t^2)}.
	\]
	Verifique que se cumple que
	\[
        a(t)=\rho'(t) T(t)+\rho(t) T' (t).
	\]
	para cada $t\in\R$.

\begin{respuesta}
	Se tiene que
	\begin{align*}
        v(t)
        &=\alpha'(t)=(\cos(t),-\sen(t),2t),
        &\rho(t)&=\|v(t)\|=\sqrt{1+4t^2},\\
        a(t)&=\alpha''(t)=(-\sen(t),-\cos(t),2),
        &\rho'(t)&=\dfrac{4t}{\sqrt{1+4t^2}}
	\end{align*}
	para cada $t\in\R$. Puesto que $\alpha'(t)\neq0$ para cada $t\in\R$, se sigue que
	\begin{align*}
		T(t)=\l(\dfrac{\cos(t)}{\sqrt{1+4t^2}},\dfrac{-\sen(t)}{\sqrt{1+4t^2}},\dfrac{2t}{\sqrt{1+4t^2}}\r)
	\end{align*}
	y además, 
	\[
		T'(t)
		=\l(\dfrac{-(1+4t^2) \sen(t) -4t\cos(t))}{(1 + 4 t^2)^\frac{3}{2}},
		\dfrac{ 4t\sen(t)-(1+4t^2) \cos(t))}{(1 + 4 t^2)^\frac{3}{2}},
		\dfrac{ 2}{(1 + 4 t^2)^\frac{3}{2}}\r)
	\]
	para todo $t\in\R$. Finalmente, para cada $t\in\R$, se tiene que
	\begin{align*}
		\rho'(t) T(t)+&\rho(t) T' (t)=\\
		&= \dfrac{4t}{\sqrt{1+4t^2}}\l(\dfrac{\cos(t)}{\sqrt{1+4t^2}},\dfrac{-\sen(t)}{\sqrt{1+4t^2}},\dfrac{2t}{\sqrt{1+4t^2}}\r)\\
	    &\phantom{=}+\sqrt{1+4t^2}\l(\dfrac{-(1+4t^2) \sen(t) -4t\cos(t))}{(1 + 4t^2)^\frac{3}{2}},
			\dfrac{ 4t\sen(t)-(1+4t^2) \cos(t))}{(1 + 4 t^2)^\frac{3}{2}},
			\dfrac{ 2}{(1 + 4 t^2)^\frac{3}{2}}\r)\\[2mm]
		&= 4t\l(\dfrac{\cos(t)}{1+4t^2},\dfrac{-\sen(t)}{1+4t^2},\dfrac{2t}{1+4t^2}\r)\\
		&\phantom{=}+\l(\dfrac{-(1+4t^2) \sen(t) -4t\cos(t))}{1 + 4t^2},
			\dfrac{ 4t\sen(t)-(1+4t^2) \cos(t))}{1 + 4t^2},
			\dfrac{ 2}{1 + 4t^2}\r)\\[2mm]
		&=\l(\dfrac{-(1+4t^2) \sen(t))}{1 + 4t^2},
			\dfrac{-(1+4t^2) \cos(t))}{1 + 4t^2},
			\dfrac{8t^2+2}{1 + 4t^2}\r)\\[1mm]
		&=(-\sen(t),-\cos(t),2)\\
		&=a(t).
	\end{align*}
	Así, se concluye que
	\[
		a(t)=\rho'(t) T(t)+\rho(t) T' (t)
	\]
	para cada $t\in\R$.
\end{respuesta}

%%%%%%%%%%%%%%%%%%%%%%%%%%%%%%%%%%%%%%%%%%%
%%%%%%%%%%%%%%%%%%%%%%%%%%%%%%%%%%%%%%%%%%%
%%%%%%%%%%%%%%%%%%%%%%%%%%%%%%%%%%%%%%%%%%%
% Tema 3: Geometria de curvas
\item 
	Calcule la longitud de $C$ representada por la trayectoria $\func{\alpha}{[0,2\pi]}{\R^3}$ donde $\alpha$ se define como en el ejercicio anterior. Además, calcular su vector curvatura y su curvatura.

\begin{respuesta}
	Utilizando los cálculos anteriores y la definición, se tiene que
	\[
		\ell(C)=\int_0^{2\pi}\rho(t)dt=\int_0^{2\pi}\sqrt{1 + 4t^2}dt=
		\pi \sqrt{1 + 16 \pi^2} + \dfrac{1}{4} \senh^{-1}(4\pi)\approx 40.4.
	\]
	Por otro lado, se tiene que
	\begin{align*}
		K(t)=\dfrac{T'(t)}{\rho(t)}
		&=\dfrac{1}{\sqrt{1+4t^2}}\l(\dfrac{-(1+4t^2) \sen(t) -4t\cos(t))}{(1 + 4 t^2)^\frac{3}{2}},
		\dfrac{ 4t\sen(t)-(1+4t^2) \cos(t))}{(1 + 4 t^2)^\frac{3}{2}},
		\dfrac{ 2}{(1 + 4 t^2)^\frac{3}{2}}\r)\\[2mm]
		&=\l(\dfrac{-(1+4t^2) \sen(t) -4t\cos(t))}{(1 + 4 t^2)^2},
		\dfrac{ 4t\sen(t)-(1+4t^2) \cos(t))}{(1 + 4 t^2)^2},
		\dfrac{ 2}{(1 + 4 t^2)^2}\r)
	\end{align*}
	y
	\begin{align*}
		\kappa(t)=\dfrac{\|T'(t)\|}{\rho(t)}
		&=\dfrac{1}{(1+4t^2)^2}
		\sqrt{(-(1+4t^2) \sen(t) -4t\cos(t)))^2
				+(4t\sen(t)-(1+4t^2) \cos(t)))^2
				+4}\\
		&=\dfrac{1}{(1+4t^2)^2}
			\sqrt{
			(1+4t^2)^2 +16t^2+4}\\
		&=\dfrac{1}{(1+4t^2)^2}
			\sqrt{
				(4 t^4 + 16 t^2 + 5}
	\end{align*}
	para todo $t\in\lcl0,2\pi\rcl$.
\end{respuesta}

%%%%%%%%%%%%%%%%%%%%%%%%%%%%%%%%%%%%%%%%%%%
%%%%%%%%%%%%%%%%%%%%%%%%%%%%%%%%%%%%%%%%%%%
%%%%%%%%%%%%%%%%%%%%%%%%%%%%%%%%%%%%%%%%%%%
% Tema 3: Geometria de curvas
\item    
	Dada la trayectoria 
 	\[
        \funcion{\alpha}{\left[0,2\pi \right]}{\R^2}{t}{a(1-\cos (t)) \veci+a(t-\sen (t)) \vecj}
    \]
	con $a>0$. Hallar 
	\begin{enumerate}
	    \item La longitud de la curva $C$ descrita por $\alpha$. (Distancia) 
	    \item Desplazamiento entre los puntos inicial y final de la curva $C$.
	\end{enumerate}
	
	\begin{respuesta}
	
	\begin{enumerate}
	  \item
	 	El vector velocidad correspondiente es
    	\[
       		v(t)=\alpha'(t)= a \sen(t) \veci+(a-a \cos (t)) \vecj,
        \]
    	para todo $t\in[0,2\pi]$ y la rapidez asociada es
    	\[
            \rho(t) = \| v(t) \| = a \sqrt{2-2\cos(t)} 
      	\]
    	para todo $t\in\left[0,2\pi\right]$. Por la definición de la longitud de arco de la curva descrita por $r$, se tiene que
    	\begin{align*}
    		\ell(C)
                & =\int_0^{2\pi} \|v(t)\|\ dt\\
                & =\int_0^{2\pi}a \sqrt{2-2\cos(t)}\  dt\\
                & =8a.\qedhere
        \end{align*}
    \item
        El desplazamiento $\Delta r$ es la diferencia entre los puntos final e inicial, que son $\alpha(2\pi)$ y $\alpha(0)$ respectivamente, entonces
        \begin{align*}
            \Delta r&=a(0,2\pi)-a(0,0)\\
            &=a(0,2\pi)
        \end{align*}
        y su norma es $\| \Delta r\|=2a\pi$.
    \end{enumerate}
\end{respuesta}



%%%%%%%%%%%%%%%%%%%%%%%%%%%%%%%%%%%%%%%%%%%
%%%%%%%%%%%%%%%%%%%%%%%%%%%%%%%%%%%%%%%%%%%
%%%%%%%%%%%%%%%%%%%%%%%%%%%%%%%%%%%%%%%%%%%
% Tema 3: Geometria de curvas
\item 
    La trayectoria definida por
	\[
        \funcion{\alpha}{[0, 2\pi]}{\R^2}{t}
        {\alpha(t) = (a\sen{t}, b\cos{t})}
	\]
    donde $0 < b < a$ son números reales, es la parametrización de una elipse. Probar que la longitud $L$ de una elipse está dada por la integral
    \[
        L = 4a\int_{0}^{\pi/2}{\sqrt{1 - e^{2}\sen^{2}{t}}\ dt,}
    \]
    donde $e = \dfrac{\sqrt{a^{2} - b^{2}}}{a}$. Al número $e$ se lo llama excentricidad de la elipse.

% \begin{respuesta} 
%     Para calcular la longitud $L$ necesitamos determinar la norma de la derivada de $\alpha$. Notemos que $\alpha$ es derivable y que su derivada $\alpha'$ está dada por:
%     \[ 
%         \alpha'(t) = (a\cos{t}, -b\sen{t}),
%     \]
%     cuya norma es
%     \[
%         \| \alpha'(t) \| = \sqrt{a^2 \cos^{2}{t} + b^2 \sen^{2}{t}},
%     \]
%     para todo $t \in [0, 2\pi]$. Ahora, dado que la elipse considerada se encuentra centrada en el origen, la longitud $L$ se determina como 4 veces la longitud en un cuadrante, lo que equivale a calcular la integral para $0 \leq t \leq \pi/2$. Entonces,
%     \begin{align*}
%         L 
%             &= 4\int_{0}^{\pi/2}{\| \alpha'(t)\|}\ dt\\
%             &= 4\int_{0}^{\pi/2}{\sqrt{a^2 \cos^{2}{t} + b^2 \sen^{2}{t}}}\ dt\\
%             &= 4\int_{0}^{\pi/2}{\sqrt{a^2(1 - \sen^{2}{t}) + b^2 \sen^{2}{t}}}\ dt\\
%             &= 4\int_{0}^{\pi/2}{\sqrt{a^2 - (a^2 - b^2)\sen^{2}{t}}}\ dt\\
%             &= 4\int_{0}^{\pi/2}{\sqrt{a^2\left(1 - \frac{a^2 - b^2}{a^2}\sen^{2}{t}\right)}}\ dt\\
%             &= 4a\int_{0}^{\pi/2}{\sqrt{1 - \frac{a^2 - b^2}{a^2}\sen^{2}{t}}}\ dt\\
%             &= 4a\int_{0}^{\pi/2}{\sqrt{1 - e^2\sen^{2}{t}}}\ dt.\qedhere
%     \end{align*}
% \end{respuesta}

%%%%%%%%%%%%%%%%%%%%%%%%%%%%%%%%%%%%%%%%%%%
%%%%%%%%%%%%%%%%%%%%%%%%%%%%%%%%%%%%%%%%%%%
%%%%%%%%%%%%%%%%%%%%%%%%%%%%%%%%%%%%%%%%%%%
% Tema 3: Geometria de curvas
\item 
    Dada la curva definida por la trayectoria
    \[
        \funcion{\alpha}{\R}{\R^3}
        {t}{\alpha(t) = (e^{t}\cos{t}, e^{t}\sen{t}, e^{t})}
	\]
    Hallar la curvatura y la ecuación del plano osculador cuando $t = 0$.
        
\begin{respuesta}
    Nótese que $\alpha$ es al menos dos veces derivable. En efecto, sus dos primeras derivadas $\alpha'$ y $\alpha''$ están dadas por
    \[
        \alpha'(t) = (e^{t}\cos{t} - e^{t}\sen{t}, e^{t}\cos{t} + e^{t}\sen{t}, e^{t}),
    \]
    y
    \[
        \alpha''(t) = (-2e^{t}\sen{t}, 2e^{t}\cos{t}, e^{t}),
    \]
    para todo $t \in \mathbb{R}$. Además,
    \[
        \| \alpha'(t) \| = \sqrt{(e^{t}\cos{t} - e^{t}\sen{t})^{2} + (e^{t}\cos{t} + e^{t}\sen{t})^2 + e^{2t}} = \sqrt{3} e^{t},
    \]
    para todo $t\in\R$. En particular, en $t = 0$, se tiene que $\alpha'(0)= (1,1,1)$, $\alpha''(0) = (0,2,1)$ y $\|\alpha'(0)\| = \sqrt{3}$, por tanto, la curvatura de $\alpha$ en $t = 0$, es
    \begin{align*}
        \kappa(0) 
            &= \frac{\|\alpha'(0) \times \alpha''(0) \|}{\| \alpha'(0) \|^{3}}\\
            &= \frac{\|(1,1,1) \times (0,2,1) \|}{\sqrt{3}^{3}}\\
            &= \frac{\|(-1,-1,2)\|}{3\sqrt{3}}\\
            &= \frac{\sqrt{2}}{3}.
    \end{align*}
    Para determinar la ecuación del plano osculador, calculemos el vector unitario tangente $T(t)$ y el vector normal principal $N(t)$. El vector unitario tangente está dado por
    \[
        T(t) = \frac{1}{\| \alpha'(t)\|}\alpha'(t) = \frac{\sqrt{3}}{3}\alpha'(t),
    \]
    para todo $t\in\R$ y cuya derivada $T'$ está dada por
    \[
        T'(t) = \frac{\sqrt{3}}{3}(-\cos{t} - \sen{t}, \cos{t} - \sen{t}, 0), 
    \]
    con norma $\| T'(t)\| = \frac{\sqrt{6}}{3}$, para todo $t \in\R$. El vector normal principal, está dado por
    \[
        N(t) = \frac{1}{\|T'(t)\|}T'(t) = \frac{\sqrt{2}}{2}T'(t).
    \]
    Nótese que un vector perpendicular al plano osculador es $B(t) = T(t) \times N(t)$ para todo $t\in\R$, por lo tanto, en $t = 0$, 
    \[
        B(0) = T(0) \times N(0) 
        = \l(\frac{\sqrt{3}}{3}\alpha'(0)\r)\times \l(\frac{\sqrt{2}}{2}T'(0)\r)
        =\frac{\sqrt{6}}{6} \l( (1,1,1) \times (-1, 1, 0) \r)
        = \frac{\sqrt{6}}{6} (-1, -1, 2).
    \]
    Con esto, se tiene que la ecuación del plano está dada por
    \[
        B(0)\cdot (x,y,z) = B(0)\cdot \alpha(0),
    \]
    es decir,
    \[
        (-1, -1, 2)\cdot(x,y,z) =(-1, -1, 2)\cdot(1, 0, 1)
    \]
    así, la ecuación del plano es
    \[
        -x - y + 2z = -1.\qedhere
    \]
\end{respuesta}


%%%%%%%%%%%%%%%%%%%%%%%%%%%%%%%%%%%%%%%%%%%
%%%%%%%%%%%%%%%%%%%%%%%%%%%%%%%%%%%%%%%%%%%
%%%%%%%%%%%%%%%%%%%%%%%%%%%%%%%%%%%%%%%%%%%
% Tema 3: Geometria de curvas
\item 
    Si un punto se mueve por una curva, parametrizada por la función
    \[
        \func{\alpha}{I}{\R^3,}
    \]
    de manera que sus vectores velocidad y aceleración tienen siempre longitud constante, probar que la curvatura de la curva es constante en todos sus puntos.

% \begin{respuesta}
%     Tenemos que $\|v(t)\|=k_1$ y $\|a(t)\|=k_2$ para todo $t\in I$, donde $k_1$ y $k_2$ son constantes reales. Así
%     \[
%         T(t)=\frac{v(t)}{k_1},
%     \]
%     por lo tanto
%     \[
%         \kappa(t)=\frac{\|T'(t)\|}{k_1} = \frac{\|a(t)\|}{k_1^2} = \frac{k_2}{k_1^2},
%     \]
%     de donde se tiene que la curvatura es constante.
% \end{respuesta}

%%%%%%%%%%%%%%%%%%%%%%%%%%%%%%%%%%%%%%%%%%%
%%%%%%%%%%%%%%%%%%%%%%%%%%%%%%%%%%%%%%%%%%%
%%%%%%%%%%%%%%%%%%%%%%%%%%%%%%%%%%%%%%%%%%%
% Tema 3: Geometria de curvas
\item
    Determinar el punto de la parábola de ecuación $y=x^2$ en el que la curvatura alcanza su valor máximo. 
    
\begin{respuesta}
    Definamos la trayectoria
    \[
        \funcion{\alpha}{\R}{\R^2}{t}{(t,t^2).}
    \]
    Tenemos que
    \[
        \kappa(t)=\frac{1}{\rho(t)}\|T'(t)\|,
    \]
    para todo $t\in\R$, donde
    \begin{align*}
        T'(t)
            & = \l(\frac{\alpha'(t)}{\|\alpha'(t)\|}\r)' \\
            & = \frac{\rho(t)\alpha''(t)-\rho'(t)\alpha'(t)}{\rho^3(t)}.
    \end{align*}
    Así,
    \[
        \kappa(t)=\frac{2}{\l(\sqrt{1+4t^2}\r)^3}
    \]
    para todo $t\in\R$.
    Ahora, buscamos $t\in\R$ de modo que $\kappa$ alcance su valor máximo. Notemos que $\kappa$ es diferenciable en $\R$, por lo tanto
    \[
        \kappa'(t)=-\frac{24t}{(1+4t^2)^\frac{5}{2}},
    \]
    para todo $t\in\R$, de donde se tiene que $\kappa'(t)=0$ si y solo si $t=0$, es decir que $t=0$ es un punto crítico para $\kappa$. Finalmente, notemos que $\kappa'(t)>0$ si $t<0$ y $\kappa'(t)<0$ si $t>0$, es decir, $\kappa$ alcanza su valor máximo cuando $t=0$.
\end{respuesta}

%%%%%%%%%%%%%%%%%%%%%%%%%%%%%%%%%%%%%%%%%%%
%%%%%%%%%%%%%%%%%%%%%%%%%%%%%%%%%%%%%%%%%%%
%%%%%%%%%%%%%%%%%%%%%%%%%%%%%%%%%%%%%%%%%%%
% Tema 3: Geometria de curvas
\item 
    Demostrar que en una circunferencia de radio 1 la curvatura es constante.

\begin{respuesta}
    Consideremos la trayectoria
    \[
        \funcion{\alpha}{[0,2\pi]}{\R^2}{t}{(\cos(t),\sin(t)).}
    \]
    Se tiene que
    \[
        \rho(t)=1
        \texty
        T(t)=(-\sen(t),\cos(t)),
    \]
    para todo $t\in[0,2\pi]$, por lo tanto
    \begin{align*}
        \kappa(t)
            & = \frac{1}{\rho(t)}\|T'(t)\| \\
            & = \|(-\cos(t),-\sen(t))\| \\
            & = 1,
    \end{align*}
    para todo $t\in[0,2\pi]$, lo que implica que la curvatura es constante.

    %%%%%%%%%%%%%%%%%%%%%%%%%%%%%%%%%%%%%%%%%%% JO
%%%%%%%%%%%%%%%%%%%%%%%%%%%%%%%%%%%%%%%%%%%
%%%%%%%%%%%%%%%%%%%%%%%%%%%%%%%%%%%%%%%%%%%
% Tema 5: Limites y diferenciabilidad
\item 
    Sea el campo escalar
    \[
        \funcion{f}{\R^2}{\R}{(x,y)}{xy^2.}
    \]
    Dado que $f$ es un polinomio, $f$ es continuo y se cumple que
    \[
        \Lim_{(x,y)\to (3,2)}f(x,y)=f(3,2)=12.
    \]
    Para $\epsilon=1$, determine explícitamente un $\delta>0$ tal que
    \[
    (\forall (x,y)\in\R^2)(\|(x,y)-(3,2)\|< \delta \Imp |xy^2-12|<1).
    \]

\begin{respuesta}    
    De la definición de norma, para cada $(x,y)\in\R^2$, se tiene que
    \[
        \|(x,y)-(3,2)\|=\sqrt{(x-3)^2+(y-2)},
    \]
    luego $\|(x,y)-(3,2)\|\geq |x-3|$ y $\|(x,y)-(3,2)\|\geq |y-2|$. Queremos acotar el término
    \[
        |xy^2-12|
    \]
    cerca del punto $(3,2)$. Con esto y dado que $(x,y)\to (3,2)$,  $|x-3|$ y $|y-2|$ deben ser términos pequeños, se tiene que para cada $(x,y)\in\R^2$, 
    \begin{equation}\label{eq:dtri}
        \begin{split}|xy^2-12|
        &=|(x-3)(y^2-4)+4x+3y^2-24|\\
        &=|(x-3)(y^2-4)+4(x-3)+3(y^2-4)|\\
        &\leq |x-3||y^2-4|+4|x-3|+3|y^2-4|\\
        &=|x-3||y-2||y+2|+4|x-3|+3|y-2||y+2|.
        \end{split}
    \end{equation}
    Así, debemos acotar cada uno de los términos anteriores. Como el $\delta$ debe ser pequeño, podemos suponer que $\delta\leq 1$. Entonces si $(x,y)\in\R^2$ tales que
    \[
        \|(x,y)-(3,2)\|<\delta\leq 1,
    \]
    se tiene que 
    \[
        |x-3|\leq \|(x,y)-(3,2)\|<\delta \yds 
        |y-2|\leq \|(x,y)-(3,2)\|<\delta.
    \]
    Ahora bien, como $|y-2|<1$, se tiene que $|y+2|<5$. Por todo esto,
    \[
        |x-3||y-2||y+2|+4|x-3|+3|y-2||y+2|
        <(\delta)( \delta) 5+4\delta +3 (\delta) 5
        =5\delta^2+19\delta. 
    \]
    De la última desigualdad y \eqref{eq:dtri}, es suficiente encontrar $\delta>0$, tal que
    \[
        5\delta^2+19\delta\leq 1
    \]
    lo qué ocurre cuando 
    \[
        0<\delta\leq \dfrac{ \sqrt{381} - 19}{10}.
    \]
    Por lo tanto, es suficiente tomar
    \[
        \delta=\dfrac{ \sqrt{381} - 19}{10}\approx 0.0519. \qedhere
    \]
\end{respuesta}

%%%%%%%%%%%%%%%%%%%%%%%%%%%%%%%%%%%%%%%%%%% JO
%%%%%%%%%%%%%%%%%%%%%%%%%%%%%%%%%%%%%%%%%%%
%%%%%%%%%%%%%%%%%%%%%%%%%%%%%%%%%%%%%%%%%%%

% Tema 5: Limites y diferenciabilidad
\item 
    Suponga que está trabajando en una empresa de productos lácteos cuyo producto estrella es un yogur, que se vende en un envase cónico de altura 30 cm y de radio 20 cm. La empresa cuenta con un sensor que llena el envase de yogur ``exactamente'' hasta la altura determinada. Al momento de elaborar el envase se podría tolerar errores en la altura o el radio, pero el error en la cantidad del líquido que se coloca no puede ser mayor a 1 decímetro cúbico. Usted está encargado de resolver este problema. Para ello determine una función que modele la aprobación o rechazo de los envases y el rango en el cual pueden estar la altura y radio de los mismos.
    
\begin{respuesta}
    Si $x$ es la altura y $y$ el radio del envase medido en decímetros, se tiene que su volumen es
    \[
        \dfrac{1}{3}\pi y^2 x.        
    \]
    Como la constante $\dfrac{\pi}{3}\approx 1$, no va a alterar los cálculos ?`por qué? No la vamos a considerar y se define la función
    \[
        \funcion{f}{\R^+\times \R^+}{\R^+}{(x,y)}{xy^2.}
    \]
    El caso ideal es cuando $x=3$ y $y=2$, ahora queremos saber que tan cerca pueden estar estas variables para que el volumen del líquido no sea muy diferente al ideal. Lo que queremos saber es que tolerancia $\delta >0$, me permite deducir
    \[
        \sqrt{(x-3)^2+(y-2)^2}<\delta \Imp |xy^2-12|<1.
    \]
    Por el ejercicio anterior, conocemos que esto se verifica para cualquier
    \[
        0<\delta \leq \dfrac{ \sqrt{381} - 19}{10}\approx 0.0519.
    \]
    Puesto que entre más precisión tenga la maquina de medición, más costosa y sensible será, se considera la tolerancia permitida mayor $0.0519$ decímetros; es decir, el radio debe estar en el rango $[1.9481;2.0519]$ y la altura en el rango $[2.9481;3.0519]$ en decímetros. En otras palabras, el radio debe estar en el rango $[19.481;20.519]$ y la altura en el rango $[29.481;30.519]$ en centímetros.   
\end{respuesta}

%%%%%%%%%%%%%%%%%%%%%%%%%%%%%%%%%%%%%%%%%%% JO
%%%%%%%%%%%%%%%%%%%%%%%%%%%%%%%%%%%%%%%%%%%
%%%%%%%%%%%%%%%%%%%%%%%%%%%%%%%%%%%%%%%%%%%
% Tema 5: Limites y diferenciabilidad
\item
	Demostrar que el campo escalar
	\[
		\funcion{f}{\R^n}{\R}{x}{\|x\|}
	\]
	es una función continua en $\R^n$.

\begin{respuesta}
	Sea $a\in\R^n$ arbitrario. Vamos a demostrar que $f$ es continua en $a$; es decir, que
	\[
		\lim_{x\to a}f(x)=f(a).
	\]
	Sea $\epsilon>0$. Vamos a demostrar que existe $\delta>0$ tal que para todo $x\in\R^n$,
	\[
		0<\|x-a\|<\delta \Imp \l|\|x\|-\|a\|\r|<\epsilon.
	\]
	Probemos que tomando $\delta =\epsilon$ se cumple la proposición anterior. Sea $x\in\R^n$ tal que
	\[
		0<\|x-a\|<\delta.
	\]
	Por la desigualdad triangular se tiene que
	\[
		\|x\|=\|(x-a)+a\|\leq\|x-a\|+\|a\|
	\]
	o equivalentemente 
	\[
		\|x\|-\|a\|\leq \|x-a\|
	\]
	y de idéntica manera se tiene que
	\[
		\|a\|-\|x\|\leq \|x-a\|
	\]
	o
	\[
		-\|x-a\|\leq \|x\|-\|a\|.
	\]
	Así,
	\[
		-\|x-a\|\leq \|x\|-\|a\|\leq \|x-a\|
	\]
	lo que es equivalente a que
	\[
		\l|\|x\|-\|a\|\r|\leq \|x-a\|
	\]
	y dado que $\|x-a\|<\delta=\epsilon$, se tiene que
	\[
		\l|\|x\|-\|a\|\r|<\epsilon.
	\]
\end{respuesta}

%%%%%%%%%%%%%%%%%%%%%%%%%%%%%%%%%%%%%%%%%%%
%%%%%%%%%%%%%%%%%%%%%%%%%%%%%%%%%%%%%%%%%%%
%%%%%%%%%%%%%%%%%%%%%%%%%%%%%%%%%%%%%%%%%%%
% Tema 5: Limites y diferenciabilidad
\item 
    Calcular, paso a paso
    \[
        \lim_{x\to 2}\lim_{y\to 0}f(x,y),
    \]
    donde
    \[
        \funcion{f}{[0,3]\times[-1,1]}{\R}
        {(x,y)}{f(x,y)=x^2+y+2.}
    \]
    
\begin{respuesta}
    Para $x\in [0,3]$, definimos
    \[
        \funcion{f_x}{[-1,1]}{\R}{y}{f_x(y)=x^2+y+2.}
    \]
    Notemos que, para $x\in[0,3]$, fijo, se tiene que
    \[
        \lim_{y\to 0}f_x(y)=\lim_{y\to 0}x^2+y+2=x^2+2.
    \]
    Por lo tanto, podemos definir la función
    \[
        \funcion{g}{[0,3]}{\R}{x}{g(x)=x^2+2.}
    \]
    Finalmente, tenemos que
    \[
        \lim_{x\to 2}g(x)=\lim_{x\to 2}x^2+2=6.
    \]
    Por lo tanto,
    \[
        \lim_{x\to 2}\lim_{y\to 0}f(x,y) = 6.\qedhere
    \]
\end{respuesta}

%%%%%%%%%%%%%%%%%%%%%%%%%%%%%%%%%%%%%%%%%%% PE
%%%%%%%%%%%%%%%%%%%%%%%%%%%%%%%%%%%%%%%%%%%
%%%%%%%%%%%%%%%%%%%%%%%%%%%%%%%%%%%%%%%%%%%
% Tema 5: Limites y diferenciabilidad
\item
    Sea $\func{f}{\Omega}{\R}$, con $\Omega \subset \R^2$. Se sabe que si 
    \[
        \lim_{(x,y)\to(a,b)}f(x,y)=L
    \]
    y existen los dos límites iterados 
    \[
              \dlim_{x\to a}\dlim_{y\to b}f(x,y) \texty \dlim_{y\to b}\dlim_{x\to a}f(x,y)
    \]
    entonces estos son iguales a $L$. 
    \begin{enumerate}
        \item Sea
            \[
                \funcion{g}{\Omega}{\R}{(x,y)}{\frac{x-y}{x+y}}
            \]
            donde $\Omega=\lbrace (x,y)\in \R^2:x+y\neq 0 \rbrace $. Verificar, cuando $(x,y)\to (0,0)$, que los límites iterados son diferentes y deducir que $g$ no tiende a un límite.
        \item Sea
            \[
                \funcion{g}{\Omega}{\R}{(x,y)}{\frac{x^2y^2}{x^2y^2+(x-y)^2}}
            \]
            donde $\Omega \in \R^2 \setminus \lbrace 0 \rbrace$. Demostrar que aunque los límites iterados son iguales cuando $(x,y)\to (0,0)$, la función $g$ no tiende a ningún límite.
        \item Sea
            \[
                \funcion{g}{\R^2}{\R}{(x,y)}
                    {\begin{cases}                    
                    x\sen\l(\dfrac{1}{y}\r) &\text{si }y\neq 0,\\[4mm]
                    0&\text{si }y= 0.
                \end{cases}}
            \]
            Verificar que 
            \[
                \lim_{(x,y)\to (0,0)}g(x,y)=0,
            \]
            sin embargo los límites iterados son diferentes.
    \end{enumerate}

\begin{respuesta}
    \begin{enumerate}
        \item 
            Se tiene que
            \[
                \dlim_{x\to 0}\dlim_{y\to 0}g(x,y)=1,
            \]
            mientras que
            \[
                \dlim_{y\to 0}\dlim_{x\to 0}g(x,y)=-1.
            \]
            Si el límite
            \[
                \dlim_{(x,y)\to (0,0)}g(x,y)
            \]
            existe, ya que los dos límites iterados existen entonces estos deberían ser iguales, lo cual es una contradicción. Finalmente, se puede decir  que \textbf{ya que los límites iterados existen y no son iguales entonces el \boldmath$\lim_{(x,y)\to (0,0)}g(x,y)$ no existe.} 
            
        \item
            Se tiene que
            \[
                \dlim_{x\to 0}\dlim_{y\to 0}g(x,y)=0,
            \]
            mientras que
            \[
                \dlim_{y\to 0}\dlim_{x\to 0}g(x,y)=0.
            \]
            Nótese que la existencia del límite 
            \[
                \dlim_{(x,y)\to (0,0)}g(x,y)
            \]
            implica la igualdad de los límites iterados si estos existen, no recíprocamente. Por esto \textbf{la igualdad de los límites iterados no nos dice nada sobre el \boldmath$\lim_{(x,y)\to (0,0)}g(x,y)$.} 
            
            En este ejemplo, si calculamos el límite de $g$ cuando $(x,y)\to (0,0)$ por la trayectoria inyectiva
            \[
                \funcion{\alpha}{\R}{\R^2}{t}{(t,t),}
            \]
            se tiene que $\alpha(0)=(0,0)$, entonces
            \begin{align*}
                \dlim_{t\to 0}f(\alpha(t))&=\dlim_{t\to 0}\dfrac{t^4}{t^4}\\
                &=1.
            \end{align*}
            De esta manera se concluye que el $\lim_{(x,y)\to (0,0)}g(x,y)$ no existe ya que debería ser único.
            
        \item   
            Se tiene que 
            \begin{align*}
                \dlim_{y\to 0}\dlim_{x\to 0}g(x,y)&=\dlim_{y\to 0}\l(\dlim_{x\to 0} x\sen \dfrac{1}{y} \r)\\
                &=\dlim_{y\to 0}0\\
                &=0.
            \end{align*}
            Por otro lado
            \begin{align*}
                \dlim_{x\to 0}\dlim_{y\to 0}g(x,y)&=\dlim_{x\to 0}\l(\dlim_{y\to 0} x\sen \dfrac{1}{y} \r)
            \end{align*}
            no existe ya que $\sen \dfrac{1}{y}$ no tiende a un límite cuando $y\to 0$.
            
            \textbf{Si al menos uno de los límites iterados no existe, entonces no se puede concluir nada sobre el \boldmath$\lim_{(x,y)\to (0,0)}g(x,y)$.} 
            
            Además, nótese que $\sen \dfrac{1}{y}$, cuando $y\to 0$, varía entre $[-1,1]$, por lo que 
            \[
                \dlim_{(x,y)\to (0,0)}x\sen \dfrac{1}{y}=0.
            \]
            \textbf{A pesar de la existencia de este límite, no se puede concluir que los límites iterados son iguales ya que uno de ellos no existe.} 
            
    \end{enumerate}
    
\end{respuesta}

%%%%%%%%%%%%%%%%%%%%%%%%%%%%%%%%%%%%%%%%%%% JO
%%%%%%%%%%%%%%%%%%%%%%%%%%%%%%%%%%%%%%%%%%%
%%%%%%%%%%%%%%%%%%%%%%%%%%%%%%%%%%%%%%%%%%%
% Tema 5: Limites y diferenciabilidad
\item
	Dadas la función 
	\[
		\funcion{f}{\R^2\setminus\{(1,0)\}}{\R}{x}%
			{\dfrac{x_1^2-2x_2-1}{\sqrt{(x_1-1)^2+x_2^2}}.}
	\]
	y la trayectoria
	\[
		\funcion{\alpha}{\R}{\R^2}{t}{(t,1-t).}
	\]
	¿Existe el límite de $f$ en $(1,0)$ por la trayectoria $\alpha$?
	

\begin{respuesta}
	Se puede observar que $\alpha$ es inyectiva. Además, para todo $t\in\R\setminus\{1\}$, se tiene que
	\begin{align*}
		f(\alpha(t))
			&=f(t,1-t)\\
			&=\dfrac{t^2-2(1-t)-1}{\sqrt{(t-1)^2+(1-t)^2}}\\
			&=\dfrac{t^2+2t-3}{\sqrt 2 |t-1|}\\
			&=\dfrac{(t+3)(t-1)}{\sqrt 2 |t-1|}\\
			&=
				\begin{cases}
					-\dfrac{t+3}{\sqrt 2}&\text{ si }t<1\\[2mm]
					\dfrac{t+3}{\sqrt 2}&\text{ si }t>1,\\
				\end{cases}
	\end{align*}
	de donde
	\[
		\lim_{t\to 1^-}f(\alpha(t))=-2\sqrt 2
		\yds 
		\lim_{t\to 1^+}f(\alpha(t))=2\sqrt 2.
	\]
	Así, el límite
	\[
		\lim_{t\to 1}f(\alpha(t))
	\]
	no existe.
\end{respuesta}

%%%%%%%%%%%%%%%%%%%%%%%%%%%%%%%%%%%%%%%%%%% FC
%%%%%%%%%%%%%%%%%%%%%%%%%%%%%%%%%%%%%%%%%%%
%%%%%%%%%%%%%%%%%%%%%%%%%%%%%%%%%%%%%%%%%%%
% Tema 5: Limites y diferenciabilidad
\item
    Dado el campo escalar:
    \[
        \funcion{f}{\Omega}{\R}{(x,y)}{\dfrac{xy+1}{x^2-y^2}}
    \]
    donde $\Omega=\lbrace (x,y)\in\R^2:x^2\neq y^2 \rbrace$.
    Determine el límite de $f$ cuando $(x,y) \to (1,-1)$ a través de las trayectorias:
    \[
        \funcion{\alpha}{\R}{\R^2}{t}{(t,-1)}
    \]
    y 
    \[
        \funcion{\beta}{\R}{\R^2}{s}{(s,-s^2)}
    \]
    ¿Qué se puede concluir?
\begin{respuesta}
    Se debe calcular el límite
    \[
        \lim_{(x,y)\to (1,-1)} f(x,y).
    \]
    Por definición de límites por trayectorias, para $\alpha(t)$
    \[
        \lim_{t\to t_0} f(\alpha(t))
    \]
    donde $t_0=1$ ya que $\alpha(1)=(1,-1)$, entonces
    \[
    \begin{aligned}
        \lim_{t\to t_0} f(\alpha(t))&=\lim_{t\to 1} \dfrac{-t+1}{t^2-1}\\
        &=\lim_{t\to 1} \dfrac{-t+1}{(t-1)(t+1)}\\
        &=\lim_{t\to 1} \dfrac{-1}{(t+1)}\\
        &=-\dfrac{1}{2}.
    \end{aligned}
    \]
    Para $\beta(t)$
    \[
        \lim_{s\to s_0} f(\beta(s))
    \]
    donde $s_0=1$ ya que $\beta(1)=(1,-1)$, entonces 
    \[
    \begin{aligned}
        \lim_{s\to s_0} f(\beta(s))&=\lim_{s\to 1} \dfrac{-s^3+1}{s^2-s^4}\\
        &=\lim_{s\to 1} \dfrac{s^2+s+1}{(s+1)(s^2+1)}\\
        &=\dfrac{3}{2}.
    \end{aligned}
    \]
El límite de $f$ cuando $(x,y) \rightarrow (1,-1)$ a través de las trayectorias $\alpha(t)$ y $\beta(t)$ son distintos, por lo tanto se puede concluir que el límite de $f$ cuando $(x,y) \to (1,-1)$ no existe.
\end{respuesta}

%%%%%%%%%%%%%%%%%%%%%%%%%%%%%%%%%%%%%%%%%%% JS
%%%%%%%%%%%%%%%%%%%%%%%%%%%%%%%%%%%%%%%%%%%
%%%%%%%%%%%%%%%%%%%%%%%%%%%%%%%%%%%%%%%%%%%
% Tema 5: Limites y diferenciabilidad
\item
	Dado el campo escalar
    \[
    	\funcion{f}{\R^2}{\R}{(x,y)}%
            	{\begin{cases}
            		\dfrac{\sen(x^2+y^2)}{x^2+y^2} &\text{ si }   (x,y)\neq(0,0), \\
             		1 &  \text{ si }  (x,y)=(0,0).
                 \end{cases}}
    \]
	Probar que $f$ es continuo en $(0,0)$.

\begin{respuesta}
    Para que el campo $f$ sea continuo en $(0,0)$, necesariamente se debe dar que
    \[
        \dlim_{(x,y)\to (0,0)}f(x,y)=f(0,0),
    \]
    para lo cual se analiza la existencia del límite de $f$ cuando $(x,y)\to (0,0)$. Se realiza el cambio de variable $\theta=x^2+y^2$, además nótese que si $(x,y)\to (0,0)$ entonces $\theta \to 0$, de aquí
    \begin{align*}
        \dlim_{(x,y)\to (0,0)}f(x,y)&=\dlim_{\theta\to 0}\dfrac{\sen \theta}{\theta}\\
        &=1\\
        &=f(0,0).
    \end{align*}
    Con esto se verifica que $f$ es continua en $(0,0)$.
\end{respuesta}

%%%%%%%%%%%%%%%%%%%%%%%%%%%%%%%%%%%%%%%%%%% FC
%%%%%%%%%%%%%%%%%%%%%%%%%%%%%%%%%%%%%%%%%%%
%%%%%%%%%%%%%%%%%%%%%%%%%%%%%%%%%%%%%%%%%%%
% Tema 5: Limites y diferenciabilidad
\item
    Dado el campo escalar:
    \[
        \funcion{f}{\R^2}{\R}{(x,y)}{f(x,y)=
    \begin{cases}
         1 & \text{si } y \geq x^4, \\
         -1 & \text{si } y\leq0, \\
         0 & \text{caso contrario.}
    \end{cases}}
    \]
    Obtenga los siguientes límites o explique porqué no existen:
    \begin{enumerate}
        \item $\dlim_{(x,y) \to (0,1)} f(x,y)$
        \item $\dlim_{(x,y) \to (2,3)} f(x,y)$
        \item $\dlim_{(x,y) \to (0,0)} f(x,y)$
    \end{enumerate}
    \begin{respuesta}
        \begin{enumerate}
            \item $\dlim_{(x,y) \to (0,1)} f(x,y)=1$ porque cualquier camino que pase por $(0,1)$ y su entorno, satisface la condición $y \geq x^4$ en dicho punto.
            \item $\dlim_{(x,y) \to (2,3)} f(x,y)=0$ porque cualquier camino que pase por $(2,3)$ y su entorno, no satisface las condiciones $y \leq x^4$ o $y \geq x^4$ en dicho punto.
            \item Tomando las trayectorias,
                \[
                    \funcion{\alpha}{\R}{\R^2}{t}{(0,t)}
                \]
                nótese que $\alpha(0)=(0,0)$, entonces
                \[
                    \dlim_{t \to 0} f\l(\alpha(t)\r)=1
                \]
                y
                \[
                    \funcion{\beta}{\R}{\R^2}{t}{(t^2,t)}
                \]
                nótese que $\beta(0)=(0,0)$, entonces
                \[
                    \dlim_{t \to 0} f\l(\beta(t)\r)=0
                \]
                
                Por lo tanto el límite no existe en $(0,0)$.
\end{enumerate}
\end{respuesta}

%%%%%%%%%%%%%%%%%%%%%%%%%%%%%%%%%%%%%%%%%%%
%%%%%%%%%%%%%%%%%%%%%%%%%%%%%%%%%%%%%%%%%%%
%%%%%%%%%%%%%%%%%%%%%%%%%%%%%%%%%%%%%%%%%%%
% Tema 5: Limites y diferenciabilidad
\item 
    Determinar el valor de $\alpha\in\R$ para que el campo escalar
    \[
        \funcion{f}{\R^2}{\R}{x}
        {\begin{cases}
            \dfrac{x^3+xy^2+2x^2+2y^2}{x^2+y^2}&\text{ si} (x,y)\neq(0,0),\\
            \alpha&\text{ si} (x,y)=(0,0)
        \end{cases}}
    \]
    sea continua en $(0,0)$.
    
\begin{respuesta}
    % Para que la función sea continua en $(0,0)$, es preciso que
    % \[
    %     \lim_{(x,y)\to(0,0)} f(x,y) = f(0,0).
    % \]
    % Se tiene que
    % \begin{align*}
    %     \dfrac{x^3+xy^2+2x^2+2y^2}{x^2+y^2}&=\dfrac{(x+2)(x^2+y^2)}{x^2+y^2}\\
    %     &=x+2,
    % \end{align*}
    % por lo tanto 
    % \begin{align*}
    %      \lim_{(x,y)\to(0,0)} f(x,y)&= \lim_{(x,y)\to(0,0)} (x+2)\\
    %      &=2,
    % \end{align*}
    entonces $f(0,0)$ debería ser 2, es decir  $\alpha=2$.
\end{respuesta}

%%%%%%%%%%%%%%%%%%%%%%%%%%%%%%%%%%%%%%%%%%% JR
%%%%%%%%%%%%%%%%%%%%%%%%%%%%%%%%%%%%%%%%%%%
%%%%%%%%%%%%%%%%%%%%%%%%%%%%%%%%%%%%%%%%%%%
% Tema 5: Limites y diferenciabilidad
\item 
    Considere el campo escalar
    \[
        \funcion{f}{\R^2}{\R}{(x,y)}{
        \begin{cases}
            \dfrac{3(x-1)^2y + (x-1)y^2 + (x-1)^3 + y^3}{(x-1)^3 + y^3} &\text{ si } (x,y)\neq(1,0),\\
            1 & \text{ si } (x,y) = (1,0).
        \end{cases}}
    \]
    Estudie la continuidad de $f$ en $(1,0)$.
    
% \begin{respuesta}
%     Para determinar si $f$ es continua en $(1,0)$ estudiemos la existencia del límite de $f$ cuando $(x,y)$ tiende a $(1,0)$.
    
%     Iniciemos analizando la existencia de los límites iterados, notemos que
%     \[
%         \dlim_{x\to 1}\dlim_{y\to 0}f(x,y)=\dlim_{y\to 0}\dlim_{x\to 1}f(x,y)=1,
%     \]
%     dado que los límites iterados existen y son iguales, si el límite existe, debe ser igual a 1.
    
%     Ahora, analicemos la existencia del límite mediante la trayectoria
%     \[
%         \funcion{\alpha}{\R}{\R^2}{t}{(t,t-1),}
%     \]
%     tenemos que $\alpha(t)=(1,0)$ en $t=1$. Entonces, el límite de $f$ a lo largo de la trayectoria $\alpha$ cuando $t$ tiende a 1 es
    
%     \[
%         \dlim_{t\to1}{f(\alpha(t))}=\dlim_{t\to1}{f(t,t-1)}=3,
%     \]
%     por lo tanto, si el límite existe, debe ser igual a 3, lo que contradice lo anterior. De aquí, concluimos que no existe el límite de $f$ cuando $(x,y)$ tiende a $(1,0)$, lo que implica que $f$ no es continua en $(1,0)$.
% \end{respuesta}

%%%%%%%%%%%%%%%%%%%%%%%%%%%%%%%%%%%%%%%%%%% JR
%%%%%%%%%%%%%%%%%%%%%%%%%%%%%%%%%%%%%%%%%%%
%%%%%%%%%%%%%%%%%%%%%%%%%%%%%%%%%%%%%%%%%%%    
% Tema 5: Limites y diferenciabilidad
\item 
    Considere el campo escalar
    \[
        \funcion{f}{\R^2}{\R}{(x,y)}{
        \begin{cases}
            \dfrac{xy^2}{x^2 + y^4} & \text{ si } x \neq 0,\\
            0 & \text{ si } x = 0.
        \end{cases}}
    \]
    Muestre que no es continuo en el origen.
    
\begin{respuesta}
    Para determinar si $f$ es continua en $(0,0)$ estudiemos la existencia del límite de $f$ cuando $(x,y)$ tiende a $(0,0)$.
    
    Iniciemos analizando la existencia de los límites iterados, notemos que
    \[
        \dlim_{x\to 0}\dlim_{y\to 0}f(x,y)=\dlim_{y\to 0}\dlim_{x\to 0}f(x,y)=0,
    \]
    dado que los límites iterados existen y son iguales, si el límite existe, debe ser igual a 0.
    
    Ahora, analicemos la existencia del límite mediante la trayectoria
    \[
        \funcion{\alpha}{\R}{\R^2}{t}{(t^2,t),}
    \]
    tenemos que $\alpha(t)=(0,0)$ en $t=0$. Entonces, el límite de $f$ a lo largo de la trayectoria $\alpha$ cuando $t$ tiende a 0 es
    
    \[
        \dlim_{t\to0}{f(\alpha(t))}=\dlim_{t\to0}{f(t^2,t)}=\dfrac{1}{2},
    \]
    por lo tanto, si el límite existe, debe ser igual a $\dfrac{1}{2}$, lo que contradice lo anterior. De aquí, concluimos que no existe el límite de $f$ cuando $(x,y)$ tiende a $(0,0)$, lo que implica que $f$ no es continua en el origen.
\end{respuesta}

%%%%%%%%%%%%%%%%%%%%%%%%%%%%%%%%%%%%%%%%%%% JR
%%%%%%%%%%%%%%%%%%%%%%%%%%%%%%%%%%%%%%%%%%%
%%%%%%%%%%%%%%%%%%%%%%%%%%%%%%%%%%%%%%%%%%%    
% Tema 5: Limites y diferenciabilidad
\item 
    Considere el campo escalar
    \[
        \funcion{f}{\R^3}{\R}{(x,y,z)}{
        \begin{cases}
            \dfrac{xy + yz^2 + xz^2}{x^2 + y^2 + z^4} & \text{ si } (x,y,z) \neq (0,0,0),\\
            0 & \text{ si } (x,y,z) = (0,0,0).
        \end{cases}}
    \]
    Muestre que es no es continuo en el origen.
    
\begin{respuesta}
    Analicemos este límite por trayectorias, tomemos $k\in\R$ y consideremos la trayectoria
    \[
        \funcion{\alpha_k}{\R}{\R^3}{t}{(kt^2,t^2,t).}
    \]
    Notemos que si $\alpha_k(t)=(0,0,0)$, entonces $t=0$. Luego, tenemos que 
    \[
        \dlim_{t\to0}{f(\alpha_k(t))}=\dlim_{t\to0}{f(kt^2,t^2,t)}=\dlim_{t\to0}{\dfrac{kt^4 + t^4 + kt^4}{k^2t^4 + t^4 + t^4}}=\dfrac{2k+1}{k^2 + 2}.
    \]
    Ya que el límite anterior toma valores distintos que dependen del valor de $k$, concluimos que no existe el límite de $f$ cuando $(x,y,z)$ tiende a $(0,0,0)$, lo que implica que $f$ no es continua en el origen.

%%%%%%%%%%%%%%%%%%%%%%%%%%%%%%%%%%%%%%%%%%%
%%%%%%%%%%%%%%%%%%%%%%%%%%%%%%%%%%%%%%%%%%%
%%%%%%%%%%%%%%%%%%%%%%%%%%%%%%%%%%%%%%%%%%%
% Tema 5: Limites y diferenciabilidad
\item 
Sea la función:
	\[
        \funcion{F}{\R^2}{\R^3}{(x,y)}
        {\l(e^x,e^{-y},xy\r)}
    \]
Hallar la derivada $F'(a;y)$ en el punto $a=(1,1)$ y la dirección $y=(1,2)$.
    \begin{respuesta}\
Por definición de derivada de $F: \R^2 \rightarrow \R^3 $ en $a \in R^2$ respecto a $ y\in R^2$.
    \[
       f'(a;y)=\l(f_1'(a;y),f_2'(a;y),f_3'(a;y)\r)
    \]
La derivada de cada campo escalar será:
	\[
    \begin{aligned}
       f_1'(a;y) & = \lim_{h\to 0} \dfrac{f_1(a+hy)-f_1(a)}{h}\\
        & =  \lim_{h\to 0} \dfrac{f_1((1,1)+h(1,2))-f_1(1,1)}{h}\\
        & =  \lim_{h\to 0} \dfrac{f_1((1+h,1+2h)-f_1(1,1)}{h}\\
        & =  \lim_{h\to 0} \dfrac{e^{1+h}-e^1}{h}\\
        & =  \lim_{h\to 0} \dfrac{e(e^h-1)}{h}\\
        & =  e
     \end{aligned}
    \]
    
	\[
    \begin{aligned}
       f_2'(a;y) & = \lim_{h\to 0} \dfrac{f_2(a+hy)-f_2(a)}{h}\\
        & =  \lim_{h\to 0} \dfrac{e^{-1-2h}-e^{-1}}{h}\\
        & =  \lim_{h\to 0} \dfrac{-2e^{-1}(e^{-2h}-1)}{-2h}\\
        & =  \dfrac{-2}{e}
     \end{aligned}
    \]
    
	\[
    \begin{aligned}
       f_3'(a;y) & = \lim_{h\to 0} \dfrac{f_3(a+hy)-f_3(a)}{h}\\
        & =  \lim_{h\to 0} \dfrac{(1+h)(1+2h)}{h}\\
        & =  \lim_{h\to 0} \dfrac{1+2h+h+2h^2-1}{h}\\
        & =  3
     \end{aligned}
    \]
Por lo tanto $F'(a;y)$ en el punto $a=(1,1)$ y la dirección $y=(1,2)$ es:
    \[
       F'(a;y)=\l(e,\dfrac{-2}{e},3\r)
    \]
\end{respuesta}

%%%%%%%%%%%%%%%%%%%%%%%%%%%%%%%%%%%%%%%%%%%
%%%%%%%%%%%%%%%%%%%%%%%%%%%%%%%%%%%%%%%%%%%
%%%%%%%%%%%%%%%%%%%%%%%%%%%%%%%%%%%%%%%%%%%
% Tema 6: Derivadas parciales y gradiente
\item
    Dado el campo escalar
    \[
        \funcion{f}{\R^n}{\R}{x}{f{(x)=\| x \|^n}}
    \]
    Calcule $f'(a,y)$
    \begin{respuesta}
    Definimos
    \[
        \funcion{g}{I}{\R}{t}{g{(t)=f(a+ty)}}
    \]
    donde $a,y\in \R^n$        
    \[
        g(t)=\| a+ty \|^n
    \]  
    \[
        g(t)=(\| a+ty \|^2)^{\frac{n}{2}}
    \]
    \[
        g(t)=[(a+ty)\cdot (a+ty)]^{\frac{n}{2}}
    \]
    Derivando     
    \[
        g'(t)=\frac{n}{2}[(a+ty)\cdot (a+ty)]^{\frac{n}{2}-1}[y\cdot(a+ty)+(a+ty)\cdot y]
    \]
    Si $t=0$
    \[
        g'(0)=\frac{n}{2}(a\cdot a)^{\frac{n}{2}-1}[ya+ay]
    \]
    \[
        g'(0)=\frac{n}{2}(\| a \|^2)^{\frac{n}{2}-1}(2ay)
    \]
    \[
        g'(0)=n\|a\|^{n-2}(ay)
    \] 
    \[
        f'(a,y)=n\|a\|^{n-2}(ay)
    \] 
\end{respuesta}

%%%%%%%%%%%%%%%%%%%%%%%%%%%%%%%%%%%%%%%%%%%
%%%%%%%%%%%%%%%%%%%%%%%%%%%%%%%%%%%%%%%%%%%
%%%%%%%%%%%%%%%%%%%%%%%%%%%%%%%%%%%%%%%%%%%
% Tema 6: Derivadas parciales y gradiente
\item 
    Dados $\Omega\subseteq \R^2$, y $\func{f}{\Omega}{\R}$ un campo escalar diferenciable tal que su derivadas parciales sean 0 en todo punto de $\Omega$, demostrar que $f$ es constante.

\begin{respuesta}
    Sean $x,y\in \Omega$, debemos demostrar que $f(x)=f(y)$. Llamemos $h=y-x$, dado que $f$ es diferenciable, existen las derivadas en todo punto y en toda dirección, por lo tanto, por el teorema del valor medio para campos escalares, tenemos que existe $\zeta\in(0,1)$ tal que
    \[
        f(x+h)-f(x) = f'(x+\zeta h;h) = \nabla f(x+\zeta h)\cdot h.
    \]
    Dado que todas las derivadas parciales de $f$ son 0, se tiene que $\nabla f(x+\zeta h)=0$, por lo tanto
    \[
        f(x+h)-f(x)=0,
    \]
    lo que equivale a
    \[
        f(y)=f(x),
    \]
    y por lo tanto la función es constante.
\end{respuesta}

%%%%%%%%%%%%%%%%%%%%%%%%%%%%%%%%%%%%%%%%%%%
%%%%%%%%%%%%%%%%%%%%%%%%%%%%%%%%%%%%%%%%%%%
%%%%%%%%%%%%%%%%%%%%%%%%%%%%%%%%%%%%%%%%%%%
% Tema 6: Derivadas parciales y gradiente
\item
    Dado el campo escalar
    \[
        \funcion{f}{\R^3}{\R}{(x,y,z)}{\cos{(xyz)}}
    \]
    Utilizando la definición determine las derivadas parciales de $f$.
    
% \begin{respuesta}
%     Sea $(x,y,z)\in\R^3$. Supongamos que $y\neq 0$ y $z\neq 0$. Tenemos que
%     \begin{align*}
%         \frac{f((x,y,z)+h(1,0,0))-f(x,y,z)}{h}
%             & = \frac{\cos((x+h)yz)-\cos(xyz))}{h} \\
%             & = \frac{\cos(xyz+hyz)-\cos(xyz))}{h} \\
%             & = \frac{\cos(xyz)\cos(hyz)-\sen(xyz)\sen(hyz)-\cos(xyz))}{h} \\
%             & = \frac{\cos(xyz)(\cos(hyz)-1)}{h} - \frac{\sen(xyz)\sen(hyz)}{h}\\
%             & = yz\cos(xyz)\frac{\cos(hyz)-1}{hyz} - yz\sen(xyz)\frac{\sen(hyz)}{hyz},
%     \end{align*}
%     por lo tanto
%     \begin{align*}
%         \lim_{h\to 0}\frac{f((x,y,z)+h(1,0,0))-f(x,y,z)}{h}
%             & = \lim_{h\to 0} \l(yz\cos(xyz)\frac{\cos(hyz)-1}{hyz} - yz\sen(xyz)\frac{\sen(hyz)}{hyz} \r)\\
%             & = - yz\sen(xyz).
%     \end{align*}
%     Así, tenemos que
%     \[
%         D_1f(x,y,z)=-yz\sen(xyz)
%     \]
%     si $y\neq 0$ y $z\neq 0$. De manera similar se calcula $D_2f(x,y,z)$ y $D_3f(x,y,z)$.
   
% \end{respuesta}

%%%%%%%%%%%%%%%%%%%%%%%%%%%%%%%%%%%%%%%%%%% JR
%%%%%%%%%%%%%%%%%%%%%%%%%%%%%%%%%%%%%%%%%%%
%%%%%%%%%%%%%%%%%%%%%%%%%%%%%%%%%%%%%%%%%%%
% Tema 6: Derivadas parciales y gradiente
\item 
    Utilizando la función
    \[
        \funcion{p}{\R^{+}\times\R^{+}}{\R}{(u,v)}{p(u,v)=1.01u^{0.75}v^{0.25},}
    \]
    Cobb y Douglas modelaron la economía americana de 1899 a 1922, en donde la primera variable representa a la cantidad de la mano de obra y la segunda a la cantidad de capital.
    \begin{enumerate}
        \item Determine $D_1 p$ y $D_2 p$.
        \item Encuentre la productividad marginal de la mano de obra y la productividad marginal del capital en el año 1920 en $(194, 407)$ (comparado con los valores asignados en 1899 de $(100,100)$). Interprete los resultados.
        \item En el año 1920, ¿qué producción tendría más beneficio, un incremento de inversión de capital o un incremento en el gasto en mano de obra?
    \end{enumerate}

\begin{respuesta}
    \begin{enumerate}
        \item 
        Notemos que $p$ es diferenciable en $\R^{+} \times \R^{+}$, por tanto, sus derivadas parciales $D_1 p$ y $D_2 p$ existen y están definidas por
        \[
            D_1 p(u,v) = 0.7575 u^{-0.25} v^{0.25}
        \]
        y
        \[
            D_2 p(u,v) = 0.2525 u^{0.75} v^{-0.75}
        \]
        para cada $(u,v)\in \R^{+} \times \R^{+}$, luego, se tiene que 
        \[
            \funcion{D_1 p}{\R^{+}\times\R^{+}}{\R}{(u,v)}{D_1 p(u,v) = 0.7575 u^{-0.25} v^{0.25}}
        \]
        y
        \[
            \funcion{D_2 p}{\R^{+}\times\R^{+}}{\R}{(u,v)}{D_2 p(u,v) = 0.2525 u^{0.75} v^{-0.75}.}
        \]
        \item
        Para determinar la productividad marginal evaluamos $D_1 p$ y $D_2 p$ en los valores asignados para los años 1920 y 1899, respectivamente.
        
        Para 1920 tenemos que,
        \[
            D_1 p (194, 407) \approx 0.912 \texty D_2 p (194, 407) \approx 0.319.
        \]
        Para 1920 tenemos que,
        \[
            D_1 p (100, 100) \approx 0.758 \texty D_2 p (100, 100) \approx 0.555.
        \]
        Como podemos notar, comparando los valores marginales obtenidos para cada año concluimos que respecto a la mano de obra, fue mas rentable invertir en este rubro en el año 1920 que en el año 1899. De forma similar, se concluye que respecto al capital, fue más rentable invertir en este rubro en 1899 que en 1920. 
        \item 
        Gracias al literal anterior concluimos que por cada unidad invertida en el gasto de mano de obra se obtuvo una ganancia de 0.912 unidades, mientras que por cada unidad gastada en inversión de capital, solamente se obtuvieron 0.319 unidades de ganancia. Por lo tanto, se obtuvo un mayor beneficio al gastar en la mano de obra que en invertir una mayor cantidad de capital.
    \end{enumerate}
\end{respuesta}

%%%%%%%%%%%%%%%%%%%%%%%%%%%%%%%%%%%%%%%%%%%
%%%%%%%%%%%%%%%%%%%%%%%%%%%%%%%%%%%%%%%%%%%
%%%%%%%%%%%%%%%%%%%%%%%%%%%%%%%%%%%%%%%%%%%
% Tema 6: Derivadas parciales y gradiente
\item 
    Dado $a\in\R^n$, determinar el vector gradiente del campo escalar definido por
    \[
        \funcion{f}{\R^n}{\R}{x}{f(x)=a\cdot x.}
    \]

\begin{respuesta}
    Empecemos calculando las derivadas parciales del campo escalar, así, sea $i\in\{1,\ldots,n\}$, tenemos que, para $x\in\R^n$,
    \begin{align*}
        D_if(x) 
            & = D_i \l( \sum_{j=1}^n a_jx_j \r)\\
            & = \sum_{j=1}^n D_i (a_jx_j)\\
            & = a_i.
    \end{align*}
    Por lo tanto, tenemos que
    \[
        \nabla f(x)
        =\l(D_1f(x), \ldots, D_nf(x)\r)
        =(a_1,\ldots,a_n) = a,
    \]
    para todo $x\in \R^n$.
\end{respuesta}

%%%%%%%%%%%%%%%%%%%%%%%%%%%%%%%%%%%%%%%%%%%
%%%%%%%%%%%%%%%%%%%%%%%%%%%%%%%%%%%%%%%%%%%
%%%%%%%%%%%%%%%%%%%%%%%%%%%%%%%%%%%%%%%%%%%
% Tema 6: Derivadas parciales y gradiente
\item 
    Dados $\Omega\subset \R^2$ y $\func{f,g}{\Omega}{\R}$ dos campos escalares derivables tal que $g(x)\neq 0$ para todo $x\in \Omega$, demuestre que
    \[
        \nabla\l(\frac{f}{g}\r)(x)=\frac{g(x)\nabla f(x)-f(x)\nabla g(x)}{g(x)^2}.
    \]

\begin{respuesta}
    Para $x\in\Omega$, tenemos que
    \begin{align*}
        D_1\l(\frac{f}{g}\r)(x)
            & = D_1 \l(\frac{f(x)}{g(x)}\r) \\
            & = \frac{g(x) D_1 f(x)  - f(x) D_1g(x)}{g(x)^2}
    \end{align*}
    y
    \begin{align*}
        D_2\l(\frac{f}{g}\r)(x)
            & = D_2\l(\frac{f(x)}{g(x)}\r) \\
            & = \frac{g(x) D_2f(x) - f(x) D_2g(x)}{g(x)^2}
    \end{align*}
    Por lo tanto,
    \begin{align*}
        \nabla\l(\frac{f}{g}\r)(x)
            & = \l(D_1\l(\frac{f}{g}\r)(x),D_2\l(\frac{f}{g}\r)(x) \r)\\[2mm]
            & = \l( \frac{g(x)D_1 f(x) + f(x)D_1g(x)}{g(x)^2} , \frac{g(x)D_2 f(x) + f(x)D_2g(x)}{g(x)^2} \r)\\[2mm]
            & = \frac{g(x)\l(D_1f(x),D_2f(x)\r) - f(x)\l(D_1g(x),D_2g(x)\r)}{g(x)^2}\\[2mm]
            & = \frac{g(x)\nabla f(x)-f(x)\nabla g(x)}{g(x)^2}.\qedhere
    \end{align*}
\end{respuesta}

%%%%%%%%%%%%%%%%%%%%%%%%%%%%%%%%%%%%%%%%%%%
%%%%%%%%%%%%%%%%%%%%%%%%%%%%%%%%%%%%%%%%%%%
%%%%%%%%%%%%%%%%%%%%%%%%%%%%%%%%%%%%%%%%%%%
% Tema 6: Derivadas parciales y gradiente
\item
    Sean $\func{f}{\R^2}{\R}$ un campo escalar diferenciable, $a=(1,2)$, $u=(2,2)$ y $v=(1,1)$. Se tiene que la derivada de la función en el punto $a$ en dirección al punto $u$ es 2 y la derivada de la función en el punto $a$ en dirección al punto $v$ es $-2$. Determinar el vector gradiente de $f$ en $a$ y calcular la derivada de $f$ en $a$ en dirección al punto $(4,6)$. 

    \begin{respuesta}
    Tenemos que $f'(a;u-a)=2$ y $f'(a;v-a)=-2$. Por otra parte, se tiene que
    \begin{align*}
        f'(a;u-a) 
            & = \nabla f(a) \cdot (u-a) \\
            & = (D_1f(a),D_2f(a))\cdot (1,0) \\
            & = D_1f(a)
    \end{align*}
    y
    \begin{align*}
        f'(a;v-a) 
            & = \nabla f(a) \cdot (v-a) \\
            & = (D_1f(a),D_2f(a))\cdot (0,-1) \\
            & = -D_2f(a),
    \end{align*}
    por lo tanto
    \[
        D_1f(a)=2
        \texty
        D_2f(a)=2,
    \]
    lo que implica que
    \[
        \nabla f(a)=(2,2). 
    \]
    Finalmente, tenemos que la dirección de (1,2) a (4,6) es (3,4), entonces la derivada direccional en esta dirección es
    \[
        \nabla f(a) \cdot \left( \frac{3}{5},\frac{4}{5} \right)=\frac{14}{5} \qedhere
    \]
\end{respuesta} 



%%%%%%%%%%%%%%%%%%%%%%%%%%%%%%%%%%%%%%%%%%%
%%%%%%%%%%%%%%%%%%%%%%%%%%%%%%%%%%%%%%%%%%%
%%%%%%%%%%%%%%%%%%%%%%%%%%%%%%%%%%%%%%%%%%%
% Tema 6: Derivadas parciales y gradiente
\item
    Sea la función
   		\[
            \funcion{f}{\R^3}{\R}{(x,y,z)}{x^2+y^2-4z^2.}
        \]
    Determinar la razón de cambio de $f$ en el punto $a=(1,3,-2)$ en la dirección del vector $v=\veci-\vecj+\veck$.
% \begin{respuesta}
%     Tenemos que
% 	\[
%         \nabla f(x,y,z)=(2x,2y,-8z),
%     \]
%     por lo tanto
%     \[
%         \nabla f(1,3,-2)=(2,6,16).
%     \]
%     Además, el vector director de v es
%     \[
%         u=\dfrac{1}{\sqrt{3}}(1,-1,1),
%     \]
%         de donde
%     \[
%         f'(a;u)=\nabla f(a)\cdot u=(2,6,16) \cdot \dfrac{1}{\sqrt{3}}(1,-1,1)=\dfrac{12}{\sqrt{3}},
%     \]
%     la cual es la razón de cambio buscada.
% \end{respuesta}

%%%%%%%%%%%%%%%%%%%%%%%%%%%%%%%%%%%%%%%%%%%
%%%%%%%%%%%%%%%%%%%%%%%%%%%%%%%%%%%%%%%%%%%
%%%%%%%%%%%%%%%%%%%%%%%%%%%%%%%%%%%%%%%%%%%
% Tema 6: Derivadas parciales y gradiente
\item 
    Dado el campo escalar
	\[
		\funcion{f}{\R^2}{\R}{x}{2x_1^2-3x_2.}
	\]
	Encuentre la dirección en el que la derivada direccional de $f$ en el punto $(1,2)$ es máxima y verifique que esta coincide con la del gradiente de $f$ en dicho punto.

\begin{respuesta}
	Como $f$ es un polinomio, $f$ tiene todas sus derivadas direccionales en todo punto y toda dirección. Así, para cada $y$ unitario de $\R^2$, se tiene que
	\begin{align*}
		f'((1,2);y)
			&=\lim_{h\to 0}\dfrac{f((1,2)+hy)-f(1,2)}{h}    \\
			&=\lim_{h\to 0}\dfrac{f(1+hy_1,2+hy_2)-f(1,2)}{h}\\	
			&=\lim_{h\to 0}\dfrac{2(1+hy_1)^2-3(2+hy_2)+4}{h}    \\
	        &=\lim_{h\to 0}\dfrac{4hy_1+2h^2y_1^2-3hy_2}{h}    \\
			&=4y_1-3y_2.
	\end{align*}
    Como $y$ es unitario, se cumple que $y_1^2+y_2^2=1$. Por lo tanto, $|y_2|=\sqrt{1-y_1^2}$. Por lo tanto,
    \[
        f'((1,2);y)
        =\begin{cases}
            4y_1-3\sqrt{1-y_1^2}& \text{si }y_2\geq0,\\
            4y_1+3\sqrt{1-y_1^2}& \text{si }y_2<0.
        \end{cases}
    \]
    Para hallar la dirección unitaria $y$ en la cual $f'((1,2);y)$ es máxima, conviene definir dos funciones. Sean
    \[
        \funcion{\alpha}{[-1,1]}{\R}{x}{4x-3\sqrt{1-x^2}}
        \yds
        \funcion{\beta}{[-1,1]}{\R}{x}{4x+3\sqrt{1-x^2}.}
    \]
    Mediante la optimización usual en variable real, se deduce que
    \[
        \max\{\alpha(x):x\in[-1,1]\}=\alpha(1)=4
        \yds
        \max\{\beta(x):x\in[-1,1]\}=\alpha(0.8)=5.
    \]
    Así, la dirección que maximiza la derivada direccional es
    \[
        y_1=0.8
        \yds
        y_2=-0.6.
    \]
    Por otra parte,
    \[
        \nabla f(x)=(4x_1,-3)
    \]
    para cada $x\in\R^2$. En particular,
    \[
        \nabla f(1,2)=(4,-3),
    \]
    de donde, su dirección es
    \[
        \dfrac{\nabla f(1,2)}{\|\nabla f(1,2)\|}
        =\dfrac{(4,-3)}{5}
        =(0.8;-0.6),
    \]
	lo que prueba que la dirección del gradiente maximiza la derivada direccional en el punto $(1,2)$.
\end{respuesta}

%%%%%%%%%%%%%%%%%%%%%%%%%%%%%%%%%%%%%%%%%%%
%%%%%%%%%%%%%%%%%%%%%%%%%%%%%%%%%%%%%%%%%%%
%%%%%%%%%%%%%%%%%%%%%%%%%%%%%%%%%%%%%%%%%%%
% Tema 6: Derivadas parciales y gradiente
\item
    Se detectó una explosión nuclear en un lugar indeterminado, el equipo especializado trata de buscar el origen de la misma, para lo cual se envía un dron con materiales de medición. A partir de un punto de referencia, el dron se desplaza 1 metro al norte y detecta que la radiación aumenta 0.1 Sv, luego de esto, regresa al punto original y se desplaza un metro al este y la radiación disminuye 0.2 Sv. Determinar, aproximadamente, la dirección en la cual ocurrió la explosión.

\begin{respuesta}
    Primero, utilizaremos las siguientes notaciones:
    \begin{itemize}
        \item $x$: distancia, en metros, desde el punto de referencia en dirección norte.
        \item $y$: distancia, en metros, desde el punto de referencia en dirección este.
        \item $f(x,y)$: nivel de radiación, en Sv, en las coordenadas $(x,y)$.
    \end{itemize}
    Así, se tiene que 
    \[
        D_2f(0,0)\approx 0.1
        \texty
        D_1f(0,0)\approx -0.2,
    \]
    por lo tanto,
    \[
        \nabla f(0,0)\approx (-0.2,0.1).
    \]
    Dado que el gradiente de una función da la dirección de máximo crecimiento, tenemos que la dirección aproximada en la cual ocurrió la explosión es: $(-0.2,0.1)$.
\end{respuesta}

%%%%%%%%%%%%%%%%%%%%%%%%%%%%%%%%%%%%%%%%%%%
%%%%%%%%%%%%%%%%%%%%%%%%%%%%%%%%%%%%%%%%%%%
%%%%%%%%%%%%%%%%%%%%%%%%%%%%%%%%%%%%%%%%%%%
% Tema 6: Derivadas parciales y gradiente
\item 
    Del ejercicios anterior, se detecta que el nivel de radiación del lugar en el que se encuentra el dron es crítico, es decir, el dron no puede avanzar a un lugar con mayor radiación que la existente en su posición. ¿En qué dirección puede ir el dron para que no exista un incremento en el nivel de la radiación?

\begin{respuesta}
    Se busca una dirección $(h,k)\in\R^2$ tal que
    \[
        f'((0,0);(h,k))=0,
    \]
    dado que $f'((0,0);(h,k))=\nabla f(0,0)\cdot (h,k)$, se tiene que se necesita que
    \[
        (-0.2,0.1)\cdot (h,k) = 0
    \]
    que equivale a $-2h+k=0$, por lo tanto, $k=2h$, así, una dirección en la que puede ir el dron para que no exista incremento en el nivel de radiación es $(1,2)$.
\end{respuesta}

%%%%%%%%%%%%%%%%%%%%%%%%%%%%%%%%%%%%%%%%%%%
%%%%%%%%%%%%%%%%%%%%%%%%%%%%%%%%%%%%%%%%%%%
%%%%%%%%%%%%%%%%%%%%%%%%%%%%%%%%%%%%%%%%%%%
% Tema 6: Derivadas parciales y gradiente
\item
    Sea el campo escalar
    \[
        \funcion{f}{D\subseteq\R^3}{\R}{(x,y,z)}{f{(x,y,z)=axy^2+byz+cz^2x^3}}
    \]
    donde $a,b$ y $c\in R$ y son constantes. Halle los valores de $a,b$ y $c$ tales que la derivada direccional de $f$ en el punto $(1,2,-1)$ tenga el valor máximo de 64 en una dirección paralela al eje $z$.
    
\begin{respuesta}
    El valor máximo de la derivada sucede cuando
    \[
        f'\l(a;\frac{\nabla f(a)}{\|\nabla f(a)\|}\r)=\|\nabla f(a)\|
    \]
    Primero debemos obtener el vector gradiente de $f$ en a. 
    \[
        \nabla f(a)=\l(D_1f(a),D_2f(a),D_3f(a)\r)
    \]
    Las derivadas parciales de $f$ en el punto $a=(1,2,-1)$
    \[
        D_1f(a)=(4a+3c),
        \qquad
        D_2f(a)=(4a-b)
        \texty
        D_3f(a)=(2b-2c)
    \]
    Entonces:
    \[
    \nabla f(a)=(4a+3c,4a-b,2b-2c)
           \]
    Por otro lado, como el vector $y$ es paralelo al eje $z$
    \[
        y=(0,0,k),
    \]
    con $k\in\R$. Para que la derivada de $f$ en el punto $a$ y en la dirección del vector $y$ sea máxima, se debe cumplir
    \[
        y=\frac{\nabla f(a)}{\| \nabla f(a) \|}=(0,0,k),
    \]
    \[
        \|\nabla f(a) \|=\sqrt{(4a+3c)^2+(4a-b)^2+(2b-2c)^2}
    \]
    Entonces:
    \[
        \frac{\nabla f(a)}{\| \nabla f(a)\|}=\frac{(4a+3c,4a-b,2b-2c)}{\sqrt{(4a+3c)^2+(4a-b)^2+(2b-2c)^2}}=(0,0,k)
    \]
    Considerando que $\|\nabla f(a) \| \neq0$, se obtienen el sistema de ecuaciones
    \[
        \begin{cases}
            4a+3c=0 \\
            4a-b=0 \\
            2b-2c=k,
        \end{cases}
    \]
    del cual se sigue que
    \[
        a=\frac{3k}{32}
        \qquad
        b=\frac{3k}{8}
        \texty
        c=-\frac{k}{8}
    \]
    Reemplazando estos valores en el módulo del vector gradiente, se tiene que
    \[
        \| \nabla f(a)\| = \sqrt{k^2} 
    \]
    Pero se considera que el valor máximo de la derivada de $f$ en $a$ en la dirección del vector $y$ es igual a 64.
    \[
        64 = |k|
    \]
    Entonces sustituyendo este valor,  obtenemos
    \[
        a=6
        \qquad
        b=24
        \texty
        c=-8
    \]
    o
    \[
        a=-6
        \qquad
        b=-24
        \texty
        c=8
    \]
    
\end{respuesta}

%%%%%%%%%%%%%%%%%%%%%%%%%%%%%%%%%%%%%%%%%%%
%%%%%%%%%%%%%%%%%%%%%%%%%%%%%%%%%%%%%%%%%%%
%%%%%%%%%%%%%%%%%%%%%%%%%%%%%%%%%%%%%%%%%%%
% Tema 6: Derivadas parciales y gradiente
\item 
    La distribución de temperatura de una placa metálica está dada por la función
	\[
        \funcion{T}{\R^2}{\R}
        {(x,y)}{T(x,y) = xe^{2y} + y^3 e,}
	\]
	Donde $(x,y)$ representa las coordenadas de un punto sobre la placa.
    \begin{enumerate}
        \item ¿En qué dirección aumenta más rápidamente la temperatura de la placa metálica en el punto $(2, 0)$?
        \item ¿En qué dirección la temperatura decrece más rápidamente?
    \end{enumerate}

% \begin{respuesta}
%     Recordemos que para un campo escalar $f$ definido sobre un abierto $D\subseteq\R^2$, en un punto $a\in D$, la dirección en donde $f'(a;w)$ alcanza su mayor valor es
%     \[
%         w = \frac{\nabla f(a)}{\|\nabla f(a)\|},
%     \]
%     siempre que $\|\nabla f(a)\|\neq 0$.

%     Calculemos el gradiente de $T$, para esto, tomemos un punto $a =(a_1, a_2)\in\R^2$ y una dirección $w=(w_1,w_2)\in\R^2$ cualesquiera. Se define la función $g$ por 
% 	\[
%         \funcion{g}{\R}{\R}
%         {t}{g(t) = T(a + tw) = (a_1 + t w_1)e^{2(a_2 + t w_2)} + (a_2 + t w_2)^3 e.}
% 	\]
%     Notemos que $g$ es derivable en todo $t\in\R$, con derivada 
%     \[
%         g'(t) = T'(a + tw; w) = w_1 e^{2(a_2 + t w_2)} + 2w_2(a_1 + tw_1) e^{2(a_2 + t w_2)} + 3w_2(a_2 + t w_2)^2 e,
%     \]
%     para todo $t\in\R$.

%     En particular, si $t = 0$, se tiene que $g'(0) = T'(a;w) = w_1 e^{2a_2} + 2w_2 a_1 e^{2a_2} + 3w_2a^2_2 e$. Entonces, si tomamos $a = (x, y)$, se tiene que el gradiente de $T$, está dado por
%     \[
%         \nabla T(x, y) = (e^{2y}, 2xe^{2y} + 3y^2e),
%     \]
%     para todo $(x, y)\in\R^2$.

%     \begin{enumerate}
%     \item 
%         Si $a = (2, 0)$ se tiene que $\nabla T(a) = (1, 4)$ con $\|\nabla T(a) \| =  \sqrt{17}$. Por tanto, la dirección $w$ en donde la temperatura de la placa metálica en el punto $a = (2, 0)$ aumenta más rápidamente es
% 	    \[
% 	        w = \l(\dfrac{1}{\sqrt{17}}, \dfrac{4}{\sqrt{17}}\r).
% 	    \]
%     \item 
%         Gracias a lo anterior, la dirección $w$ en donde la temperatura de la placa metálica en el punto $a = (2, 0)$ decrece más rápidamente es
% 	    \[
% 	        w = \l(-\dfrac{1}{\sqrt{17}}, -\dfrac{4}{\sqrt{17}}\r).
% 	    \]
%     \end{enumerate}
% \end{respuesta}

%%%%%%%%%%%%%%%%%%%%%%%%%%%%%%%%%%%%%%%%%%%
%%%%%%%%%%%%%%%%%%%%%%%%%%%%%%%%%%%%%%%%%%%
%%%%%%%%%%%%%%%%%%%%%%%%%%%%%%%%%%%%%%%%%%%
% Tema 6: Derivadas parciales y gradiente
\item 
    La ecuación de la superficie de una colina es $z = 1200 - 3x^2 - 2y^2$, donde la distancia se mide en metros, el eje $x$ apunta al Este y el eje $y$ apunta al Norte. Un hombre se encuentra en el punto $(-10, 5, 850)$.
    \begin{enumerate}
        \item ¿Cuál es la dirección de la ladera más pronunciada?
        \item Si el hombre se mueve en la dirección Este. ¿Está ascendiendo o descendiendo? ¿Cuál es la pendiente?
        \item Si el hombre se mueve en la dirección S-O. ¿Está ascendiendo o descendiendo? ¿Cuál es la pendiente?
    \end{enumerate}

\begin{respuesta}
    Supongamos que la variable $z$, que representa la altura a la cual se encuentra el individuo, es una función que depende de las variables $x$ e $y$, las coordenadas cartesianas de la ubicación de este, es decir, a $z$ le asociamos la función $f$ definida por
    \[
        \funcion{f}{\R^2}{\R}
        {(x, y)}{f(x, y) = 1200 - 3x^2 - 2y^2.}
    \]

    \begin{enumerate}
    \item 
        Recordemos que para un campo escalar $h$ definido sobre un abierto $D\subseteq\R^2$, en un punto $a\in D$, la dirección en donde $h'(a;w)$ alcanza su mayor valor es
        \[
            w = \frac{\nabla h(a)}{\|\nabla h(a)\|},
        \]
        siempre que $\|\nabla h(a)\|\neq 0$.

        Sean $a = (a_1, a_2)\in\R^2$ y $w = (w_1, w_2)\in\R^2$ un punto y una dirección cualesquiera, calculemos la derivada de $f$ en el punto $a$ respecto a la dirección $w$, para esto definimos la función $g$ por
        \[
            \funcion{g}{\R}{\R}
            {t}{g(t) = f(a + tw) = 1200 - 3(a_1 + tw_1)^2 - 2(a_2 + tw_2)^2.}
        \]
        Notemos que $g$ es derivable para todo $t\in\R$ con derivada
        \[
            g'(t) = f'(a + tw; w) = - 6w_1(a_1 + tw_1) - 4w_2(a_2 + tw_2),
        \]
        para todo $t\in\R$. En particular, si $t=0$ se tiene que $g'(0) = f'(a;w) = - 6w_1 a_1 - 4w_2 a_2$.

        Entonces, el gradiente de $f$ está dado por
        \[
            \nabla f(x, y) = (- 6x , - 4y),
        \]
        para todo $(x, y)\in\R^2$. 

        En particular, si $(x, y) = (-10, 5)$ metros, se tiene que $\nabla f(-10, 5) = (60 , - 20)$, con $\|\nabla f(-10, 5)\| = \|(60 , - 20)\| = 20\sqrt{10}$. 

        Por tanto, la dirección $w$ de la ladera más pronunciada es
        \[
            w = \l(\dfrac{3}{\sqrt{10}}, -\dfrac{1}{\sqrt{10}}\r).
        \]
    \item 
        En este caso, la dirección que el individuo toma es $w = (1, 0)$, entonces, con $a = (-10, 5)$ metros se tiene que $f'(a;w) = 60$, lo que implica que el individuo asciende por la colina con una pendiente de $60$.
    \item 
        En este caso, la dirección que el individuo toma es $w = \l(-\dfrac{\sqrt{2}}{2},-\dfrac{\sqrt{2}}{2}\r)$, entonces, con $a = (-10, 5)$ metros se tiene que $f'(a;w) = -20\sqrt{2}$, lo que implica que el individuo desciende por la colina de pendiente $-20\sqrt{2}$.
    \end{enumerate}
\end{respuesta}


%%%%%%%%%%%%%%%%%%%%%%%%%%%%%%%%%%%%%%%%%%%
%%%%%%%%%%%%%%%%%%%%%%%%%%%%%%%%%%%%%%%%%%%
%%%%%%%%%%%%%%%%%%%%%%%%%%%%%%%%%%%%%%%%%%%
% Tema 5: Limites y diferenciabilidad
\item 
    Dado el campo escalar
    \[
        \funcion{f}{\R^2}{\R}{x}{x_1^2+x_1x_2,}
    \]
    comprobar que es diferenciable en todo $\R^2$.
    
\begin{respuesta}
    Sea $x\in\R^2$, para determinar su diferenciabilidad, debemos estudiar el límite
    \[
        \lim_{h\to 0}\frac{f(x+h)-f(x)-\nabla f(x)\cdot h}{\|h\|}.
    \]
    Se tiene que el gradiente de la función es
    \[
        \nabla f(x) = (2x_1+x_2,x_1),
    \]
    así, tenemos que
    \begin{align*}
        \lim_{h\to 0}\frac{f(x+h)-f(x)-\nabla f(x)\cdot h}{\|h\|}
            & =  \lim_{h\to 0}\frac{f(x_1+h_1,x_2+h_2)-f(x_1,x_2) - (2x_1+x_2,x_1)\cdot(h_1,h_2)}{\|h\|}\\
            & =  \lim_{h\to 0}\frac{(x_1+h_1)^2+(x_1+h_1)(x_2+h_2)-x_1^2-x_1x_2 - 2x_1h_1-x_2h_1-x_1h_2}{\|h\|}\\
            & =  \lim_{h\to 0}\frac{h_1h_2+h_1^2}{\|h\|}.
    \end{align*}
    Notemos que, para $h\in\R^2$, con $h\neq 0$,
    \[
        \l|\frac{h_1h_2+h_2^2}{\|h\|}\r| 
        = \frac{\l|(h_1,h_2)\cdot(h_2,h_2)\r|}{\|h\|}
        \leq \frac{\|(h_1,h_2)\|\|(h_2,h_2)\|}{\|h\|}
        =\|(h_2,h_2)\|=|h_2|\sqrt{2},
    \]
    Por lo tanto,
    \[
        \lim_{h\to 0}\l|\frac{h_1h_2+h_2^2}{\|h\|}\r|\leq \lim_{h\to 0}|h_2|\sqrt{2}  =0.
    \]
    Así, se tiene que
    \[
        \lim_{h\to 0}\frac{f(x+h)-f(x)-\nabla f(x)\cdot h}{\|h\|}=0,
    \]
    por lo tanto, se tiene que la función es diferenciable.
\end{respuesta}

%%%%%%%%%%%%%%%%%%%%%%%%%%%%%%%%%%%%%%%%%%%
%%%%%%%%%%%%%%%%%%%%%%%%%%%%%%%%%%%%%%%%%%%
%%%%%%%%%%%%%%%%%%%%%%%%%%%%%%%%%%%%%%%%%%%
% Tema 5: Limites y diferenciabilidad
\item
    Dado el campo escalar
    \[
        \funcion{f}{\R^2}{\R}{(x,y)}{\sen(xy),}
    \]
    comprobar que es diferenciable con continuidad en todo $\R^2$. Además, hallar su derivada en el punto $(\pi,1)$.
    
\begin{respuesta}
    Para determinar si la función es diferenciable con continuidad, hallemos primero sus derivadas parciales, las cuales son:
    \[
        D_1f(x,y)=y\cos(xy)
        \texty
        D_2f(x,y)=x\cos(xy),
    \]
    para todo $(x,y)\in\R^2$. Dado que sus derivadas parciales son continuas, se tiene que $f$ es diferenciable con continuidad, además, 
    \[
        \nabla f(x,y)=(y\cos(xy),x\cos(xy)),
    \]
    para todo $(x,y)\in\R^2$, por lo tanto,
    \[
        \funcion{f'(x,y)}{\R^2}{\R}{(h,k)}{\nabla f(x,y)\cdot(h,k)=y\cos(xy)h+x\cos(xy)k),}
    \]
    para todo $(x,y)\in\R^2$. Con esto, se tiene $f'(\pi,1)$ es la función dada por
    \[
        \funcion{f'(\pi,1)}{\R^2}{\R}{(h,k)}{f'(\pi,1)(h,k)=-h-\pi k.}\qedhere
    \]
\end{respuesta}

%%%%%%%%%%%%%%%%%%%%%%%%%%%%%%%%%%%%%%%%%%%
%%%%%%%%%%%%%%%%%%%%%%%%%%%%%%%%%%%%%%%%%%%
%%%%%%%%%%%%%%%%%%%%%%%%%%%%%%%%%%%%%%%%%%%
% Tema 5: Limites y diferenciabilidad
\item
    Dado el campo escalar 
    \[
        \funcion{f}{\R^2}{\R}{(x,y)}
            {\begin{cases}
                \dfrac{xy}{x^2+y^2} & \text{si }(x,y)\neq(0,0), \\
                0 & \text{si }(x,y)=(0,0).
            \end{cases}}
    \]
    Determinar las derivadas parciales de $f$ en $(0,0)$ y comprobar que $f$ no es diferenciable en $(0,0)$.
    
\begin{respuesta}
    Para las derivadas parciales, se tiene que
    \[
        D_1f(0,0)
        = f'((0,0);(1,0))
        = \lim_{t\to 0} \frac{f((0,0)+t(1,0))-f(0,0)}{t}
        = \lim_{t\to 0} \frac{f(t,0)-0}{t} = 0.
    \]
    De manera análoga, se obtiene que
    \[
        D_2f(0,0) = 0,
    \]
    de donde
    \[
        \nabla f(0,0)=(0,0)
    \]
    
    Ahora, para analizar su diferenciablidad en $(0,0)$, analicemos 
    \[
        \lim_{(h,k)\to (0,0)}\frac{f((0,0))+(h,k))-f(0,0)-\nabla f(0,0)\cdot (h,k)} {\norma{(h,k)}},
    \]
    así, tenemos que
    \begin{align*}
        \lim_{(h,k)\to (0,0)}\frac{f((0,0))+(h,k))-f(0,0)-\nabla f(0,0)\cdot (h,k)} {\norma{(h,k)}}
            & =  \lim_{(h,k)\to (0,0)}\frac{\dfrac{hk}{h^2+k^2}-(0,0)\cdot (h,k)}{\norma{(h,k)}}\\
            & =  \lim_{(h,k)\to (0,0)}\frac{hk}{(\sqrt{h^2+k^2})^3}.
    \end{align*}
    Si consideramos el límite mediante la curva de ecuación $y=x$, tenemos que este límite no existe, por lo tanto, se tiene que la función no es diferenciable.
\end{respuesta}



%%%%%%%%%%%%%%%%%%%%%%%%%%%%%%%%%%%%%%%%%%%
%%%%%%%%%%%%%%%%%%%%%%%%%%%%%%%%%%%%%%%%%%%
%%%%%%%%%%%%%%%%%%%%%%%%%%%%%%%%%%%%%%%%%%%
% Tema 6: Derivadas parciales y gradiente
\item 
    Sean el campo escalar
	\[
		\funcion{f}{\R^2}{\R}{x}{\|x\|^2}
	\]
	y la trayectoria
	\[
		\funcion{r}{\lcl0,1\rcl}{\R}{t}{(t,t^2).}
	\]
	Muestre que se cumple que
	\[
		(f\circ r)'(t)=\nabla f(r(t))\cdot r'(t)
	\]
	para todo $t\in \lcl0,1\rcl$.

% \begin{respuesta}
% 	Por la definición de derivada, de una trayectoria, se tiene que
% 	\[
% 		r'(t)=(1,2t)
% 	\]
% 	para cada $t\in\R$. Además,
% 	\begin{align*}
% 		(f\circ r)(t)
% 			&=f(r(t))\\
% 			&=f(t,t^2)\\
% 			&=t^2 +t^4,
% 	\end{align*}
% 	de donde
% 	\[
% 		(f\circ r)'(t)=2t+4t^3
% 	\]
% 	para todo $t\in \lcl0,1\rcl$. Por otro lado,
% 	\[
% 		\nabla f(x)=2(x_1,x_2)
% 	\]
% 	para todo $x\in\R^2$. Finalmente, se tiene que
% 	\[
% 		\nabla f(r(t))\cdot r'(t)
% 		=\nabla f(t,t^2)\cdot r'(t)
% 		=(2(t,t^2))\cdot (1,2t)
% 		=2t+4t^3
% 	\]
% 	para cada $t\in\R$. Así,
% 	\[
% 		(f\circ r)'(t)=\nabla f(r(t))\cdot r'(t)
% 	\]
% 	para todo $t\in \lcl0,1\rcl$.
% \end{respuesta}

%%%%%%%%%%%%%%%%%%%%%%%%%%%%%%%%%%%%%%%%%%%
%%%%%%%%%%%%%%%%%%%%%%%%%%%%%%%%%%%%%%%%%%%
%%%%%%%%%%%%%%%%%%%%%%%%%%%%%%%%%%%%%%%%%%%
% Tema 6: Derivadas parciales y gradiente
\item Sean 
	\[
        \funcion{f}{\R^2}{\R}
        {(x,y)}{f(x,y) = x^2 y - y^2,}
	\]
    y 
    \[
        \funcion{r}{\R}{\R^2}
        {t}{r(t) = (\sen(t), e^t).}
	\]
    Hallar $(f \circ r)'(t)$ para todo $t\in\R$.

\begin{respuesta}
    Notemos que $f$ es derivable en $\R^2$ y $r$ es derivable en $\R$, por lo tanto $f\circ r$ es derivable en $\R$. Tenemos además que
    \[
        r'(t)=(\cos(t), e^t)
    \]
    para todo $t\in\R$ y
    \[
        \nabla f(x,y) = (2xy, x^2 - 2y),
    \]
    para todo $(x,y)\in\R^2$, por lo tanto
    \begin{align*}
        (f \circ r)'(t) 
            &= \nabla f(r(t)) \cdot r'(t)\\
    		&= (2\sen(t) e^t, \sen^2(t) - 2e^t)\cdot (\cos(t), e^t)\\
            &= e^t(\sen(2t) + \sen^2(t) - 2e^t),
    \end{align*}
    para todo $t\in\R$.
\end{respuesta}

%%%%%%%%%%%%%%%%%%%%%%%%%%%%%%%%%%%%%%%%%%%
%%%%%%%%%%%%%%%%%%%%%%%%%%%%%%%%%%%%%%%%%%%
%%%%%%%%%%%%%%%%%%%%%%%%%%%%%%%%%%%%%%%%%%%
% Tema 6: Derivadas parciales y gradiente
\item
Hallar la derivada direccional de la función:
	\[
        \funcion{f}{\R^3}{\R}{(x,y,z)}
        {x+y+z^2}
    \]
en $a=(1,1,-2)$ a lo largo de la curva C que resulta de la intersección de la superficie $S_1:2x^2-2y^2+z^2=4$ con el plano $S_2:x+y-z=4$  

    \begin{respuesta}
Aplicando la definición de la regla de la cadena para la derivada direccional a lo largo de una trayectoria $\alpha$,
	\[
        (f\circ \alpha)'(t)=\nabla f(r(t)) \cdot T(t)
    \]
conociendo que:
	\[
        T(t)=\dfrac{\alpha'(t)}{\|\alpha'(t)\|}
    \]
En un punto dado , el vector gradiente es normal a la curva de nivel que pasa por ese punto. Ahora si se tiene un punto en una superficie, hay infinitas curvas de nivel que pasan por ese punto, y todas comparten el mismo vector normal, de donde se puede deducir que el vector gradiente de una superficie, es ortogonal al plano tangente en ese punto. Es decir, que si se tienen 2 superficies distintas que se intersecan en un mismo punto, se tendrán dos vectores linealmente independientes que a su vez resultan ortogonales al vector tangente de la curva de intersección de las dos superficies. 
	\[
    \begin{aligned}
        r'(t) & = \nabla S_1 \times \nabla S_1\\
        & = \begin{bmatrix}
    	4 & -4 & -4\\
   		1 & 1 & -1
		\end{bmatrix}\\
        & = (8,0,8)
    \end{aligned}
    \]
El vector tangente unitario será:
	\[
        T(t)=\dfrac{(8,0,8)}{\sqrt{128}}
    \]
Por otro lado,
	\[
    \begin{aligned}
	\nabla f(r(t)) & = \nabla f(x,y,z)\\
    & = (1,1,2z)
    \end{aligned}
    \]
Por lo tanto,
	\[
    \begin{aligned}
        (f\circ r)'(t) & = (1,1,-4) \cdot \dfrac{(8,0,8)}{\sqrt{128}}\\
        & = \dfrac{-24}{\sqrt{128}}
        \end{aligned}
    \]
\end{respuesta}

%%%%%%%%%%%%%%%%%%%%%%%%%%%%%%%%%%%%%%%%%%%
%%%%%%%%%%%%%%%%%%%%%%%%%%%%%%%%%%%%%%%%%%%
%%%%%%%%%%%%%%%%%%%%%%%%%%%%%%%%%%%%%%%%%%%
% Tema 6: Derivadas parciales y gradiente
\item 
    Sea 
    \[
        \funcion{f}{\R^3}{\R}{(x,y,z)}
        {\sqrt[2]{x^2+y^2}+(x^2+y^2)^\frac{3}{2}-z.}
    \]
    Encontrar un vector $V \in \R^3$ normal a la superficie de nivel $L_0(f)$ en cualquier punto $(x,y,z) \neq 0$ de la superficie. Hallar el coseno del ángulo $\theta$ formado por el vector $V$ y el eje $z$ y determinar el límite de $\cos(\theta)$ cuando $(x,y,z) \to (0,0,0)$  
    
\begin{respuesta}
    Sea $(x,y,z)\in L_0(f)$ tal que $(x,y,z)\neq 0$. Tenemos que $\nabla f(x,y,z)$ es normal a $L_0(f)$ en el punto $(x,y,z)$. Por lo tanto, definimos
    \begin{align*}
        V
            & =\nabla f(x,y,z)\\
            & =\left(\frac{3x^3+3xy^2+x}{\sqrt{x^2+y^2}},\frac{3y^3+3x^2y+y}{\sqrt{x^2+y^2}},-1\right).
    \end{align*}
    Así, el coseno del ángulo $\theta$ que forma el vector $V$ con el eje $z$ está dado por
    \begin{align*}
        \cos (\theta)
            & = \left( \frac{V \cdot (0,0,1)}{\|V\| \|(0,0,1)\|} \right) \\
            & = -\frac{1}{9(x^2+y^2)^2+6(x^2+y^2)+2},
    \end{align*}
    de donde tenemos que
    \begin{align*}
        \lim_{(x,y,z) \to (0,0,0)}\cos (\theta)
            & = \lim_{(x,y,z) \to (0,0,0)} -\frac{1}{9(x^2+y^2)^2+6(x^2+y^2)+2} \\
            & = -\frac{1}{2}.
    \end{align*}
\end{respuesta}

%%%%%%%%%%%%%%%%%%%%%%%%%%%%%%%%%%%%%%%%%%%
%%%%%%%%%%%%%%%%%%%%%%%%%%%%%%%%%%%%%%%%%%%
%%%%%%%%%%%%%%%%%%%%%%%%%%%%%%%%%%%%%%%%%%%
% Tema 6: Derivadas parciales y gradiente
\item
    Dado $a\in\R$, hallar la ecuación cartesiana del plano tangente a la superficie 
    \[
       \{(x,y,z)\in\R^3 : xyz=a^3\}
    \]
    en cualquiera de sus puntos.
	
\begin{respuesta}
    La superficie es el conjunto de nivel $L_{a^3}(f)$ de la función 
    \[ 
        \funcion {f}{\R^3}{\R}{(x,y,z)}{xyz}
    \]
    y el gradiente de $f$ en cualquier punto $(x_0,y_0,z_0)$ de la superficie es
    \[
        \nabla f(x_0,y_0,z_0)=(y_0z_0,x_0z_0,x_0y_0).
    \]
    Por otra parte, tenemos que $\nabla f(x_0,y_0,z_0)$ es normal a la superficie en $(x_0,y_0,z_0)$ por lo tanto normal al plano tangente a la superficie en dicho punto, entonces, para todo $(x,y,z)$ que pertenece al plano, se tiene que
    \begin{align*}
        0
            & = ((x,y,z)-(x_0,y_0,z_0)) \cdot \nabla f(x_0,y_0,z_0) \\
            & = y_0z_0x+x_0z_0y+x_0y_0z-3x_0y_0z_0,
    \end{align*}
    por lo tanto, la ecuación cartesiana del plano es
    \[
        \frac{x}{3x_0}+\frac{y}{3y_0}+\frac{z}{3z_0}=1.
    \]
\end{respuesta}

%%%%%%%%%%%%%%%%%%%%%%%%%%%%%%%%%%%%%%%%%%%
%%%%%%%%%%%%%%%%%%%%%%%%%%%%%%%%%%%%%%%%%%%
%%%%%%%%%%%%%%%%%%%%%%%%%%%%%%%%%%%%%%%%%%%
% Tema 5: Limites y diferenciabilidad
\item 
    Dado el campo vectorial
	\[
		\funcion{F}{\R^2\setminus\{0\}}{\R^2}{x}{\l(\dfrac{x_1x_2}{x_1^2+x_2^2},\dfrac{x_1^2x_2^2}{x_1^2+x_2^2}\r).}
	\]
	Estudie la existencia del límite de $F$ cuando $x$ tiende a $(0,0)$.

%\begin{respuesta}
% 	Definamos las componentes de $F$:
% 	\[
% 	    \funcion{f_1}{\R^2}{\R}{x}{\dfrac{x_1x_2}{x_1^2+x_2^2}}
% 	    \texty
% 	    \funcion{f_2}{\R^2}{\R}{x}{\dfrac{x_1^2x_2^2}{x_1^2+x_2^2}}
% 	\]
% 	Para determinar si existe
% 	\[
% 		\Lim_{x\to 0} f_1(x)=\Lim_{x\to 0} \dfrac{x_1x_2}{x_1^2+x_2^2}
% 	\]
% 	analicemos sus límites por dos trayectorias. Sean la trayectorias
% 	\[
% 	    \funcion{\alpha_1}{\R\setminus\{0\}}{\R^2}{t}{(t,t)}
% 	    \yds
% 		\funcion{\alpha_2}{\R\setminus\{0\}}{\R^2}{t}{(t,-t).}
% 	\]
% 	Así,
% 	\[
% 		\Lim_{t\to 0}f_1(r_1(t))
% 			=\Lim_{t\to 0}f_1(t,t)
% 			=\Lim_{t\to 0}\dfrac{t^2}{2t^2}\\
% 			=\Lim_{t\to 0}\dfrac{1}{2} 
% 			=\dfrac{1}{2}\\
% 	\]
% 	y
% 	\[
% 		\Lim_{t\to 0}f_1(r_2(t))
% 			=\Lim_{t\to 0}f_1(t,-t)\\
% 			=\Lim_{t\to 0}\dfrac{-t^2}{2t^2}\\
% 			=\Lim_{t\to 0}-\dfrac{1}{2}\\
% 			=-\dfrac{1}{2}
% 	\]
% 	de donde
% 	\[
% 		\Lim_{x\to 0} f_1(x)=\Lim_{x\to 0} \dfrac{x_1x_2}{x_1^2+x_2^2}
% 	\]
% 	no existe. Como la primera componente no tiene límite en $(0,0)$, no hace falta estudiar la existencia del límite de la segunda componente, y se deduce que $\Lim_{x\to 0} F(x)$ no existe.
%\end{respuesta}

%%%%%%%%%%%%%%%%%%%%%%%%%%%%%%%%%%%%%%%%%%%
%%%%%%%%%%%%%%%%%%%%%%%%%%%%%%%%%%%%%%%%%%%
%%%%%%%%%%%%%%%%%%%%%%%%%%%%%%%%%%%%%%%%%%%
% Tema 5: Limites y diferenciabilidad
\item 
	Sea el campo vectorial $\func{f}{\R^n}{\R^m}$, donde $n$ y $m$ son enteros positivos. Demuestre que si $f$ es una aplicación lineal, entonces $f$ es una función continua en todo su dominio.

\begin{respuesta}
	Sea $a\in\R^n$. Vamos a demostrar que $f$ es continua en $a$. Sea $\vepsilon>0$ debemos hallar $\delta>0$ tal que para todo $x\in\R^n$ se cumpla que
	\[
	    0<\|x-a\|<\delta \Imp \|f(x)-f(a)\|<\vepsilon.
	\]
	Para cualquier $x\in\R^n$, se tiene que
	\begin{align*}
		f(x)-f(a)
		&=f\l(\sum_{i=1}^nx_ie^i\r)-f\l(\sum_{i=1}^na_ie^i\r)\\
		&=\sum_{i=1}^nx_if\l(e^i\r)-\sum_{i=1}^na_if\l(e^i\r)\\
		&=\sum_{i=1}^n(x_i-a_i)f\l(e^i\r),
	\end{align*}
	sea $k=\max\{\|f(e^i)\|: i=1,2,\ldots,n\}$; se sigue que
	\begin{align*}
		\|f(x)-f(a)\|
		&=\l\|\sum_{i=1}^n(x_i-a_i)f\l(e^i\r)\r\|\\
		&\leq\sum_{i=1}^n|x_i-a_i| \l\|f\l(e^i\r)\r\|\\
		&\leq k\sum_{i=1}^n|x_i-a_i|\\
		&\leq k \sqrt{n}\|x-a\|,
	\end{align*}
	donde en el ultimo paso utilizamos la desigualdad de Cauchy-Schwarz con los vectores de $\R^n$: $(|x_1-a_1|,|x_2-a_2|,\ldots,|x_n-a_n|)$ y $(1,1,\ldots,1)$. Finalmente, se tiene que si tomamos $\delta=\dfrac{\vepsilon}{k \sqrt{n}}>0$ se tiene que para todo $x\in\R^n$ tal que
	\[
	    0<\|x-a\|<\delta
	\]
	se cumple que
	\[
	    \|f(x)-f(a)\|\leq k \sqrt{n}\|x-a\|<  k \delta\sqrt{n}  =\vepsilon.
	\]
	Por la definición de continuidad, se tiene que $f$ es continua en $a$.
\end{respuesta}

%%%%%%%%%%%%%%%%%%%%%%%%%%%%%%%%%%%%%%%%%%%
%%%%%%%%%%%%%%%%%%%%%%%%%%%%%%%%%%%%%%%%%%%
%%%%%%%%%%%%%%%%%%%%%%%%%%%%%%%%%%%%%%%%%%%
% Tema 5: Limites y diferenciabilidad
\item 
	Dado el campo vectorial
	\[
		\funcion{F}{\R^2}{\R^2}{(x,y)}{(x^2-y,xy).}
	\]
	Suponiendo que $F$ es diferenciable, calcular $F'(1,1)$.

\begin{respuesta}
    Suponiendo que $F$ es diferenciable, se tiene que
    \[
        [F'(1,1)(h,k)] = J_F(1,1) [(h,k)]
    \]
    para todo $(h,k)\in\R^2$. Por lo tanto, empecemos calculando la matiz jacobiana de $F$,
    \[
        J_F(x,y)=
            \begin{bmatrix}
                2x  & -1\\
                y   & x
            \end{bmatrix},
    \]
    de donde
    \[
        J_F(1,1)=
            \begin{bmatrix}
                2  & -1\\
                1   & 1
            \end{bmatrix}.
    \]
    Así, tenemos que
    \[
        [F'(1,1)(h,k)] 
            = J_F(1,1) [(h,k)]
            = \begin{bmatrix}
                2  & -1\\
                1   & 1
            \end{bmatrix}
            \begin{bmatrix}
                h \\
                k 
            \end{bmatrix}
            = \begin{bmatrix}
                2h - k \\
                h + k 
            \end{bmatrix},
    \]
    para todo $(h,k)\in\R^2$, por lo tanto
    \[
        F'(1,1)(h,k) = (2h - k,h + k),
    \]
    para todo $(h,k)\in\R^2$. Con esto, se tiene que
    \[
        \funcion{F'(1,1)}{\R^2}{\R}{(h,k)}{F'(1,1)(h,k) = (2h - k,h + k).}
        \qedhere
    \]
\end{respuesta}



%%%%%%%%%%%%%%%%%%%%%%%%%%%%%%%%%%%%%%%%%%%
%%%%%%%%%%%%%%%%%%%%%%%%%%%%%%%%%%%%%%%%%%%
%%%%%%%%%%%%%%%%%%%%%%%%%%%%%%%%%%%%%%%%%%%
% Tema 5: Limites y diferenciabilidad
\item
    Dada la función vectorial 
    \[
        \funcion{F}{D}{\R^3}{(x_1,x_2)}{f(x_1,x_2)=\l(e^{x_1+x_2},\frac{1}{x_1},\sen(x_1x_2)\r)}
    \]
    donde $D=\{(x_1,x_2)\in\R^2:x_1\neq 0\}$. Para $a=(-1,0)$ y $y=\in\R^2$, determine la derivada de $F$ respecto al vector $y$ en el punto $a$.

\begin{respuesta} 
	Por definición
	\[
        [F'(a;y)]=J_f(a)[y].
    \]
	Se sabe que
    \[
        J_f(a)=\left( \begin{array}{lcr}
        D_{1}f_1(a)    &D_{2}f_1(a) \\
        D_{1}f_2(a)    &D_{2}f_2(a)\\
        D_{1}f_3(a)    &D_{2}f_3(a)
        \end{array}\right) 
    \]
    Se tiene que    
    \[
        J_f(x)= \begin{pmatrix}
        & e^{x_1+x_2}    &e^{x_1+x_2} \\
        & \frac{1}{x_2}  &0\\
        & x_2\cos(x_1x_2)    &x_1\cos(x_1x_2)
        \end{pmatrix} 
    \]
    Evaluando en a=(-1,0)
    \[
        J_f(a)= \begin{pmatrix}
        e^{-1} &e^{-1} \\
        -1     &0\\
        0      &-1
        \end{pmatrix}  
    \]
    Entonces la derivada de $F$ en $a$ en la dirección de cualquier vector $y$  es
    \[
        F'(a;y)= \begin{pmatrix}
        e^{-1} &e^{-1} \\
        -1     &0\\
        0      &-1
        \end{pmatrix}  
        \begin{pmatrix}
        y_1  \\
        y_2  
        \end{pmatrix}= 
        \begin{pmatrix}
        e^{-1}y_1+e^{-1}y_2\\
        -y_1     \\
        -y_2     
        \end{pmatrix} 
    \]
\end{respuesta}


%%%%%%%%%%%%%%%%%%%%%%%%%%%%%%%%%%%%%%%%%%%
%%%%%%%%%%%%%%%%%%%%%%%%%%%%%%%%%%%%%%%%%%%
%%%%%%%%%%%%%%%%%%%%%%%%%%%%%%%%%%%%%%%%%%%
% Tema 5: Limites y diferenciabilidad
\item 
	Dado el campo vectorial
	\[
		\funcion{F}{\R^2}{\R^2}{x}{(x_1^2-x_2,x_1x_2^2).}
	\]
	Mostrar que $F$ es diferenciable en todo $\R^2$.

\begin{respuesta}
	Primero calculemos su matriz Jacobiana:
	\[
		J_F(x)=
		\begin{pmatrix}
			2x_1&-1\\
			x_2^2&2x_1x_2
		\end{pmatrix}
	\]
	para todo $x\in\R^2$.
    
    Ahora, sea $a\in\R^2$, vamos a demostrar que $F$ es diferenciable en $a$. Procediendo como en el ejercicio anterior, se tiene que
    \[
        F'(a_1,a_2)(h_1,h_2)=(2a_1h_1-h_2, a_2^2h_1+2a_1a_2h_2),
    \]
    para todo $h\in\R^2$. Por otro lado, para cada $h\in\R^2$ no nulo, se tiene que
	\begin{align*}
		(a+h)-F(a)-F'(a)(h)
			&=((a_1+h_1)^2-(a_2+h_2),(a_1+h_1)(a_2+h_2)^2)
			-(a_1^2-a_2,a_1 a_2^2) \\
			&\phantom{=}-(2a_1h_1-h_2, a_2^2h_1+2a_1a_2h_2) \\
			&=\l(h_1^2, a_1h_2^2+2a_2h_1h_2+h_1h_2^2\r),
	\end{align*}
	por lo tanto
	\[
	    \frac{F(a+h)-F(a)-F'(a)(h)}{\|h\|}=\l(\frac{h_1^2}{\|h\|},\frac{a_1h_2^2+2a_2h_1h_2+h_1h_2^2}{\|h\|}\r).
	\]
	Dado que
	\[
		|h_i|\leq \|h\|
	\]
	para $i\in\{1, 2\}$, se sigue que
	\[
		0\leq \dfrac{h_1^2}{\|h\|}\leq \|h\|, 
		\qquad
		0\leq \dfrac{h_2^2}{\|h\|}\leq \|h\|
		\yds
		0\leq \dfrac{|h_1h_2|}{\|h\|}\leq \frac{\|h\|}{2}, 
	\]
	de donde, por el teorema del Sanduche, 
	\[
		\lim_{h\to 0}\dfrac{h_1^2}{\|h\|}=0
		\yds
		\lim_{h\to 0}\dfrac{a_1h_2^2+2 a_2h_1h_2+h_1h_2^2}{\|h\|}=0.
	\]
	De esta manera, se tiene que
	\[
		\lim_{h\to 0}\dfrac{F(a+h)-F(a)-F'(a)(h)}{\|h\|}=0
	\]
	lo que significa que la función $f$ es diferenciable en $a$.
\end{respuesta}

%%%%%%%%%%%%%%%%%%%%%%%%%%%%%%%%%%%%%%%%%%%
%%%%%%%%%%%%%%%%%%%%%%%%%%%%%%%%%%%%%%%%%%%
%%%%%%%%%%%%%%%%%%%%%%%%%%%%%%%%%%%%%%%%%%%
% Tema 5: Limites y diferenciabilidad
\item
    Dado el campo vectorial
    \[
        \funcion{F}{\R^2}{\R^2}{(x,y)}{f(x,y)=(x+y^2,y-2xy)}
    \]
    Verifique que $F$ es diferenciable en el punto $a=(2,1)$.

%%%%%%%%%%%%%%%%%%%%%%%%%%%%%%%%%%%%%%%%%%%
%%%%%%%%%%%%%%%%%%%%%%%%%%%%%%%%%%%%%%%%%%%
%%%%%%%%%%%%%%%%%%%%%%%%%%%%%%%%%%%%%%%%%%%
% Tema 6: Derivadas parciales y gradiente
\item
	Dados los campos vectoriales
	\[
 	    \funcion{G}{\R^2}{\R^3}{(r,s)}{(r+s,r-s,rs)}
 	    \texty
 	    \funcion{F}{\R^3}{\R^3}{(x,y,z)}{(x^2-y^2,xyz,x^2+y^2+z)}
	\]
	Determine $J_{F\circ G}(a)$, para $a\in\R^2$.
	
\begin{respuesta} 
	Si $G$ es diferenciable en $a$ y $F$ es diferenciable en $G(a)$, entonces
    \[ 
        J_{F\circ G} (a)= J_F (G(a))J_G (a)
    \]
    Se debe plantear la jacobiana de $F$ y de $G$
    \begin{align*}
        &J_F(x,y,z) = \begin{pmatrix}
        2x & -2y & 0 \\
        yz & xz & xy\\
        2x & 2y & 1
        \end{pmatrix}\\
        &J_g(r,s) = \begin{pmatrix}
        1 & 1  \\
        1 & -1 \\
        s & r
        \end{pmatrix}
    \end{align*}
    Evaluando la jacobiana de $F$ en $G(a)$
    \[ 
        J_F(G(r,s)) = \begin{pmatrix}
        2(r+s) & -2(r-s) & 0 \\
        (r-s)rs & (r+s)rs & (r+s)(r-s)\\
        2(r+s) & 2(r-s) & 1
        \end{pmatrix}
    \]
    Operando
    \[ 
        J_F(r,s) = \begin{pmatrix}
        2r+2s & -2r+2s & 0 \\
        r^2s-rs^2 & r^2s+rs^2 & r^2-s^2\\
        2r+2s & 2r-2s & 1
        \end{pmatrix}
	\]
    Entonces calculamos $(F\circ G)'(r,s)$ 
    \begin{align*}
        (F\circ G)'(r,s)= \begin{pmatrix}
        2r+2s & -2r+2s & 0 \\
        r^2s-rs^2 & r^2s+rs^2 & r^2-s^2\\
        2r+2s & 2r-2s & 1
        \end{pmatrix} \begin{pmatrix}
        1 & 1  \\
        1 & -1 \\
        s & r
        \end{pmatrix}\\
	    (F\circ G)'(r,s)= \begin{pmatrix}
        4s & 4r \\
        3r^2s-s^3 & r^3-3rs^2\\
        4r+s & 4s+r
        \end{pmatrix}
    \end{align*} 
    Si $a=(1,2)$, entonces
    \[ 
        (F\circ G)'(a)= \begin{pmatrix}
        8 & 4 \\
        -2 & -11\\
        6 & 9
        \end{pmatrix}
    \]
\end{respuesta}

%%%%%%%%%%%%%%%%%%%%%%%%%%%%%%%%%%%%%%%%%%%
%%%%%%%%%%%%%%%%%%%%%%%%%%%%%%%%%%%%%%%%%%%
%%%%%%%%%%%%%%%%%%%%%%%%%%%%%%%%%%%%%%%%%%%
% Tema 6: Derivadas parciales y gradiente
\item
	Dados los campos vectoriales
    \[
        \funcion{F}{\R^3}{\R^2}{(x,y,z)}{(x^2+y+z,2x+y+z^2)}
        \texty
        \funcion{G}{\R^3}{\R^3}{(u,v,w)}{(uv^2w^2,w^2\sen(v),u^2e^v)}
    \]
    Determinar la matriz jacobiana de $F\circ G$ en $(u,v,w)$.

% \begin{respuesta}
%     Sea,
%     \[
%         J_{F\circ G}(u,v,w)=J_F(G(u,v,w))J_G(u,v,w).
%     \]
%     De tal forma
%     \[
%         J_{F\circ G}(u,v,w)=
%         \left[
%         \begin{array}{ccc}
% 		0&1&1\\
% 		2&1&2u^2\\
% 		\end{array} 
%         \right]
%         \left[
%         \begin{array}{ccc}
% 		0&0&0\\
% 		0&w^2&0\\
%         2u&u^2&0\\
% 		\end{array} 
%         \right]=
%         \left[
%         \begin{array}{ccc}
% 		2u&u^2+w^2&0\\
% 		4u^3&w^2+2u^4&0\\
% 		\end{array} 
%         \right]
%     \]
% \end{respuesta}

%%%%%%%%%%%%%%%%%%%%%%%%%%%%%%%%%%%%%%%%%%%
%%%%%%%%%%%%%%%%%%%%%%%%%%%%%%%%%%%%%%%%%%%
%%%%%%%%%%%%%%%%%%%%%%%%%%%%%%%%%%%%%%%%%%%
% Tema 6: Derivadas parciales y gradiente
\item
    Considere las funciones
	\[
 	    \funcion{f}{\R^2}{\R}{(x,y)}{xy}
 	    \texty 
 	    \funcion{G}{\R^2}{\R^2}{(u,v)}{(\cos(v+u),\sen(v-u)).}
    \]
	Determine $D_1(f\circ G)$ y $D_2(f\circ G)$ en el punto $(u,v)\in\R^2$.
     
\begin{respuesta} 
	Dado $(u,v)\in\R^2$, tenemos que para $D_1(f\circ G)(u,v)$
	\[
	D_1(f\circ G)(u,v)=D_1f(G(u,v))D_1g_1(u,v)+D_2f(G(u,v)D_1g_2(u,v)
	        \]
		\[
	D_1g_1(u,v)=- \sen (v+u)\qquad  D_1g_2(u,v)=-\cos(v-u)
	        \]
    y para $(x,y)\in \R^2$
    	\[
	D_1f(x,y)=y\qquad D_2f(x,y)=x
	        \]
	por lo tanto:
		\[
	D_1f(G(u,v))=D_1f(\cos(v+u),\sen(v-u))=\sen(v-u)
	        \]
	y
	\[
	D_2f(G((u,v))=D_2f(\cos(v+u),\sen(v-u))=\cos(v+u)
	        \]
	así, se tiene que
    \begin{align*}
    D_1(f\circ G)(u,v)&=D_1f(G(u,v))D_1g_1(u,v)+D_2f(G(u,v))D_1g_2(u,v),\\
    &=(\sen(v-u))(-\sen(v+u))+(\cos(v+u))(-\cos(v-u)),\\
    &=-[\cos(v-u)\cos(v+u)+\sen(v-u)\sen(v+u)],\\
    &=-[\cos(v-u-v-u)],\\
    D_1(f\circ G)(u,v)&=-\cos(2u).
    \end{align*}
    Para hallar $D_2(f\circ G)(u,v)$
    \[
	D_2(f\circ G)((u,v))=D_1f(G(u,v))D_2g_1(u,v)+D_2f(G(u,v))D_2g_2(u,v)
	        \]
	En este caso necesitamos         
	\[
	D_2g_1(u,v)=-\sen(v+u)\qquad D_2g_2(u,v)=\cos(v-u)
	        \]
    por tanto
	 \begin{align*}
    D_2(f\circ G)(u,v)&=(\sen(v-u))(-\sen(v+u))+(\cos(v+u))(\cos(v-u)),\\
    &=\cos(v-u)\cos(v+u)-\sen(v-u)\sen(v+u),\\
    &=\cos(v-u+v+u),\\
    D_2(f\circ G)(u,v)&=\cos(2v).
    \end{align*}        
\end{respuesta}

%%%%%%%%%%%%%%%%%%%%%%%%%%%%%%%%%%%%%%%%%%%
%%%%%%%%%%%%%%%%%%%%%%%%%%%%%%%%%%%%%%%%%%%
%%%%%%%%%%%%%%%%%%%%%%%%%%%%%%%%%%%%%%%%%%%
% Tema 6: Derivadas parciales y gradiente
\item
    Supongamos que:
	\[
        \funcion{f}{\Omega}{\R}{(x,y,z)}
        {xy+yz+\phi \l( \dfrac{x}{y}\r)} \texty \funcion{G}{\R^3}{\R^3}{(r,s,t)}
        {\l(1+rse^t,rs^2e^{-t},r^2s \sen t\r)}
    \]
    Donde $\phi: \R \rightarrow \R$ es una función derivable y $\Omega=\{(x,y,z)\in\R^3:y\neq 0\} $. Determinar $D_2(f \circ G) $ en el punto $(r,s,t)=(2,1,0)$ sabiendo que $\phi' \l(\dfrac{3}{2}\r)=-1 $

\begin{respuesta}
    Para $(r,s,t) \in \R^3$ se tiene que:
	\[
	    D_2(f \circ G)(r,s,t)=D_1f(G(r,s,t))D_2g_1(r,s,t)+D_2f(G(r,s,t))D_2g_2(r,s,t)+D_3f(G(r,s,t))D_2g_3(r,s,t).
    \]
    Además
	\[
    	D_2g_1(r,s,t)=re^t,
    \]

    \[
	    D_2g_2(r,s,t)=2rse^{-t}
    \]
    y
    \[
	    D_2g_3(r,s,t)=r^2 \sen t.
    \]
    Para $(x,y,z) \in \R^3$,
    \[
	    D_1f(x,y,z)=\l(y + \phi' \l( \dfrac{x}{y} \r) \dfrac{1}{y}\r),
    \]
    
    \[
	    D_2f(x,y,z)=\l(x+z+ \phi' \l( \dfrac{x}{y} \r) \dfrac{-x}{y^2}\r)
    \]
    y
    \[
	    D_3f(x,y,z)=y.
    \]
    Por lo tanto,
    \[
        \begin{aligned}
        D_1f(G(r,s,t))&= D_1f\l(1+rse^t,rs^2e^{-t},r^2s \sen t\r)\\&=rs^2e^{-t} + \phi' \l( \dfrac{1+rse^t}{rs^2e^{-t}} \r) \dfrac{1}{rs^2e^{-t}},
        \end{aligned}
    \]
    
	\[
        \begin{aligned}
	    D_2f(G(r,s,t))&= D_2f\l(1+rse^t,rs^2e^{-t},r^2s \sen t\r)\\&=1+rse^t + r^2s \sen t + \phi' \l( \dfrac{1+rse^t}{rs^2e^{-t}} \r) \dfrac{-1-rse^t}{\left(rs^2e^{-t}\right)^2}
        \end{aligned}
    \]
    y
	\[
        \begin{aligned}
	    D_3f(G(r,s,t))&= D_3f\l(1+rse^t,rs^2e^{-t},r^2s \sen t\r)\\&=rs^2e^{-t}.
        \end{aligned}
    \]
    Se tiene entonces que,
	\[
	    D_2(f \circ G)(r,s,t)=D_1f(G(r,s,t))D_2g_1(r,s,t)+D_2f(G(r,s,t))D_2g_2(r,s,t)+D_3f(G(r,s,t))D_2g_3(r,s,t)
    \]
    \[
        \begin{aligned}
	    D_2(f \circ G)(r,s,t)&=\l(rs^2e^{-t} + \phi' \l( \dfrac{1+rse^t}{rs^2e^{-t}} \r) \dfrac{1}{rs^2e^{-t}}\r) \l(re^t\r) \\ &+ \l(1+rse^t + r^2s \sen t + \phi' \l( \dfrac{1+rse^t}{rs^2e^{-t}} \r) \dfrac{-1-rse^t}{\left(rs^2e^{-t}\right)^2}\r) \l(2rse^{-t}\r) \\ &+ \l(rs^2e^{-t}\r) \l(r^2 \sen t\r).
        \end{aligned}
    \]
    Evaluando en el punto $(r,s,t)=(2,1,0)$ y conociendo que $\phi' \l(\dfrac{3}{2}\r)=-1 $,
    \[
        \begin{aligned}
	    D_2(f \circ G)(2,1,0)&=\l(2+\phi' \l( \dfrac{3}{2}\r) \l(\dfrac{1}{2}\r)\r)(2) + \l(3+ \phi' \l( \dfrac{3}{2}\r)\l(\dfrac{-3}{4}\r)\r)(4) + (2)(0)\\
        &=\l(2- \dfrac{1}{2}\r)(2) + \l(3+ \dfrac{3}{4}\r)(4) + 0\\
        &=3+15+0\\
        &=18
        \end{aligned}
    \]
\end{respuesta}

%%%%%%%%%%%%%%%%%%%%%%%%%%%%%%%%%%%%%%%%%%%
%%%%%%%%%%%%%%%%%%%%%%%%%%%%%%%%%%%%%%%%%%%
%%%%%%%%%%%%%%%%%%%%%%%%%%%%%%%%%%%%%%%%%%%
% Tema 6: Derivadas parciales y gradiente
\item
    Supongamos que:
	\[
        \funcion{f}{\R^2}{\R}{(x,y)}
        {x^2+y^2}
    \]
    y
    \[
        \funcion{G}{\R^2}{\R^2}{(u,v)}
        {\l(e^{u+v},u^2+v\r)}
    \]
    Determinar $D_1(f \circ G) $ en el punto $(u,v)=(1,1) \in \R^2$.

\begin{respuesta}
    Para $(u,v) \in \R^2$ se tiene que:
	\[
	D_1(f \circ G)(u.v)=D_1f(G(u,v)D_1g_1(u,v)+D_2f(G(u,v))D_1g_2(u,v)
    \]
    Además
	\[
	D_1g_1(u,v)=e^{u+v}
    \]
    y
    \[
	D_1g_2(u,v)=2u
    \]
    y para $(x,y) \in \R^2$,
    \[
	D_1f(x,y)=2x
    \]
    y
    \[
	D_2f(x,y)=2y
    \]

    Por lo tanto,
    \[
    D_1f(G(r,s,t))=D_1f\l(e^{u+v},u^2+v\r)=2\l(e^{u+v}\r)
    \]
    y
	\[
	D_2f(G(r,s,t))= D_2f\l(e^{u+v},u^2+v\r)=2\l(u^2+v\r)
    \]

    Se tiene entonces que,
    \[
    \begin{aligned}
	D_1(f \circ G)(u.v)&=D_1f(G(u,v)D_1g_1(u,v)+D_2f(G(u,v))D_1g_2(u,v)\\ &=2\l(e^{u+v}\r)\l(e^{u+v}\r) + 2\l(u^2+v\r)\l(2u\r)
    \end{aligned}
    \]
    Evaluando en el punto $(u,v)=(1,1) \in \R^2$,
    \[
    \begin{aligned}
	D_1(f \circ G)(u.v)&=2\l(e^{1+1}\r)\l(e^{1+1}\r) + 2\l(1^2+1\r)\l(2(1)\r)\\ &=2e^4+8
    \end{aligned}
    \]
\end{respuesta}

%%%%%%%%%%%%%%%%%%%%%%%%%%%%%%%%%%%%%%%%%%%
%%%%%%%%%%%%%%%%%%%%%%%%%%%%%%%%%%%%%%%%%%%
%%%%%%%%%%%%%%%%%%%%%%%%%%%%%%%%%%%%%%%%%%%
% Tema 7: Funcion implicita y Hessiana
\item Sea 
	\[
        \funcion{f}{\R^2}{\R}
        {(x,y)}{f(x,y) = 
            \begin{cases}
            e^x + e^y + \dfrac{xy}{x^2 + y^2}& \text{si } (x, y) \neq (0, 0),\\
            2 & \text{si } (x, y) = (0, 0).
            \end{cases}}
	\]
    Hallar $D_{1,1}f(0, 0)$ y $D_{2,1}f(0, 0)$.

\begin{respuesta}
    Calculemos en primer lugar $D_1f$. Nótese que para todo $(x, y)\in\R^2\setminus\{(0,0)\}$
    \[
        D_1f(x,y) = e^x + \dfrac{y^3 - x^2 y}{(x^2 + y^2)^2},
    \]
    mientras que si $(x, y) = (0, 0)$, se tiene que
    \[
        \lim_{h \to 0}{\dfrac{f(0 + h, 0) - f(0, 0)}{h}} = \lim_{h \to 0}{\dfrac{f(h, 0) - f(0, 0)}{h}} = \lim_{h \to 0}{\dfrac{e^h - 1}{h}} = 1.
    \]
    Por tanto, 
    \[
        \funcion{D_1f}{\R^2}{\R^2}
        {(x,y)}{D_1f(x,y) = 
            \begin{cases}
            e^x + \dfrac{y^3 - x^2 y}{(x^2 + y^2)^2}& \text{si } (x, y) \neq (0, 0),\\
            1 & \text{si } (x, y) = (0, 0).
            \end{cases}}
	\]
    Calculemos ahora $D_{1,1}(f)(0,0)$,
    \begin{align*}
        D_{1,1}f(0,0) 
            & = \lim_{h \to 0}{\dfrac{D_1f(0+h, 0) - D_1f(0, 0)}{h}}\\
    		& = \lim_{h \to 0}{\dfrac{D_1f(h, 0) - D_1f(0, 0)}{h}}\\
            & = \lim_{h \to 0}{\dfrac{e^h - 1}{h}} = 1,
    \end{align*}
    es decir, $D_{1,1}f(0,0) = 1$.
    Finalmente, tenemos que
    \begin{align*}
        \lim_{h \to 0}{\dfrac{D_1f(0, 0+ h) - D_1f(0, 0)}{h}}
    	    &= \lim_{h \to 0}{\dfrac{D_1f(0, h) - D_1f(0, 0)}{h}}\\
            &= \lim_{h \to 0}{\dfrac{1}{h^2}} \\
            &= +\infty,
    \end{align*}
    es decir, $D_{2,1}f(0,0)$ no existe.
\end{respuesta}

%%%%%%%%%%%%%%%%%%%%%%%%%%%%%%%%%%%%%%%%%%% FC
%%%%%%%%%%%%%%%%%%%%%%%%%%%%%%%%%%%%%%%%%%%
%%%%%%%%%%%%%%%%%%%%%%%%%%%%%%%%%%%%%%%%%%%
% Tema 7: Funcion implicita y Hessiana
\item
Dado el campo escalar:
    \[
        \funcion{f}{\R^3}{\R}{(x,y,z)}{\dfrac{x+y}{y-z}}
    \]
Obtenga su matriz Hessiana en el punto $a=(0,1,-1)$:
\begin{respuesta}
    Por definición de la matriz Hessiana en $\R^3$.
    \[
        H_f(x,y,z)=\l(J_{\nabla f}(x,y,z)\r)^T= \begin{pmatrix}
        D_{1,1}f(x,y,z) & D_{1,2}f(x,y,z) & D_{1,3}f(x,y,z)\\
        D_{2,1}f(x,y,z) & D_{2,2}f(x,y,z) & D_{2,3}f(x,y,z)\\
        D_{3,1}f(x,y,z) & D_{3,2}f(x,y,z) & D_{3,3}f(x,y,z)
        \end{pmatrix}
    \]
Se tiene entonces que las segundas derivadas parciales de $f$ respecto a $x$,
    \[
        D_{1,1}f(x,y,z)=0
        \yds
        D_{2,1}f(x,y,z)=-\dfrac{1}{(y-z)^2}
        \yds
        D_{3,1}f(x,y,z)=\dfrac{1}{(y-z)^2},
    \]
para $y-z \neq 0$. Respecto a $y$,
    \[
        D_{1,2}f(x,y,z)=-\dfrac{1}{(y-z)^2}
        \yds
        D_{2,2}f(x,y,z)=\dfrac{2(x+z)}{(y-z)^3}
        \yds
        D_{3,2}f(x,y,z)=-\dfrac{(y-z)+2(x+z)}{(y-z)^3},
    \]
para $y-z \neq 0$. Respecto a $z$,
    \[
        D_{1,3}f(x,y,z)=\dfrac{1}{(y-z)^2}
        \yds
        D_{2,3}f(x,y,z)=\dfrac{(y-z)-2(x+z)}{(y-z)^3}
        \yds
        D_{3,3}f(x,y,z)=\dfrac{2(y+x)}{(y-z)^3},
    \]
para $y-z \neq 0$. Evaluando sus segundas derivadas parciales en $a=(0,1,-1)$, se obtiene la matriz Hessiana de $f$ en $a$,
    \[
        H_f(a)= \begin{pmatrix}
        0 & -\dfrac{1}{4} & \dfrac{1}{4}\\
        -\dfrac{1}{4} & -\dfrac{1}{4} & 0\\
        \dfrac{1}{4} & 0 & \dfrac{1}{4}
        \end{pmatrix}
    \]
\end{respuesta}

%%%%%%%%%%%%%%%%%%%%%%%%%%%%%%%%%%%%%%%%%%%
%%%%%%%%%%%%%%%%%%%%%%%%%%%%%%%%%%%%%%%%%%%
%%%%%%%%%%%%%%%%%%%%%%%%%%%%%%%%%%%%%%%%%%%
% Tema 7: Funcion implicita y Hessiana
\item  
    Dada la función
    \[
        \funcion{f}{\R^3}{\R}{(x,y,z)}{f(x,y,z)=\sen(xyz),}
    \]
    determinar $\nabla f(x,y,z)$ y $H_f(x,y,z)$ para todo $(x,y,z)\in\R^3$.

% \begin{respuesta}
%     Primero calculemos el gradiente de la función, así, se tiene que
%     \[
%         \nabla f(x,y,z)
%         = (yz\sen(xyz),xz\sen(xyz),xy\sen(xyz)),
%     \]
%     para $(x,y,z)\in\R^3$. Ahora, determinemos la matriz jacobiana del gradiente:
%     \[
%         J_{\nabla f}(x,y,z)
%         = \begin{pmatrix}
%         (y^2z^2\cos(xyz) & xy\cos(xyz) z^2 + \sen(xyz) z & x z \cos(xyz) y^2 + \sen(xyz) y\\
%         x y\cos(xyz) z^2 + \sen(xyz) z & x^2 z^2 \cos(xyz) & y z \cos(xyz) x^2 + \sen(xyz) x\\
%         x z \cos(xyz) y^2 + \sen(xyz) y & y z \cos(xyz) x^2 + \sen(xyz) x & x^2 y^2 \cos(xyz))
%         \end{pmatrix},
%     \]
%     para $(x,y,z)\in\R^3$. Por lo tanto, transponemos esta matriz par hallar la matriz hessiana; tenemos que
%     \[
%         H_{f}(x,y,z)
%         = \begin{pmatrix}
%         (y^2z^2\cos(xyz) & xy\cos(xyz) z^2 + \sen(xyz) z & x z \cos(xyz) y^2 + \sen(xyz) y\\
%         x y\cos(xyz) z^2 + \sen(xyz) z & x^2 z^2 \cos(xyz) & y z \cos(xyz) x^2 + \sen(xyz) x\\
%         x z \cos(xyz) y^2 + \sen(xyz) y & y z \cos(xyz) x^2 + \sen(xyz) x & x^2 y^2 \cos(xyz))
%         \end{pmatrix},
%     \]
%      para $(x,y,z)\in\R^3$.
% \end{respuesta}

%%%%%%%%%%%%%%%%%%%%%%%%%%%%%%%%%%%%%%%%%%% JO
%%%%%%%%%%%%%%%%%%%%%%%%%%%%%%%%%%%%%%%%%%%
%%%%%%%%%%%%%%%%%%%%%%%%%%%%%%%%%%%%%%%%%%%
% Tema 7: Funcion implicita y Hessiana
\item 
    Sea el campo escalar
    \[
        \funcion{f}{\R^2}{\R}{(x,y)}{x^2e^{x-y}.}
    \]
\begin{enumerate}
    \item Encuentre la aproximación lineal en el punto $(1,0)$.
    \item Encuentre la aproximación cuadrática en el punto $(1,0)$.
    \item Calcule con ambas aproximaciones estime $f(1,0)$ y $f(1,1)$; y compare con su valor real.
\end{enumerate}

\begin{respuesta}
    $f$ es diferenciable pues es la multiplicación de un polinomio con una función exponencial. Así,
    \begin{align*}
        \nabla f(a)&= \l(a_1 (a_1 + 2) e^{a_1 - a_2}, a_1^2 (-e^{a_1 - a_2})\r)\\
        H_f(a)&=
        \begin{pmatrix}
            (a_1^2 e^{a_1 - a_2} + 4 a_1 e^{a_1 - a_2} + 2 e^{a_1 - a_2} & a_1^2 (-e^{a_1 - a_2}) - 2 a_1 e^{a_1 - a_2}\\
            a_1^2 (-e^{a_1 - a_2}) - 2 a_1 e^{a_1 - a_2} & a_1^2 e^{a_1 - a_2})  
        \end{pmatrix}
    \end{align*}
    para todo $a\in\R^2$. Por lo tanto la aproximación lineal de $f$ en el punto $a=(1,0)$ para $(x,y)\in\R^2$ cerca de $a$ es
    \begin{align*}
        f(x,y)
        &\approx f(1,0)+\nabla f (a) \cdot ((x,y)-a)\\
        &\approx e+(3e,-e) \cdot (x-1,y)\\
        &\approx e(1+3x-3-y)\\
        &\approx e(1+3x-y-3).
    \end{align*}
    Además, la aproximación cuadrática de $f$ en el punto $a$ para $(x,y)\in\R^2$ cerca de $a$ es
    \begin{align*}
        f(x,y)
        &\approx f(1,0)+\nabla f (a) \cdot ((x,y)-a)+\dfrac1 2 ((x,y)-a) H_f(a) ((x,y)-a)^t\\
        &\approx e(1+3x-y-3)+\dfrac1 2 (x-1,y) 
        \begin{pmatrix}
            7e & -3e\\ 
            -3e &e
        \end{pmatrix}
        (x-1,y)^t\\
        &\approx e(1+3x-y-3)+\dfrac{e}2(7 x^2 - 6 x y - 14 x + y^2 + 6 y + 7)\\
        &\approx \dfrac{e}{2}(7 x^2 - 6 x y - 8 x + y^2 + 4 y + 3).
    \end{align*}
    Por lo tanto, la aproximación lineal de $f$ en $(1,0)$ para $(1,0)$ es
    \[
        e(1+3(1)-(0)-3)=e
    \]
    y la cuadrática es
    \[
        \dfrac{1}{2}(7-8+3)=e.
    \]
    Lo que indica que ambas aproximaciones coinciden con $f(1,0)=e$, lo cual es lo esperado ya que estás aproximaciones entre más cerca al punto $a=(1,0)$ son mejor. La ventaja es que para números cercanos a $(1,0)$ también funcionan. En efecto, ambas aproximaciones para $(1,1)$ son
    \[
        e(1-1)= 0 \yds \dfrac{e}2\approx 1.36,
    \]
    respectivamente. Comparado con el valor $f(1,1)=1$ la aproximación cuadrática es mejor. En la práctica es más cómodo trabajar con aproximaciones pues son polinomios aunque se pierda un poco de exactitud.
\end{respuesta}

%%%%%%%%%%%%%%%%%%%%%%%%%%%%%%%%%%%%%%%%%%%
%%%%%%%%%%%%%%%%%%%%%%%%%%%%%%%%%%%%%%%%%%%
%%%%%%%%%%%%%%%%%%%%%%%%%%%%%%%%%%%%%%%%%%%
% Tema 7: Funcion implicita y Hessiana
\item  
    Dada la función
    \[
        \funcion{f}{\R^2}{\R}{(x,y)}{f(x,y)=\sen(x)\cos(y),}
    \]
    determinar la aproximación lineal y la aproximación cuadrática de la función alrededor del punto $(0,0)$.

\begin{respuesta}
    Llamemos $F_1$ a la aproximación lineal de $f$ en $(0,0)$, tenemos que
    \[
        F_1(x,y) = f(0,0) + \nabla f(0,0) \cdot ((x,y)-(0,0)),
    \]
    para todo $(x,y)\in\R ^2$, por lo tanto, empecemos calculando el gradiente de $f$,
    \[
        \nabla f(x,y) = (\cos(x) \cos(y), -\sen(x) \sen(y) ),
    \]
    para $(x,y)\in\R^2$. Con esto, se tiene que
    \[
        \nabla f(0,0) = (1,0).
    \]
    Por lo tanto, la aproximación lineal de $f$ en $(0,0)$ es
    \begin{align*}
        F_1(x,y)
            & = f(0,0) + \nabla f(0,0) \cdot ((x,y)-(0,0))\\
            & = 0 + (1,0)\cdot (x,y)\\
            & = x,
    \end{align*}
    para $(x,y)\in\R^2$.
    
    Ahora, llamemos $F_2$ a la aproximación cuadrática de $f$ en $(0,0)$, tenemos que
    \[
        F_2(x,y) = f(0,0) + \nabla f(0,0) \cdot ((x,y)-(0,0))  + \frac{1}{2!}[(x,y)-(0,0)]^T Hf(0,0) [(x,y)-(0,0)],
    \]
    para todo $(x,y)\in\R ^2$, por lo tanto, empecemos calculando la jacobiana de $f$,
    \[
        H_f(x,y) = 
        \begin{pmatrix}
            -\sen(x)\cos(y) & -\cos(x)\sen(y)\\
            -\cos(x) \sen(y) & \sen(x) (-\cos(y)))
        \end{pmatrix},
    \]
    para $(x,y)\in\R^2$. Con esto, se tiene que
    \[
        H_f(0,0) = \begin{pmatrix}
            0 & 0\\
            0 & 0
        \end{pmatrix},
    \]
    Por lo tanto, la aproximación cuadrática de $f$ en $(0,0)$ es
    \begin{align*}
        F_2(x,y)
            & = f(0,0) + \nabla f(0,0) \cdot ((x,y)-(0,0))+ \frac{1}{2!}[(x,y)-(0,0)]^T Hf(0,0) [(x,y)-(0,0)]\\
            & = 0 + (1,0)\cdot (x,y) +
            \frac 1 2 
            \begin{pmatrix}
            x&y
            \end{pmatrix}
            \begin{pmatrix}
            0 & 0\\
            0 & 0
            \end{pmatrix}
            \begin{pmatrix}
            x\\y
            \end{pmatrix}\\
            & = x,
    \end{align*}
    para $(x,y)\in\R^2$.
\end{respuesta}

%%%%%%%%%%%%%%%%%%%%%%%%%%%%%%%%%%%%%%%%%%%
%%%%%%%%%%%%%%%%%%%%%%%%%%%%%%%%%%%%%%%%%%%
%%%%%%%%%%%%%%%%%%%%%%%%%%%%%%%%%%%%%%%%%%%
% Tema 7: Funcion implicita y Hessiana
\item  
    Utilizando el ejercicio anterior, encuentre una estimación del valor de $\sen(0.2)\cos(0.1)$. Para esto, utilice la aproximación lineal y la aproximación cuadrática de la función.

% \begin{respuesta}
%     Tomemos la aproximación lineal de la función del ejercicio anterior, así, se tiene que
%     \[
%         \sen(0.2)\cos(0.1)=f(0.2,0.1)\approx F_1(0.2,0.1) = 0.2.
%     \]
% \end{respuesta}

%%%%%%%%%%%%%%%%%%%%%%%%%%%%%%%%%%%%%%%%%%%
%%%%%%%%%%%%%%%%%%%%%%%%%%%%%%%%%%%%%%%%%%%
%%%%%%%%%%%%%%%%%%%%%%%%%%%%%%%%%%%%%%%%%%%
% Tema 7: Funcion implicita y Hessiana
\item
	Considere $\func{g}{\R^2}{\R}$ una función diferenciable tal que $g(x,f(x))=0$ y una función diferenciable $\func{f}{\R}{\R}$. Exprese la derivada $f$ en términos de las derivadas parciales de $g$.
        
% \begin{respuesta}
%     Se sabe que
%     \begin{align*}
%         \frac{\partial g}{\partial x}(x,f(x))=0,\\
%     \end{align*}
%     aplicando la regla de la cadena se tiene que
%     \begin{align*}
%          \frac{\partial g}{\partial x}(x,f(x))=D_1g(x,f(x))\dfrac{\partial x}{\partial x}(x)+D_2g(x,f(x))\dfrac{\partial f}{\partial x}(x)=0,
%     \end{align*}
%     de donde
%     \[
%         \dfrac{\partial f}{\partial x}(x)=-\dfrac{D_1g}{D_2g}(x,f(x)).
%     \]
% \end{respuesta}

%%%%%%%%%%%%%%%%%%%%%%%%%%%%%%%%%%%%%%%%%%%
%%%%%%%%%%%%%%%%%%%%%%%%%%%%%%%%%%%%%%%%%%%
%%%%%%%%%%%%%%%%%%%%%%%%%%%%%%%%%%%%%%%%%%%
% Tema 7: Funcion implicita y Hessiana
\item
    Suponga que la función $f:\R^2 \rightarrow \R$ cumple que
    \[
        e^{xy}-e^{yf(x,y)}+f(x,y)e^x-1=0
    \]
    para todo $(x,y)\in \R^2$. Determine $D_1f$ y $D_2f$ en el punto $(0,0)$.
    
\begin{respuesta} 
    Definamos la función
    \[
        \funcion{g}{\R^3}{\R}{(x,y,z)}{e^{xy}-e^{yz}+ze^x-1,}
    \]
    tenemos que cumple
    \[
        g(x,y,f(x,y))=0
    \]
	para todo $(x,y)\in \R^2$. Notemos que, en el punto $(0,0)$, tenemos que
	\[
	    e^{0}-e^{0}+f(0,0)e^0-1=0,
    \]	
    de donde, $f(0,0)=1$. Por tanto,
	\[
        D_1f(x,y)=-\frac{D_1g(x,y,f(x,y))}{D_3g(x,y,f(x,y))}	
    \]
	para todo $(x,y)\in \R^2$ tal que $D_3g(x,y,f(x,y))\neq 0$. Con esto,  se calcula, para $(x,y,z)\in \R^3$,
	\[
        D_1g(x,y,z)=ye^{xy}+ze^x,
	\]
	y
	\[
	    D_3g(x,y,z)=-ye^{yz}+e^x.
	\]      
	Entonces 
	\begin{align*}
	    D_1f(x,y)
	    & =-\frac{D_1g(x,y,f(x,y))}{D_3g(x,y,f(x,y))}\\
		& =-\frac{ye^{xy}+f(x,y)e^x}{e^x-ye^{yf(x,y)}}.
	\end{align*}
	Evaluando en el punto $(0,0)$
	\[
        D_1f(0,0)=-1.
    \]
	Para la derivada $D_2f(x,y)$
	\[
	    D_2f(x,y)=-\frac{D_2g(x,y,f(x,y))}{D_3g(x,y,f(x,y))},
    \]
	derivando $g$ con respecto a $y$
	\[
        D_2g(x,y,z)=xe^{xy}-ze^{yz},
    \]
	por lo tanto 
	\[
        D_2f(x,y)=-\frac{xe^{xy}-f(x,y)e^{yf(x,y)}}{e^x-ye^{yf(x,y)}},
    \]
	evaluando en el punto $(0,0)$
	\[
        D_2f(0,0)=-1.
    \]
\end{respuesta}

%%%%%%%%%%%%%%%%%%%%%%%%%%%%%%%%%%%%%%%%%%%
%%%%%%%%%%%%%%%%%%%%%%%%%%%%%%%%%%%%%%%%%%%
%%%%%%%%%%%%%%%%%%%%%%%%%%%%%%%%%%%%%%%%%%%
% Tema 7: Funcion implicita y Hessiana
\item 
    Suponga que la función $\func{f}{\R^2}{\R}$, verifica que
    \[
        x^2 f^2(x, y) + x y^2 - f^3(x, y) + 4 y f(x, y) = 5,
    \]
    para todo $(x, y)\in \R^2$. Determinar $D_1 f(x, y)$ y $D_2 f(x, y)$ para $(x, y)\in \R^2$.

\begin{respuesta}
    Se define la función
    \[
    \funcion{g}{\R^3}{\R}
    {(x, y, z)}{g(x, y, z) = x^2 z^2 + x y^2 - z^3 + 4yz - 5,}
	\]
    entonces, se verifica que
    \[
    g(x, y, f(x, y)) = 0,
    \]
    para todo $(x, y)\in D$. Luego, se tiene que
    \[
    D_1 f(x, y) = -\dfrac{D_1 g(x, y, f(x, y))}{D_3 g(x, y, f(x, y))},
    \]
    y
    \[
    D_2 f(x, y) = -\dfrac{D_2 g(x, y, f(x, y))}{D_3 g(x, y, f(x, y))},
    \]
    para todo $(x, y) \in D$ tal que $D_3 g(x, y, f(x, y)) \neq 0$.
    Calculemos las derivadas buscadas:
    \begin{align*}
    D_1 g(x, y, f(x, y)) &= 2xz^2 + y^2,\\
    D_2 g(x, y, f(x, y)) &= 2xy + 4z, \mbox{ y}\\
    D_3 g(x, y, f(x, y)) &= 2x^2 z - 3z^2 + 4y,
    \end{align*}
    para todo $(x, y, z)\in\R^3$, por lo cual:
    \begin{align*}
        D_1 f(x, y) &= -\dfrac{D_1 g(x, y, f(x, y))}{D_3 g(x, y, f(x, y))}\\
        &= \dfrac{2xf^2(x, y) + y^2}{3f^2(x, y) - 2x^2 f(x, y) - 4y},
    \end{align*}
    \begin{align*}
        D_2 f(x, y) &= -\dfrac{D_2 g(x, y, f(x, y))}{D_3 g(x, y, f(x, y))}\\
        &= \dfrac{2xy + 4f(x, y)}{3f^2(x, y) - 2x^2 f(x, y) - 4y},
    \end{align*}
    para todo $(x, y)\in D$ tal que $3f^2(x, y) - 2x^2 f(x, y) - 4y \neq 0$.
\end{respuesta}

%%%%%%%%%%%%%%%%%%%%%%%%%%%%%%%%%%%%%%%%%%%
%%%%%%%%%%%%%%%%%%%%%%%%%%%%%%%%%%%%%%%%%%%
%%%%%%%%%%%%%%%%%%%%%%%%%%%%%%%%%%%%%%%%%%%
% Tema 7: Funcion implicita y Hessiana
\item
    Consideremos las funciones
    \[
        \funcion{f}{\R^2}{\R}{(x,y)}{F(x,y),}
    \]
    y 
    \[
        \funcion{h}{\R^2}{\R}{(x,y)}{h(x,y).}
    \]
    Si se cumple que
    \[
        f(x+y+h(x,y),x^2+y^2+(h(x,y))^2)=0
    \]
    para todo $(x,y)\in\R^2$, determinar las derivadas parciales de $h$ para $(x,y) \in\R^2$ en función de las derivadas parciales de $f$. 

\begin{respuesta}
    Sean
    \[
    \funcion{g_1}{D \subseteq \R^2}{\R}{(x,y)}{x+y+h(x,y)}  \texty  \funcion{g_2}{D \subseteq \R^2}{\R}{(x,y)}{x^2+y^2+(h(x,y))^2}
    \]
    entonces
    \[
    \frac{\partial}{\partial x}f(g_1(x,y),g_2(x,y))=D_1f(g_1(x,y),g_2(x,y))\frac{\partial g_1}{\partial x}(x,y)+D_2f(g_1(x,y),g_2(x,y))\frac{\partial g_2}{\partial x}(x,y)
    \]
    para todo $(x,y) \in D$. Para lo cual se calcula  
    \begin{align*}
    \frac{\partial}{\partial x}f(g_1(x,y),g_2(x,y))&=0\\
    \frac{\partial g_1}{\partial x}(x,y)&=1+\frac{\partial h}{\partial x}(x,y),\\
    \frac{\partial g_2}{\partial x}(x,y)&=2x+2h(x,y) \frac{\partial h}{\partial x}(x,y)
    \end{align*}
    para todo $(x,y) \in D$. Por lo tanto
    \[
    D_1f(g_1(x,y),g_2(x,y))\left(1+\frac{\partial h}{\partial x}(x,y)\right)+D_2f(g_1(x,y),g_2(x,y))\left(2x+2h(x,y) \frac{\partial h}{\partial x}(x,y)\right)=0
    \]
    y
    \[
    \frac{\partial h}{\partial x}(x,y)=-\frac{D_1f(g_1(x,y),g_2(x,y))+2xD_2f(g_1(x,y),g_2(x,y))}{D_1f(g_1(x,y),g_2(x,y))+2h(x,y)D_2f(g_1(x,y),g_2(x,y))}
    \]
    para todo $(x,y) \in D$. 
    
    De manera similar se calcula
    \[
    \frac{\partial h}{\partial y}(x,y)=-\frac{D_1f(g_1(x,y),g_2(x,y))+2yD_2f(g_1(x,y),g_2(x,y))}{D_1f(g_1(x,y),g_2(x,y))+2h(x,y)D_2f(g_1(x,y),g_2(x,y))}
    \]
    para todo $(x,y) \in D$.
\end{respuesta} 

%%%%%%%%%%%%%%%%%%%%%%%%%%%%%%%%%%%%%%%%%%%
%%%%%%%%%%%%%%%%%%%%%%%%%%%%%%%%%%%%%%%%%%%
%%%%%%%%%%%%%%%%%%%%%%%%%%%%%%%%%%%%%%%%%%%
% Tema 8: Campos vectoriales
\item
    Verificar que en el campo vectorial
    \[
    \funcion{F}{\R^3}{\R^3}{(x,y,z)}{(xy^2z^2,z^2\sen y,x^2e^y)}
    \]
    la $\dive(F)=y^2z^2+z^2 \cos y$.
    
% \begin{respuesta}
%   	La divergencia de $F$
%     \[
%     \dive(F)(x,y,z)=\triangledown \cdot F (x,y,z) = \dfrac{\partial(xy^2z^2)}{\partial x}+\dfrac{\partial(z^2\sen y)}{\partial y}+\dfrac{\partial(x^2e^y)}{\partial z}=y^2z^2+z^2\cos y+0.
%     \]
% \end{respuesta}

%%%%%%%%%%%%%%%%%%%%%%%%%%%%%%%%%%%%%%%%%%%
%%%%%%%%%%%%%%%%%%%%%%%%%%%%%%%%%%%%%%%%%%%
%%%%%%%%%%%%%%%%%%%%%%%%%%%%%%%%%%%%%%%%%%%
% Tema 8: Campos vectoriales
\item 
    Dado $\func{F}{\R^3}{\R^3}$ un campo vectorial y $\func{g}{\R^3}{\R}$ un campo escalar, ambos diferenciable, demostrar que
    \[
        \dive(gF) = \nabla g\cdot F + g\dive(F) 
    \]

\begin{respuesta}
    Tomemos 
    \[
        F(x,y,z) = (f_1(x,y,z),f_2(x,y,z),f_3(x,y,z)) 
    \]
    para $(x,y,z)\in\R^3$, se tiene que (omitiendo por un momento las evaluaciones)
    \begin{align*}
        \dive(gF) & = \dive(gf_1,gf_2,gf_3) \\
            & = \parcial{gf_1}{x} + \parcial{gf_1}{y} + \parcial{gf_1}{z} \\
            & = \parcial{g}{x} f_1 + g\parcial{f_1}{x} 
            + \parcial{g}{y} f_2 + g\parcial{f_2}{y} 
            + \parcial{g}{z} f_3 + g\parcial{f_3}{z} \\
            & = \l( \parcial{g}{x},\parcial{g}{y},\parcial{g}{z} \r) \cdot (f_1,f_2,f_3) + g\l(\parcial{f_1}{x} +\parcial{f_1}{y} +\parcial{f_1}{z} \r)\\
            & = \nabla g\cdot f + g\dive(f).
    \end{align*}
\end{respuesta}

%%%%%%%%%%%%%%%%%%%%%%%%%%%%%%%%%%%%%%%%%%%
%%%%%%%%%%%%%%%%%%%%%%%%%%%%%%%%%%%%%%%%%%%
%%%%%%%%%%%%%%%%%%%%%%%%%%%%%%%%%%%%%%%%%%%
% Tema 8: Campos vectoriales
\item
    Sean
    \[
        \funcion{F}{\R^3}{\R^3}{(x,y,z)}{(x,y,z)}
    \]
    y $a \in \R^3$ un vector constante. Demostrar que $\rot(a \times F)=2a.$

\begin{respuesta}
    Sea $a=(a_1,a_2,a_3)$ entonces
    \[
        (a \times F){(x,y,z)}=(a_2z-a_3y ,a_3 x-a_1z, a_1y-a_2x)
    \]
    de aquí se tiene que 
    \begin{align*}
        \rot(a \times F)(x,y,z)
            &=\left( \frac{\partial}{\partial y} (a_1y-a_2x)-\frac{\partial}{\partial z}(a_3 x-a_1z),\frac{\partial}{\partial z} (a_2z-a_3y)-\frac{\partial}{\partial x}(a_1y-a_2x),\frac{\partial}{\partial x} (a_3 x-a_1z)-\frac{\partial}{\partial y}(a_2z-a_3y) \right)\\
            &=(2a_1,2a_2,2a_3)\\
            &=2(a_1,a_2,a_3),
    \end{align*}
    con lo cual se demuestra que $\rot(a \times F)=2a$
\end{respuesta}

%%%%%%%%%%%%%%%%%%%%%%%%%%%%%%%%%%%%%%%%%%%
%%%%%%%%%%%%%%%%%%%%%%%%%%%%%%%%%%%%%%%%%%%
%%%%%%%%%%%%%%%%%%%%%%%%%%%%%%%%%%%%%%%%%%%
% Tema 8: Campos vectoriales
\item
    Dado un campo escalar $\phi:\R^3 \longrightarrow \R$, determinar la divergencia y el rotacional de $\nabla \phi$.
    
\begin{respuesta}
    Se supone que $F:\R^3 \longrightarrow \R^3$ es un gradiente, es decir $F=\nabla \phi $, donde $\phi:\R^3 \longrightarrow \R$ es un campo escalar.  
    
    Entonces, omitiendo las notaciones por un momento,
    \[
    F=\left( \frac{\partial \phi}{\partial x},\frac{\partial \phi}{\partial y},\frac{\partial \phi}{\partial z} \right)
    \]
    y la divergencia de $F$ es
    \[
    \dive(F)=\frac{\partial^2 \phi}{\partial x^2}+\frac{\partial^2 \phi}{\partial y^2}+\frac{\partial^2 \phi}{\partial z^2} .
    \]
    Así pues, la divergencia de un gradiente $F=\nabla \phi$ es el laplaciano de $\phi$ y simbólicamente se expresa como
    \[
    \dive(F)=\dive(\nabla \phi)=\nabla ^2 \phi.
    \]
    Por otro lado se tiene que el rotacional de $F$ viene dado por
    \[
    \rot(F)=\left( \frac{\partial^2 \phi}{\partial y \partial z}-\frac{\partial^2 \phi}{\partial z \partial y},\frac{\partial^2 \phi}{\partial z \partial x}-\frac{\partial^2 \phi}{\partial x \partial z},\frac{\partial^2 \phi}{\partial x \partial y}-\frac{\partial^2 \phi}{\partial y \partial z} \right), 
    \]
    cuando estas derivadas parciales mixtas son continuas entonces 
    \[
    \rot(F)=\rot(\nabla \phi)=0.
    \]
    Esto demuestra que $rot(F)=0$ es condición necesaria para que un campo vectorial $F$  derivable con continuidad sea un gradiente.
\end{respuesta}



%%%%%%%%%%%%%%%%%%%%%%%%%%%%%%%%%%%%%%%%%%%
%%%%%%%%%%%%%%%%%%%%%%%%%%%%%%%%%%%%%%%%%%%
%%%%%%%%%%%%%%%%%%%%%%%%%%%%%%%%%%%%%%%%%%%
% Tema 8: Campos vectoriales
\item
	Calcular el vector rotacional del campo vectorial
    \[
        \funcion{F}{\R^3}{\R^3}{(x,y,z)}{(x^2-y,xe^z,xy)}
    \]
    y evalué $\rot(F)(2,5,9)$.

\begin{respuesta}
    Se conoce que
    \[
    \rot F= \triangledown \times F,
    \]
    es decir,
    \[
        \rot F=
        \left[
        \begin{array}{ccc}
		i & j & k \\
		\dfrac{\partial}{\partial x} & \dfrac{\partial}{\partial y} & \dfrac{\partial}{\partial z}\\
        x^2-y & xe^z & xy
        \end{array} 
        \right]
        =(x-xe^z,-1-y,e^z)
    \]
    evaluando,
    \[
    \rot F(2,5,9)=(2(1-e^9), -1-5, e^9).
    \]
\end{respuesta}

%%%%%%%%%%%%%%%%%%%%%%%%%%%%%%%%%%%%%%%%%%% PE
%%%%%%%%%%%%%%%%%%%%%%%%%%%%%%%%%%%%%%%%%%%
%%%%%%%%%%%%%%%%%%%%%%%%%%%%%%%%%%%%%%%%%%%
% Tema 8: Campos vectoriales
\item
    \begin{enumerate}
        \item Dado el campo vectorial
            \[
                \funcion{G}{\R^2}{\R^2}{(x,y)}{(x^2,y^2)}
            \]
            calcular la divergencia de $G$ en los puntos $(2,2)$ y $(-2,-2)$, interpretar el resultado.
    
        \item Dado el campo vectorial 
            \[
                \funcion{F}{\R^2}{\R^2}{(x,y)}{(-y,x)}
            \]
            calcular su rotacional en cualquier punto e interpretar el resultado.
    \end{enumerate}
\begin{respuesta}
    \begin{enumerate}
        \item Se tiene que
        \[
            \dive (G)(x,y)=D_1g_1(x,y)+D_2g_2(x,y)=2x+2y,
        \]
        de donde $\dive (G)(2,2)=8$ y $\dive (G)(-2,-2)=-8$. Con esto se concluye que el punto $(2,2)$ es una fuente y el punto $(-2,-2)$ un sumidero.
    
        \item Se tiene que
        \[
            \rot(F)(x,y)=D_1f_2(x,y)-D_2f_1(x,y)=2
        \]
        para cualquier $(x,y)\in\R^2$, con lo que se concluye que el campo rota en sentido antihorario en todo punto $(x,y)\in\R^2$.
    \end{enumerate}
    
%%%%%%%%%%%%%%%%%%%%%%%%%%%%%%%%%%%%%%%%%%%
%%%%%%%%%%%%%%%%%%%%%%%%%%%%%%%%%%%%%%%%%%%
%%%%%%%%%%%%%%%%%%%%%%%%%%%%%%%%%%%%%%%%%%%
% Tema 9: Campos conservativos
\item 
	Identifique si la función dada es un campo conservativo. Si lo es determine su potencial.
	\begin{enumerate}
	\item 
		\[
			\funcion{F}{\R^2}{\R^2}{x}{(3x_1^2-2x_2^2,4x_1x_2+3).}
		\]

\begin{respuesta}
	Dado que el campo es diferenciable, se tiene que 
	\[
	    \parcial{f_1}{x_2}(x_1,x_2) = -4x_2
	\]
	y
	\[
	    \parcial{f_2}{x_1}(x_1,x_2) = 4x_2
	\]
	para todo $(x_1,x_2)\in\R^2$. Dado que estos no son iguales, se tiene que $F$ no es conservativo.
\end{respuesta}

	\item 
		\[
		    \funcion{F}{\R^2}{\R^2}{x}{(x_2e^{x_1}+\sen(x_2),e^{x_1}+x_1\cos(x_2)).}
		\]
	
\begin{respuesta}
    Dado que el campo es diferenciable, se tiene que 
	\[
	    \parcial{f_1}{x_2}(x_1,x_2) = e^{x_1} + \cos(x_2)
	\]
	y
	\[
	    \parcial{f_2}{x_1}(x_1,x_2) = e^{x_1} + \cos(x_2)
	\]
	para todo $(x_1,x_2)\in\R^2$. Dado que estos son iguales y el dominio es $\R^2$, el cual es simplemente conexo, se tiene que $F$ es conservativo.
	
	Ahora, debemos hallar un campo escalar
	\[
		\func{f}{\R^2}{\R}
	\]
	tal que
	\[
		D_1 f(x)=x_2e^{x_1}+\sen(x_2)
		\yds
		D_2 f(x)=e^{x_1}+x_1\cos(x_2)
	\]
	para todo $x\in\R^2$. Al integrar $D_1 f(x)$ respecto a $x_1$, se tiene que
	\[
		f(x)=x_2e^{x_1}+x_1\sen(x_2) + g(x_2),
	\]
	para todo $x\in\R^2$ donde $g$ es una función real y derivable. Así,
	\[
		D_2 f(x)= e^{x_1}+x_1\cos(x_2) + g'(x_2)
	\]
	de donde
	\[
		e^{x_1}+x_1\cos(x_2) + g'(x_2)=e^{x_1}+x_1\cos(x_2)
	\]
	o también
	\[
		g'(x_2)=0
	\]
	para todo $x\in\R^2$. Finalmente, se deduce que existe $k\in\R$ tal que
	\[
		g(x_2)=k
	\]
	para todo $x_2\in\R^2$. Por lo tanto,
	\[
		\funcion{f}{\R^2}{\R}{x}{x_2e^{x_1}+x_1\sen(x_2) + k}
	\]
    es el potencial de $F$.
\end{respuesta}



%%%%%%%%%%%%%%%%%%%%%%%%%%%%%%%%%%%%%%%%%%%
%%%%%%%%%%%%%%%%%%%%%%%%%%%%%%%%%%%%%%%%%%%
%%%%%%%%%%%%%%%%%%%%%%%%%%%%%%%%%%%%%%%%%%%
% Tema 9: Campos conservativos
\item 
	Sea el campo vectorial
	\[
		\funcion{F}{\R^3}{\R^3}{x}{(P(x),Q(x),R(x))} 
	\]
	donde $P$, $Q$ y $R$ tienen derivadas de primer orden continuas. Demuestre que si $F$ es conservativo, entonces
	\[
		D_2 P =D_1 Q,\qquad D_3 P =D_1 Q \yds D_3 Q =D_2 R.
	\]


% \begin{respuesta}
% 	Como $F$ es conservativo, existe un campo escalar
% 	\[
% 		\func{f}{\R^3}{\R}
% 	\]
% 	tal que
% 	\[
% 		D_1 f =P, \qquad D_2 f =Q \yds D_3 f =R.
% 	\]
% 	Dado que $P$, $Q$ y $R$ tienen derivadas de primer orden continuas, entonces las derivadas parciales mixtas de $f$ coinciden:
% 	\[
% 		D_i D_j f = D_j D_i f
% 	\]
% 	para $i, j = 1, 2 ,3$. Así,
% 	\begin{align*}
% 		D_2 P&= D_2 D_1 f=D_1 D_2 f=D_1 Q \\
% 		D_3 P&= D_3 D_1 f=D_1 D_3 f=D_1 R \\
% 		D_3 Q&= D_3 D_2 f=D_2 D_3 f=D_2 R. \qedhere
% 	\end{align*}
% \end{respuesta}

%%%%%%%%%%%%%%%%%%%%%%%%%%%%%%%%%%%%%%%%%%%
%%%%%%%%%%%%%%%%%%%%%%%%%%%%%%%%%%%%%%%%%%%
%%%%%%%%%%%%%%%%%%%%%%%%%%%%%%%%%%%%%%%%%%%
% Tema 9: Campos conservativos
\item
    Dado el campo vectorial
    \[
        \funcion{F}{\R^2}{\R^2}{(x,y)}{F(x,y)=(2xy^3,3x^2 y^2)}
    \]
    Determine el campo escalar $f:\R^2\rightarrow \R$, tal que $F(x,y)=\nabla f(x,y)$ para $(x,y)\in\R^2$.
    
\begin{respuesta}
    Se calcula
    \[
        D_1F_2=6xy^2 \texty D_2F_1=6xy^2,
    \]
    con lo cual se verifica que $F$ es un gradiente ya que $\R^2$ es simplemente conexo. Se sabe entonces que existe una función $\func{f}{\R^2}{\R}$ tal que \[
        F(x,y)=\nabla f
    \]
    donde
    \begin{align*}
         F_1(x,y)&=\dfrac{\partial f}{\partial x}(x,y)=2xy^3,\\
         F_2(x,y)&=\dfrac{\partial f}{\partial y}(x,y)=3x^2y^2,
    \end{align*}
    integrando la primera ecuación con respecto $x$ se tiene
    \[
        f(x,y)=x^2y^3+g(y),
    \]
    derivando esta expresión con respecto a $y$
    \[
        \dfrac{\partial f}{\partial y}(x,y)=3x^2y^2+\dfrac{dg}{dy}(y)=3x^2y^2
    \]
    de donde 
    \[
        \dfrac{dg}{dy}(y)=0
    \]
    entonces $g(y)=c$, con $c\in\R$.
    Finalmente, la función potencial del campo vectorial $F$ es
    \[
        \funcion{f}{\R^2}{\R}{(x,y)}{x^2y^3+c}
    \]
    con $c\in\R$.
\end{respuesta}
        

%%%%%%%%%%%%%%%%%%%%%%%%%%%%%%%%%%%%%%%%%%%
%%%%%%%%%%%%%%%%%%%%%%%%%%%%%%%%%%%%%%%%%%%
%%%%%%%%%%%%%%%%%%%%%%%%%%%%%%%%%%%%%%%%%%%
% Tema 10: Optimizacion
\item 
	Determinar el mínimo de la función
	\[
	    \funcion{C}{\R^2}{\R}{x}
	    {3x_1^2+2x_1x_2+4x_2^2-12x_1-26x_2+50.}
	\]

\begin{respuesta}
	Las derivadas parciales del campo escalar $C$ son
	\[
		D_1 C(x)=6x_1+2x_2-12\yds 
		D_2 C(x)=8x_2+2x_1-26
	\]
	para todo $x\in\R^2$. Ahora, para encontrar los puntos críticos del campo escalar debemos resolver el sistema
	\[
		\begin{cases}
		6x_1+2x_2=12\\
		8x_2+2x_1=26,
		\end{cases}
	\]
	que es equivalente a $\nabla C(x)=0$. La solución del sistema anterior es $x_1= 1$ y $x_2=3$. Para determinar si el punto crítico $(1,3)$ es un mínimo o un máximo debemos hallar la Hessiana. Así,
	\[
		H_f(x)=
			\begin{pmatrix}
				6&2\\
				2&8
			\end{pmatrix}
	\]
	para cada $x\in\R^2$. La matriz Hessiana es constante y 
	\[
		\det(H_f(x))=48-4=44>0.
	\]
	Como los dos menores de la matriz Hessiana son siempre positivos, por el criterio de Sylvester, se deduce que $f$ alcanza su mínimo en $(1,3)$ y este mínimo es
	\[
		f(1,3)=3+2(3)+4(9)-12-26(3)+50=5.
	\]
\end{respuesta}




%%%%%%%%%%%%%%%%%%%%%%%%%%%%%%%%%%%%%%%%%%%
%%%%%%%%%%%%%%%%%%%%%%%%%%%%%%%%%%%%%%%%%%%
%%%%%%%%%%%%%%%%%%%%%%%%%%%%%%%%%%%%%%%%%%%
% Tema 10: Optimizacion
\item 
    En funciones de una sola variable es imposible, en el caso de funciones continuas, tener dos máximos relativos y ningún mínimo relativo. Pero si las funciones son de dos variables, sí existen esas funciones. Demuestre que la función
    \[
    \funcion{f}{\R^2}{\R}
    {(x,y)}{f(x,y)=-(x^2-1)^2-(x^2y-x-1)^2}
	\]
    tiene sólo dos puntos críticos y tiene máximos relativos en ambos puntos.
    
\begin{respuesta}
    Notemos que $f$ es derivable para todo $(x, y)\in\R^2$, con gradiente 
    \[
    \nabla (x,y) = (-4x(x^2-1) - 2(2xy-1)(x^2y-x-1), -2x^2(x^2y-x-1)),
    \]
    para todo $(x, y)\in\R^2$, además, sus dos primeras derivadas parciales $D_1 f$ y $D_2 f$ son continuas para todo $(x, y)\in\R^2$, por lo tanto $f$ es diferenciable con continuidad en todo $\R^2$.
    
    Encontremos los puntos críticos de $f$, para obtenerlos igualemos el gradiente a cero, es decir
    \[
    \nabla f(x,y) = (0, 0),
    \]
    
    de aquí se obtiene el siguiente sistema 
    \begin{align}
    -4x(x^2-1) - 2(2xy-1)(x^2y-x-1) = 0,\\
    -2x^2(x^2y-x-1) = 0.
    \end{align}
    A partir de este sistema se tiene que 
    \[    
    \begin{array}{ccccc}
	x^2 = 1 & & \texty & & x^2y-x-1 = 0,
	\end{array}
    \]
    de donde se obtiene los dos puntos críticos de $f$
     \[    
    \begin{array}{ccccc}
	a_1 = (1,2), & & y & & a_2 = (-1, 0).
	\end{array}
    \]
    La matriz hessiana de $f$ está dada por 
        \[
        H_f(x,y)= (J_{\nabla f} (x,y))^T = 
        \begin{bmatrix}
         -12x^2-12x^2y^2 + 12xy +4y +2 & -4x(x^2y-x-1)-2x^2(2xy-1) \\
         -4x(x^2y-x-1)-2x^2(2xy-1) & -2x^4
        \end{bmatrix},
    	\]
    para todo $(x, y)\in\R^2$.
    \begin{enumerate}
    \item Ahora, si analizamos la matriz hessiana de $f$ en $a_1$, obtenemos 
        \[
        H_f(a_1)= (J_{\nabla f} (a_1))^T = 
        \begin{bmatrix}
         -26 & -6 \\
         -6 & -2
        \end{bmatrix},
    	\]
    dado que su determinante es $|H_f(a_1)| = 16 > 0$, junto con $D_{1,1}f(a_1) = -26 < 0$, se concluye que $f$ en $a_1$ alcanza un máximo con valor igual a 0.
    \item Y si analizamos la matriz hessiana de $f$ en $a_2$, obtenemos 
        \[
        H_f(a_2)= (J_{\nabla f} (a_2))^T = 
        \begin{bmatrix}
         -10 & 2 \\
         2 & -2
        \end{bmatrix},
    	\]
    dado que su determinante es $|H_f(a_2)| = 16 > 0$, junto con $D_{1,1}f(a_2) = -10 < 0$, por tanto, se concluye que $f$ en $a_2$ alcanza un máximo con valor igual a 0.
    \end{enumerate}    
    Entonces, se ha mostrado que $f$ tiene dos puntos críticos: $a_1=(1,2)$ y $a_2=(-1,0)$, además, en estos dos puntos $f$ alcanza un máximo relativo igual a 0.
\end{respuesta}

%%%%%%%%%%%%%%%%%%%%%%%%%%%%%%%%%%%%%%%%%%%
%%%%%%%%%%%%%%%%%%%%%%%%%%%%%%%%%%%%%%%%%%%
%%%%%%%%%%%%%%%%%%%%%%%%%%%%%%%%%%%%%%%%%%%
% Tema 10: Optimizacion
\item
    Sea 
    \[
    \funcion{f}{\R^2}{\R}
    {(x,y)}{f(x,y)=4+x^3+y^3-3xy,}
	\]
    clasificar sus puntos críticos.
    
\begin{respuesta}
    Notemos que $f$ es derivable para todo $(x, y)\in\R^2$, con gradiente 
    \[
    \nabla (x,y) = (3x^2-3y,3y^2-3x),
    \]
    para todo $(x, y)\in\R^2$, además, sus dos primeras derivadas parciales $D_1 f$ y $D_2 f$ son continuas para todo $(x, y)\in\R^2$, por lo tanto $f$ es diferenciable con continuidad en todo $\R^2$.
    
    Encontremos los puntos críticos de $f$, para obtenerlos igualemos el gradiente a cero, es decir
    \[
    \nabla f(x,y) = (0, 0).
    \]
    
    A partir de lo cual se obtiene el siguiente sistema 
    \begin{align*}
    3x^2-3y = 0,\\
    3y^2-3x = 0.
    \end{align*}
    Resolviendo el sistema anterior se tienen los puntos críticos de $f$
     \[    
    \begin{array}{ccccc}
	a_1 = (0,0), & & y & & a_2 = (1,1).
	\end{array}
    \]
    La matriz hessiana de $f$ está dada por 
        \[
        H_f(x,y)= (J_{\nabla f} (x,y))^T = 
        \begin{bmatrix}
         6x & -3 \\
         -3 & 6y
        \end{bmatrix}.
    	\]
    para todo $(x, y)\in\R^2$.
    \begin{enumerate}
    \item Ahora, si analizamos la matriz hessiana de $f$ en $a_1$, obtenemos 
        \[
        H_f(a_1)= (J_{\nabla f} (a_1))^T = 
        \begin{bmatrix}
         0 & -3 \\
         -3 & 0
        \end{bmatrix},
    	\]
    dado que su determinante es $|H_f(a_1)| = -9 < 0$, por lo cual se concluye que $a_1$ es un punto de ensilladura de $f$.
    \item Y si analizamos la matriz hessiana de $f$ en $a_2$, obtenemos 
        \[
        H_f(a_2)= (J_{\nabla f} (a_2)^T = 
        \begin{bmatrix}
         6 & -3 \\
         -3 & 6
        \end{bmatrix},
    	\]
    dado que su determinante es $|H_f(a_2)| = 27 > 0$, junto con $D_{1,1}f(a_2) = 6 > 0$, por tanto, se concluye que $f$ en $a_2$ alcanza un mínimo relativo con valor igual a 3.
    \end{enumerate}
    Finalmente, se concluye que $a_1=(0,0)$ es un punto de ensilladura de $f$ y que en $a_2=(1,1)$ $f$ alcanza un mínimo relativo con valor igual a 3.
\end{respuesta}

%%%%%%%%%%%%%%%%%%%%%%%%%%%%%%%%%%%%%%%%%%%
%%%%%%%%%%%%%%%%%%%%%%%%%%%%%%%%%%%%%%%%%%%
%%%%%%%%%%%%%%%%%%%%%%%%%%%%%%%%%%%%%%%%%%%
% Tema 10: Optimizacion
\item
    Sea 
    \[
        \funcion{f}{I}{\R}{(x,y)}{x^2+y^2-x-y+1} 
    \]
    donde $I$ es el disco definido por $\lbrace (x,y) \in \R^2:x^2+y^2 \leqslant 1 \rbrace$. Hallar los máximos y mínimos de la función $f$.
	
\begin{respuesta}
    Para hallar los puntos críticos se analiza el gradiente de la función $f$
    \[
        \nabla f(x,y)=(2x-1,2y-1)
    \]
    el cual es igual a (0,0) para $(x,y)=(\frac{1}{2},\frac{1}{2})$.
    La frontera se puede expresar como
    \[
        \funcion{r}{[0,2\pi]}{\R^2}{t}{(\sen t, \cos t)}
    \]
    entonces $f(r(t))=2-\cos t-\sen t$, $t \in [0,2\pi]$, los puntos críticos de la frontera serán los puntos para los que $f'(r(t))=0$.
    \[
        f'(r(t))=\sen t - \cos t,
    \]
    así los puntos críticos son $t=\dfrac{\pi}{4}$, $t=\dfrac{5\pi}{4}$ y los extremos $t=0$ y $t=2\pi$.
        
    Los valores de la función $f$ en los puntos críticos son
       
    \begin{align*}     
        f\left (\frac{1}{2},\frac{1}{2} \right )&=\frac{1}{2}\\
        f\left(r\left(\frac{\pi}{4}\right)\right)&=f\left (\frac{\sqrt[]{2}}{2},\frac{\sqrt[]{2}}{2} \right )=2-\sqrt[]{2}\\
        f\left(r\left(\frac{5\pi}{4}\right)\right)&=f\left (-\frac{\sqrt[]{2}}{2},-\frac{\sqrt[]{2}}{2} \right )=2+\sqrt[]{2}\\
        f(r(0))&=f(r(2\pi))=f(0,1)=1.
    \end{align*}
    Comparando todos los valores se puede ver que el mínimo absoluto esta en $\left (\frac{1}{2},\frac{1}{2} \right )$ y el máximo absoluto en $\left (-\frac{\sqrt[]{2}}{2},-\frac{\sqrt[]{2}}{2} \right )$.
\end{respuesta}

%%%%%%%%%%%%%%%%%%%%%%%%%%%%%%%%%%%%%%%%%%%
%%%%%%%%%%%%%%%%%%%%%%%%%%%%%%%%%%%%%%%%%%%
%%%%%%%%%%%%%%%%%%%%%%%%%%%%%%%%%%%%%%%%%%%

% Tema 10: Optimizacion
\item 
    Este ejercicio indica que si la hessiana de un campo escalar se anula en un punto crítico, en este podría haber un máximo, un mínimo o un punto de silla. Considere los 3 campos escalares
    \[
        \funcion{f}{\R^2}{\R}{(x,y)}{x^4+x^2+y^4,}
        \qquad
        \funcion{g}{\R^2}{\R}{(x,y)}{-x^4-x^2-y^4}
        \yds
        \funcion{h}{\R^2}{\R}{(x,y)}{x^4-x^2+y^4.}
    \]
    Muestre que $\nabla f(0,0)=\nabla g(0,0)=\nabla h(0,0)=0$ y $\l|H_f(0,0)\r||=\l|H_g(0,0)\r|=\l|H_h(0,0)\r|=0$, sin embargo $(0,0)$ es un mínimo absoluto de $f$, es un máximo absoluto de $g$ y es un punto silla de $h$. 

\begin{respuesta}
    Por calculo directo, se tiene que para cada $(x,y)\in\R^2$,
    \[
        \nabla f(x,y)=(4x^3+2x,4y^3),
        \qquad
        \nabla g(x,y)=(-4x^3-2x,-4y^3)
        \yds
        \nabla h(x,y)=(4x^3-2x,4y^3)
    \]
    y, también
    \[
        H_f(x,y)=
        \begin{pmatrix}
        12x^2+2&0\\
        0&12y^2\\
    \end{pmatrix},
    \qquad
    H_g(x,y)=
    \begin{pmatrix}
        -12x^2-2&0\\
        0&-12y^2\\
    \end{pmatrix}
    \yds
    H_h(x,y)=
    \begin{pmatrix}
        12x^2-2&0\\
        0&12y^2
    \end{pmatrix}.
    \]
    En particular,
    \[
        \nabla f(0,0)=\nabla g(0,0)=\nabla h(0,0)=0
    \]
    y
    \[
        H_f(0,0)=
        \begin{pmatrix}
        2&0\\
        0&0\\
    \end{pmatrix},
    \qquad
    H_g(0,0)=
    \begin{pmatrix}
        -2&0\\
        0&0\\
    \end{pmatrix}
    \yds
    H_h(0,0)=
    \begin{pmatrix}
        -2&0\\
        0&0
    \end{pmatrix},
    \]
    de donde
    \[
        \l|H_f(0,0)\r||=\l|H_g(0,0)\r|=\l|H_h(0,0)\r|=0.
    \]
    Es claro que $f(x,y)\geq 0$ para todo $(x,y)\in\R^2$ y dado que $f(0,0)=0$, se deduce que $(0,0)$ es el mínimo absoluto de $f$. Además, notemos que $g(x,y)\leq 0$ para todo $(x,y)\in\R^2$ y dado que $g(0,0)=0$, se deduce que $(0,0)$ es el máximo absoluto de $g$. Finalmente, dado que para todo $y\in\R*$,
    \[
        h(0,y)=y^4>0
    \]
    y para todo $x\in\lop-1,1\rop\setminus\{0\}$,
    \[
        h(x,0)=x^4-x^2=x^2(x+1)(x-1)<0
    \]
    y $h(0,0)=0$, se deduce que $(0,0)$ es un punto silla pues hay direcciones en que el campo escalar $h$ crece y direcciones en el que decrece. Esto nos indica que si el determinante de la hessiana de un campo escalar es 0 en un punto crítico, no se puede determinar su naturaleza.
\end{respuesta}

%%%%%%%%%%%%%%%%%%%%%%%%%%%%%%%%%%%%%%%%%%%
%%%%%%%%%%%%%%%%%%%%%%%%%%%%%%%%%%%%%%%%%%%
%%%%%%%%%%%%%%%%%%%%%%%%%%%%%%%%%%%%%%%%%%%

% Tema 10: Optimizacion
\item 
    Se conoce que el campo escalar
    \[
        \funcion{f}{\R^4}{\R}{(w,x,y,z)}{\dfrac{4y^3+2w^2-y^2-w^3}{1+x^2+2z^2}}
    \]
    tiene dos puntos críticos
    \[
        a=\l(\dfrac{4}{3},0,0,0\r)
        \yds
        b=\l(0,0,\dfrac{1}{6},0\r).
    \]
    Además, se conoce que
    \[
        H_f(w,0,y,0)=
        \begin{pmatrix}
        4-6w&0&0&0\\
        0&2w^3-4w^2-8y^3+2y^2&0&0\\
        0&0&24y-2&0\\
        0&0&0&4w^3-8w^2-16y^3+4y^2        
        \end{pmatrix}
    \]
    para cada $w,y\in\R$. Determine si $a$ y $b$ son extremos de $f$ y cual es su naturaleza.
\begin{respuesta}
    Por las hipótesis, se tiene que
    \[
        H_f(a)=
        \begin{pmatrix}
        -4&0&0&0\\
        0&-\dfrac{64}{27}&0&0\\
        0&0&-2&0\\
        0&0&0&-\dfrac{128}{27}        
        \end{pmatrix}
        \yds
        H_f(b)=
        \begin{pmatrix}
        4&0&0&0\\
        0&\dfrac{1}{54}&0&0\\
        0&0&2&0\\
        0&0&0&\dfrac{1}{27}        
        \end{pmatrix}.
    \]
    Por Álgebra Lineal sabemos directamente que
    \[
        -4,
        \qquad 
        -\dfrac{64}{27},
        \qquad
        -2
        \yds
        -\dfrac{128}{27}        
    \]
    son los 4 valores propios de $H_f(a)$; luego, por el criterio de Sylvester, $H_f(a)$ es definida negativa, de donde $f$ alcanza un máximo relativo en $a$. De igual forma, 
    \[
        4,
        \qquad
        \dfrac{1}{54},
        \qquad
        2
        \yds
        \dfrac{1}{27}        
    \]
    son los 4 valores propios de $H_f(b)$; luego, por el criterio de Sylvester, $H_f(b)$ es definida positiva, de donde $f$ alcanza un mínimo relativo en $b$. Por lo tanto,
    \[
        f(a)=\dfrac{32}{27}
        \yds
        f(b)=-\dfrac{1}{108}
    \]
    son un máximo y un mínimo local, respectivamente.
\end{respuesta}

%%%%%%%%%%%%%%%%%%%%%%%%%%%%%%%%%%%%%%%%%%% FC
%%%%%%%%%%%%%%%%%%%%%%%%%%%%%%%%%%%%%%%%%%%
%%%%%%%%%%%%%%%%%%%%%%%%%%%%%%%%%%%%%%%%%%%
% Tema 10: Optimizacion
\item 
    Dada la función
    \[
    \funcion{u}{\R^3}{R}{(x,y,z)}{-x^2-2y^2-z^3+xy+75z+16x+27y+120}
    \]
    Determine sus puntos críticos y su naturaleza
\begin{respuesta}
    Se encuentran los puntos críticos de $u$, para ello debe cumplirse que $\nabla u (x,y,z)=0$. 

    Se tiene que:
        \[
        \begin{aligned}
        D_1u(x,y,z)&=-2x+y+16=0\\
        D_2u(x,y,z)&=-4y+x+27=0\\
        D_3u(x,y,z)&=-3z^2+75=0.
        \end{aligned}
        \]
    Resolviendo el sistema de ecuaciones se obtiene dos puntos críticos $(13,10,5)$ y $(13,10,-5)$, para determinar si es un máximo o mínimo aplicamos el criterio de Sylvester.
    Entonces la matriz Hessiana de $u$ estará dada por
    \[
        H_u(x,y,z)= \begin{pmatrix}
        -2 & 1 & 0\\
        1 & -4 & 0\\
        0 & 0 & -6z
        \end{pmatrix}
    \]
    En $(x,y,z)=(13,10,5)$,
    \[
        H_u(13,10,5)= \begin{pmatrix}
        -2 & 1 & 0\\
        1 & -4 & 0\\
        0 & 0 & -30
        \end{pmatrix}
    \]
    de donde se obtiene que sus determinantes menores son
    \[
        \begin{aligned}
        H_1(13,10,5)&=-2\\
        H_2(13,10,5)&=7\\
        H_3(13,10,5)&=-210,
        \end{aligned}
    \]
    la matriz Hessiana se define negativa, por lo que en el punto $(x,y,z)=(13,10,5)$ hay un máximo.

    Mientras que en $(x,y,z)=(13,10,-5)$,
    \[
    H_u(13,10,-5)= \begin{pmatrix}
        -2 & 1 & 0\\
        1 & -4 & 0\\
        0 & 0 & 30
        \end{pmatrix}
    \]
     de donde se obtiene que sus determinantes menores son
    \[
        \begin{aligned}
        H_1(13,10,5)&=-2\\
        H_2(13,10,5)&=7\\
        H_3(13,10,5)&=210,
        \end{aligned}
    \]
    por lo que la naturaleza de este punto es un punto silla; es decir, no es ni máximo ni mínimo.
\end{respuesta}

\item
    Dada la función
    \[
    \funcion{f}{I}{R}{(x,y)}{x^2+y^2+\frac{1}{x^2y^2}}
    \]
    con $I=\lbrace (x,y)\in\R^2:x\neq 0 \land y\neq 0 \rbrace$ determinar la naturaleza de sus puntos críticos.

\begin{respuesta}
    Primero se buscan los puntos críticos, estos son los puntos en donde el gradiente es cero
    \[
    \nabla f(x,y)=\left(2x-\frac{2}{x^3y^2},2y-\frac{2}{x^2y^3}\right).
    \]
    Las soluciones del sistema
    
    \[
    2x-\frac{2}{x^3y^2}=0
    \]
    \[
    2y-\frac{2}{x^2y^3}=0
	\]
    son $(1,1)$, $(1,-1)$, $(-1,1)$ y $(-1,-1)$.
    Para determinar si estos puntos son máximos, mínimos o puntos de silla, se analiza la matriz hesiana
    \[
    H_f(x,y)=\left( {\begin{array}{cc}
   2+\dfrac{6}{x^4y^2} & \dfrac{4}{x^3y^3} \\
   \dfrac{4}{x^3y^3} & 2+\dfrac{6}{x^2y^4} \\
  \end{array} } \right)
    \]
    Se verifica que $\det H_f(a,b)>0$ y $\dfrac{\partial^2f}{\partial x^2 }(a,b)>0$ donde $(a,b)$ es cualquiera de los puntos críticos. Entonces en cada punto crítico hay un mínimo y $f(a,b)=3$ para todos los puntos críticos.
\end{respuesta}

%%%%%%%%%%%%%%%%%%%%%%%%%%%%%%%%%%%%%%%%%%%
%%%%%%%%%%%%%%%%%%%%%%%%%%%%%%%%%%%%%%%%%%%
%%%%%%%%%%%%%%%%%%%%%%%%%%%%%%%%%%%%%%%%%%%
% Tema 10: Optimizacion
\item
    Dada la función
    \[
     \funcion{u}{\R^+ \times \R^+}{\R}{(x,y)}{655x+468y+\frac{xy}{30}-\frac{x^2}{40}-\frac{y^2}{80}}
    \]
    determinar sus puntos críticos y naturaleza.

% \begin{respuesta}
%     Para maximizar la función, calculamos primero su gradiente y obtenemos
%     \[
%     \nabla u(x,y)=\left(655+\frac{y}{30}-\frac{x}{20},
%     468+\frac{x}{30}-\frac{y}{40}\right)
%     \]
%     Se calcula los puntos para los cuales $\nabla u(x,y)=(0,0)$ y se determina que su punto crítico es $(230220,325680)$. 
%     Para determinar su naturaleza, calculamos la matriz hessiana \[
%         H_u(x,y)=\begin{pmatrix}
%         -\frac{1}{20} &\frac{1}{30}\\
%         \frac{1}{30} & -\frac{1}{40}
%         \end{pmatrix}
%     \]
%     Se verifica que para todo $(x,y)\in\R^+ \times \R^+$
%     \[
%     d_1=-\frac{1}{20}<0 \texty  
%     d_2=\frac{1}{800}-\frac{1}{900}>0,
%     \]
%     se concluye que la función $u$ tiene un valor máximo relativo en el punto $a=(230220,325680)$.
% \end{respuesta}

%%%%%%%%%%%%%%%%%%%%%%%%%%%%%%%%%%%%%%%%%%%
%%%%%%%%%%%%%%%%%%%%%%%%%%%%%%%%%%%%%%%%%%%
%%%%%%%%%%%%%%%%%%%%%%%%%%%%%%%%%%%%%%%%%%%
% Tema 10: Optimizacion
\item 
    Dado los campos escalares
    \[
        \funcion{d}{\Omega}{\R}{(x,y,z)}{x^2+y^2+z^2} \texty \funcion{g}{\Omega}{\R}{(x,y,z)}{z-\dfrac{1}{xy}}
    \]
    con $\Omega=\lbrace (x,y,z)\in\R^3:x\neq 0 \land y\neq 0 \rbrace$. Resolver el problema
    \[
        \begin{cases}
            \max d(x,y,z)\\
            \text{sujeto  a:}\\
            g(x,y,z)=0.
        \end {cases}
    \]
    
\begin{respuesta}
    Para resolver el problema, se define la función $\func{L}{\R^4}{\R}$
    \[
        L(\lambda,x,y,z)=d(x,y,z)-\lambda(g(x,y,z)-0)=x^2+y^2+z^2-\lambda \l(z-\dfrac{1}{xy}\r)
    \]
    cuyo gradiente es
    \[
        \nabla L(\lambda,x,y,z)=\l(-z+\dfrac{1}{xy},2x-\dfrac{\lambda}{x^2y},2y-\dfrac{\lambda}{xy^2},2z-\lambda \r)
    \]
    para $(\lambda,x,y,z)\in\R^4$. Los puntos en los que $\nabla L(\lambda,x,y,z)=0$, es decir, los puntos críticos de $L$ son:
    \begin{align*}
        a_1&=(2,1,1,1),\\
        a_2&=(2,-1,-1,1),\\ 
        a_3&=(-2,1,-1,-1) \texty \\ 
        a_4&=(-2,-1,1,-1).
    \end{align*}
    Para determinar la naturaleza de los puntos críticos, se calcula la matriz hesiana
    \[
            H_L(\lambda,x,y,z) = 
            \begin{pmatrix}
                0 & -\dfrac{1}{x^2y} & -\dfrac{1}{xy^2} & -1\\ \\
                -\dfrac{1}{x^2y} & 2+\dfrac{2\lambda}{x^3y} & \dfrac{\lambda}{x^2y^2} & 0\\ \\
                -\dfrac{1}{xy^2} & \dfrac{\lambda}{x^2y^2} & 2+\dfrac{2\lambda}{xy^3} & 0\\ \\
                -1 & 0 & 0 & 2
            \end{pmatrix}.
    \]
    Con esto se tiene
    \[
        H_L(a_1)=
        \begin{pmatrix}
            0 & -1 & -1 & -1\\
            -1 & 6 & 2 & 0\\
            -1 & 2 & 6 & 0\\
            -1 & 0 & 0 & 2
        \end{pmatrix},
    \]
    \[
        H_L(a_2)=
        \begin{pmatrix}
            0 & 1 & 1 & -1\\
            1 & 6 & 2 & 0\\
            1 & 2 & 6 & 0\\
            -1 & 0 & 0 & 2
        \end{pmatrix},
    \]
    \[
        H_L(a_3)=
        \begin{pmatrix}
            0 & 1 & -1 & -1\\
            1 & 6 & -2 & 0\\
            -1 & -2 & 6 & 0\\
            -1 & 0 & 0 & 2 
        \end{pmatrix}
    \]
    y
    \[
        H_L(a_4)=
        \begin{pmatrix}
            0 & -1 & 1 & -1\\
            -1 & 6 & -2 & 0\\
            1 & -2 & 6 & 0\\
            -1 & 0 & 0 & 2
        \end{pmatrix}.
    \]
    Por otro lado se tiene que $n=3$ y $m=1$ entonces se calcula para cada punto crítico $(-1)^mdet(H_{2m+1})$ y $(-1)^mdet(H_{m+n})$, entonces
    \[
        -det(H_{L_3}(a_1))=8 \texty -det(H_{L_4}(a_1))=48,
    \]
    \[
        -det(H_{L_3}(a_2))=8 \texty -det(H_{L_4}(a_2))=48,
    \]
    \[
        -det(H_{L_3}(a_3))=8 \texty -det(H_{L_4}(a_3))=48,
    \]
    \[
        -det(H_{L_3}(a_4))=8 \texty -det(H_{L_4}(a_4))=48.
    \]
    Se puede notar que, para todos los puntos críticos, la secuencia de determinantes consta solamente de número positivos, entonces $d$ tiene un mínimo, con un valor de 3, en todos los puntos críticos.
\end{respuesta}

%%%%%%%%%%%%%%%%%%%%%%%%%%%%%%%%%%%%%%%%%%%
%%%%%%%%%%%%%%%%%%%%%%%%%%%%%%%%%%%%%%%%%%%
%%%%%%%%%%%%%%%%%%%%%%%%%%%%%%%%%%%%%%%%%%%
% Tema 10: Optimizacion
\item 
    Dadas la funciones
    \[
         \funcion{f}{\R^3}{\R}{(x,y,z)}{xyz}
    \]
  	y,
	\[
        \funcion{g}{\R^3}{\R}{(x,y,z)}{x + y+ z.}
    \] 
    Resolver el problema 
    \[
        \begin{cases}
        \begin{array}{c}
		\max f(x,y,z)\\
		\text{sujeto  a:}\\
			g(x,y,z)=a\\		
		\end{array}
        \end{cases} 
    \]
    donde $a\in\N^{*}$, es constante.
    
\begin{respuesta}
        Para resolver el problema definimos la función que involucra a $f$ y a $g$ con la condición impuesta por el problema, siendo 
        \[
            \funcion{h}{\R^4}{\R}{(\lambda,x,y,z)}{f(x,y,z) - \lambda(g(x,y,z) - a).}
        \]
        Notemos que se tiene que
        \[
            h(\lambda,x,y,z) = xyz - \lambda(x+y+z-a)
        \]
        para todo $(\lambda,x,y,z)\in\R^4$, además, $h$ es una función diferenciable en $\R^4$, con vector gradiente 
        \[
            \nabla h(x,y,z,\lambda) = 
            \begin{pmatrix}
                -(x+y+z-a) \\
                yz - \lambda\\
                xz - \lambda\\
                xy - \lambda
            \end{pmatrix}
        \]
        y matriz hessina dada por
        \[
            H_h(x,y,z,\lambda) = 
            \begin{pmatrix}
                0 & -1 & -1 & -1\\
                -1 & 0 & z & y\\
                -1 & z & 0 & x\\
                -1 & y & x & 0
            \end{pmatrix}
        \]
        para cada $(\lambda,x,y,z)\in\R^4$.
        
        Determinemos ahora, los puntos crpiticos de $h$, para esto igualamos el vector gradiente a 0, es decir, hacemos $\nabla h(\lambda,x,y,z) = (0,0,0,0)$, obteniendo el siguiente sistema:
        \[
            \begin{array}{cc}
                yz = \lambda,\\
                xz = \lambda,\\
                xy = \lambda,\\
                x+y+z = a,
            \end{array}
        \]
        cuya solución es $x=y=z=\dfrac{a}{3}$ y $\lambda=\dfrac{a^2}{9}$. Analizando la matriz hessiana en $w=\l(\dfrac{a^2}{9},\dfrac{a}{3},\dfrac{a}{3},\dfrac{a}{3}\r)$ obtenemos
        \[
            H_h(w) = 
            \begin{pmatrix}
                0 & -1 & -1 & -1 \\
                -1 & 0 & \dfrac{a}{3} & \dfrac{a}{3} \\
                -1 & \dfrac{a}{3} & 0 & \dfrac{a}{3} \\
                -1 & \dfrac{a}{3} & \dfrac{a}{3} & 0 
            \end{pmatrix}
        \]
        Para determinar si en $\l(\dfrac{a}{3},\dfrac{a}{3},\dfrac{a}{3}\r)$, el único punto crítico restringido, $f$ alcanza un máximo local, utilicemos el criterio del hessiano, para esto, notemos que $m=1$ y $n=3$, luego la sucesión de determinantes a considerar es
        \[
            (-1)^{m}\text{det H}_{2m+1} =  -\text{det H}_{3} = 
            -\begin{vmatrix}
                0 & -1 & -1\\
                -1 & 0 & \dfrac{a}{3}\\
                -1 & \dfrac{a}{3} & 0 
            \end{vmatrix} = -\dfrac{2a}{3}
        \]
        y
        \[
            (-1)^{m}\text{det H}_{m+n} =  -\text{det H}_{4} = 
            -\begin{vmatrix}
                0 & -1 & -1 & -1 \\
                -1 & 0 & a/3 & a/3 \\
                -1 & a/3 & 0 & a/3 \\
                -1 & a/3 & a/3 & 0 
            \end{vmatrix} = \dfrac{a^2}{3},
        \]
        como podemos observar se inicia con un valor negativo y luego con un positivo, por lo cual, el valor que resuelve al problema dado es  $(x^{*},y^{*},z^{*})=\l(\dfrac{a}{3},\dfrac{a}{3},\dfrac{a}{3}\r)$ en donde el máximo de $f$ es $f(x^{*},y^{*},z^{*}) = \dfrac{a^3}{27}$.
\end{respuesta}
    
%%%%%%%%%%%%%%%%%%%%%%%%%%%%%%%%%%%%%%%%%%%
%%%%%%%%%%%%%%%%%%%%%%%%%%%%%%%%%%%%%%%%%%%
%%%%%%%%%%%%%%%%%%%%%%%%%%%%%%%%%%%%%%%%%%%
% Tema 10: Optimizacion
\item 
    Dado $\Omega=\{(x,y,z)\in\R^3:x+y-z=1\ \wedge \ x=2y \}$, encuentre los extremos relativos del campo escalar
    \[
        \funcion{f}{\Omega}{\R}{(x,y,z)}{x^2-xy+xyz^2.}
    \]
    
\begin{respuesta}
    Dado que el $\dom(f)\neq \R^3$, debemos utilizar el multiplicador de Lagrange para determinar los extremos relativos de $f$. Definamos
    \[
        \funcion{L}{\R^2\times \Omega}{\R}{(\lambda_1,\lambda_2 x,y,z)}{f(x,y,z)-\lambda_1(x+y-z-1)-\lambda_2(x-2y).}
    \]
    Por definición de gradiente y hessiana, se tiene que para todo $(\lambda_1,\lambda_2,x,y,z)\in \R^2\times \Omega$,
    \[
        \nabla L(\lambda_1,\lambda_2,x,y,z)
        =(
        1+z-x-y,2y-x,2x+y(z^2-1)-\lambda_1-\lambda_2,
        x(z^2-1)-\lambda_1+2\lambda_2,
        2xyz+\lambda_1
        )
    \]
    y,
    \[
        H_L(\lambda_1,\lambda_2, x,y,z)=
        \begin{pmatrix}
            0&0&-1&-1&1\\
            0&0&-1&2&0\\
            -1&-1&2&z^2-1&2yz\\
            -1&2&z^2-1&0&2xz\\
            1&0&2yz&2xz&2xy
        \end{pmatrix}.        
    \]
    Es fácil ver que,
    \begin{align*}
        \nabla L(\lambda_1,\lambda_2,x,y,z)=0 
        \Di x+y=z+1 &\andm x=2y\andm 2x+y(z^2-1)=\lambda_1+\lambda_2\\
        &\andm x(z^2-1)=\lambda_1-2\lambda_2
        \andm 2xyz=-\lambda_1.
    \end{align*}
    El sistema anterior tiene una única solución
    \[
        a=(0,0,0,0,-1).
    \]
    Por definición
    \[
        H_L(a)=
        \begin{pmatrix}
            0&0&-1&-1&1\\
            0&0&-1&2&0\\
            -1&-1&2&0&0\\
            -1&2&0&0&0\\
            1&0&0&0&0
        \end{pmatrix}.        
    \]
    Mediante cálculo directo, se tiene que
    \[
        \det(H_L(a))=8
    \]
    así, para $m=2$,
    \[
        (-1)^m\det(H_L(a))_{2m+1}>0
    \]
    de donde $f$ alcanza un mínimo relativo en $a$. Así el valor, $(0,0,-1)$ minimiza al campo escalar $f$.
    
    Un ejercicio recomendado para el lector es reemplazar las variables y verificar que al optimizar una función real adecuada, se obtiene la misma respuesta.
\end{respuesta}



%%%%%%%%%%%%%%%%%%%%%%%%%%%%%%%%%%%%%%%%%%%
%%%%%%%%%%%%%%%%%%%%%%%%%%%%%%%%%%%%%%%%%%%
%%%%%%%%%%%%%%%%%%%%%%%%%%%%%%%%%%%%%%%%%%%
% Tema 10: Optimizacion
\item Dado los campos escalares
    \[
        \funcion{m}{\R^3}{\R}{(x,y,z)}{6x-y^2+xz+60} \texty \funcion{p}{\R^3}{\R}{(x,y,z)}{x^2+y^2+z^2}
    \]
    Resolver el problema
    \[
        \begin{cases}
            \max m(x,y,z)\\
            \text{sujeto  a:}\\
            p(x,y,z)=36.
        \end {cases}
    \]
    
% \begin{respuesta} 
%     Definimos la función que involucra a $m$ y a $p$.
%         \[
%             \funcion{h}{\R^4}{\R}{(\lambda,x,y,z,)}{m(x,y,z) - \lambda(p(x,y,z)-36).}
%         \]
%         Notemos que se tiene que
%         \[
%             h(\lambda,x,y,z) = (6x-y^2+xz+60) - \lambda\l(x^2+y^2+z^2-36\r)
%         \]
%         para todo $(x,y,z,\lambda)\in\R^4$, además, $h$ es una función diferenciable en $\R^4$, con vector gradiente 
%         \[
%             \nabla h(\lambda,x,y,z) = 
%             \begin{pmatrix}
%                 -(x^2+y^2+z^2-36)\\
%                 6+z - 2\lambda x\\
%                 -2y - 2\lambda y\\
%                 x - 2\lambda z
%             \end{pmatrix}
%         \]
%         Determinemos ahora, los puntos críticos de $h$, para esto igualamos el vector gradiente a 0, es decir, hacemos $\nabla h(\lambda,x,y,z) = (0,0,0,0)$, obteniendo el siguiente sistema:
%         \[
%             \begin{array}{cc}
%                 x^2+y^2+z^2=36\\
%                 6+z = 2\lambda x\\
%                 -2y = 2\lambda y\\
%                 x = 2\lambda z
%             \end{array}
%         \]
%         Las soluciones de este sistema son: $(0,0,0,-6)$, $\l(\sqrt{3}/2,3\sqrt{3},0,3\r)$, $\l(-\sqrt{3}/2,-3\sqrt{3},0,3\r)$ y $\l(-1,-4,\pm 4,2\r)$
%         Analizamos la matriz Hessiana de $h$ para determinar la naturaleza de los puntos críticos.
%         \[
%             H_h(\lambda,x,y,z) = 
%             \begin{pmatrix}
%                 0 & -2x & -2y & -2z\\
%                 -2x & -2\lambda & 0 & 1\\
%                 -2y & 0 & -2-2\lambda & 0\\
%                 -2z & 1 & 0 & -2\lambda
%             \end{pmatrix}
%         \]
%         para cada $(\lambda,x,y,z)\in\R^4$. 
        
%         De donde se obtiene que el campo más débil está en los puntos: $\l(-4,\pm 4,2\r)$
% \end{respuesta}
    
%%%%%%%%%%%%%%%%%%%%%%%%%%%%%%%%%%%%%%%%%%%
%%%%%%%%%%%%%%%%%%%%%%%%%%%%%%%%%%%%%%%%%%%
%%%%%%%%%%%%%%%%%%%%%%%%%%%%%%%%%%%%%%%%%%%
% Tema 10: Optimizacion
\item
    Dado los campos escalares
    \[
        \funcion{f}{\R^3}{\R}{(x,y,z)}{xy+zy}\\ \mbox{, } \funcion{g}{\R^3}{\R}{(x,y,z)}{x^2+y^2} \texty \funcion{h}{\R^3}{\R}{(x,y,z)}{yz}
    \]
    Resolver el problema
    \[
        \begin{cases}
            \max f(x,y,z)\\
            \text{sujeto  a:}\\
            g(x,y,z)=2 \texty\\
            h(x,y,z)=2.
        \end {cases}
    \]
    
\begin{respuesta}
    Para resolver el problema, se define la función $\func{L}{\R^5}{\R}$
    \[
        L(\lambda_1,\lambda_2,x,y,z)=f(x,y,z)-\lambda_1(g(x,y,z)-2)-\lambda_2(h(x,y,z)-2)=xy+zy-\lambda_1 \l(x^2+y^2-2\r)-\lambda_2 \l(yz-2\r)
    \]
    cuyo gradiente es
    \[
        \nabla L(\lambda_1,\lambda_1,x,y,z)=\l( \begin{matrix} -\l(x^2+y^2-2\r)\\-\l(yz-2\r)\\y-2\lambda_1 x\\x+z-2\lambda_1 y-\lambda_2 z\\y-\lambda_2 y \end{matrix} \r)^T.
    \]
    Los puntos críticos se obtiene igualando el gradiente a $(0,0,0,0,0)$. Resolviendo el sistema de ecuaciones se obtienen los siguientes puntos críticos
    \begin{align*}
        a_1=&\l(\dfrac{1}{2},1,1,1,2\r),\\
        a_2=&\l(-\dfrac{1}{2},1,1,-1,-2\r),\\
        a_3=&\l(-\dfrac{1}{2},1,-1,1,2\r),\\
        a_4=&\l(\dfrac{1}{2},1,-1,-1,-2\r).
	\end{align*}
	Para determinar la naturaleza de estos puntos se calcula la matriz hesiana
	\[
        H_L(\lambda_1,\lambda_2, x,y,z)=
        \begin{pmatrix}
            0 & 0 & -2x & -2y & 0 \\
            0 & 0 & 0 & -z & -y \\
            -2x & 0 & -2\lambda_1 & 1 & 0 \\
            -2y & -z & 1 & -2\lambda_1 & 1-\lambda_2 \\
            0 & -y & 0 & 1-\lambda_2 & 0
        \end{pmatrix},        
    \]
    de donde
    \[
        H_L(a_1)=
        \begin{pmatrix}
            0 & 0 & -2 & -2 & 0 \\
            0 & 0 & 0 & -2 & -1 \\
            -2 & 0 & -1 & 1 & 0 \\
            -2 & -2 & 1 & -1 & 0 \\
            0 & -1 & 0 & 0 & 0
        \end{pmatrix},        
    \]
    \[
        H_L(a_2)=
        \begin{pmatrix}
            0 & 0 & -2 & 2 & 0 \\
            0 & 0 & 0 & 2 & 1 \\
            -2 & 0 & 1 & 1 & 0 \\
            2 & 2 & 1 & 1 & 0 \\
            0 & 1 & 0 & 0 & 0
        \end{pmatrix},        
    \]
    \[
        H_L(a_3)=
        \begin{pmatrix}
            0 & 0 & 2 & -2 & 0 \\
            0 & 0 & 0 & -2 & -1 \\
            2 & 0 & 1 & 1 & 0 \\
            -2 & -2 & 1 & 1 & 0 \\
            0 & -1 & 0 & 0 & 0
        \end{pmatrix},        
    \]
    y
    \[
        H_L(a_4)=
        \begin{pmatrix}
            0 & 0 & 2 & 2 & 0 \\
            0 & 0 & 0 & 2 & 1 \\
            2 & 0 & -1 & 1 & 0 \\
            2 & 2 & 1 & -1 & 0 \\
            0 & 1 & 0 & 0 & 0
        \end{pmatrix}.        
    \]
    Nótese que en este caso $m=2$ y $n=3$ entonces para cada punto crítico se busca el valor de $(-1)^m\det(H_{2m+1})$, se tiene
    \begin{align*}
        \det H_{L_5}(a_1)=-16,\\
        \det H_{L_5}(a_2)=16,\\
        \det H_{L_5}(a_3)=16,\\
        \det H_{L_5}(a_4)=-16.
    \end{align*}
    Analizando la secuencia de gradientes para cada punto crítico, se puede ver que en los puntos $a_1$ y $a_4$ se tiene un máximo de valor 3 y en los puntos $a_2$ y $a_3$ se tiene un mínimo de valor 1.
\end{respuesta}
    
%%%%%%%%%%%%%%%%%%%%%%%%%%%%%%%%%%%%%%%%%%%
%%%%%%%%%%%%%%%%%%%%%%%%%%%%%%%%%%%%%%%%%%%
%%%%%%%%%%%%%%%%%%%%%%%%%%%%%%%%%%%%%%%%%%%
% Tema 10: Optimizacion
\item
    Dadas la funciones
    \[
         \funcion{V}{\R^3}{\R}{(x,y,z)}{xyz}
    \]
  	y,
	\[
        \funcion{A}{\R^3}{\R}{(x,y,z)}{2xy+2xz+2zy.}
    \] 
    Resolver el problema 
    \[
        \begin{cases}
        \begin{array}{c}
		\max V(x,y,z)\\
		\text{sujeto  a:}\\
			A(x,y,z)=21600.\\		
		\end{array}
        \end{cases} 
    \]

\begin{respuesta}

    Definamos la función $L:\R^4 \rightarrow \R$
    \[
        L(\lambda,x,y,z)=V(x,y,z)-\lambda (A(x,y,z)-21600)=xyz-\lambda ((2xy+2xz+2zy)-c)
    \]
    para $\lambda, x,y,z \in \R$. Procedemos a optimizar la función, para esto, hallamos el gradiente:
    \[
        \triangledown L(\lambda,x,y,z)=
        \left[
        \begin{array}{c}
		-(2xy+2xz+2zy-21600)\\
		yz-2\lambda y-2 \lambda z\\
        xz-2\lambda x-2\lambda z\\
        xy-2\lambda x-2\lambda y
		\end{array} 
        \right]^T
    \]
    para $\lambda, x,y,z \in \R$. Para obtener los puntos críticos, igualamos el gradiente a $(0,0,0,0)$ y obtenemos
    \[
        \left[
        \begin{array}{c}
		-(2xy+2xz+2zy-21600)\\
		yz-2\lambda y-2 \lambda z\\
        xz-2\lambda x-2\lambda z\\
        xy-2\lambda x-2\lambda y
		\end{array} 
        \right]=
        \left[
        \begin{array}{c}
        0\\
        0\\
        0\\
        0
        \end{array} 
        \right]
    \]
    de donde, se obtiene el sistema de ecuaciones
    \begin{align*}
        -(2xy+2xz+2zy-21600)&=0\\
        yz-2\lambda y-2 \lambda z&=0\\
        xz-2\lambda x-2\lambda z&=0\\
        xy-2\lambda x-2\lambda y&=0.
    \end{align*}
    Del sistema anterior se obtienen los puntos críticos
    \begin{align*}
        a_1=(15,60,60,60) \texty\\
        a_2=(-15,-60,-60,-60),
    \end{align*}
    para analizar la naturaleza de estos puntos se analiza la matriz hesiana
    \[
        H_L(\lambda, x,y,z)=
        \begin{pmatrix}
            0 & -2y-2z & -2x-2z & -2x-2y\\
            -2y-2z & 0 & z-2\lambda & y-2\lambda \\
            -2x-2z & z-2\lambda & 0 & x-2\lambda\\
            -2x-2y & y-2\lambda & x-2\lambda & 0
        \end{pmatrix}.        
    \]
    Se sabe que $m=1$ y $n=2$ entonces se calculará $(-1)^m\det H_{2m+1}$ y $(-1)^m\det H_{m+n}$ para cada punto crítico y se tiene
    \begin{align*}
        -\det H_{L_3}(a_1)<0 \texty -\det H_{L_4}(a_1)>0,\\
        -\det H_{L_3}(a_2)>0 \texty -\det H_{L_4}(a_2)<0,
    \end{align*}
    de donde, por la secuencia de los determinantes, se concluye que se tiene un máximo en $a_1$ y un punto de silla en $a_2$
\end{respuesta}

%%%%%%%%%%%%%%%%%%%%%%%%%%%%%%%%%%%%%%%%%%%
%%%%%%%%%%%%%%%%%%%%%%%%%%%%%%%%%%%%%%%%%%%
%%%%%%%%%%%%%%%%%%%%%%%%%%%%%%%%%%%%%%%%%%%
% Tema 10: Optimizacion
\item 
    Una empresa de Quito produce principalmente arroz y café. Imagine que es contratado para determinar la cantidad que debe producirse de cada uno para que el costo de producción sea lo menor posible. Empíricamente se conoce que los costos de producción del arroz y café están dados respectivamente por
	\[
		3x^2-12x+15\yds 2y^2-13y+5
	\]
	donde $x$ y $y$ representan la cantidad en toneladas de arroz y café producidos. Por otro lado, por la época de sequía, la producción de café es el doble de costosa que lo normal. Finalmente, el costo del transporte de ambos productos está dada por
	\[
		25 + 2xy
	\]
	donde $x$ y $y$ representan la cantidad en toneladas de arroz y café producidos, respectivamente. Determine la cantidad de arroz y café que debe producir la empresa para que el costo sea mínimo.


\begin{respuesta}
	Se utilizará la siguiente notación
	\begin{itemize}
	    \item $x:$ cantidad producida de toneladas de arroz
	    \item $y:$ cantidad producida de toneladas de café
	    \item $C(x,y):$ costo de producción.
	\end{itemize}
	Se sabe que el costo de producción resulta de la suma del costo de producción de arroz, costo de producción de café y costo de transporte. Además, durante el mes en análisis, el costo de producción del café es el doble de lo normal, entonces el costo de producción total viene dado por 
	\[
	\funcion{C}{\R^+\times \R^+}{\R}{(x,y)}{(3x^2-12x+15)+2(2y^2-13y+5)+(25+2xy)=3x^2+2xy+4y^2-12x-26y+50.}
	\]
	Es importante tomar en cuenta que como no existe producción negativa se debe considerar el dominio de $C$ a $\R^+\times \R^+$. 
	
	La función $C$ es la que se quiere optimizar, encontrar el mínimo en este caso. De acuerdo al ejercicio 4 de la hoja de ejercicios 8, la producción de arroz y café, para que el costo de producción sea mínimo, debe ser 1 y 3 toneladas respectivamente.
\end{respuesta}

%%%%%%%%%%%%%%%%%%%%%%%%%%%%%%%%%%%%%%%%%%%%%%
%%%%%%%%%%%%%%%%%%%%%%%%%%%%%%%%%%%%%%%%%%%%%%
%%%%%%%%%%%%%%%%%%%%%%%%%%%%%%%%%%%%%%%%%%%%%%

% Tema 10: Optimizacion
\item
	Una fábrica produce dispositivos de dos tipos, A y B, cuyos precios por unidad son 700 y 500 dólares, respectivamente. El costo de producir $x$ dispositivos tipo A y $y$ dispositivos tipo B está dado por la siguiente función
	\[
        \funcion{C}{\R^+ \times \R^+ }{\R}{(x,y)}{45x+32y-\frac{xy}{30}+\frac{x^2}{40}+\frac{y^2}{80}.}
    \]
	Determine los valores de $x$ y $y$ para que la utilidad sea máxima.

\begin{respuesta}
    Primero, utilizaremos las siguientes notaciones:
    	\begin{itemize}
    	    \item $x:$ cantidad de objetos $A$ producidos,
    	    \item $y:$ cantidad de objetos $B$ producidos,
        	\item $U(x,y)$: Utilidad, en dólares, al vender $x$ dispositivos tipo A y $y$ dispositivos tipo B,
            \item $V(x,y)$: ingreso, en dólares, producido al vender $x$ dispositivos tipo A y $y$ dispositivos tipo B,
            \item $C(x,y)$: costo de producción, en dólares, de $x$ dispositivos tipo A y $y$ dispositivos tipo B.
        \end{itemize}
    Se sabe que la utilidad es la diferencia entre el ingreso y el costo de producción. El ingreso viene dado por
    \[
        \funcion{V}{\R^+\times \R^+}{\R}{(x,y)}{700x+500y,}
    \]
    por tanto la función a maximizar es
	\[
     \funcion{U}{\R^+\times \R^+}{\R}{(x,y)}{655x+468y+\frac{xy}{30}-\frac{x^2}{80}+\frac{y^2}{40}.}
    \]
    Del ejercicio 12 de la hoja de ejercicios 8 se sabe que si se producen 230220 dispositivos de tipo A y 25680 dispositivos de tipo B, la utilidad será máxima.
\end{respuesta}    

%%%%%%%%%%%%%%%%%%%%%%%%%%%%%%%%%%%%%%%%%%%%%%%%
%%%%%%%%%%%%%%%%%%%%%%%%%%%%%%%%%%%%%%%%%%%%%%%%
%%%%%%%%%%%%%%%%%%%%%%%%%%%%%%%%%%%%%%%%%%%%%%%%
% Tema 10: Optimizacion
\item
    Se tiene una superficie de ecuación $z=\dfrac{1}{xy}$, con $(x,y) \in \R^2$ tal que $xy\neq 0$. Hallar los puntos de la superficie más cercanos al origen. 

\begin{respuesta}
    Se utilizará la siguiente notación
    \begin{itemize}
        \item $(x,y):$ punto sobre la superficie
        \item $d(x,y):$ distancia del punto $(x,y)$ al origen, en unidades de longitud
        \item $D(x,y):$ cuadrado de la distancia del punto $(x,y)$ al origen.
        
    \end{itemize}
   Se sabe que la distancia de un punto al origen viene definida por $d=\sqrt{x^2+y^2+z^2}$, como $d \geq 0$ entonces la distancia será mínima cuando $x^2+y^2+z^2$ sea mínimo. Finalmente se definen las funciones
   \[
        \funcion{D}{\R^3}{\R}{(x,y,z)}{x^2+y^2+z^2}
   \]
   y
   \[
         \funcion{g}{\R^* \times \R^* \times \R }{\R}{\R}{z-\dfrac{1}{xy}.}
   \]
   Con esto se debe optimizar, 
   \[
        \begin{cases}
            \max D(x,y,z)\\
            \text{sujeto  a:}\\
            g(x,y,z)=0.
        \end {cases}
    \]
   Por el ejercicio 1 de la hoja de ejercicios 9 se sabe que la distancia es mínima en los puntos $(1,1,1)$, $(1,-1,-1)$, $(-1,-1,1)$ y $(-1,1,-1)$ y será $d=\sqrt{3}$.
   
   Otra forma de optimizar este problema se detalla en el ejercicio 11 de la hoja de ejercicios 8.

\end{respuesta}

%%%%%%%%%%%%%%%%%%%%%%%%%%%%%%%%%%%%%%%%%%%%%
%%%%%%%%%%%%%%%%%%%%%%%%%%%%%%%%%%%%%%%%%%%%%
%%%%%%%%%%%%%%%%%%%%%%%%%%%%%%%%%%%%%%%%%%%%%

% Tema 10: Optimizacion
    \item 
    Dado un número natural mayor que cero, dividirlo en tres sumandos de manera que su producto sea el máximo.

\begin{respuesta}
    Sea $a\in\N^{*}$, supongamos que:
    \begin{table}[h]
        \begin{tabular}{cl}
            $x$: &el primer sumando,\\
            $y$: &el segundo sumando,\\
            $z$: &el tercer sumando.
        \end{tabular}
    \end{table}
    
    Notemos que se debe verificar que $x+y+z=a$ siendo su producto $xyz$ el máximo. Se definen entonces:
    \[
         \funcion{f}{\R^3}{\R}{(x,y,z)}{xyz}
    \]
  	y,
	\[
        \funcion{g}{\R^3}{\R}{(x,y,z)}{x + y+ z.}
    \] 
    A partir de lo cual, buscamos resolver el siguiente problema 
    \[
        \begin{cases}
        \begin{array}{c}
		\max f(x,y,z)\\
		\text{sujeto  a:}\\
			g(x,y,z)=a.\\		
		\end{array}
        \end{cases} 
    \]
    En base al ejercicio 2 de la hoja de ejercicios 9, se puede concluir que los sumandos de un número natural $a> 0$, cuyo producto $xyz$ es máximo son $x=y=z=\dfrac{a}{3}$.
\end{respuesta}

%%%%%%%%%%%%%%%%%%%%%%%%%%%%%%%%%%%%%%%%%%%%%%%%%%
%%%%%%%%%%%%%%%%%%%%%%%%%%%%%%%%%%%%%%%%%%%%%%%%%%
%%%%%%%%%%%%%%%%%%%%%%%%%%%%%%%%%%%%%%%%%%%%%%%%%%

% Tema 10: Optimizacion
\item Se desea instalar un radiotelescopio sobre la superficie de un planeta recién descubierto. Para minimizar la interferencia, se quiere colocar el radiotelescopio donde el campo magnético del planeta sea más débil. El planeta es esférico con un radio de 36 unidades. Con base en un sistema de coordenadas cuyo origen es el centro del planeta, la fuerza del campo magnético está dada por $m(x,y,z)=6x-y^2+xz+60$. ¿Dónde se debe ubicar el radiotelescopio?

\begin{respuesta}
    Se usará la siguiente notación
    \begin{itemize}
        \item $(x,y,z):$ punto sobre la superficie del planeta,
        \item $m(x,y,z):$ fuerza del campo magnético, unidades de fuerza.
    \end{itemize}
    Se quiere encontrar el punto donde la fuerza ejercida por el campo magnético sea mínima. Además los puntos donde se quiere medir este campo está sobre la superficie del planeta, entonces se definen la funciones
    \[
        \funcion{m}{\R^3}{\R}{(x,y,z)}{6x-y^2+xz+60} \texty \funcion{p}{\R^3}{\R}{(x,y,z)}{x^2+y^2+z^2.}
    \]
    Se quiere optimizar
    \[
        \begin{cases}
            \max m(x,y,z)\\
            \text{sujeto  a:}\\
            p(x,y,z)=36.
        \end {cases}
    \]
    En base al ejercicio 4 de la hoja de ejercicios 9, los puntos donde la intensidad del campo magnético es mínima, donde se debería instalar el radiotelescopio, son $(-4,4,2)$ y $(-4,-4,2)$.
    
    Los puntos críticos del lagrangiano de este problema de optimización son $(0,0,0,-6)$, $(\sqrt{3}/2,3\sqrt{3},0,3)$, $(-\sqrt{3}/2,-3\sqrt{3},0,3)$ y $\l(-1,-4,\pm 4,2\r)$.
    
\end{respuesta}

%%%%%%%%%%%%%%%%%%%%%%%%%%%%%%%%%%%%%%%%%%%%%%%%%%
%%%%%%%%%%%%%%%%%%%%%%%%%%%%%%%%%%%%%%%%%%%%%%%%%%
%%%%%%%%%%%%%%%%%%%%%%%%%%%%%%%%%%%%%%%%%%%%%%%%%%


% Tema 10: Optimizacion
\item
	La temperatura en un punto $(x,y,z)$ del espacio tridimensional está dada por la función 
    \[
        \funcion{f}{\R^3}{\R}{(x,y,z)}{xy+zy.}
    \] 
    Determinar los puntos donde la temperatura es máxima y mínima, sobre la intersección de las superficies de ecuación $x^2+y^2=2$ y $yz=2$.
    
\begin{respuesta}
    \begin{itemize}
        \item $(x,y,z):$ punto del espacio $\R^3$ que pertenece a las superficies de ecuación $x^2+y^2=2$ y $yz=2$,
        \item $t(x,y,z):$ temperatura en el punto $(x,y,z)$, en unidades de temperatura. 
    \end{itemize}
    Se definen las funciones 
     \[
        \funcion{t}{\R^3}{\R}{(x,y,z)}{xy+zy}\\ \mbox{, } \funcion{g}{\R^3}{\R}{(x,y,z)}{x^2+y^2} \texty \funcion{h}{\R^3}{\R}{(x,y,z)}{yz.}
    \]
    Se quieren encontrar los puntos donde la temperatura es máxima y mínima, entonces se resolverá el siguiente problema de optimización
    \[
        \begin{cases}
            \max t(x,y,z)\\
            \text{sujeto  a:}\\
            g(x,y,z)=2 \texty\\
            h(x,y,z)=2.
        \end {cases}
    \]
    Del ejercicio 5 de la hoja de ejercicios 9, se tiene que en los puntos $(1,1,2)$ y $(-1,-1,-2)$ la temperatura es máxima e igual a 3, mientras que en los puntos $(1,-1,-2)$ y $(-1,1,2)$ la temperatura es mínima e igual a 1.
\end{respuesta}

%%%%%%%%%%%%%%%%%%%%%%%%%%%%%%%%%%%%%%%%%%%
%%%%%%%%%%%%%%%%%%%%%%%%%%%%%%%%%%%%%%%%%%%
%%%%%%%%%%%%%%%%%%%%%%%%%%%%%%%%%%%%%%%%%%%

% Tema 10: Optimizacion
\item
    Se desea construir una cisterna hermética, en forma de un paralelepípedo regular de área 21600 $m^2$, de tal forma que su el volumen sea máximo. Encontrar las dimensiones que debe tener la cisterna.
    
\begin{respuesta}

    Primero utilizaremos las siguientes notaciones:
    \begin{itemize}
        \item $x$: ancho de la cisterna en metros,
        \item $y$: largo de la cisterna en metros,
        \item $z$: altura de la cisterna en metros,
        \item $V(x,y,z)$: volumen de la cisterna en metros cúbicos,
        \item $A(x,y,z)$: área de la cisterna en metros cuadrados.
    \end{itemize}
    Nótese que se quiere encontrar el valos máximo de $V$ sabiendo que tiene una limitación (restricción), el área. Por otro lado, sabemos que las dimensiones de la cisterna serán reales positivos. Se definen las funciones
    \[
         \funcion{V}{\R^+ \times \R^+ \times \R^+}{\R}{(x,y,z)}{xyz}
    \]
  	y,
	\[
        \funcion{A}{\R^+ \times \R^+ \times \R^+}{\R}{(x,y,z)}{2xy+2xz+2zy.}
    \] 
    Se quiere resolver el problema
    \[
        \begin{cases}
        \begin{array}{c}
		\max V(x,y,z)\\
		\text{sujeto  a:}\\
			A(x,y,z)=21600.\\		
		\end{array}
        \end{cases} 
    \]

    En base al ejercicio 6 de la hoja de ejercicios 9 se tiene que las dimensiones de la cisterna, para que el volumen de ésta sea  máximo, son $(60,60,60)$. Nótese que el punto crítico $(-60,-60,-60)$ no es considerado ya que no son parte del dominio de $V$ en este problema.  
\end{respuesta}

%%%%%%%%%%%%%%%%%%%%%%%%%%%%%%%%%%%%%%%%%%%
%%%%%%%%%%%%%%%%%%%%%%%%%%%%%%%%%%%%%%%%%%%
%%%%%%%%%%%%%%%%%%%%%%%%%%%%%%%%%%%%%%%%%%%
% Tema 10: Optimizacion
\item
    Dada una superficie representada por la función 
    \[
        \funcion{f}{I}{\R}{(x,y)}{x^2+y^2-x-y+1} 
    \]
    donde $I$ es el disco definido por $\lbrace (x,y) \in \R^2:x^2+y^2 \leq 1 \rbrace$, determinar los puntos de la misma que tienen mayor y menor altura. 

%%%%%%%%%%%%%%%%%%%%%%%%%%%%%%%%%%%%%%%%%%%
%%%%%%%%%%%%%%%%%%%%%%%%%%%%%%%%%%%%%%%%%%%
%%%%%%%%%%%%%%%%%%%%%%%%%%%%%%%%%%%%%%%%%%%
% Tema 10: Optimizacion
\item
    La temperatura en cualquier punto de coordenadas $(x,y,z)$ ubicado en la esfera de ecuación $x^2+y^2+z^2=11$ viene dada por $T(x,y,z)=20 + 2x +2y + z^2$. Determinar las temperaturas extremas sobre la curva de intersección de la esfera con el plano $P$ de ecuación $x+y+z=3$.
    
% \begin{respuesta}
%     Objetivo: Maximizar y minimizar la temperatura T:
% 	\[
%       T(x,y,z)=20 + 2x +2y + z^22
%     \]
%     para $(x,y,z) \in \R^3$. Se tienen dos restricciones, la esfera,
% 	\[
%       x^2+y^2+z^2=11,
%     \]
%     y el plano,
% 	\[
%       x+y+z=3.
%     \]
%     Por el método de multiplicadores de Lagrange, tenemos\comentario{no podemos sacar el gradiente a $E$, pues esta es una esfera, no una función, debemos definir una función cuya imagen sea la esfera y a esa sí le podemos sacar el gradiente.}:
% 	\[
%     \begin{aligned}
%       \nabla T(x,y,z) &= \lambda \nabla E(x,y,z)+ \mu \nabla P(x,y,z)\\
%       (x,y,2z)&=\lambda(2x,2y,2z)+ \mu(1,1,1)
%     \end{aligned}
%     \]
%     Se obtiene el sistema de ecuaciones:
% 	\[
%         \begin{cases}
%             2 = 2 \lambda x+\mu \\
%             2 = 2 \lambda y+\mu \\
%             2z =2 \lambda z+\mu \\
%             x^2+y^2+z^2=11\\
%             x+y+z=3
%         \end{cases}
%     \]
%     Resolviéndolo se tiene que $\lambda=0$ y $x=y$.\\
%     Para $\lambda=0$, los puntos críticos son:
% 	\[
%         A(x,y,z)=( -1,3,1)
%     \]
%     \[
%         B(x,y,z)=(3,-1,1)
%     \]
%     Para $x=y$, los puntos críticos son:
% 	\[
%         C(x,y,z)=\l( 1+\dfrac{2\sqrt{3}}{3},1+\dfrac{2\sqrt{3}}{3},1-\dfrac{4\sqrt{3}}{3}\r)
%     \]
%     \[
%         D(x,y,z)=\l( 1-\dfrac{2\sqrt{3}}{3},1-\dfrac{2\sqrt{3}}{3},1+\dfrac{4\sqrt{3}}{3}\r)
%     \]
%     Evaluando $A,B,C,D \in \R^3$ en $T(x,y,z)$,
% 	\[
%     \begin{aligned}
%       T(A) &= 25\\
%       T(B) &= 25\\
%       T(C) &= \dfrac{75}{16}\\
%       T(D) &= \dfrac{75}{16}
%     \end{aligned}
%     \]
%     Por lo tanto $A,B \in \R^3$ son mínimos y $C,D \in \R^3$ son máximos.
% \end{respuesta}


%%%%%%%%%%%%%%%%%%%%%%%%%%%%%%%%%%%%%%%%%%%
%%%%%%%%%%%%%%%%%%%%%%%%%%%%%%%%%%%%%%%%%%%
%%%%%%%%%%%%%%%%%%%%%%%%%%%%%%%%%%%%%%%%%%%
% Tema 10: Optimizacion
\item
    En $\R^3$, hallar el punto del plano $P$ de ecuación $2x-2y+z=4$, que está más próximo al origen de coordenadas.

% \begin{respuesta}
%     Objetivo: Minimizar la distancia punto-plano, esta distancia está dada por la ecuación \comentario{no se da por ecuaciones, se da por funciones}:
% 	\[
%       d^2=x^2+y^2+z^2
%     \]
%     \textcolor{red}{
%     \[
%     \funcion{d}{\R^3}{\R}{(x,y,z)}{x^2+y^2+z^2}
%     \], de esta forma, toma sentido sacar el gradiente a $d$}
%     para todo $(x,y,z) \in \R^3$. Se tiene una sola restricción que es el plano $2x-2y+z=4$.\\

%     Por el método de multiplicadores de Lagrange, tenemos:
% 	\[
%     \begin{aligned}
%       \nabla d(x,y,z) &= \lambda \nabla P(x,y,z)\\
%       (2x,2y,2z)&=\lambda(2,-2,1)
%     \end{aligned}
%     \]
%     Se obtiene el sistema de ecuaciones:
% 	\[
%         \begin{cases}
%             2x = 2 \lambda \\
%             2y = -2 \lambda \\
%             2z = \lambda \\
%             2x-2y+z = 4
%         \end{cases}
%     \]
%     Se obtiene un único extremo:
% 	\[
%         (x,y,z)=\l( \frac{8}{9},-\frac{8}{9},\dfrac{4}{9}\r),
%     \]
%     que es un mínimo en la frontera de $\R^3$.
% \end{respuesta}

%%%%%%%%%%%%%%%%%%%%%%%%%%%%%%%%%%%%%%%%%%%
%%%%%%%%%%%%%%%%%%%%%%%%%%%%%%%%%%%%%%%%%%%
%%%%%%%%%%%%%%%%%%%%%%%%%%%%%%%%%%%%%%%%%%%
% Tema 11: Integrales dobles
\item 
	Dado el campo escalar 
	\[
		\funcion{f}{\R^2}{\R}{(x,y)}{x^2-2xy}
	\] 
	y la región
	\[
		D=\{(x,y)\in\R^2: x \geq 0, \ y \geq 0, \ 3x + y \leq 3\},
	\]
	calcule la integral doble
	\[
		\iint_D f dA.
	\]

\begin{respuesta}
	Dado que la función $f$ es un polinomio, $f$ es continua en todo $\R^2$. En particular $f$ es continua en $D$, de donde 
	\[
		\iint_D f dA
	\]
	existe. Por definición de integral doble se tiene que
	\[
		\iint_D f dA=\int_0^3\int_{0}^{3-3x}f(x,y) dy dx.
	\]
	Entonces
	\begin{align*}
		\int_{0}^1\int_{0}^{3-3x}f(x,y) dy dx
			&=\int_0^1\int_{0}^{3-3x}(x^2-2xy) dy dx\\
			&=\int_0^1\left.(x^2y-xy^2)\right|_0^{3-3x} dy dx\\
			&=\int_0^1(x^2(3-3x)-x(3-3x)^2) dy dx\\
			&=\int_0^1(-12 x^3 + 21 x^2 - 9 x) dx\\
			&=\left.-3 x^4 + 7 x^3 - \dfrac{9}2 x^2\right|_0^1\\
			&=-\dfrac{1}{2}.
	\end{align*}
	Por lo tanto,
	\[
		\iint_D f dA=-\dfrac{1}{2}.\qedhere
	\]
\end{respuesta}

%%%%%%%%%%%%%%%%%%%%%%%%%%%%%%%%%%%%%%%%%%%
%%%%%%%%%%%%%%%%%%%%%%%%%%%%%%%%%%%%%%%%%%%
%%%%%%%%%%%%%%%%%%%%%%%%%%%%%%%%%%%%%%%%%%%
% Tema 11: Integrales dobles
\item 
	Dados el campo escalar 
	\[
		\funcion{f}{\R^2}{\R}{(x,y)}{x}
	\] 
	y 
	\[
		\funcion{C_1}{\lcl0,1\rcl}{\lcl0,2\rcl}{x}{4x-2x^2}
		\yds
		\funcion{C_2}{\lcl1,2\rcl}{\lcl0,2\rcl}{x}{2(x-2)^2}
	\]
	se define por $D$ a la región limitada por las gráficas de eje $x$, $C_1$ y $C_2$. Exprese la integral doble
	\[
		\iint_D f dA.
	\]
	como integrales iteradas de dos formas distintas. Calcule una de ellas.

\begin{respuesta}
	Dado que la función $f$ es un polinomio, $f$ es continua en todo $\R^2$. En particular $f$ es continua en $D$, de donde 
	\[
		\iint_D f dA
	\]
	existe. La región $D$ se puede observar en el siguiente gráfico
	\begin{center}
		\includegraphics[scale=0.7]{Graficos/fig09:02.eps}
	\end{center}

	Por otra parte, $C_1$ y $C_2$ son funciones biyectivas. Se tiene que
	\[
		y= C_1(x)\Di x=\dfrac{2-\sqrt{4-2y}}{2}
		\yds
		v= C_2(u)\Di u=2-\sqrt{\dfrac{v}{2}}
	\]
	donde $x\in\lcl0,1\rcl$, $u\in\lcl1,2\rcl$ y $y,v\in\lcl0,2\rcl$. 

	Para calcular la integral doble tenemos dos posibilidades:
	\begin{itemize}
	\item 
    	Respecto a $dx dy$, se tiene que
        \[
			\iint_D f dA=
			\int_{0}^1\int_{0}^{4x-2x^2}x dy dx+
			\int_{1}^2\int_{0}^{2(x-2)^2}x dy dx
		\]
		
	\item 
    	Respecto a $dy dx$, se tiene que
		\[
			\iint_D f dA=
			\int_{0}^2\int_{\frac{2-\sqrt{4-2y}}{2}}^{2-\sqrt{\frac{2}{y}}} x dx dy.
		\]
	\end{itemize}
% 	Como se puede ver la segunda integral es más simple de calcular:
% 	\begin{align*}
% 		\iint_D f dA
% 		&=\int_{0}^2\int_{\frac{2-\sqrt{4-2y}}{2}}^{2-\sqrt{\frac{y}{2}}}f(x,y) dx dy\\
% 		&=\int_{0}^2\int_{\frac{2-\sqrt{4-2y}}{2}}^{2-\sqrt{\frac{y}{2}}} y dx dy\\
% 		&=\int_{0}^2 x y\left. \right|_{\frac{2-\sqrt{4-2y}}{2}}^{2-\sqrt{\frac{y}{2}}} dy\\
% 		&=\int_{0}^2 y\left[2-\sqrt{\frac{y}{2}} - \dfrac{2-\sqrt{4-2y}}{2}\right]dy\\
% 		&=\int_{0}^2 y\left[1-\sqrt{\frac{y}{2}} + \dfrac{\sqrt{4-2y}}{2}\right]dy\\
% 		&=\int_{0}^2 \left[y-\frac{y^{3/2}}{\sqrt{2}} + \dfrac{y\sqrt{4-2y}}{2}\right]dy\\
% 		&=\dfrac{22}{15}.
% 	\end{align*}
% 	Por lo tanto,
% 	\[
% 		\iint_D f dA=\dfrac{22}{15}.\qedhere
% 	\]
\end{respuesta}

%%%%%%%%%%%%%%%%%%%%%%%%%%%%%%%%%%%%%%%%%%%
%%%%%%%%%%%%%%%%%%%%%%%%%%%%%%%%%%%%%%%%%%%
%%%%%%%%%%%%%%%%%%%%%%%%%%%%%%%%%%%%%%%%%%%
% Tema 11: Integrales dobles
\item 
    Dados el campo escalar
    \[
        \funcion{f}{\R^2}{\R}{(x,y)}{f(x,y)=2xy}
    \]
    y la región $D$ definida como el trapecio cuyos vértices son: $A=(1,1)$, $B=(4,1)$, $C=(2,4)$ y $D=(3,4)$; determine:
    \[
        \iint_D f(x,y) \,dA
    \]

\begin{respuesta} 
    Primero debemos definir la región D. La región D limitada por este trapecio se puede ver de manera natural como una región tipo 2. 
    \begin{center}
        \includegraphics[scale=.7]{Graficos/fig09:03.eps}
    \end{center}
    La recta $l_1$ que pasa por A y C es
    \[
        y=3x-2
            \]
    o bien 
    \[
        x=\alpha(y)=\frac{(y+2)}{3}
            \]
    La recta $l_2$ que pasa por B y D es 
    \[
        y=-3x+13
            \]
    o bien 
    \[
        x=\beta(y)=\frac{(13-y)}{3}
            \]
    de este modo la región que el trapecio encierra se puede ver limitada por \begin{itemize}
    \item la recta horizontal $y=1$, por debajo, 
    \item la recta horizontal $y=4$, por encima, 
    \item la recta $x=\alpha(y)=(y+2)/3$, por la izquierda y
    \item la recta $x=\beta(y)=(13-y)/3$, por la derecha.
    \end{itemize}
    por tanto la región D queda descrita como 
    \[
        D=\{(x,y)\in\R^2:(y+2)/3\leq x \leq (13-y)/3; 1\leq y \leq 4\}
            \]
    entonces
    \[
        \iint_D 2xy \,dx\,dy=\int_{1}^{4}\int_{(y+2)/3}^{(13-y)/3}2xy\,dx\,dy
            \]
    donde primero se hace la integral respecto de $x$ (evaluándose entre $(13-y)/3$ y $(y+2)/3$) y luego respecto a $y$ (evaluándose entre 1 y 4).
    \begin{align*}
        \iint_D 2xy \,dx\,dy&=\int_{1}^{4}\left(x^2y\right)_{(y+2)/3}^{(13-y)/3}\,dx\\
        &=\int_{1}^{4}y \left[\left(\frac{13-y}{3}\right)^2-\left(\frac{y+2}{3}\right)^2\right]\,dx\\
        &=\frac{1}{9}\int_{1}^{4}(-30y^2+165y)\,dy\\
        &=\frac{1}{9}\left(-10y^3+\frac{165}{2}y^2\right) _1^{4} \\
        &=\frac{135}{2}
    \end{align*}
\end{respuesta}

%%%%%%%%%%%%%%%%%%%%%%%%%%%%%%%%%%%%%%%%%%%
%%%%%%%%%%%%%%%%%%%%%%%%%%%%%%%%%%%%%%%%%%%
%%%%%%%%%%%%%%%%%%%%%%%%%%%%%%%%%%%%%%%%%%%
% Tema 11: Integrales dobles
\item 
    Dados el campo escalar
    \[
        \funcion{f}{\R^2}{\R}{(x,y)}{f(x,y)=x+y}
    \]
	y $D$ la región que corresponde la cuarta parte, situada en el primer cuadrante, del anillo circular limitado por las circunferencias de ecuaciones
	\[
        x^2+y^2=1
    \]
	y        
	\[
        x^2+y^2=4.
    \]
	Determinar         
	\[
	    \iint_D f(x,y) \,dA
    \]
	           

\begin{respuesta} 
	Primero definimos la región D, esta región no se puede ver globalmente como una región de tipo 1 o de tipo 2. 
	\begin{center} 
        \includegraphics[width=8cm]{Graficos/fig09:04.eps}
    \end{center}
	Se debe partirla en dos subregiones, ya sea con la recta $x=1$ (en cuyo caso las subregiones resueltas son de tipo 1), o con la recta $y=1$ (en cuyo caso las regiones resueltas son de tipo 2). Tomando la primera opición.
	
	\[
	    D=D_1\cup D_2
	        \]
	entonces
	\[
	    D_1=\{(x,y)\in \R^2:\sqrt{1-x^2}\leq y \leq \sqrt{4-x^2}; 0\leq x \leq 1\}
	        \]
	y
	\[
	    D_2=\{(x,y)\in \R^2:0 \leq y \leq \sqrt{4-x^2}; 1\leq x \leq 2\}.
	        \]
	por tanto se tiene 
    \begin{align*}
        \iint_D (x+y) \,dx\,dy=&\iint_{D_1}(x,y)\,dx \,dy+\iint_{D_2}(x,y)\,dx \,dy\\
        =&\int_{0}^{1}\int_{\sqrt{1-x^2}}^{\sqrt{4-x^2}}(x+y)\,dy\,dx+\int_{1}^{2}\int_{0}^{\sqrt{4-x^2}}(x+y)\,dy\,dx\\
        =&\int_0^1x(\sqrt{4-x^2}-\sqrt{1-x^2})\,dx+\int_0^1\frac{1}{2}(4-x^2-(1-x^2))\,dx\\
       & +\int_1^2x\sqrt{4-x^2}\,dx+\int_1^2\frac{1}{2}(4-x^2)\,dx\\
       =&\left[-\frac{1}{3}(4-x^2)^{3/2}+\frac{1}{3}(1-x^2)^{3/2}+\frac{3}{2}x\right]_0^1+\left[-\frac{1}{3}(4-x^2)^{3/2}+2x-\frac{1}{6}x^3\right]_1^2\\
       =&\frac{14}{3}
     \end{align*}             
\end{respuesta}
    

%%%%%%%%%%%%%%%%%%%%%%%%%%%%%%%%%%%%%%%%%%%
%%%%%%%%%%%%%%%%%%%%%%%%%%%%%%%%%%%%%%%%%%%
%%%%%%%%%%%%%%%%%%%%%%%%%%%%%%%%%%%%%%%%%%%
% Tema 11: Integrales dobles
\item
    Dibujar la región de integración y calcular la integral doble $I=\iint_R x \cos(x+y) dxdy$, donde $R$ es el triángulo de vértices $(0,0), (\pi,0)$, y $(\pi,\pi)$

\begin{respuesta}
    \begin{align*}
    I=&\int \int_R x \cos(x+y) dxdy\\
    =&\int _{0}^{\pi} \left[\int _{0}^{x}  x \cos(x+y) dy \right]dx\\
    =&\int _{0}^{\pi} x \sen(x+y) \mid_{0}^{\pi}\ dx\\
    =&-\dfrac{3\pi}{2}
    \end{align*}
\end{respuesta}

%%%%%%%%%%%%%%%%%%%%%%%%%%%%%%%%%%%%%%%%%%%
%%%%%%%%%%%%%%%%%%%%%%%%%%%%%%%%%%%%%%%%%%%
%%%%%%%%%%%%%%%%%%%%%%%%%%%%%%%%%%%%%%%%%%%
% Tema 11: Integrales dobles
\item 
    Dado el campo escalar
    \[
        \funcion{f}{\R^2}{\R}
        {(x,y)}{f(x,y)=x^3}
	\]
    y la región
	\[
        D=\{(x,y)\in\R^2 \, : \, 0 \leq y \leq 4 - x^2, \, -2 \leq x \leq 2 \},
    \]
    calcular $\iint_D f(x,y)\, dxdy.$

\begin{respuesta}
\end{respuesta}
    
%%%%%%%%%%%%%%%%%%%%%%%%%%%%%%%%%%%%%%%%%%%
%%%%%%%%%%%%%%%%%%%%%%%%%%%%%%%%%%%%%%%%%%%
%%%%%%%%%%%%%%%%%%%%%%%%%%%%%%%%%%%%%%%%%%%
% Tema 11: Integrales dobles
\item 
    Dado el campo escalar
    \[
        \funcion{f}{\R^2}{\R}
        {(x,y)}{f(x,y)=3}
	\]
    y la región
	\[
        D=\{(x,y)\in\R^2 \, : \, - \sqrt{1-x^2} \leq y \leq \sqrt{1-x^2}, \, 0 \leq x \leq 1 \},
    \]
    calcular $\iint_D f(x,y)\, dydx.$

\begin{respuesta}
\end{respuesta}

%%%%%%%%%%%%%%%%%%%%%%%%%%%%%%%%%%%%%%%%%%%
%%%%%%%%%%%%%%%%%%%%%%%%%%%%%%%%%%%%%%%%%%%
%%%%%%%%%%%%%%%%%%%%%%%%%%%%%%%%%%%%%%%%%%%
% Tema 12: Cambio de variable
\item
    Calcular
    \[
        \iint_D \sqrt{1-\frac{x^2}{a^2}-\frac{y^2}{b^2}}\,dx\,dy
    \]
    donde $D$ es la región interior de la elipse de ecuación $\dfrac{x^2}{a^2}+\dfrac{y^2}{b^2}=1$ 
    
\begin{respuesta} 
    Primero identifiquemos la región $D$
	\begin{center}
        \includegraphics[width=8cm]{Graficos/fig10:05.eps}
    \end{center}
	Para resolver esta integral tomaremos el cambio de variable a coordenadas polares, es decir el cambio de variable
	\[
        x=ar\cos\theta \qquad y=br\sen \theta
    \]
    es decir, tenemos la función definida por 
    \[
    P(r,\theta)=(ar\cos\theta,br\sen\theta)
    \]
    para $(r,\theta)\in\R^2$. Ahora necesitamos encontrar $D^*$ tal que $P(D^*)=D$   
    \begin{itemize}
        \item Si $x=a$; $\theta=0 \Imp r=1$
        \item $0\leq\theta<2\pi$
    \end{itemize}
    entonces la región $D^*=[0,1]\times[0,2\pi]$
    por lo tanto 
    \begin{align*}
        \iint_D \sqrt{1-\frac{x^2}{a^2}-\frac{y^2}{b^2}}\,dx\,dy&=\iint_{D^*}\sqrt{1-r^2}abr\,dr\,d\theta\\    
        &=ab\int_0^1\int_0^{2\pi}\sqrt{1-r^2}r\,dr\,d\theta\\
        &=2\pi ab\int_0^1\sqrt{1-r^2}r\,dr\\
        &=2\pi ab \left(-\frac{1}{3}(1-r^2)^{3/2}\right)\mid_{0}^{1}\\
        &=\frac{2}{3}\pi ab\qedhere
    \end{align*}
\end{respuesta}

%%%%%%%%%%%%%%%%%%%%%%%%%%%%%%%%%%%%%%%%%%%
%%%%%%%%%%%%%%%%%%%%%%%%%%%%%%%%%%%%%%%%%%%
%%%%%%%%%%%%%%%%%%%%%%%%%%%%%%%%%%%%%%%%%%%
% Tema 12: Cambio de variable
\item
	Compruebe la siguiente igualdad mediante la introducción de un conveniente cambio de variables,
    \[
        \iint_R f(xy) \,dx dy= \ln(2) \int_{1}^{2} f(u)\, du, 
    \]
    donde $R$ es la región del primer cuadrante limitada por las curvas de ecuaciones $xy=1, xy=2, y=x, y=4x$

\begin{respuesta}
    Se resuelve utilizando el cambio de variable,
    \begin{align*}
        xy=&u\\
        y=&v
    \end{align*}
    y
    \begin{align*}
        x=\dfrac{u}{v}\\
        y=v.\\
    \end{align*}
    Por lo que la imagen de cada curva es,
    \begin{align*}
        u&=1\\
        u&=2\\
        u&=v^2\\
        v^2&=4u.
    \end{align*}
    Como $x$ e $y$ son positivos, $u$ y $v$ también son positivos. Así
    \[
        J(u,v)=
        \left[
        \begin{array}{cc}
		\dfrac{1}{v} & -\dfrac{u}{v^2}\\
		0 & 1
        \end{array} 
        \right]
        =\dfrac{1}{v}>0,
    \]
    entonces,
    \[
    	\int_1^2 
    	\left[
        \int_{\sqrt{u}}^{\sqrt{4u}} \dfrac{1}{v} f(u) dv
        \right] du 
        = \ln(2) \int_{1}^{2} f(u) du, 
    \]
\end{respuesta}

%%%%%%%%%%%%%%%%%%%%%%%%%%%%%%%%%%%%%%%%%%%
%%%%%%%%%%%%%%%%%%%%%%%%%%%%%%%%%%%%%%%%%%%
%%%%%%%%%%%%%%%%%%%%%%%%%%%%%%%%%%%%%%%%%%%
% Tema 12: Cambio de variable
\item
    Calcular la integral 
    \[
        \iint_R e \ ^{\frac{x-y+1}{x+y-2}} \,dx dy, 
    \]
    donde $R$ es el triángulo de vértices $(0,0), (2,0), (0,2)$.
    
 \begin{respuesta}
%     Sea el cambio de variable,
%     \begin{align*}
%         u=x-y+1\\
%         v=x+y-2,
%     \end{align*}
%     de donde,
%     \begin{align*}
%         x=\dfrac{u+v+1}{2}\\
%         y=\dfrac{v-u+3}{2}.
%     \end{align*}
%     Se obtiene la matriz Jacoviana,
%     \[
%         J(u,v)=
%         \left[
%         \begin{array}{cc}
% 		\dfrac{1}{2} & \dfrac{1}{2}\\
% 		-\dfrac{1}{2} & \dfrac{1}{2}
%         \end{array} 
%         \right]=
%         \dfrac{1}{2}.
%     \]
%     Si,
%     \begin{align*}
%         (x,y)=(0,0)\\
%         (x,y)=(2,0)\\
%         (x,y)=(0,2)
%     \end{align*}
%     para $(u,v)$ se tiene,
%     \begin{align*}
%         (u,v)=(1,-2)\\
%         (u,v)=(3,0)\\
%         (u,v)=(-1,0).
%     \end{align*}
%     Resolviendo la integral doble
%     \[
%         \int\int_R e \ ^{\frac{x-y+1}{x+y-2}} dx dy=
%         \dfrac{1}{2} \int_{-2}^{0}\int_{-v-1}^{v+3} e^{\frac{u}{v}}\  du dv=
%         +\infty
%     \]
 \end{respuesta}

%%%%%%%%%%%%%%%%%%%%%%%%%%%%%%%%%%%%%%%%%%%
%%%%%%%%%%%%%%%%%%%%%%%%%%%%%%%%%%%%%%%%%%%
%%%%%%%%%%%%%%%%%%%%%%%%%%%%%%%%%%%%%%%%%%%
% Tema 12: Cambio de variable
\item Evalúe la integral
	\[
        \iint_D (x+y)\,dA
    \]
    Siendo $D$ la región que se encuentra fuera del círculo de ecuación $x^2+y^2=1$ y dentro del cardioide de ecuación $r=1 + \cos \theta$.

\begin{respuesta}
    Las ecuaciones de las curvas dadas deben estar en el mismo sistema de coordenadas. Si, 
 	\[
        x=r \cos \theta 
        \yds 
        y=r \sen \theta
    \]
    La ecuación del círculo será $r=1$.
    Es decir las 2 funciones quedan definidas por:
    \[
        P(r,\theta)=(r \cos \theta,r \sen \theta),
    \]
    para  $(r,\theta) \in \R^2$, con lo cual
    \[
        \det(J_P(r,\theta))=r
    \]
    para  $(r,\theta) \in \R^2$, ahora encontramos $D^*$ tal que $T(D^*)=D$, vemos que:
    \[
        D^*=\l[\frac{3 \pi}{2}, \frac{5 \pi}{2}\r] \times \l[1, 1+\cos \theta \r]
    \]
    para  $(x,y) \in \R^2$, tenemos que:
	\[
 	\begin{aligned}
      \iint_D f(x,y)dydx &=\iint_{D^*} f(P(r,\theta)) | det(J_P(r,\theta))  | drd \theta\\
      &=\iint_{D^*} (r \cos \theta + r \sen \theta)r drd \theta \\
      &= \int_{\frac{3 \pi}{2}}^{\frac{5 \pi}{2}} \int_{1}^{1+ \cos \theta} (r \cos \theta + r \sen \theta)r drd \theta\\
       &=\dfrac{15 \pi+32}{24}
     \end{aligned}
    \]
\end{respuesta}

%%%%%%%%%%%%%%%%%%%%%%%%%%%%%%%%%%%%%%%%%%%
%%%%%%%%%%%%%%%%%%%%%%%%%%%%%%%%%%%%%%%%%%%
%%%%%%%%%%%%%%%%%%%%%%%%%%%%%%%%%%%%%%%%%%%
% Tema 12: Cambio de variable
\item 
    Use coordenadas polares para expresar
	\[
		\int_{1/\sqrt{2}}^1 \int_{\sqrt{1-x^2}}^x xy \,dy \,dx +
		\int_{1}^{\sqrt{2}} \int_{0}^x xy \,dy \,dx +
		\int_{\sqrt{2}}^2 \int_{0}^{\sqrt{4-x^2}} xy \,dy \,dx
	\]
	en una integral doble y luego halle su valor.

\begin{respuesta}
	Como se puede observar en el dibujo. Se está integrando sobre el área rayada. 
    \begin{center}
		\includegraphics[scale=0.6]{Graficos/fig10:03.eps}
	\end{center}

	Así se tiene que la suma de las integrales anteriores coincide con
	\[
		\int_{0}^{\frac{\pi}{4}} \int_1^2 (r^2 \sen(t)\cos(t)) r \,dr\, dt.
	\]
	Además, se tiene que
	\[
		\int_{0}^{\frac{\pi}{4}} \int_1^2  r^3 \dfrac{\sen(2t)}2 \,dr\, dt=\dfrac{15}{16}.
	\]
	El valor de la suma de las integrales es
	\[
		\int_{1/\sqrt{2}}^1 \int_{\sqrt{1-x^2}}^x xy\, dy\, dx +
		\int_{1}^{\sqrt{2}} \int_{0}^x xy \,dy\, dx +
		\int_{\sqrt{2}}^2 \int_{0}^{\sqrt{4-x^2}} xy \,dy\, dx
		=\dfrac{15}{16}.\qedhere
	\]
\end{respuesta}

%%%%%%%%%%%%%%%%%%%%%%%%%%%%%%%%%%%%%%%%%%%
%%%%%%%%%%%%%%%%%%%%%%%%%%%%%%%%%%%%%%%%%%%
%%%%%%%%%%%%%%%%%%%%%%%%%%%%%%%%%%%%%%%%%%%
% Tema 13: Integrales triples
\item
    Un sólido está limitado por debajo mediante las superficies de ecuación $z=3y^2$, por arriba mediante el plano de ecuación $z=3$ y en los extremos por los planos de ecuaciones  $x=-1$ y $x=2$.
\begin{enumerate}
    \item 
        Encuentre el centro de masas de este sólido, si este es un sólido homogéneo.
    \item 
        Ahora suponga que la densidad del sólido es dada por la función definida por $\delta(x,y,z)=z+x^2$ para $(x,y,z)\in\R^3$. Encuentre el centro de masa del sólido 
\end{enumerate}

\begin{respuesta} 
	Primero visualizaremos el sólido:
	\begin{center}
        \includegraphics[width=8cm]{Graficos/fig10:06.eps}
    \end{center}
	\begin{itemize}
	    \item Si $z=3$, entonces \\
	     $-1\leq y \leq 1$
	    \item Para el caso de $x$\\
	     $-1\leq x \leq 2$
	    \item Para $z$\\
	     $ 3y^2\leq z \leq 3$
	\end{itemize}
	Con esto podemos definir la región en el espacio, de la siguiente manera
	\[
	w=\{ (x,y,z)\in \R^3: 3y^2\leq z \leq3, -1\leq y \leq 1, -1\leq x \leq 2\}
	        \]
	El centro de la masas $(\overline{x}, \overline{y}, \overline{z})$ está dado por
	\[
	\overline{x}=\frac{\iiint_w x\,\delta(x,y,z)\,dx\,dy\,dz}{\iiint_w \,\delta(x,y,z)\,dx\,dy\,dz}
	        \]
	\[
	\overline{y}=\frac{\iiint_w y\,\delta(x,y,z)\,dx\,dy\,dz}{\iiint_w \,\delta(x,y,z)\,dx\,dy\,dz}
	        \]        
	\[
	\overline{z}=\frac{\iiint_w z\,\delta(x,y,z)\,dx\,dy\,dz}{\iiint_w \,\delta(x,y,z)\,dx\,dy\,dz}
	        \]  
	donde $\delta(x,y,z)$ representa la desidad del sólido en el punto $(x,y,z)$.\\
	Para el literal (a), esta densidad $\delta(x,y,z)=1$(cuerpo homogéneo), entonces primero evaluaremos 
	
	\begin{align*}
	    \iiint_w\,\delta(x,y,z)\,dx\,dy\,dz &= \iiint_w\,dx\,dy\,dz\\
	    &=\int_{-1}^2\int_{-1}^1\int_{3y^2}^3\,dz\,dy\,dz\\
	    &=12
	\end{align*}
	a continuación evaluaremos las integrales 
	\begin{itemize}
	    \item [i)]
        \begin{align*}
    	    \iiint_w x\,\delta(x,y,z)\,dx\,dy\,dz &= \iiint_w x\,dx\,dy\,dz\\
    	    &=\int_{-1}^2\int_{-1}^1\int_{3y^2}^3x\,dz\,dy\,dx\\
    	    &=6
    	\end{align*}
    	\item [ii)]
        \begin{align*}
    	    \iiint_w y\,\delta(x,y,z)\,dx\,dy\,dz &= \iiint_w y\,dx\,dy\,dz\\
    	    &=\int_{-1}^2\int_{-1}^1\int_{3y^2}^3 y\,dz\,dy\,dx\\
    	    &=0
    	\end{align*}
    	\item [iii)]
        \begin{align*}
    	    \iiint_w z\,\delta(x,y,z)\,dx\,dy\,dz &= \iiint_w z\,dx\,dy\,dz\\
    	    &=\int_{-1}^2\int_{-1}^1\int_{3y^2}^3 z\,dz\,dy\,dx\\
    	    &=\frac{108}{5}
    	\end{align*}
    \end{itemize}
	con esto podemos concluir que
	\[
	(\overline{x}, \overline{y}, \overline{z})=\left(\frac{1}{2}, 0, \frac{9}{5}\right)
	        \]
    Para el literal (b), la densidad $\delta(x,y,z)=z+x^2$ (cuerpo no homogéneo), entonces	 
    	\begin{align*}
	    \iiint_w\,\delta(x,y,z)\,dx\,dy\,dz &= \iiint_w(z+x^2)\,dx\,dy\,dz\\
	    &=\int_{-1}^2\int_{-1}^1\int_{3y^2}^3 (z+x^2)\,dz\,dy\,dz\\
	    &=\frac{168}{5}
	\end{align*}
	para el caso de las otras integrales  
	\begin{itemize}
	    \item [i)]
        \begin{align*}
    	    \iiint_w x\,\delta(x,y,z)\,dx\,dy\,dz &=\int_{-1}^2\int_{-1}^1\int_{3y^2}^3x(z+x^2)\,dz\,dy\,dx\\
    	    &=\frac{129}{5}
    	\end{align*}
    	\item [ii)]
        \begin{align*}
    	    \iiint_w y\,d(x,y,z)\,dx\,dy\,dz&=\int_{-1}^2\int_{-1}^1\int_{3y^2}^3 y(z+x^2)\,dz\,dy\,dx\\
    	    &=0
    	\end{align*}
    	\item [iii)]
        \begin{align*}
    	    \iiint_w z\,\delta(x,y,z)\,dx\,dy\,dz&=\int_{-1}^2\int_{-1}^1\int_{3y^2}^3 z(z+x^2)\,dz\,dy\,dx\\
    	    &=\frac{2376}{35}
    	\end{align*}
    \end{itemize}
    con esto podemos concluir que
	\[
        (\overline{x}, \overline{y}, \overline{z})=\left(\frac{43}{56}, 0, \frac{99}{49}\right)
    \]
\end{respuesta}

%%%%%%%%%%%%%%%%%%%%%%%%%%%%%%%%%%%%%%%%%%%
%%%%%%%%%%%%%%%%%%%%%%%%%%%%%%%%%%%%%%%%%%%
%%%%%%%%%%%%%%%%%%%%%%%%%%%%%%%%%%%%%%%%%%%
% Tema 11: Integrales dobles
\item
    Sea $D$ una región cuya frontera esta constituida por las curvas de ecuaciones $y=\dfrac{1}{x}$, $y=\dfrac{x}{4}$ y $y=\sqrt[]{x}$, donde $x$ y $y$ están en metros. Determinar el área $A$ de la región $D$ considerándola como una región de tipo I y después de tipo II.

\begin{respuesta}
    \begin{center}
        \includegraphics[scale=.3]{Graficos/IM2.jpeg}
    \end{center}

    Primero se encuentra la intersección de las curvas para dividir la región $D$ en dos regiones de tipo I. 
    Estos puntos son (0,0), (1,1) y $\l(2,\dfrac{1}{2}\r)$.
    
    Se definen entonces las regiones
    \[
        D_1=\lbrace (x,y)\in \R^2: 0\leqslant x\leqslant 1,f(x)\leqslant y\leqslant g(x) \rbrace
    \]
    y
    \[
        D_2=\lbrace (x,y)\in \R^2: 1\leqslant x\leqslant 2,h(x)\leqslant y\leqslant k(x) \rbrace
    \]
    donde 
    \[
        \funcion{f}{[0,1]}{\R}{x}{\frac{x}{4}},
    \]
    \[
        \funcion{g}{[0,1]}{\R}{x}{x^\frac{1}{2}}, 
    \]
    \[
        \funcion{h}{[1,2]}{\R}{x}{\frac{x}{4}},
    \]
    y
    \[
        \funcion{k}{[1,2]}{\R}{x}{\dfrac{1}{x}.}
    \]
    El área de la región $D$ es igual a la suma de las áreas de las regiones $D_1$ y $D_2$. En general, el área de una región $R$ viene dada por $\iint_R dA$.
    
    Entonces
    \begin{align*}
        A&=\iint_{D_1} dA + \iint_{D_2} dA\\
            &=\int_0^1\left(\int_{f(x)}^{g(x)}dy \right)dx +\int_1^2\left(\int_{h(x)}^{k(x)}dy \right)dx\\
            &=\int_0^1\left[x^{\frac{1}{2}}-\frac{x}{4} \right]dx + \int_1^2 \left[\frac{1}{x}-\frac{x}{4} \right]dx
    \end{align*}
    Finalmente, al resolver la última expresión se encuentra que el área de la región $D$ es $0.86m^2$.

    Ahora, para resolver este problema con la región $D$ como una región de tipo II, primero se encuentra la intersección de las curvas para dividir la región $D$ en dos regiones de tipo II. 
    Estos puntos son (0,0), (1,1) y $\l(2,\dfrac{1}{2}\r)$.
    
    Se definen entonces las regiones
    \[
        D_1=\lbrace (x,y)\in \R^2: 0\leqslant y\leqslant \frac{1}{2},f(y)\leqslant x\leqslant g(y) \rbrace
    \]
    y
    \[
        D_2=\lbrace (x,y)\in \R^2: \frac{1}{2} \leqslant y\leqslant 1,h(y)\leqslant x\leqslant k(y) \rbrace
    \]
    donde 
    \[
        \funcion{f}{[0,\frac{1}{2}]}{\R}{y}{y^2},
    \]
    \[
        \funcion{g}{[0,\frac{1}{2}]}{\R}{y}{4y}, 
    \]
    \[
        \funcion{h}{[\frac{1}{2},1]}{\R}{y}{y^2},
    \]
    y
    \[
        \funcion{k}{[\frac{1}{2},1]}{\R}{y}{\frac{1}{y}}.
    \]
    El área de la región $D$ es igual a la suma de las áreas de las regiones $D_1$ y $D_2$. En general, el área de una región $R$ viene dada por $\iint_R dA$.
    
    Entonces
    \begin{align*}
        A&=\iint_{D_1} dA + \iint_{D_2} dA\\
            &=\int_0^{\frac{1}{2}}\left(\int_{f(y)}^{g(y)}dx \right)dy +\int_{\frac{1}{2}}^1\left(\int_{h(y)}^{k(y)}dx \right)dy\\
            &=\int_0^\frac{1}{2}\left[4y-y^2 \right]dx + \int_\frac{1}{2}^1 \left[\frac{1}{y}-y^2 \right]dx
    \end{align*}
    Finalmente, al resolver la última expresión se encuentra que el área de la región $D$ es $0.86m^2$.
\end{respuesta}

%%%%%%%%%%%%%%%%%%%%%%%%%%%%%%%%%%%%%%%%%%%
%%%%%%%%%%%%%%%%%%%%%%%%%%%%%%%%%%%%%%%%%%%
%%%%%%%%%%%%%%%%%%%%%%%%%%%%%%%%%%%%%%%%%%%
% Tema 12: Cambio de variable
\item
    Calcular el volumen del sólido limitado superiormente por el paraboloide de ecuación $z=1-x^2-y^2$ e inferiormente por el plano de ecuación $z=1-y$.

\begin{respuesta}
    \begin{center}
        \includegraphics[scale=.4]{Graficos/IM1.jpeg}
    \end{center}
    Definimos las funciones:
    \[
        \funcion{f_1}{\R^2}{\R}{(x,y)}{1-x^2-y^2} 
    \]
    y
    \[
        \funcion{f_2}{\R^2}{\R}{(x,y)}{1-y} 
    \]
      
    Identificamos $D$, que es la región que se proyecta sobre el plano $xy$ producto de la intersección del paraboloide con el plano, ésta es el círculo de ecuación:
    \begin{equation}\label{regD}
        x^2+\l(y-\dfrac{1}{2}\r)^2=\dfrac{1}{4}.
    \end{equation}
    Se realiza el cambio de variable a coordenadas cilíndricas, es decir, el cambio de variable    
    \[
        \funcion{T}{\Omega}{D}{(r,\theta,z)}{(r \cos \theta,r \sen \theta,z),}
    \]
    nótese que 
    \[
        f_1(r,\theta)=1-r^2
    \]
    y
    \[
        f_2(r,\theta)=1-r\sen \theta,
    \]
    además, si se aplica el cambio de variable en (\ref{regD}) se obtiene    $r=\sen \theta$. 
    
    De lo anterior
    \[
        \Omega=\l\{(r,\theta,z)\in\R^3:0\leq \theta \leq \pi, 0\leq r \leq \sen \theta, 1-r\sen \theta \leq z \leq 1-r^2 \r\}
    \]
    y
    \[
        \det(J_T(r,\theta,z))=r
    \]
    para  $(r,\theta,z) \in \Omega$. Entonces
    
	\[
    \begin{aligned}
      \iiint_D dV&=\iiint_\Omega \det(J_T(r,\theta,z)) dzdrd\theta\\
      &=\int_0^\pi \int_0^{\sen\theta} \int_{f_2(r,\theta)}^{f_1(r,\theta)}rdzdrd\theta\\
      &=\int_0^\pi \int_0^{\sen\theta}r(1-r^2-1+r\sen \theta )drd \theta\\
      &= \int_{0}^{\pi} \dfrac{\sen^4 \theta}{12}d \theta\\
      &=\dfrac{ \pi}{32}.
     \end{aligned}
    \]
    
    De manera similar, se puede realizar el cambio de variable
    \[
        \funcion{T}{\Omega}{D}{(r,\theta,z)}{\l(r \cos \theta,r \sen \theta+\dfrac{1}{2},z\r)}
    \]
    donde
    \[
        \Omega=\l\{(r,\theta,z)\in\R^3:0\leq \theta \leq 2\pi, 0\leq r \leq \dfrac{1}{2}, \dfrac{1}{2}-r\sen \theta \leq z \leq \dfrac{3}{4}-r^2-r\sen \theta \r\}
    \]
    y
    \[
        \det(J_T(r,\theta,z))=r
    \]
    para  $(r,\theta,z) \in \Omega$. Se sugiere al estudiante calcular el volumen utilizando este cambio de variable y comparar con el resultado anterior.   
    
\end{respuesta}

%%%%%%%%%%%%%%%%%%%%%%%%%%%%%%%%%%%%%%%%%%%
%%%%%%%%%%%%%%%%%%%%%%%%%%%%%%%%%%%%%%%%%%%
%%%%%%%%%%%%%%%%%%%%%%%%%%%%%%%%%%%%%%%%%%%
% Tema 12: Cambio de variable
\item
	Calcular el área de la región limitada por la elipse de ecuación
    \[
        \dfrac{x^2}{a^2}+\dfrac{y^2}{b^2}=1.
    \]
    
\begin{respuesta}
    % La elipse es canónica, es decir, su centro se encuentra en el centro del sistema coordenado. Se encuentra los puntos de intersección con los ejes. Si, 
    % \begin{align*}
    % x=&0,\\
    % y=&\pm b,\\
    % \end{align*}
    % si 
    % \begin{align*}
    % y=&0,\\
    % x=&\pm a.\\
    % \end{align*}
    % El área de integración es
    % \[
    % R=\left\lbrace (x,y)\in\R^2:\dfrac{x^2}{a^2} + \dfrac{y^2}{b^2} \ \leq 1 \right\rbrace.
    % \]
    % Entonces, el área A por definición
    % \begin{align*}
    % A=&\int \int_R dxdy\\
    % =&4 \int _{0}^{a} \int _{0}^{b\sqrt{1-\dfrac{x^2}{a^2}}} \ dydx\\
    % =&4 \int _{0}^{a} b\sqrt{1-\dfrac{x^2}{a^2}} \ dx,
    % \end{align*}
    % sea
    % \begin{align*}
    % x=&a \sen(t)\\
    % dx=&a \cos(t) dt, 
    % \end{align*}
    % para todo $t \in \left[0,\dfrac{\pi}{2}\right]$. Por lo que
    % \begin{align*}
    % A=&4 ab \int _{0}^{a} \cos(t) \cos(t) dt\\
    % =& 4 ab \dfrac{\pi}{4}\\
    % =& \pi ab
    % \end{align*}
    $A=\pi ab$ %respuesta
\end{respuesta} 

%%%%%%%%%%%%%%%%%%%%%%%%%%%%%%%%%%%%%%%%%%%
%%%%%%%%%%%%%%%%%%%%%%%%%%%%%%%%%%%%%%%%%%%
%%%%%%%%%%%%%%%%%%%%%%%%%%%%%%%%%%%%%%%%%%%
% Tema 12: Cambio de variable
\item 
    Sea $D$ el semi anillo limitado por los círculos centrados en el origen de radio 1 y 4 unidades; al lado derecho del eje $y$. Determine el área de $D$ utilizando integración doble y coordenadas polares.

\begin{respuesta}
%     Se utilizan coordenadas polares, entonces
%     \[
%         \funcion{T}{\Omega}{D}{(r,\theta)}{(r\cos \theta,r\sen \theta)}
%     \]
% 	donde
% 	\[
% 	    \Omega=\l\{(r,\theta)\in\R^2:-\dfrac{\pi}{2} \leq \theta \leq \dfrac{\pi}{2}, 1\leq r \leq 4 \r\}
% 	\]
% 	y
% 	\[
% 	    J_T(r,\theta)=r.
% 	\]
% 	Entonces
% 	\[
% 		A=\iint_D dA=\iint_\Omega J_T(r,\theta) dr d\theta
% 		=\int_{-\frac{\pi}{2}}^{\frac{\pi}{2}}\int_{1}^4 r \ dr dt
% 		=\pi \int_{1}^4 r dr
% 		=\dfrac{15\pi}{2}.
% 	\]
	Así, el área buscada es $\dfrac{15\pi}{2}$ unidades cuadradas.
\end{respuesta}

%%%%%%%%%%%%%%%%%%%%%%%%%%%%%%%%%%%%%%%%%%%
%%%%%%%%%%%%%%%%%%%%%%%%%%%%%%%%%%%%%%%%%%%
%%%%%%%%%%%%%%%%%%%%%%%%%%%%%%%%%%%%%%%%%%%
% Tema 11: Integrales dobles
\item
    Calcular el área de la región $R$ limitada por las funciones:
	\[
        \funcion{f}{\R}{\R}{x}
        {\sen(x)}
        \texty
        \funcion{g}{\R}{\R}{x}
        {\cos(x)}
    \]
    en el intervalo $\l[0, \pi \r]$.

\begin{respuesta}  
    \begin{center}
        \includegraphics[scale=.3]{Graficos/IM3.jpeg}
    \end{center}
    La región $R$ solicitada es tipo 3, y se la resolverá como tipo 1 por intervalos.
    La región $R \subseteq \R^2$ está compuesta por $R_1 \cup R_2$. El punto de intersección es $x=\dfrac{\pi}{4}$. Sean
    \[
    \begin{aligned}
      R_1&=\l\lbrace(x,y) \in \R^2: f(x) \leq y \leq g(x)\ , 0 \leq x \leq \dfrac{\pi}{4}\r\rbrace\\
      &=\l\lbrace(x,y) \in \R^2: \sen(x) \leq y \leq \cos(x)\ , 0 \leq x \leq \dfrac{\pi}{4}\r\rbrace
    \end{aligned}
    \]
    y
    \[
    \begin{aligned}
      R_2&=\l\lbrace(x,y) \in \R^2: g(x) \leq y \leq f(x)\ , \dfrac{\pi}{4} \leq x \leq \pi \r\rbrace\\
      &=\l\lbrace(x,y) \in \R^2: \cos(x) \leq y \leq \sen(x)\ , \dfrac{\pi}{4} \leq x \leq \pi \r\rbrace
    \end{aligned}
    \]
    El área se calcula con la integral:
    \[
    \begin{aligned}
      A&=\iint_R dA\\
      &=\iint_{R_1} dA_1 + \iint_{R_2} dA_2\\
      &=\int_{0}^{\frac{\pi}{4}} \int_{\sen(x)}^{\cos(x)} dydx +\int_{\frac{\pi}{4}}^{\pi} \int_{\cos(x)}^{\sen(x)} dydx\\
      &=\l((\sqrt{2}-1)+(\sqrt{2}+1)\r)\\
      &=2\sqrt{2}
    \end{aligned}
    \]
\end{respuesta}

%%%%%%%%%%%%%%%%%%%%%%%%%%%%%%%%%%%%%%%%%%%
%%%%%%%%%%%%%%%%%%%%%%%%%%%%%%%%%%%%%%%%%%%
%%%%%%%%%%%%%%%%%%%%%%%%%%%%%%%%%%%%%%%%%%%
% Tema 11: Integrales dobles
\item
    Calcular el volumen del sólido del primer octante limitado por el cilindro de ecuación $x^2+y^2=1$, el plano de ecuación $z=0$ y superiormente por el plano de ecuación $x+y+z=3$.

\begin{respuesta}
    \begin{center}
        \includegraphics[scale=.5]{Graficos/IM4.jpeg}
    \end{center}
    La región $R$ es la proyección sobre el plano $xy$, y corresponde al primer cuadrante del círculo de ecuación:
    \[
      x^2+y^2=1.
    \]
    La región R es tipo 3, y la acotamos como tipo 1:
	\[
      R=\l\lbrace(x,y) \in \R^2: 0 \leq y \leq \sqrt{1-x^2} , 0 \leq x \leq 1 \r\rbrace
    \]
    El volumen solicitado es:
	\[
      V=\iint_R f(x,y)dA      
    \]
    donde $f(x,y)$ es la superficie que limita superiormente al sólido que se proyecta desde las región $R$, es decir el plano $x+y+z=3$
	\[
    \begin{aligned}
      V&=\int_{0}^{1} \int_{0}^{\sqrt{1-x^2}} (3-x-y)dydx\\
      &\approx1,69u^3
    \end{aligned}
    \]
    
    Se puede calcular este volumen con cambio de variables a coordenadas cilíndricas
    \[
        \funcion{T}{\Omega}{\R^3}{(r,\theta,z)}{(r\cos\theta,r\sen\theta,z),}
    \]
    donde
    \[
        \Omega=\l\{(r,\theta,z)\in\R^3:0\leq r\leq 1,0\leq \theta \leq \dfrac{\pi}{2},0\leq z \leq 3-r\sen\theta-r\cos\theta \r\}
    \]
    y 
    \[
        J_T(r,\theta,z)=r
    \]
    entonces
    \begin{align}
        V=\iiint_\Omega rdzdrd\theta&= \int_0^{\dfrac{\pi}{2}}\int_0^1\int_0^{3-r\sen\theta-r\cos\theta}rdzdrd\theta\\
        &=\int_0^{\dfrac{\pi}{2}}\int_0^1 r(3-r\sen\theta-r\cos\theta)drd\theta\\
        &=\int_0^{\dfrac{\pi}{2}} \l( \dfrac{3}{2}-\dfrac{1}{3}(\sen\theta+\cos\theta) \r) d\theta\\
        &=\dfrac{3\pi}{4}-\dfrac{2}{3}
    \end{align}
    de aquí en volumen del sólido es $V\approx 1.69$

    
\end{respuesta}

%%%%%%%%%%%%%%%%%%%%%%%%%%%%%%%%%%%%%%%%%%%
%%%%%%%%%%%%%%%%%%%%%%%%%%%%%%%%%%%%%%%%%%%
%%%%%%%%%%%%%%%%%%%%%%%%%%%%%%%%%%%%%%%%%%%
% Tema 14: Cambio de variable en integrales triples
\item 
    Una piscina circular de diámetro 50 metros tiene una profundidad constante a lo largo de este a oeste (eje $x$) y crece linealmente de 2 metros al sur a 5 metros al finalizar el norte (eje $y$). Encuentre el volumen del agua que llenaría la piscina.

\begin{respuesta}
	La piscina se puede modelar como un cilindro limitado por dos planos. El cilindro está definido por la ecuación
	\[
		x^2+(y-25)^2=25^2
	\]
	y los planos están definidos por la ecuaciones
	\[
		z=-\dfrac{3y}{50} \texty z=2.
	\]
	En el siguiente gráfico se puede ver la proyección de la piscina en el plano $yz$

	\begin{center}
		\includegraphics[scale=0.65]{Graficos/fig10:02.eps}
	\end{center}
	
	Se utilizan coordenadas cilíndricas
	\[
	    \funcion{T}{\Omega}{\R^3}{(r,\theta,z)}{(r\cos\theta,r\sen\theta+25,z)}
	\]
	donde
	\[
	    \Omega=\l\{0\leq r\leq 25,0\leq \theta\leq 2\pi, f_1=-\dfrac{3(r\sen\theta+25)}{50}\leq z \leq 2 \r\}
	\]
	y
	\[
	    J_T((r,\theta,z))=r.
	\]
	Entonces el volumen de la piscina viene definido por
	\begin{align*}
	    V&=\iiint_\Omega J_T((r,\theta,z))dzdrd\theta=\int_0^{25} \int_0^{2\pi} \int_{f_1}^{2}rdzd\theta dr\\
	    &=\int_0^{25} \int_0^{2\pi} \l(\dfrac{7r}{2}+\dfrac{3r^2\sen\theta}{50} \r)d\theta dr\\
	    &=\int_0^{25}  7\pi r dr,
	\end{align*}
	de aquí el volumen de la piscina es $V\approx 6872.23$
\end{respuesta}


%%%%%%%%%%%%%%%%%%%%%%%%%%%%%%%%%%%%%%%%%%%
%%%%%%%%%%%%%%%%%%%%%%%%%%%%%%%%%%%%%%%%%%%
%%%%%%%%%%%%%%%%%%%%%%%%%%%%%%%%%%%%%%%%%%%
% Tema 11: Integrales dobles
\item 
    Encuentre el centro de masa de una lamina de forma de triángulo rectángulo isósceles si su densidad en cualquier punto es proporcional al cuadrado de la distancia del vértice opuesto de la hipotenusa.


\begin{respuesta}
	Sea $c>0$ la longitud del cateto del triángulo rectángulo. Sea la densidad
	\[
		\funcion{\rho}{\R^2}{\R}{(x,y)}{k(x^2+y^2)}
	\]
	donde $k>0$. En el siguiente gráfico se puede observar el triángulo rectángulo
	\begin{center}
		\includegraphics[scale=0.7]{Graficos/fig10:04.eps}
	\end{center}
    donde su región está definida por
    \[
        R=\l\{(x,y)\in\R^2:0\leq x \leq c,0\leq y \leq c-x \r\}
    \]
	y su masa es
	\[
		m=\int_0^c \int_{0}^{c-x}\rho(x,y) \ dy dx=
		k \int_0^c \int_{0}^{c-x}(x^2+y^2) \ dy dx= \dfrac{k c^4}{6}.
	\]
	Además, si el centro de masa es $(\overline{x},\overline{y})$, se tiene que
	\[
		\overline{x}
		=\dfrac{1}{m} \ \int_0^c \int_{0}^{c-x}x\rho(x,y) \ dy dx
		=\dfrac{k}{m} \ \int_0^c \int_{0}^{c-x} (x^3+xy^2) \ dy dx
		=\dfrac{k}{m} \dfrac{c^5}{15}
		=\dfrac{2c}{5}
	\]
	y
	\[
		\overline{y}
		=\dfrac{1}{m} \ \int_0^c \int_{0}^{c-x}y\rho(x,y) \ dy dx
		=\dfrac{k}{m} \ \int_0^c \int_{0}^{c-x} (yx^2+y^3) \ dy dx
		=\dfrac{k}{m} \dfrac{c^5}{15}
		=\dfrac{2c}{5}.
	\]
	Por lo tanto, el centro de masa del triángulo rectángulo isósceles de cateto $c$ cuyo vértice opuesto a la hipotenusa está en el origen es
	\[
		\l(\dfrac{2c}{5},\dfrac{2c}{5}\r).\qedhere
	\]
\end{respuesta}

\item Calcular el volumen del casquete esférico limitado por las superficies de ecuaciones:
\[
    \begin{array}{cccc}
        x^2 + y^2 + z^2 = a^2, &  x^2 + y^2 + z^2 = b^2 & \text{y} & x^2 + y^2 = z^2,\\
    \end{array}
\]
con $z\geq0$, donde $a$ y $b$ son números reales tales que $0 < a < b$.

\begin{respuesta}\hspace{0pt}

Notemos que la representación gráfica del casquete esférico es
\begin{center}
	\includegraphics[scale=0.6]{Graficos/fig12_00.png}
\end{center}
Utilizando coordenadas esféricas tenemos que para este volumen:
\[
    \begin{array}{ccc}
        \l\{
        \begin{array}{ccl}
        x&=& r\cos(\theta)\sen(\varphi),  \\
        y&=& r\sen(\theta)\sen(\varphi),  \\
        z&=& r\cos(\varphi),
    \end{array}
    \r. & 
    \text{donde} & 
    \l\{
    \begin{array}{ccccc}
        a&\leq&r&\leq&b,  \\
        0&\leq&\theta&\leq&2\pi,  \\
        0&\leq&\varphi&\leq&\pi/4.
    \end{array}
    \r.
    \end{array}
\]
Dado que el jacobiano de la transformación es $r^{2}\sen(\varphi)$, si $V$ representa el valor del volumen, se tiene que:
\begin{align*}
    V &= \int_{a}^{b}{\int_{0}^{2\pi}{\int_{0}^{\pi/4}{r^{2}\sen(\varphi)\,d\varphi}\,d\theta}\,dr}\\
      &= \int_{a}^{b}{r^{2}\,dr}\int_{0}^{2\pi}{\,d\theta}\int_{0}^{\pi/4}{\sen(\varphi)\,d\varphi}\\
      &= \l(\dfrac{b^{3} - a^{3}}{3}\r)(2\pi)\l(1 - \dfrac{\sqrt{2}}{2}\r)\\
      &= \dfrac{\pi}{3}(2 - 2\sqrt{2})(b^{3} - a^{3}).
\end{align*}
Luego, el volumen del casquete esférico es $\dfrac{\pi}{3}(2 - 2\sqrt{2})(b^{3} - a^{3})$ unidades cúbicas.\qedhere
\end{respuesta}
\item Calcular el momento de inercia alrededor del eje $z$ del sólido de densidad constante limitado por el elipsoide de ecuación $36x^{2} + 9y^{2} + 4z^{2} = 36$.

\begin{respuesta}\hspace{0pt}

Por hipótesis, la densidad $\rho$ es tal que $\rho(x,y,z)=k$ para cada $(x,y,z)\in\R^{3}$, donde $k\in\R^+$. Se define $\Omega=\l\{(x,y,z)\in\R^{3}\,:\,x^{2} + \dfrac{1}{4}y^{2} + \dfrac{1}{9}z^{2} = 1\r\}$, entonces vamos a determinar el valor de
\[
    M_{Z} = \iiint_{\Omega}{k(x^{2} + y^{2})\,dxdydz}
\]
Para determinar el valor de la integral anterior, utilizamos coordenadas esféricas, sean 
\[
    \l\{
        \begin{array}{ccl}
        x&=& r\cos(\theta)\sen(\varphi),  \\
        y&=& 2r\sen(\theta)\sen(\varphi),  \\
        z&=& 3r\cos(\varphi),
    \end{array}
    \r.
\]
a partir de lo cual se obtiene que
\[
    \l\{
    \begin{array}{ccccc}
        0&\leq&r&\leq&1,  \\
        0&\leq&\theta&\leq&2\pi,  \\
        0&\leq&\varphi&\leq&\pi.
    \end{array}
    \r.
\]
Se define entonces la transformación
\[
    \funcion{P}{[0, 1]\times[0, 2\pi]\times[0, \pi]}{\R^{3}}{(r, \theta, \varphi)}{(r\cos(\theta)\sen(\varphi), 2r\sen(\theta)\sen(\varphi), 3r\cos(\varphi))}
\]
cuyo jacobiano es $J_{P}(t, \theta, \varphi) = -6r^{2}\sen(\varphi)$ para cada $(r, \theta, \varphi)\in[0, 1]\times[0, 2\pi]\times[0, \pi]$. Por lo tanto, se tiene que
\begin{align*}
    M_{z} &= \iiint_{\Omega}{k(x^{2} + y^{2})\,dxdydz}\\
          &= \int_{0}^{1}{\int_{0}^{2\pi}{\int_{0}^{\pi}{k(r^{2}\cos^{2}(\theta)\sen^{2}(\varphi) + 4r^{2}\sen^{2}(\theta)\sen^{2}(\varphi))(6r^{2}\sen^{2}(\varphi))\,d\varphi}\,d\theta}\,dr}\\
          &= 8k\pi.
\end{align*}
Por lo cual, el momento de inercia alrededor del eje $z$ de este sólido es $8k\pi$.\qedhere
\end{respuesta}

%%%%%%%%%%%%%%%%%%%%%%%%%%%%%%%%%%%%%%%%%%%
%%%%%%%%%%%%%%%%%%%%%%%%%%%%%%%%%%%%%%%%%%%
%%%%%%%%%%%%%%%%%%%%%%%%%%%%%%%%%%%%%%%%%%%
% Tema 15: Integrales de linea
\item Dados el campo escalar
    \[
    \funcion{f}{\R^3}{\R}
    {x}{=\|x\|^2}
	\]
    y la trayectoria
	\[
    \funcion{r}{[0,1]}{\R^3}
    {t}{(\cos{4t},\sen{4t},3t),}
	\]
    calcular $\int_{r}{f\,ds}$.

\begin{respuesta}
    Notemos que $r$ es derivable  con derivada
    \[
    r'(t) = (-4\sen{4t},4\cos{4t},3),
    \]
    con
    \[
    \|r'(t)\| = \sqrt{16\sen^{2}{4t}+16\cos^{2}{4t}+9} = \sqrt{16+9}=5,
    \]
    para todo $t\in[0,1]$, además, 
    \[
    f(r(t)) = f(\cos{4t},\sen{4t},3t) = \l(\sqrt{\cos^2{4t}+\sen^2{4t}+9t^2}\r)^2 = 1+9t^2.
    \]
    Entonces,
    \begin{align*}
    \int_{r}{f\,ds} &= \int_{0}^{1}{f(r(t))\|r'(t)\|\,dt}\\
    &=5\int_{0}^{1}{1+9t^2\, dt} = 20.
    \end{align*}
    Por lo tanto, $\int_{r}{f\,ds} = 20$.
\end{respuesta}

%%%%%%%%%%%%%%%%%%%%%%%%%%%%%%%%%%%%%%%%%%%
%%%%%%%%%%%%%%%%%%%%%%%%%%%%%%%%%%%%%%%%%%%
%%%%%%%%%%%%%%%%%%%%%%%%%%%%%%%%%%%%%%%%%%%
% Tema 15: Integrales de linea
\item
    Se tiene la trayectoria
    \[
        \funcion{\rho}{[0,\frac{7\pi}{2}]}{\R^3}{\theta}{(\cos^3 \theta, \sen^3 \theta, \theta)}
    \]
    y el campo de vectorial
    \[
        \funcion{F}{\R^3}{\R^3}{(x,y,z)}{(\sen z, \cos z, -(xy)^\frac{1}{3}).}
    \]
    Determinar la integral $\int_{\rho}F(x,y,z)ds$

\begin{respuesta}
    Se sabe que
    \[
    \int_{\rho}F(x,y,z)ds=\int_0^{\frac{7\pi}{2}}F(\rho(\theta))\cdot \rho'(\theta)d\theta
    \]
    donde
    \[
    \funcion{F\circ \rho}{[0,\frac{7\pi}{2}]}{\R^3}{\theta}{(\sen \theta,\cos \theta,-\sen \theta \cos \theta)}
    \]
    y 
    \[
    \rho'(\theta)=(-3\cos^2 \theta \sen \theta,3\sen^2 \theta \cos \theta,1).
    \]
    Entonces 
    \begin{align*}
    \int_{\rho}F(x,y,z)ds&=\int_0^{\frac{7\pi}{2}}(\sen \theta,\cos \theta,-\sen \theta \cos \theta)\cdot (-3\cos^2 \theta \sen \theta,3\sen^2 \theta \cos \theta,1) d\theta\\ 
    &=-\int_0^{\frac{7\pi}{2}}\sen \theta \cos \theta d\theta\\
    &=-\frac{1}{2}
    \end{align*}
\end{respuesta}

%%%%%%%%%%%%%%%%%%%%%%%%%%%%%%%%%%%%%%%%%%%
%%%%%%%%%%%%%%%%%%%%%%%%%%%%%%%%%%%%%%%%%%%
%%%%%%%%%%%%%%%%%%%%%%%%%%%%%%%%%%%%%%%%%%%
% Tema 15: Integrales de linea
\item
    Calcular la integral de lineal del campo escalar  
    \[
        \funcion{f}{\R^3}{\R}{(x,y,z)}{x+y+z.}
    \]
    sobre el segmento de recta que va desde $(1, 2, 3)$ hasta $(0, -1, 1)$. Luego, calcularlo en el sentido contrario.
    
\begin{respuesta}
    Una parametrización de la recta es
    \[
        r(t)=(1-t,2-3t,3-2t)
    \]
    para $0 \leq t \leq 1$. Además
    \begin{align*}
        r'(t)=&(-1,-3,-2),\\
        \| r'(t) \| =& \sqrt[]{14}.
    \end{align*}
   	Aplicando la defición, se tiene que
  	\begin{align*}
        \int_C f(x,y,z) ds = \int_0^{1} (6-6t) \sqrt[]{14}\ dt =3 \sqrt[]{14}
    \end{align*}
\end{respuesta}

%%%%%%%%%%%%%%%%%%%%%%%%%%%%%%%%%%%%%%%%%%%
%%%%%%%%%%%%%%%%%%%%%%%%%%%%%%%%%%%%%%%%%%%
%%%%%%%%%%%%%%%%%%%%%%%%%%%%%%%%%%%%%%%%%%%
% Tema 15: Integrales de linea
\item
	Calcular la integral del campo vectorial 
    \[
        \funcion{F}{\R^2}{\R^2}{(x,y)}{(3x+4y,2x+3y^2).}
    \]
    a lo largo de la trayectoria
    \[
        \funcion{r}{[0,2 \pi ]}{\R^2}{(t)}{(2 \cos (t), 2 \sen (t))}. 
    \]

% \begin{respuesta}
%     Se evalúa el campo vectorial en la trayectoria y se deriva la trayectoria, obtieniendo
%     \begin{align*}
%         r'(t)=&(-2 \sen(t),2 \cos(t))\\
%         F(r(t))=&(6 \cos(t)+ 8 \sen(t),4 \cos(t)+12\sen^2(t))
%     \end{align*}
%     para todo $t\in[0,2 \pi ]$.  Por definición de integral de línea, se tiene que
%     \begin{align*}
%         \int_C F\cdot ds &= \int_0^{2\pi} (6 \cos(t)+ 8 \sen(t),4 \cos(t)+12\sen^2(t)) \cdot (-2 \sen(t),2 \cos(t))\  dt\\
%         &\approx 46.15
%     \end{align*}
% \end{respuesta}

%%%%%%%%%%%%%%%%%%%%%%%%%%%%%%%%%%%%%%%%%%%
%%%%%%%%%%%%%%%%%%%%%%%%%%%%%%%%%%%%%%%%%%%
%%%%%%%%%%%%%%%%%%%%%%%%%%%%%%%%%%%%%%%%%%%
% Tema 15: Integrales de linea
\item Dados el campo escalar
    \[
        \funcion{f}{\Omega}{\R}
        {(x,y,z)}{\dfrac{z}{x^2 + y^2}}
	\]
	con $\Omega=\{(x,y,z)\in\R^3:(x,y)\neq (0,0)\}$
    y la trayectoria
	\[
        \funcion{r}{[0, 5]}{\R}
        {t}{(e^{2t}\cos{3t}, e^{2t}\sen{3t}, e^{2t}),}
	\]
    calcular $\displaystyle\int_{r}{f\, ds}$.
    
\begin{respuesta}
    Notemos que $r$ es derivable  con derivada
    \[
    r'(t) = (2e^{2t}\cos{3t} - 3e^{2t}\sen{3t}, 2e^{2t}\sen{3t} + 3e^{2t}\cos{3t}, 2e^{2t}),
    \]
    con
    \begin{align*}
    \|r'(t)\| &= \sqrt{13e^{4t}\cos^{2}{3t} + 13e^{4t}\sen^{2}{3t} + 4e^{4t}},\\
    					&= \sqrt{13e^{4t} + 4e^{4t}},\\
                        &= \sqrt{17e^{4t}},\\
                        &= e^{2t}\sqrt{17}.
    \end{align*}
    para todo $t\in[0, 5]$, además, 
    \begin{align*}
    f(r(t)) &= f(e^{2t}\cos{3t}, e^{2t}\sen{3t}, e^{2t}),\\
    		&= \dfrac{e^{2t}}{e^{4t}\cos^{2}{3t} + e^{4t}\sen^{2}{3t}},\\
            &= \dfrac{e^{2t}}{e^{4t}},\\
            &= e^{-2t}.
    \end{align*}
    Entonces,
    \begin{align*}
    \int_{0}^{5}{f(r(t))\|r'(t)\| \, dt} &= \sqrt{17}\int_{0}^{5}{e^{-2t} e^{2t} \, dt},\\
    		&= \sqrt{17}\int_{0}^{5}{dt},\\
            &= 5\sqrt{17}.
    \end{align*}
    Por lo tanto, $\int_{0}^{5}{f(r(t))\|r'(t)\| \, dt} = 5\sqrt{17}$.
\end{respuesta}

%%%%%%%%%%%%%%%%%%%%%%%%%%%%%%%%%%%%%%%%%%%
%%%%%%%%%%%%%%%%%%%%%%%%%%%%%%%%%%%%%%%%%%%
%%%%%%%%%%%%%%%%%%%%%%%%%%%%%%%%%%%%%%%%%%%
% Tema 15: Integrales de linea
\item Dado el campo vectorial
    \[
        \funcion{F}{\R^2}{\R^2}
        {(x,y)}{(-y, x)}
	\]
    y la función vectorial
    \[
        \funcion{r}{[0, \pi]}{\R}
        {t}{(\cos{3t}, \sen{3t}),}
	\]
    calcular $\displaystyle\int_{r}{F\, ds}$.
    
\begin{respuesta}
    Notemos que $r$ es derivable  con derivada
    \[
    r'(t) = (-3\sen{3t}, 3\cos{3t}),
    \]
    para todo $t\in[0, \pi]$, además, 
    \begin{align*}
    F(r(t)) &= F(\cos{3t}, \sen{3t}),\\
    		&= (-\sen{3t}, \cos{3t}).
    \end{align*}
    Entonces, 
    \begin{align*}
    \int_{0}^{\pi}{F(r(t))\cdot r'(t) \, dt} &=  \int_{0}^{\pi}{(-\sen{3t}, \cos{3t})\cdot (-3\sen{3t}, 3\cos{3t})\, dt},\\
    		&= \int_{0}^{\pi}{(3\sen^{2}{3t} + 3\cos^{2}{3t})\, dt},\\
            &= 3\pi.
    \end{align*}
    Por lo tanto, $\int_{0}^{\pi}{F(r(t))\cdot r'(t) \, dt} = 3\pi$.
\end{respuesta}
%%%%%%%%%%%%%%%%%%%%%%%%%%%%%%%%%%%%%%%%%%%
%%%%%%%%%%%%%%%%%%%%%%%%%%%%%%%%%%%%%%%%%%%
%%%%%%%%%%%%%%%%%%%%%%%%%%%%%%%%%%%%%%%%%%%
% Tema 15: Integrales de linea
\item
    Encuentre el centro de masa en $x$ de un alambre que tiene la forma de un triángulo isósceles que une los puntos $(-2,1)$, $(4,1)$ y $(1,10)$, donde su densidad está dada por $\delta (x,y)=x^2$, para $(x,y)\in\R^2$.

\begin{respuesta}
    El centro de masa está dado por:
	\[
    \begin{aligned}
     \bar{x}&=\dfrac{\int_C x \delta(x,y) ds}{m}\\
     \bar{y}&=\dfrac{\int_C y \delta(x,y) ds}{m}
    \end{aligned}
    \]
    donde la masa:
	\[
     m=\int_C \delta(x,y) ds.
    \]
    La curva $C$ está compuesta por 3 segmentos de recta de ecuaciones:
	\[
    \begin{aligned}
     C_1&: y=-3x+13\\
     C_2&: y=1\\
     C_3&: y=3x+7
    \end{aligned}
    \]
    sus respectivas parametrizaciones:
 	\[
    \begin{aligned}
    \funcion{\alpha_1}{[1,4]}{\R}{t}{(t,-3t+13)}&
    \funcion{\alpha_2}{[-2,4]}{\R}{t}{(t,1)} \texty
    \funcion{\alpha_3}{[-2,1]}{\R}{t}{(t,3t+7)} 
    \end{aligned}
    \]
    y sus derivadas:
 	\[
    \begin{aligned}
     C_1'(t)&=(1,-3)\\
     C_2'(t)&=(1,0)\\
     C_3'(t)&=(1,3).
    \end{aligned}
    \]
    La masa será entonces:
     \[
     \begin{aligned}
     m&=\int_{C_1} \delta(x,y) ds+\int_{C_2} \delta(x,y) ds+\int_{C_3} \delta(x,y) ds\\
      &=\int_{C_1} x^2 ds+\int_{C_2} x^2 ds+\int_{C_3} x^2 ds\\
      &=\int_{1}^{4} t^2\l(\sqrt{10}\r) dt+\int_{-2}^{4} t^2 dt+\int_{-2}^{1} t^2\l(\sqrt{10}\r) dt\\
      &=21\sqrt{10}+24+3\sqrt{10}\\
      &\approx 100
     \end{aligned}
    \]
    entonces el centro de masa en $x$ es: 
	\[
    \begin{aligned}
     \bar{x}&=\dfrac{\int_C x \delta(x,y) ds}{m}\\
     &=\dfrac{\int_{C_1} x \delta(x,y) ds+\int_{C_2} x \delta(x,y) ds+\int_{C_3} x \delta(x,y) ds}{m}\\
     &=\dfrac{\int_{C_1} x^3 ds+\int_{C_2} x^3 ds+\int_{C_3} x^3 ds}{m}
     \\
     &=\dfrac{\int_{1}^{4} t^3 \l(\sqrt{10}\r) dt+\int_{-2}^{4} t^3 dt+\int_{-2}^{1} t^3 \l(\sqrt{10}\r) dt}{m}\\
     &=\dfrac{63.75\sqrt{10}+60-3.75\sqrt{10}}{m}\\
     &\approx \dfrac{250}{100}\\
     &=2.5
    \end{aligned}
    \]
    y el centro de masa en $y$:
 
	 \[
    \begin{aligned}
     \bar{y}&=\dfrac{\int_C y \delta(x,y) ds}{m}\\
     &=\dfrac{\int_{C_1} y \delta(x,y) ds+\int_{C_2} y \delta(x,y) ds+\int_{C_3} y \delta(x,y) ds}{m}\\
     &=\dfrac{\int_{C_1} x^2y ds+\int_{C_2} x^2y ds+\int_{C_3} x^2y ds}{m}
     \\
     &=\dfrac{\int_{1}^{4} t^2(-3t+13) \l(\sqrt{10}\r) dt+\int_{-2}^{4} t^2 dt+\int_{-2}^{1} t^2(3t+7) \l(\sqrt{10}\r) dt}{m}\\
     &=\dfrac{81.75\sqrt{10}+24+9.75\sqrt{10}}{m}\\
     &=\dfrac{313}{100}\\
     &=3.13
    \end{aligned}
    \]
\end{respuesta}

%%%%%%%%%%%%%%%%%%%%%%%%%%%%%%%%%%%%%%%%%%%
%%%%%%%%%%%%%%%%%%%%%%%%%%%%%%%%%%%%%%%%%%%
%%%%%%%%%%%%%%%%%%%%%%%%%%%%%%%%%%%%%%%%%%%
% Tema 15: Integrales de linea
\item
    Se tiene una pared cuya altura está dada por
    \[
        \funcion{h}{\R^2}{\R}{(x,y)}{1+\frac{y}{3}}
    \]
    y su base por
    \[
        \funcion{\rho}{[0,\pi]}{\R^2}{t}{(30\cos^3(t),30\sen^3(t)).}
    \]
    Las medidas están en metros. Determinar el área de la pared.  
    
\begin{respuesta}
    En la siguiente figura se muestra la pared.
        
    %\includegraphics{Pared}
    
   Primero se calcula el área de la sección del primer cuadrante. 
   En este cuadrante la base viene dada por 
   \[
        \funcion{\rho}{[0,\frac{\pi}{2}]}{\R^2}{t}{(30\cos^3t,30\sen^3t).}
    \]
   Entonces el área de la mitad de la pared es
   \[
        \int_{\rho}f(x,y)ds=\int_0^{\frac{\pi}{2}}f(\rho(t))||\rho'(t)||dt,
   \]
   donde 
   \[
    \rho'(t)=(-90\cos^2t\sen t,90\sen^2t\cos t)
   \]
   y
   \[
   ||\rho'(t)||=90\sen t\cos t.
   \]
   Finalmente
   \begin{align*}
   \int_{\rho}f(x,y)ds&=\int_0^{\frac{\pi}{2}}\left(1+\frac{30\sen^3t}{3} \right)90\sen t\cos tdt\\
   &=90\int_0^\frac{\pi}{2}(\sen t +10\sen^4t)\cos tdt\\
   &=225
   \end{align*}
   y el área total de la pared es $450m^2$.
\end{respuesta}

%%%%%%%%%%%%%%%%%%%%%%%%%%%%%%%%%%%%%%%%%%%
%%%%%%%%%%%%%%%%%%%%%%%%%%%%%%%%%%%%%%%%%%%
%%%%%%%%%%%%%%%%%%%%%%%%%%%%%%%%%%%%%%%%%%%
% Tema 15: Integrales de linea
\item
    Calcule el trabajo efectuado por el campo de fuerzas definido por $f(x,y)=(y,x)$, para $(x,y)\in\R^2$, al mover una partícula desde $(0,0)$ hasta $(1,1)$ sobre cualquier trayectoria definida por $\alpha(t)=(t^n,t)$, para $t\in [0,1]$ y $n\in\N$.	
    
\begin{respuesta} 
    \[
        \alpha(t)=(t^n,t)\quad ;\quad 0\leq t \leq 1
    \]
    \[
        \alpha'(t)=(nt^{n-1},1)
    \]
    \[
        f(\alpha(t))=(t,t^n)
    \]
    Como $f(x,y)$ es un campo de fuerzas, entonces la integral de línea representa el trabajo realizado para mover la partícula
    \[
        w=\int_{\alpha_1}f\cdot ds=\int_0^1(t,t^n)\cdot(nt^{n-1},1)dt
    \]
    \[
        w=\int_0^1 nt^n+t^ndt=\left(\frac{nt^{n+1}}{n+1}+\frac{t^{n+1}}{n+1}\right)_0^1
    \]
    \[
        \frac{t^{n+1}{(n+1)}}{{n+1}}_0^1
    \]
    \[
        w=1
    \]
\end{respuesta}
%%%%%%%%%%%%%%%%%%%%%%%%%%%%%%%%%%%%%%%%%%% FC
%%%%%%%%%%%%%%%%%%%%%%%%%%%%%%%%%%%%%%%%%%%
\item Un campo de fuerza definido por
    \[
    \funcion{F}{\R^2}{\R^2}{(x,y)}{\l(2xy^2,2x^2y\r)}
    \]
Calcular el trabajo que realiza $F$ para mover una partícula a lo largo de la trayectoria de ecuación $x^2+y^2=9$ en sentido antihorario partiendo del punto $(3,0).$
\begin{enumerate}
    \item A lo largo de $\dfrac{1}{8}$ de vuelta.
    \item A lo largo de $\dfrac{1}{4}$ de vuelta.
\end{enumerate}
\begin{respuesta}
    La trayectoria es una circunferencia, y una de sus parametrizaciones es:
    \[
    \funcion{r}{[0,2\pi]}{\R^2}{(r,\theta)}{\l(3\cos (t), 3\sen (t)\r)}
    \]
    Se evalúa el campo vectorial en la trayectoria y se deriva la trayectoria, obtieniendo
    \begin{align*}
        r'(t)=&(-3 \sen(t),3 \cos(t))\\
        f(r(t))=&\l(54 \cos(t) \sen^2(t),54 \sen(t) \cos^2(t)\r)
    \end{align*}
    para todo $t\in[0,2 \pi ]$.
    \begin{enumerate}
        \item Por definición de integral de línea, se tiene que:
            \begin{align*}
            \int_C F\cdot ds &= \int_0^{\frac{\pi}{4}} \l(54 \cos(t) \sen^2(t),54 \sen(t) \cos^2(t)\r) \cdot \l(-3 \sen(t),3 \cos(t) \r)  dt\\
            &= \int_0^{\frac{\pi}{4}} \l(162 \cos^3(t) \sen(t)-162 \sen^3(t) \cos(t)\r) dt\\
            &=\dfrac{81}{4}\\
            &= 20.25
            \end{align*}
          \item La integral de línea es:
            \begin{align*}
            \int_C F\cdot ds &= \int_0^{\frac{\pi}{2}} \l(54 \cos(t) \sen^2(t),54 \sen(t) \cos^2(t)\r) \cdot \l(-3 \sen(t),3 \cos(t) \r)  dt\\
            &= \int_0^{\frac{\pi}{2}} \l(162 \cos^3(t) \sen(t)-162 \sen^3(t) \cos(t)\r) dt\\
            &=0
            \end{align*}  
    \end{enumerate}
\end{respuesta}
%%%%%%%%%%%%%%%%%%%%%%%%%%%%%%%%%%%%%%%%%%%%%%%%%%%%
%%%%%%%%%%%%%%%%%%%%%%%%%%%%%%%%%%%%%%%%%%%%%%%%%%%%
\item
Encuentre una parametrización de cada una de las siguientes curvas que resultan de la intersección de las superficies de ecuación
\begin{enumerate}
    \item $x^2 + z^2 = 1$ , $x + y = 1$.
    \item $x^2 + y^2 = 9$, $z = 0$.
    \item $y = 3x^2$, $z = 1$.
    \item $(x - 1)^2 + 4(y - 2)^2 = 4$, $z = -1$.
    \item $x^2 + y^2 + z^2 = a(x + 1)$, $x + y = a$, donde $a>0$.
\end{enumerate}

\begin{respuesta}\hspace{0pt}
    
\begin{enumerate}
    \item De reemplazar $x = 1 - y$ en $x^2 + z^2 = 1$ se obtiene 
    \[
        (y - 1)^2 + z^2 = 1.
    \]
    Sean $y = \cos(t) + 1$ y $z = \sen(t)$, entonces $x = -\cos(t)$. Se define la trayectoria
    \[
        \funcion{\alpha}{[0, 2\pi]}{\R}{t}{(-\cos(t), \cos(t) + 1, \sen(t)).}
    \]
    \item Sean $x = 3\cos(t)$ y $y = 3\sen(t)$, entonces se define
    \[
        \funcion{\alpha}{[0, 2\pi]}{\R}{t}{(3\cos(t), 3\sen(t), 0).}
    \]
    \item Sea $x = t$, se tiene entonces que $y = 3t^2$ y se define
    \[
        \funcion{\alpha}{\R}{\R}{t}{(t, 3t^2, 1).}
    \]
    \item Sean $x = 2\cos(t) + 1$ y $y = \sen(t) + 2$, entonces se define
    \[
        \funcion{\alpha}{[0, 2\pi]}{\R}{t}{(2\cos(t) + 1, \sen(t) + 2, -1).}
    \]
    \item De reemplazar $x + y = a$ en $x^2 + y^2 + z^2 = a(x + y)$ se obtiene
    \[
        2\l(y - \dfrac{a}{2}\r)^2 + z^2 = \dfrac{a^2}{2}.
    \]
    Sean $y = \dfrac{a}{2}(1 - \cos(t))$ y $z = \dfrac{a}{\sqrt{2}}\sen(t)$, junto con $x + y = a$ se obtiene que $x = \dfrac{a}{2}(1 + \cos(t))$, se define entonces
    \[
        \funcion{\alpha}{[0, 2\pi]}{\R}{t}{\l(\dfrac{a}{2}(1 + \cos(t)), \dfrac{a}{2}(1 - \cos(t)), \dfrac{a}{\sqrt{2}}\sen(t) \r).}\qedhere
    \]
\end{enumerate}
\end{respuesta}
%%%%%%%%%%%%%%%%%%%%%%%%%%%%%%%%%%%%%%%%%%%%%%%%%%%%
%%%%%%%%%%%%%%%%%%%%%%%%%%%%%%%%%%%%%%%%%%%%%%%%%%%%
\item
Considere la trayectoria
\[
    \funcion{\alpha}{\l[0, \dfrac{5}{4}\r]}{\R^3}{t}{\l(t, \dfrac{4}{3}t^{3/2}, \dfrac{1}{2}t\r)}
\]
que parametriza a una curva $\cal{C}$, halle la longitud de curva dada por la parametrización $\alpha$.

\begin{respuesta}\hspace{0pt}

% Determinemos la derivada de $\alpha$ la cual está dada por
% \[
%     \alpha'(t) = \l(1, 2t^{1/2}, \dfrac{1}{2}\r)
% \]
% con 
% \[
%     \| \alpha'(t) \| = \dfrac{1}{2}\sqrt{5 + 16t}
% \]
% para cada $t\in\l[0, \dfrac{5}{4}\r]$. Luego, la longitud de curva $L$ es
% \begin{align*}
%     L &= \int_{0}^{5/4}{\| \alpha'(t) \|\, dt}\\
%       &= \int_{0}^{5/4}{\dfrac{1}{2}\sqrt{5 + 16t}\, dt}\\
%       &= \dfrac{1}{48}\l(125 - 5\sqrt{5}\r).
% \end{align*}
Luego, la longitud de curva es $\dfrac{1}{48}\l(125 - 5\sqrt{5}\r)$.\qedhere
\end{respuesta}

%%%%%%%%%%%%%%%%%%%%%%%%%%%%%%%%%%%%%%%%%%%%%%%%%%%%
%%%%%%%%%%%%%%%%%%%%%%%%%%%%%%%%%%%%%%%%%%%%%%%%%%%%
\item
Considere la curva de ecuación $y^2 = x^3$. Halle la longitud de arco que al punto $(1, -1 )$ con $(1, 1)$.

\begin{respuesta}\hspace{0pt}
    
Notemos que la representación gráfica de la longitud recorrida es
\begin{center}
	\includegraphics[scale=2.3]{Graficos/fig13_00.png}
\end{center}
Se define la trayectoria
\[
    \funcion{\alpha}{[-1, 1]}{\R}{t}{(t^2, t^3)}
\]
que parametriza a la curva, además, se tiene que 
\[
     \begin{array}{cccc}
        \alpha(-1) = (1, -1),   &   \alpha(0) = (0, 0)  &   \text{y}  &   \alpha(1) = (1, 1).
     \end{array}
\]
Además, la derivada de $\alpha$ está dada por
\[
    \alpha'(t) = (2t, 3t^2)
\]
con 
\[
    \| \alpha'(t) \| = |t| \sqrt{4 + 9t^2}
\]
para cada $t\in[-1, 1]$. Finalmente, se tiene que
\begin{align*}
    L &= \int_{-1}^{1}{\| \alpha'(t) \|\, dt}\\
      &= \int_{-1}^{1}{|t|\sqrt{4 + 9t^2}\, dt}\\
      &= -\int_{-1}^{0}{t\sqrt{4 + 9t^2}\, dt} + \int_{0}^{1}{t\sqrt{4 + 9t^2}\, dt}\\
      &= \dfrac{1}{27}\l(26\sqrt{13} - 16\r).
\end{align*}
Por tanto, la longitud de arco $L$ es $\dfrac{1}{27}\l(26\sqrt{13} - 16\r)$.\qedhere
\end{respuesta}


%%%%%%%%%%%%%%%%%%%%%%%%%%%%%%%%%%%%%%%%%%%
%%%%%%%%%%%%%%%%%%%%%%%%%%%%%%%%%%%%%%%%%%%
%%%%%%%%%%%%%%%%%%%%%%%%%%%%%%%%%%%%%%%%%%%
% Tema 15: Integrales de linea
\item 
	Sean el campo vectorial
	\[
		\funcion{F}{\R^2}{\R^2}{x}{(x_2e^{x_1}+\sen(x_2),e^{x_1}+x_1\cos(x_2))}
	\]
	y la trayectoria
	\[
		\funcion{r}{\lcl0,\dfrac{\pi}2\rcl}{\R}{t}{(\cos(t),2\sen(t)).}
	\]
	Encuentre un potencial de $F$ y calcule el valor de
	\[
		\int_C F\cdot \ dr.
	\]

% \begin{respuesta}
% 	Por el ejercicio $1b$ de la hoja de ejercicios 8 se sabe que 
% 	\[
% 		\funcion{f}{\R^2}{\R}{x}{x_2e^{x_1}+x_1\sen(x_2) + k}
% 	\]
% 	es un potencial de $F$ donde $k\in\R$, entonces $F$ es un campo conservativo y se tiene que
% 	\begin{align*}
% 		\int_C F\cdot \ dr
% 			& = \int_C \nabla f \cdot \ dr\\
% 			&= f\l(r\l(\dfrac{\pi}{2}\r)\r)-f\l(r\l(0\r)\r)\\
% 			&= f\l(0,2\r)-f\l(1,0\r)\\
% 			&= (2+k)-k\\
% 			&= 2.\qedhere
% 	\end{align*}
% \end{respuesta}

%%%%%%%%%%%%%%%%%%%%%%%%%%%%%%%%%%%%%%%%%%%
%%%%%%%%%%%%%%%%%%%%%%%%%%%%%%%%%%%%%%%%%%%
%%%%%%%%%%%%%%%%%%%%%%%%%%%%%%%%%%%%%%%%%%%
% Tema 15: Integrales de linea
\item
    Dado el campo vectorial
    \[
        \funcion{F}{\R^2}{\R^2}{(x,y)}{(2xy^3,3x^2 y^2)}
    \]
    Determine:
    \begin{enumerate}
        \item El campo escalar $f:\R^2\rightarrow \R$, tal que $F(x,y)=\nabla f(x,y)$ para $(x,y)\in\R^2$.
        \item La integral $\int_\alpha F\cdot ds$ a lo larga de la curva parametrizada por $\alpha:[0,1]\rightarrow \R^2$, con $\alpha(t)=(t,t^2+1)$ para $t\in[0,1]$.
    \end{enumerate}

\begin{respuesta}
\begin{enumerate}
	\item Si 
	\[
	    F(x)=\nabla f(x)
	\]
	entonces
	
    \begin{align*}
        \frac{df}{dx}(x,y)&=2xy^3\\
        \frac{df}{dy}(x,y)&=3x^2y^2,
    \end{align*}
    se integra la primera ecuación con respecto a $x$
    \[
        f(x,y)=\int2xy^3dx
    \]
    \[
        f(x,y)=x^2y^3+h(y)
    \]
    para determinar $h(y)$, se deriva la función calculada con respecto a $y$
    \[
        \frac{df}{dy}(x,y)=3x^2y^2+h'(y)    
    \]
    y se resuelve
    \[
        3x^2y^2=3x^2y^2+h'(y)
    \]
    de donde
    \[
    h'(y)=0 \rightarrow h(y)=k;\quad k \in \R
    \]
    con lo cual la función potencial de $F$ es
    \[
        f(x,y)=x^2y^3+k.
    \]
    \item Para evaluar la $\int_\alpha F\cdot ds$ a lo largo de la curva $\alpha(t)=(t,t^2+1); 0\leq t \leq 1$, vamos a aplicar el segundo teorema fundamental para integrales de línea
    \[
        \int_\alpha F\cdot ds=\int_\alpha \nabla f\cdot ds=f(b)-f(a)
    \]
    donde $f$, potencial de $F$, es
    \[
        f(x,y)=x^2y^3+k.
    \]
    Por otro lado, siendo la ecuación paramétrica de la curva $\alpha(t)=(t,t^2+1)$ para  $0\leq t \leq 1$, determinamos el punto inicial $a$ y el final $b$
    \[
    a=\alpha(0)=(0,1) \quad; \quad b=\alpha(1)=(1,2)
    \]
    entonces
    \[
        \int_\alpha F\cdot ds=\int_\alpha \nabla f\cdot ds=f(1,2)-f(0,1)=8.
    \]
    
    Otra alternativa para el cálculo de esta integral es utilizar la definición de integral de línea de un campo vectorial a lo largo de una curva, entonces
    \[
        \alpha'(t)=(1,2t)
    \]
    \[
        F(\alpha(t))=(2t(t^2+1)^3,3t^2(t^2+1)^2)
    \]
    por tanto
    \begin{align*}
        \int_\alpha F\cdot ds&=\int_0^1(2t(t^2+1)^3,3t^2(t^2+1)^2)\cdot(1,2t)dt\\
        &=\int_0^1(2t(t^2+1)^3+6t^3(t^2+1)^2)dt\\
        &=8
    \end{align*}
\end{enumerate}
\end{respuesta}

%%%%%%%%%%%%%%%%%%%%%%%%%%%%%%%%%%%%%%%%%%%
%%%%%%%%%%%%%%%%%%%%%%%%%%%%%%%%%%%%%%%%%%%
%%%%%%%%%%%%%%%%%%%%%%%%%%%%%%%%%%%%%%%%%%%
% Tema 15: Integrales de linea
\item 
    Dado el campo vectorial:
		\[
        	\funcion{F}{\R^2}{\R^2}{(x,y)}{(2xy,x^2-y),} 
        \]
    verifique que es un campo conservativo y evalúe la integral
	\[
     \int_\alpha F\cdot ds
    \]
    donde $\alpha$ describe cualquier camino que une los puntos $a=(1,4)$ y $b=\l(5,\frac{1}{2}\r)$.

\begin{respuesta}
%     El $\rot F=0$, entonces se verifica que el campo vectorial $F$ es un campo conservativo y existe entonces un campo escalar $f$ tal que
%     \[
%         \nabla f=F,
%     \]
%     de donde
%     \[
%         \dfrac{\partial f}{\partial x}=2xy
%     \]
%     entonces
% 	\[
%     \begin{aligned}
%         f&=\int 2xy dx = x^2y+c(y) \\
%         \dfrac{\partial f}{\partial y}&=x^2+c'(y)
%     \end{aligned}
%     \]
%     además
%     \[
%          \dfrac{\partial f}{\partial y}=x^2-y
%     \]
%     entonces
%     \[
%         c(y)=-\dfrac{y^2}{2}+k
%     \]
%     y
%     \[
%         f(x,y)=x^2y-\dfrac{y^2}{2}+k,
%     \]
%     con $k\in\R$.

% El campo $F$ es conservativo, por lo tanto la integral de línea no depende de la trayectoria, entonces
% 	\[
%      \int_\alpha F\cdot ds=f(b)-f(a)
%     \]
%     donde
% 	\[
%      f(b)=\dfrac{99}{8}
%     \]
% y
% 	\[
%      f(a)=-\dfrac{32}{8},
%     \]
%     finalmente
	\[
     \int_\alpha F\cdot ds=\dfrac{131}{8}.
    \]
\end{respuesta}



%%%%%%%%%%%%%%%%%%%%%%%%%%%%%%%%%%%%%%%%%%%
%%%%%%%%%%%%%%%%%%%%%%%%%%%%%%%%%%%%%%%%%%%
%%%%%%%%%%%%%%%%%%%%%%%%%%%%%%%%%%%%%%%%%%%
% Tema 15: Integrales de linea
\item 
    Dado el campo de fuerzas
    \[
        \funcion{F}{\R^2}{\R^2}
        {(x,y)}{(e^x\cos y,-e^x\sen y)}
	\]
    y $a_1=(0,0)$ y $a_2=(a,b)$ con $a>0$ y $b>0$.Verificar que $F$ es un campo conservativo y calcular el trabajo realizado para mover una partícula entre estos dos puntos. 
\begin{respuesta}
    Notemos que $\rot F=0$ y el campo escalar $f$ definido por
    \[
    \funcion{f}{\R^2}{\R}
    {(x,y)}{f(x,y)=e^x\cos y}
	\]
	es tal que 
    \[
    \nabla f(x,y) = F(x,y),
    \]
    es decir, es su potencial. Se sugiere al estudiante verificar este resultado. 
    Entonces el trabajo realizado viene dado por
    \[
        \int_\alpha Fds
    \]
    donde $\alpha$ describe cualquier curva que une los puntos $a_1$ y $a_2$.
    
    Finalmente, la integral de línea de $F$ es independiente de la trayectoria y
    \begin{align*}
    \int_{\alpha}{F\cdot ds} &=  f(a_2) - f(a_1)\\
    		&=  f(a,b) - f(0,0)\\
            &= e^a\cos b - 1.
    \end{align*}
    Por tanto, el trabajo realizado por el campo para mover la particula entre $a_1$ y $a_2$ es
    \[
    \int_{\alpha}{F\cdot ds} = e^a\cos b - 1.
    \]
\end{respuesta}

%%%%%%%%%%%%%%%%%%%%%%%%%%%%%%%%%%%%%%%%%%%
%%%%%%%%%%%%%%%%%%%%%%%%%%%%%%%%%%%%%%%%%%%
%%%%%%%%%%%%%%%%%%%%%%%%%%%%%%%%%%%%%%%%%%%
% Tema 15: Integrales de linea
\item Dado el campo vectorial
    \[
    \funcion{F}{\R^3}{\R^3}
    {(x,y,z)}{(x,y^2,-z^3)}
	\]
    y $a_1=(1,1,1)$ y $a_2=(2,3,-4)$.
    Verificar si $F$ es conservativo y calcular $\int_{\alpha}{F\cdot ds}$, donde $\alpha$ es cualquier curva que une los puntos $a_1$ y $a_2$.
% \begin{respuesta}
%     Notemos que $\rot F=0$ y el campo escalar $f$ definido por
%     \[
%     \funcion{f}{\R^3}{\R}
%     {(x,y,z)}{f(x,y,z)=\dfrac{x^2}{2}+\dfrac{y^3}{3}-\dfrac{z^4}{4}}
% 	\]
% 	es tal que 
%     \[
%     \nabla f(x,y) = F(x,y).
%     \]
%     Entonces, las integrales de línea de $F$ serán independientes de la trayectoria, es decir si $\alpha$ es cualquier curva de $a_1$ a $a_2$, tenemos
%     \begin{align*}
%     \int_{\alpha}{F\cdot ds} &=  f(a_2) - f(a_1)\\
%     		&=  f(2,3,-4) - f(1,1,1)\\
%             &= -\dfrac{643}{12}.
%     \end{align*}
% \end{respuesta}
%%%%%%%%%%%%%%%%%%%%%%%%%%%%%%%%%%%%%%%%%%%%%%%%%%% 
%%%%%%%%%%%%%%%%%%%%%%%%%%%%%%%%%%%%%%%%%%%%%%%%%%%
\item Suponga que la fuerza necesaria para mover un objeto $A$ en el espacio, está dada por
\[
    \funcion{F}{\R^3}{\R^3}{(x,y,z)}{(2xz-y,-x,x^2).}
\]
Determine el trabajo producido al mover el objeto $A$ por la trayectoria
\[
    \funcion{\alpha}{[-\pi,\pi]}{\R^3}{t}{(t,\cos(t),\sen(t)).}
\]

\begin{respuesta}
    Notemos que $\dom(F)=\R^3$ es conexo por caminos y además,
    \[
        \rot F(x,y,z)=0,
    \]
    de donde $F$ es un campo conservativo y su integral es independiente de la trayectoria. Por esta razón, elijamos una parametrización más sencilla. El segmento de recta que unen los puntos $\alpha(-\pi)=(-\pi,-1,0)$ y $\alpha(\pi)=(\pi,-1,0)$ está parametrizada por
    \[
        \funcion{\beta}{[0,1]}{\R^3}{t}{(-\pi +2\pi t,-1,0).}
    \]
    Por todo esto, el trabajo está definido por
    \begin{align*}
        \int_\alpha F ds
        &=\int_\beta F ds\\
        &=\int_0^1 F(\beta(t)) \cdot \beta'(t) dt\\
        &=\int_0^1 F(-\pi +2\pi t,-1,0) \cdot (2\pi,0,0) dt\\
        &=\int_0^1 2\pi dt\\
        &=2\pi        
    \end{align*}
    Por lo tanto, el trabajo que realiza el campo al mover al pbjeto $A$ por la trayectoria $\alpha$ es $2\pi$ unidades de trabajo. Como proceso alternativo el estudiante podría hallar el potencial de $F$ y utilizar el segundo teorema fundamental del cálculo para integrales de línea.
\end{respuesta}

%%%%%%%%%%%%%%%%%%%%%%%%%%%%%%%%%%%%%%%%%%%
%%%%%%%%%%%%%%%%%%%%%%%%%%%%%%%%%%%%%%%%%%%
%%%%%%%%%%%%%%%%%%%%%%%%%%%%%%%%%%%%%%%%%%%
% Tema 16: Integrales de superficie
\item 
    Sea la parametrización
	\[
		\funcion{\alpha}{\R^2}{\R^3}{(u,v)}{(u+v,u-v,u^2+v^2)}
	\]
	y sea $S=\img(\alpha)$. Determine el vector normal estándar en el punto $(1,3)$. Luego determine si $S$ es suave.

\begin{respuesta}
	Por definición de los vectores tangentes a las curvas coordenadas se tiene que
	\[
		\funcion{T_x}{\R^2}{\R^3}{(u,v)}{(1,1,2u)}
		\yds
		\funcion{T_y}{\R^2}{\R^3}{(u,v)}{(1,-1,2v)}.
	\]
	Así, por la definición del vector normal estándar en el punto $(1,3)$ tenemos que
	\begin{align*}
		N(1,3)&= T_x(1,3) \times T_y(1,3)\\
			&= (1,1,2) \times (1,-1,6)\\
			&= (8,-4,-2)
	\end{align*}
	Por lo tanto,
	\[
		N(1,3)= (8,-4,-2).
	\]
	Finalmente, para cualquier $u,v\in\R^2$, se tiene que
	\[
	    N(u,v)= T_x(u,v) \times T_y(u,v)
		= (1,1,2u) \times (1,-1,2v)
		= (2 u + 2 v, 2 u - 2 v, -2) \neq 0 
	\]
	de donde $S$ es suave.
\end{respuesta}

%%%%%%%%%%%%%%%%%%%%%%%%%%%%%%%%%%%%%%%%%%%
%%%%%%%%%%%%%%%%%%%%%%%%%%%%%%%%%%%%%%%%%%%
%%%%%%%%%%%%%%%%%%%%%%%%%%%%%%%%%%%%%%%%%%%
% Tema 16: Integrales de superficie
\item 
    Encuentre el área superficial de la helicoidal con parametrización 
    \[
        \funcion{\alpha}{[0,1]\times[0,2\pi n]}{\R^3}{(r,\theta)}{\alpha(r,\theta)=(r\cos(\theta),r\sen(\theta),\theta),}
    \]
    donde $n$ es un entero positivo.
    \begin{figure}[h]
        \centering
        % \includegraphics[width=6cm]{Graficos/fig12:04.eps}
    \end{figure}	
	
\begin{respuesta} 
	Primero determinamos los vectores tangentes a la superficie 
	\[
	    T_r(r,\theta)=\frac{d\alpha}{dr}(r,\theta)=(\cos(\theta), \sen(\theta), 0)
	\]
	\[
	    T_{\theta}(r,\theta)=\frac{d\alpha}{d\theta}(r,\theta)=(-r\sen(\theta), r\cos(\theta), 1)
	\]
    ahora podemos encontrar el vector normal estándar	
    \[
        N(r,\theta)=T_r(r,\theta)\times T_\theta(r,\theta)
    \]
	\[
	    N(r,\theta)=\left|\begin{matrix}
	    i &j &k\\
	    \cos(\theta) &\sen(\theta) &0\\
	    -r\sen(\theta) &r\cos(\theta) &1
	    \end{matrix}\right|=(\sen(\theta), -\cos(\theta), r)
	\]
    Podemos garantizar que no existe un valor de $r$ y de $\theta$ que anule al normal, por lo tanto es una superficie suave. Por tanto calculamos el área de la superficie de la helicoidal de $n$ vueltas. Sea $A$ el área, se tiene que 
    \begin{align*}
        A&=\iint_D\|\ N(r,\theta)\|dr \,d\theta\\
        &=\int_0^{2\pi n}\int_0^{1} \sqrt{1+r^2}dr \,d\theta\\
        &=\int_0^{2\pi n}\frac{1}{2}[\sqrt{2}+\ln{(\sqrt{2}+1)}]\,d\theta\\
        &=[\sqrt{2}+\ln{(\sqrt{2}+1)}]\pi n.
    \end{align*}
\end{respuesta}
%%%%%%%%%%%%%%%%%%%%%%%%%%%%%%%%%%%%%%%%%%%
%%%%%%%%%%%%%%%%%%%%%%%%%%%%%%%%%%%%%%%%%%%
%%%%%%%%%%%%%%%%%%%%%%%%%%%%%%%%%%%%%%%%%%%
% Tema 16: Integrales de superficie
\item Obtenga una parametrización de la superficie plana limitada por la curva $\mathcal{C}$ definida por
\[
    \mathcal{C}: 
    \begin{cases}
        x^2 + y^2 = \dfrac{z^2}{2},\\
        z = y + 1.
    \end{cases}
\]

\begin{respuesta}\hspace{0pt}

Notemos que la curva $\mathcal{C}$, producto de la intersección del cono con el plano es una elipse y su respectiva proyección sobre el plano de ecuación $z = 0$ es la curva de ecuación
\[
    x^2 + \dfrac{(y - 1)^2}{2} = 1
\]
la cual es una elipse, a continuación se muestra la representación gráfica de la curva
\begin{center}
	\includegraphics[scale=0.6]{Graficos/HE14_fig00.eps}
\end{center}
Se define el dominio $\Omega$ como
\[
    \Omega = \l\{(x, y)\in\R^{2}\,:\,x^2 + \dfrac{(y - 1)^{2}}{2} \leq 1\r\}
\]
y se parametriza la curva $\mathcal{C}$ por medio de la función
\[
    \funcion{\beta}{\Omega}{\R^{3}}{(x,y)}{(x, y, y+1).}\qedhere
\]
\end{respuesta}
%%%%%%%%%%%%%%%%%%%%%%%%%%%%%%%%%%%%%%%%%%%
%%%%%%%%%%%%%%%%%%%%%%%%%%%%%%%%%%%%%%%%%%%
%%%%%%%%%%%%%%%%%%%%%%%%%%%%%%%%%%%%%%%%%%%
% Tema 16: Integrales de superficie
\item Calcule el área de la porción de superficie de ecuación $x^2 + y^2 = z^2$ situada por encima del plano de ecuación $z = 0$ y limitada por la esfera de ecuación $x^2 + y^2 + z^2 = 2ax$ con $a>0$.

\begin{respuesta}\hspace{0pt}

Notemos que el sólido $S$ cuya área vamos a determinar comprende la hoja superior del cono de ecuación $x^2 + y^2  = z^2$ junto con la parte superior de la esfera considerada. Notemos que la proyección de la intersección del cono con la esfera en el plano de ecuación $z = 0$ es la circunferencia de ecuación
\[
    \l(x - \dfrac{a}{2}\r)^2 + y^2 = \dfrac{a^2}{4}.
\]
Se definen, el domino
\[
    \Omega = \l\{(x, y)\in\R^2\,:\,\l(x - \dfrac{a}{2}\r)^2 + y^2 \leq \dfrac{a^2}{4}\r\}
\]
y la parametrización 
\[
    \funcion{\beta}{\Omega}{\R^3}{(x, y)}{\l(x, y, \sqrt{x^2 + y^2}\r)}
\]
que es tal que $S = \beta(\Omega)$. Notemos que un vector normal dado por está parametrización es 
\[
    N(x, y) = \l(-\dfrac{x}{\sqrt{x^2 + y^2}}, -\dfrac{y}{\sqrt{x^2 + y^2}}, 1\r)
\]
con norma $\|N(x, y)\| = \sqrt{2}$ para cada $(x, y)\in\Omega$, Por tanto, si $A$ representa el área buscada, entonces
\[
    A = \iint_{\Omega}{\|N(x, y)\|\,dx dy} = \iint_{\Omega}{\sqrt{2}\,dx dy} = \dfrac{\sqrt{2}}{4}a^2 \pi.\qedhere
\]
\end{respuesta}
%%%%%%%%%%%%%%%%%%%%%%%%%%%%%%%%%%%%%%%%%%%
%%%%%%%%%%%%%%%%%%%%%%%%%%%%%%%%%%%%%%%%%%%
%%%%%%%%%%%%%%%%%%%%%%%%%%%%%%%%%%%%%%%%%%%
% Tema 16: Integrales de superficie
\item Dado el recinto limitado por los planos de ecuaciones $z = y$, $z = 0$ y el cilindro de ecuación $x^2 + y^2 = a^2$, con $a>0$. Calcule el área de la porción de superficie cilíndrica comprendida entre los dos planos.

\begin{respuesta}\hspace{0pt}

Dado que consideramos un cilindro, sean 
\[
    \l\{
    \begin{array}{rcl}
        x & = & a\cos(u), \\
        y & = & a\sen(u), \\
        z & = & v,
    \end{array}
    \r.
\]
considerando solamente la parte del cilindro tal que $z\geq0$, tenemos que
\[
    0 \leq u \leq \pi 
    \texty
    0 \leq v \leq a\sen(u).
\]
Se definen, el domino $\Omega$ como
\[
    \Omega = \l\{(u, v)\in\R^2\,:\,0 \leq u \leq \pi,\,0 \leq v \leq a\sen(u)\r\}
\]
y la parametrización 
\[
    \funcion{\beta}{\Omega}{\R^3}{(u, v)}{(a\cos(u), a\sen(u), v)}
\]
Por tanto, sea $S = \beta(\Omega)$ la mitad de la superficie estudiada, dado que hemos considerado la porción del cilindro tal que $z \geq 0$. Para está parametrización tenemos que un vector normal a la superficie es
\[
    N(u, v) = (a\cos(u), a\sen(u), 0)
\]
con norma $\|N(u, v)\| = a$ para cada $(u, v)\in\Omega$. Sea $A$ el área de $S$, entonces
\[
    A = \iint_{\Omega}{\|N(u,v)\|\,du dv} = \iint_{\Omega}{a\,du dv} = 2a^2.
\]
Por lo tanto, el área del sólido del enunciado es dos veces el área de $S$, es decir, el área buscada es igual a $4a^2$.\qedhere
\end{respuesta}
%%%%%%%%%%%%%%%%%%%%%%%%%%%%%%%%%%%%%%%%%%%
%%%%%%%%%%%%%%%%%%%%%%%%%%%%%%%%%%%%%%%%%%%
%%%%%%%%%%%%%%%%%%%%%%%%%%%%%%%%%%%%%%%%%%%
% Tema 16: Integrales de superficie
\item
	Calcular el área de la sección de un paraboloide de ecuación $x^2+z^2=2ay$, con $a>0$, delimitada entre el plano $xz$ y el plano de ecuación $y=a$. 

\begin{respuesta}
    La intersección del paraboloide y el plano es
    \[
        x^2+z^2=2a^2,
    \]
    esto representa la región de integración T sobre el plano xz.  la superficie
    \[
        \funcion{f}{\R^2}{\R}{(x,z)}{\dfrac{1}{2a}(x^2+z^2)}.
    \]
    Como la superficie $S$ está dada en forma explícita $y=f(x,z)$
    \[
        \parallel \dfrac{\partial r}{\partial x} \times \dfrac{\partial r}{\partial z}\parallel = \sqrt[]{1+(\dfrac{\partial f}{\partial x})^2+(\dfrac{\partial f}{\partial z})^2}=\sqrt[]{1+\dfrac{x^2+z^2}{a^2}}.
    \]
    Sea el cambio de variable,
    \begin{align*}
        x&=r \cos \theta\\
        y&=r \sen \theta,
    \end{align*}
    donde
    \begin{align*}
        \theta\leq r\leq a \sqrt[]{2}\\
        0 \leq \theta \leq 2\pi.
    \end{align*}
    \[
        a(S)=\int \int_T \parallel \dfrac{\partial r}{\partial x} X \dfrac{\partial r}{\partial z}\parallel dx dy = \int_0^{2 \pi} \left[ \int_0^{a\sqrt[]{2}}\sqrt[]{1+\dfrac{r^2}{a^2}}\ r dr \right] d\theta=\dfrac{2}{3} \pi a^2 (3\sqrt[]{3} -1 ).
    \]
\end{respuesta}
%%%%%%%%%%%%%%%%%%%%%%%%%%%%%%%%%%%%%%%%%%%
%%%%%%%%%%%%%%%%%%%%%%%%%%%%%%%%%%%%%%%%%%%
%%%%%%%%%%%%%%%%%%%%%%%%%%%%%%%%%%%%%%%%%%%
% Tema 16: Integrales de superficie
\item 
    Sabiendo que la ecuación $(\sqrt{x^2 + y^2} -a)^2 + z^2 = b^2$ con $a>b>0$ define un toro en $\R^3$, compruebe que la función
    \[
        \funcion{\alpha}{[0,2\pi]\times[0,2\pi]}{\R^3}
        {(u,v)}{f(u,v)=((a +b\cos{v})\cos{u},(a +b\cos{v})\sen{u},b\sen{v})}
	\]
    es una paramerización del toro, con esto, demuestre que el toro es una superficie suave y calcule su área superficial.
    
\begin{respuesta}
    Notemos que si tomamos 
    \[
        \l\{
        \begin{array}{l}
            x(u,v) = (a +b\cos{v})\cos{u},\\
            y(u,v) = (a +b\cos{v})\sen{u},\\
            z(u,v) = b\sen{v}.
        \end{array}
        \r.
    \]
    Entonces:
    \begin{align*}
        \l(\sqrt{x^2 + y^2} -a\r)^2 + z^2 &= \l(\sqrt{((a +b\cos{v})\cos{u})^2+((a +b\cos{v})\sen{u})^2}-a\r)^2+(b\sen{v})^2\\
            &= \l(\sqrt{(a +b\cos{v})^2}-a\r)^2+b^2\sen^2{v}\\
            &= \l(a+b\cos{v}-a\r)^2+b^2\sen^2{v}\\
            &= b^2\cos^2{v}+b^2\sen^2{v}=b^2,
    \end{align*}
    dado que $a+b\cos{v}>0$ ya que $a>b>0$, es decir, el toro está parametrizado por $\alpha$.
    
    Notemos además que se tienen:
    \begin{align*}
        \dfrac{\partial \alpha}{\partial u}(u,v) &= (-(a+b\cos{v})\sen{u},(a+b\cos{v})\cos{u},0),\\
        \dfrac{\partial \alpha}{\partial v}(u,v) &= (-b\sen{v}\cos{u},-b\sen{v}\sen{u},b\cos{v}),
    \end{align*}
    para todo $(u,v)\in[0,2\pi]\times[0,2\pi]$. A partir de lo anterior tenemos que
    \begin{align*}
        N(u,v) &= \dfrac{\partial \alpha}{\partial u}(u,v) \times \dfrac{\partial \alpha}{\partial v}(u,v)\\
            &= b(a+b\cos{v})(\cos{v}\cos{u},\cos{v}\sen{u},\sen{v}).
    \end{align*}
    De igual forma, dado que $a+b\cos{v}>0$ ya que $a>b>0$, $N(u,v)$ nunca es cero, por lo tanto, el toro es una superficie parametrizada suave.
    
    EL área superficial está dada por
    \begin{align*}
        \iint_{\alpha}{f\cdot dS} &= \iint_{D}{\|N(u,v)\|dudv}\\
            &= \int_{0}^{2\pi}{\int_{0}^{2\pi}{\sqrt{b^2(a+b\cos{v})^2[\cos^2{v}\cos^2{u}+\cos^2{v}\sen^2{u}+\sen^2{v}]}}du dv}\\
            &= \int_{0}^{2\pi}{\int_{0}^{2\pi}{\sqrt{b^2(a+b\cos{v})^2}}du dv}\\
            &= b\int_{0}^{2\pi}{\int_{0}^{2\pi}{(a+b\cos{v})}du dv}\\
            &= 2b\pi\int_{0}^{2\pi}{(a+b\cos{v})dv}\\
            &= 4\pi^2 ab.
    \end{align*}
    Por lo tanto se ha mostrado que el toro es una superficie parametrizada suave con área superficial igual a $4\pi^2 ab\,[u^2]$, con $a>b>0$.
\end{respuesta}
%%%%%%%%%%%%%%%%%%%%%%%%%%%%%%%%%%%%%%%%%%%
%%%%%%%%%%%%%%%%%%%%%%%%%%%%%%%%%%%%%%%%%%%
%%%%%%%%%%%%%%%%%%%%%%%%%%%%%%%%%%%%%%%%%%%
% Tema 16: Integrales de superficie
\item 
    Verifique que la imagen de la superficie parametrizada por
    \[
        \funcion{\alpha}{[0,\pi]\times[0,2\pi]}{\R^3}
        {(u,v)}{f(u,v)=(2\sen{u}\cos{v},3\sen{u}\sen{v},\cos{u})}
	\]
    es un elipsoide.

\begin{respuesta}
    Notemos que si tomamos 
    \[
        \l\{
        \begin{array}{l}
            x(u,v) = 2\sen{u}\cos{v},\\
            y(u,v) = 3\sen{u}\sen{v},\\
            z(u,v) = \cos{u}.
        \end{array}
        \r.
    \]
    Entonces:
    \begin{align*}
        \dfrac{x^2(u,v)}{4} + \dfrac{y^2(u,v)}{9} + z^2(u,v) &= \dfrac{(2\sen{u}\cos{v})^2}{4} + \dfrac{(3\sen{u}\sen{v})^2}{9} + (\cos{u})^2\\
        &= \sen^2{u}\cos^2{v} + \sen^2{u}\sen^2{v} + \cos^2{u}\\
        &= 1.
    \end{align*}
    Es decir, la respectiva superficie verifica la ecuación
    \[
        \dfrac{x^2}{4} + \dfrac{y^2}{9} + z^2 = 1
    \]
    esto es, la ecuación de una elipsoide.
\end{respuesta}

%%%%%%%%%%%%%%%%%%%%%%%%%%%%%%%%%%%%%%%%%%%
%%%%%%%%%%%%%%%%%%%%%%%%%%%%%%%%%%%%%%%%%%%
%%%%%%%%%%%%%%%%%%%%%%%%%%%%%%%%%%%%%%%%%%%
% Tema 16: Integrales de superficie
\item 
    Sabiendo que la ecuación $z^2 = x^2 + y^2$ define un cono en $\R^3$. Muestre que el cono no es suave en su vértice.

\begin{respuesta}
    Notemos que la función 
    \[
        \funcion{\alpha}{\R\times[0,2\pi]}{\R^3}
        {(u,v)}{f(u,v)=(u\cos{v},u\sen{v},u)}
	\]
    parametriza al cono. 
    
    \noindent	Notemos además que se tienen:
    \begin{align*}
        \dfrac{\partial \alpha}{\partial u}(u,v) &= (\cos{v},\sen{v},1),\\
        \dfrac{\partial \alpha}{\partial v}(u,v) &= (-u\sen{v},u\cos{v},0),
    \end{align*}
    para todo $(u,v)\in\R\times[0,2\pi]$. A partir de lo anterior tenemos que
    \begin{align*}
        N(u,v) &= \dfrac{\partial \alpha}{\partial u}(u,v) \times \dfrac{\partial \alpha}{\partial v}(u,v)\\
            &= (-u\cos{v},-u\sen{v},u).
    \end{align*}
    Además, si tomamos $u=0$, entonces $N(u,v)=(0,0,0)$, por lo tanto, el cono no es $\l(\dint_a^b f\r)$ suave en su vértice.
\end{respuesta}


%%%%%%%%%%%%%%%%%%%%%%%%%%%%%%%%%%%%%%%%%%%
%%%%%%%%%%%%%%%%%%%%%%%%%%%%%%%%%%%%%%%%%%%
%%%%%%%%%%%%%%%%%%%%%%%%%%%%%%%%%%%%%%%%%%%
% Tema 16: Integrales de superficie
\item 
	Dadas
	\[
		\funcion{\alpha}{\lcl0,1\rcl \times \lcl0,2\rcl}{\R^3}{(u,v)}{(1-u,1-v,u+v)}
	\]
	y
	\[
		\funcion{f}{\R^3}{\R}{x}{x_1 x_2 x_3}
	\]
	Evalúe la integral
	\[
		\int_\alpha f\, dS.
	\]

\begin{respuesta}
	Se tiene que el vector normal estándar de $\alpha$ es
	\[
		N(u,v)=(-1,0,1)\times (0,-1,1)= (1,1,1)
	\]
	para todo $u\in \lcl0,1\rcl$ y $v\in\lcl0,2\rcl$. Con esto vemos que $S$ es suave. 

	Si $D=\lcl0,1\rcl \times \lcl0,2\rcl$, entonces
	\begin{align*}
		\int_\alpha f dS
			&=\iint_D f(\alpha(u,v))||N(u,v)|| du dv\\
			&=\sqrt 3\iint_D f(1-u,1-v,u+v) du dv\\
			&=\sqrt 3\int_0 ^1 \int_0 ^2  (1-u)(1-v)(u+v) dv du\\
			&=-\dfrac{\sqrt 3}{3}.
	\end{align*}
	Por lo tanto,
	\[
		\int_\alpha f dS
		=-\dfrac{\sqrt 3}{3}. \qedhere
	\]
\end{respuesta}

%%%%%%%%%%%%%%%%%%%%%%%%%%%%%%%%%%%%%%%%%%%
%%%%%%%%%%%%%%%%%%%%%%%%%%%%%%%%%%%%%%%%%%%
%%%%%%%%%%%%%%%%%%%%%%%%%%%%%%%%%%%%%%%%%%%
% Tema 16: Integrales de superficie
\item 
	Realice el ejercicio anterior si 
	\[
		\funcion{\alpha}{D}{\R^3}{(u,v)}{(u^2,u+v,v)}
	\]
	donde $D=\{(u,v)\in\R^2: u\geq0\wedge \ v\geq 0 \ \wedge v\leq 2-u\}$.
	
\begin{respuesta}
	Se tiene que el vector normal estándar de $\alpha$ es
	\[
		N(u,v)=(2u,1,0)\times (0,1,1)= (1, -2 u, 2 u)
	\]
	para todo $(u,v)\in D$. Con esto vemos que $S$ es suave. 
		
	Se tiene que
	\begin{align*}
		\int_\alpha f dS
			&=\iint_D f(\alpha(u,v))||N(u,v)|| du dv\\
			&=\iint_D f(u^2,u+v,v) \sqrt{1+8u^2}du dv\\
			&=\int_0^2 \int_0^{2-u} u^2 v(u+v) \sqrt{1+8u^2}du dv\\
			&\approx 2.76.
	\end{align*}
	Por lo tanto,
	\[
		\int_\alpha f dS
		\approx 2.76.
		\qedhere
	\]
\end{respuesta}



%%%%%%%%%%%%%%%%%%%%%%%%%%%%%%%%%%%%%%%%%%%
%%%%%%%%%%%%%%%%%%%%%%%%%%%%%%%%%%%%%%%%%%%
%%%%%%%%%%%%%%%%%%%%%%%%%%%%%%%%%%%%%%%%%%%
% Tema 16: Integrales de superficie
\item
    Evalúe la integral $\iint_S x^2\,dS$ donde $S$ es la superficie del cubo definido en $[-2,2]\times [-2,2]\times[-2,2]$

\begin{respuesta} 
    En primer lugar debemos parametrizar las seis caras del cubo de la siguiente manera
    \begin{table}[htbp]
    \begin{center}
    \begin{tabular}{l|l|l}
    i & $\alpha (u,v)$ para $S_i$ & cara \\
    \hline 
    1 & $(u,v,2)$ & superior\\ \hline
    2 & $(u,v,-2)$ & fondo \\ \hline
    3 & $(u,2,v)$ & derecha \\ \hline
    4 & $(u,-2,v)$ & izquierda \\ \hline
    5 & $(2,u,v)$ & frente \\ \hline
    6 & $(-2,u,v)$ & posterior \\
    \end{tabular}
    \end{center}
    \end{table}
    \\para cada $(u,v)\in[-2,2]\times[-2,2]$. $S_i$ representa la superficie de cada cara del cubo.
    
   Los vectores tangentes de la cara $S_1$ son
    \[
        T_u=(u,v)=\frac{d\alpha}{du}(u,v)=(1,0,0)
    \]
    \[
        T_v=(u,v)=\frac{d\alpha}{dv}(u,v)=(0,1,0)
    \]
    por tanto el vector normal estándar $N(u,v)=T_u(u,v)\times T_v(u,v)$ es
    \[
        N(u,v)=(0,0,1).
    \]
    Note que en cada cara del cubo se cumple que $\|N(u,v)\|=1$, entonces
    \[
        \iint_{S_i}[x(u,v)]^2\|N(u,v)\|du\,dv=\iint_{S_i}[x(u,v)]^2du\,dv
    \]
    para $1\leq i \leq 6$.
        
    Por tanto
    \begin{align*}
        \iint_{S_i}[x(u,v)]^2du\,dv=\int_{-2}^2\int_{-2}^2u^2du\,dv \quad&; \quad para\quad  i=1,2,3,4 \\
        \iint_{S_i}[x(u,v)]^2du\,dv=\int_{-2}^2\int_{-2}^24du\,dv \quad&; \quad para\quad  i=5,6 \\ 
    \end{align*}
    entonces
    \begin{align*}
        \iint_Sx^2\,ds&=\sum_{i=1}^6\iint_{S_i}[x(u,v)]^2\|N(u,v)\|du\,dv\\
        &=4\int_{-2}^2\int_{-2}^2u^2du\,dv+2\int_{-2}^2\int_{-2}^24du\,dv\\
        &=4\int_{-2}^2\frac{u^3}{3}|_{-2}^2\,dv+8\int_{-2}^2u|_{-2}^2\,dv\\
        &=\frac{256}{3}+128\\
        &=\frac{640}{3}
    \end{align*}
\end{respuesta}

%%%%%%%%%%%%%%%%%%%%%%%%%%%%%%%%%%%%%%%%%%%
%%%%%%%%%%%%%%%%%%%%%%%%%%%%%%%%%%%%%%%%%%%
%%%%%%%%%%%%%%%%%%%%%%%%%%%%%%%%%%%%%%%%%%%
% Tema 16: Integrales de superficie
\item 
    Dado el campo escalar
	\[
    	\funcion{f}{\R^3}{\R^2}{(x,y,z)}{x+z,} 
    \]
    calcular su integral sobre la superficie encerrada en la curva resultante de la intersección del cilindro de ecuación $x^2+y^2=9$ y el plano de ecuación $z=4$.

\begin{respuesta}
    La superficie es parametrizada por
    \[
        \funcion{\alpha}{D}{\R^3}{(u,v)}{(u \cos v,u \sen v,4)}
    \]
    donde $D=\{(u,v)\in\R^2:0\leq u \leq 3,0\leq v \leq 2\pi \}.$
    
    Entonces
	\[
        \iint_\alpha f dS=\iint_D f\l(\alpha (u,v) \r) \|N(u,v)\|dudv.
    \]
    Se calcula
	\[
        \begin{aligned}
            T_u(u,v)&=\dfrac{\partial \alpha}{\partial u}(u,v)=(\cos v,\sen v,0)\\
            T_v(u,v)&=\dfrac{\partial \alpha}{\partial v}(u,v)=(-u\sen v, u\cos v,0)
        \end{aligned}
    \]
    entonces
    \[
        \begin{aligned}
            \|N(u,v)\| &= \|(\cos v,\sen v,0)\times(-u\sen v, u\cos v,0)\| \\
            \|N(u,v)\| &= u.
        \end{aligned}
    \]
    Finalmente
    \[
        \begin{aligned}
            \iint_\alpha f dS&=\int_0^{2\pi} \int_0^3 \l(u \cos v + 4\r) u dudv\\
                &=36 \pi
        \end{aligned}
    \]
\end{respuesta}
%%%%%%%%%%%%%%%%%%%%%%%%%%%%%%%%%%%%%%%%%%%%%%%%%
%%%%%%%%%%%%%%%%%%%%%%%%%%%%%%%%%%%%%%%%%%%%%%%%%
%%%%%%%%%%%%%%%%%%%%%%%%%%%%%%%%%%%%%%%%%%%%%%%%%
% Tema 16: Integrales de superficie
\item 
    Dado el campo escalar
    \[
        \funcion{f}{\R^3}{\R}{(x,y,z)}{y,}
    \]
    calcule la integral de superficie de $f$ sobre cada frontera de la región $E$ acotada por el cilindro parabólico y el plano definidos por las ecuaciones: $z=1-x^2$, $y+z=2$; y los planos cartesianos.

\begin{center}
% \includegraphics[scale=0.7]{Graficos/Int_Superficie.eps}
\end{center}

\begin{respuesta}
    Vamos a calcular las 4 integrales por separado:
    \begin{enumerate}
        \item Sea $S_1$ la superficie frontera de $E$, determinada por las ecuaciones $y+z=2$ y $z=1-x^2$. Una parametrización de $S_1$ es
        \[
            \funcion{\alpha_1}{\Omega_1}{\R^3}{(x,z)}{(x,2-z,z),}
        \]
        donde 
        \[
            \Omega_1=\{(x,z): -1\leq x\leq 1\ \wedge\ 0\leq z\leq 1-x^2 \}.
        \]
        Por cálculo directo, se tiene que para cada $(x,z)\in\Omega_1$,
        \[
            D_1\alpha_1(x,z)=(1,0,0)
            \yds 
            D_2\alpha_1(x,z)=(0,-1,1);
        \]
        luego
        \[
            D_1\alpha_1(x,z)\times D_2\alpha_1(x,z)
            =(0,-1,-1).
        \]
        Por definición, se tiene que
        \begin{align*}
            \iint_{S_1} f dS
            &= \iint_{\Omega_1} f(\alpha_1(x,z)) \|D_1\alpha_1(x,z)\times D_2\alpha_1(x,z)\| \ dA\\
            &= \iint_{\Omega_1} f(x,2-z,z) \|(0,-1,1)\| \ dA\\
            &= \sqrt 2\int_{-1}^1 \int_0^{1-x^2} (2-z) \ dz dx\\
            &= \dfrac{32\sqrt 2}{15}.
        \end{align*}
        
        \item Sea $S_2$ la superficie frontera de $E$, determinada por la ecuación $z=1-x^2$. Una parametrización de $S_2$ es
        \[
            \funcion{\alpha_2}{\Omega_2}{\R^3}{(x,y)}{(x,y,1-x^2),}
        \]
        donde 
        \[
            \Omega_2=\{(x,y): -1\leq x\leq 1 \ \wedge \ 0\leq y\leq 1+x^2\}.
        \]
        En efecto, dado que la intersección del cono y del plano es $y=1+x^2$. Por cálculo directo, se tiene que para cada $(x,y)\in\Omega_2$,
        \[
            D_1\alpha_2(x,y)=(1,0,-2x)
            \yds 
            D_2\alpha_2(x,y)=(0,1,0);
        \]
        luego
        \[
            D_1\alpha_2(x,y)\times D_2\alpha_2(x,y)
            =(2x,0,1).
        \]
        Por definición, se tiene que
        \begin{align*}
            \iint_{S_2} f dS
            &= \iint_{\Omega_2} f(\alpha_2(x,y)) \|D_1\alpha_2(x,y)\times D_2\alpha_2(x,y)\| \ dA\\
            &= \iint_{\Omega_2} f(x,y,1-x^2) \|(2x,0,1)\| \ dA\\
            &= \int_{-1}^1 \int_0^{1+x^2} y \sqrt{4x^2+1} \ dy dx\\
            &= \frac{1}{2}\int_{-1}^1 (1+x^2)^2\sqrt{4x^2+1} dx\\
            &= \int_{0}^1 (1+x^2)^2\sqrt{4x^2+1} dx\\
            &\approx 3.08.
        \end{align*}
        
        \item Sea $S_3$ la superficie frontera de $E$ sobre el plano $x\ y$. Una parametrización de $S_3$ es
        \[
            \funcion{\alpha_3}{\Omega_3}{\R^3}{(x,y)}{(x,y,0),}
        \]
        donde 
        \[
            \Omega_3=[-1,1]\times [0,2].
        \]
        Por cálculo directo, se tiene que para cada $(x,y)\in\Omega_3$,
        \[
            D_1\alpha_3(x,y)=(1,0,0)
            \yds 
            D_2\alpha_3(x,y)=(0,1,0);
        \]
        luego
        \[
            D_1\alpha_3(x,y)\times D_2\alpha_3(x,y)
            =(0,0,1).
        \]
        Por definición, se tiene que
        \begin{align*}
            \iint_{S_3} f dS
            &= \iint_{\Omega_3} f(\alpha_3(x,y)) \|D_1\alpha_3(x,y)\times D_2\alpha_3(x,y)\| \ dA\\
            &= \iint_{\Omega_3} f(x,y,0) \|(0,0,1)\| \ dA\\
            &= \int_{-1}^1 \int_0^{2} y \ dy dx\\
            &= 4.
        \end{align*}
        
        \item Sea $S_4$ la superficie frontera de $E$ sobre el plano $x\ z$. Una parametrización de $S_4$ es
        \[
            \funcion{\alpha_4}{\Omega_4}{\R^3}{(x,z)}{(x,0,z),}
        \]
        donde 
        \[
            \Omega_4=\{(x,z): -1\leq x\leq 1 \ \wedge \ 0\leq z\leq 1-x^2\}.
        \]
        Por cálculo directo, se tiene que para cada $(x,z)\in\Omega_4$,
        \[
            D_1\alpha_4(x,z)=(1,0,0)
            \yds 
            D_2\alpha_4(x,z)=(0,0,1);
        \]
        luego
        \[
            D_1\alpha_4(x,z)\times D_2\alpha_4(x,z)
            =(0,1,0).
        \]
        Por definición, se tiene que
        \begin{align*}
            \iint_{S_4} f dS
            &= \iint_{\Omega_4} f(\alpha_4(x,z)) \|D_1\alpha_4(x,z)\times D_2\alpha_4(x,z)\| \ dA\\
            &= \iint_{\Omega_4} f(x,0,z) \|(0,1,0)\| \ dA\\
            &= \int_{-1}^1 \int_0^{1-x^2} 0 \ dz dx\\
            &= 0.\qedhere
        \end{align*}        
    \end{enumerate}
\end{respuesta}

%%%%%%%%%%%%%%%%%%%%%%%%%%%%%%%%%%%%%%%%%%%%%%%%
%%%%%%%%%%%%%%%%%%%%%%%%%%%%%%%%%%%%%%%%%%%%%%%%
%%%%%%%%%%%%%%%%%%%%%%%%%%%%%%%%%%%%%%%%%%%%%%%%

% Tema 16: Integrales de superficie
\item
    Encuentre la masa del cono centrado en el origen con radio de la base 2 y altura 3, si la densidad está dada por $\delta(x, y, z) = \sqrt{x^{2} + y^{2}}$ para cada $(x, y)\in\R^{2}$.
    
\begin{respuesta}\hspace{0pt}

El cono que trabajamos se representa a continuación
\begin{center}
	\includegraphics[scale=0.5]{Graficos/HE15_fig00.eps}
\end{center}
Para obtener una parametrización del cono, tomemos a $u$, como la altura respecto de la base, y a $v$, como el ángulo respecto al eje $x$. Luego, cada punto de la superficie se representa por
\[
    \begin{cases}
        x = r\cos(v),\\
        y = r\sen(v),\\
        z = u,
    \end{cases}
\]
con
\[
    0 \leq u \leq 3
    \texty
    0 \leq v \leq 2\pi.
\]
Notemos que la letra $r$ representa al radio, a continuación la eliminaremos, puesto que es un tercer parámetro. En el cono estudiado tenemos que la altura y el radio están relacionados por una función lineal. En este caso, cuando el radio es 2 la altura es 0, mientras que cuando el radio es 0 la altura es 3, por lo cual, 
\[
    r = 2 - \dfrac{2u}{3}.
\]
Por lo tanto, se define $\Omega = \{(u, v)\in\R^{2}\,:\,0 \leq u \leq 3,\,0 \leq v \leq 2\pi\}$ y la parametrización
\[
    \funcion{\alpha}{\Omega}{\R^{3}}{(u, v)}{\l(\l(2 - \dfrac{2u}{3}\r)\cos(v), \l(2 - \dfrac{2u}{3}\r)\sen(v), u\r).}
\]
El vector normal estándar a la superfice es
\[
    N(u, v) = \l(-\l(2 - \dfrac{2u}{3}\r)\cos(v), -\l(2 - \dfrac{2u}{3}\r)\sen(v), -\dfrac{2}{3}\l(2 - \dfrac{2u}{3}\r)\r)
\]
con $\|N(u, v)\| = \dfrac{\sqrt{13}}{3}\l(2 - \dfrac{2u}{3}\r)$ para cada $(u, v)\in\Omega$.

Notemos que la densidad se escribe, en función de las variables $u$ y $v$, como
\[
    \delta(\alpha(u, v)) = \l(2 - \dfrac{2u}{3}\r) 
\]
para cada $(u, v)\in\Omega$. Luego, sea $M$ la masa buscada, entonces
\begin{align*}
    M &= \iint_{\Omega}{\delta(\alpha(u, v))\|N(u, v)\|\, dA}\\
      &= \dfrac{\sqrt{13}}{3}\dint_{0}^{3}{\dint_{0}^{2\pi}{\l(2 - \dfrac{2u}{3}\r)^{2}\,dv},du}\\
      &= \dfrac{8\sqrt{13}}{3}\pi.\qedhere
\end{align*}

    De manera alternativa se puede parametrizar la superficie con
    \[
        \funcion{\beta}{[0,3]\times[0,2\pi]}{\R^3}{(r,t)}{\l(\dfrac{2}{3}r\cos t,\dfrac{2}{3}r\sen t,3-r \r),}
    \]
    de aquí
    \[
        \dfrac{\partial \beta}{\partial r}=\l(\dfrac{2}{3}\cos t,\dfrac{2}{3}\sen t,-1 \r) 
        \texty
        \dfrac{\partial \beta}{\partial t}=\l(-\dfrac{2}{3}r\sen t,\dfrac{2}{3}r\cos t,0 \r)
    \]
    entonces
    \[
        N(r,t)=\l(\dfrac{2}{3}r\cos t,\dfrac{2}{3}r\sen t,\dfrac{4}{9}r  \r)
        \texty
        \|N(r,t)\|=\dfrac{2\sqrt{13}}{9}r.
    \]
    Además
    \[
        \delta(\beta(r,t))=\dfrac{2}{3}r,
    \]
    finalmente se recomienda al estudiante verificar que
    \[
        M=\int_\beta \delta dS=\dfrac{4\sqrt{13}}{27} \int_0^{2\pi}\int_0^3 r^2drdt=\dfrac{8\sqrt{13}}{3}\pi.
    \]
\end{respuesta}

%%%%%%%%%%%%%%%%%%%%%%%%%%%%%%%%%%%%%%%%%%%
%%%%%%%%%%%%%%%%%%%%%%%%%%%%%%%%%%%%%%%%%%%
%%%%%%%%%%%%%%%%%%%%%%%%%%%%%%%%%%%%%%%%%%%
% Tema 16: Integrales de superficie
\item
    Demostrar que el momento de inercia de una esfera hueca alrededor de cualquier diámetro es $\frac{2}{3}ma^2$, donde $m$ es la masa de la esfera y $a$ el radio
    
\begin{respuesta}
    La representación paramétrica de la superficie es
    \[
        \funcion{\gamma}{[0,2\pi]\times [0,\pi]}{\R^3}{(\theta,\phi)}{(a\cos \theta \sen \phi,a\sen \theta \sen \phi,a \cos \phi)}
    \]
    de aquí se tiene que
    \[
        \frac{\partial \gamma}{\partial \theta}(\theta,\phi)=(-a\sen \theta \sen \phi,a\cos \theta \sen \phi,0)
    \]
    y
    \[
        \frac{\partial \gamma}{\partial \phi}(\theta,\phi)=(a\cos \theta \cos \phi,a\sen \theta \cos \phi,-a \sen \phi),
    \]
    entonces
    \[
        ||N(\theta,\phi)||=\left|\left| \frac{\partial \gamma}{\partial \theta}(\theta,\phi) \times \frac{\partial \gamma}{\partial \phi}(\theta,\phi) \right|\right|=a^2\sen \phi.
    \]
    Se entiende que la esfera es simétrica por lo que su densidad es constante y viene dada por $\frac{m}{A}$, donde $m$ es la masa y $A$ la superficie de la esfera. Por otro lado, al ser simétrica se puede calcular el momento de inercia con respecto a cualquier diámetro, en este caso se realizará el cálculo con respecto al eje $z$.
    
    La distancia de cualquier punto de la superficie al eje $z$ es
    \[
        d(x,y)=\sqrt[]{x^2+y^2}
    \]
    entonces
    \[
        d^2(\theta,\phi)=a^2\sen^2 \phi.
    \]
    Finalmente el momento de inercia con respecto a cualquier diámetro $l$ es
    \begin{align*}
        I&=\iint_\gamma d^2 \frac{m}{A} dS\\
            &=\iint_\gamma d^2(\gamma) \frac{m}{A}||N(\theta,\phi)||d\phi d\theta 
    \end{align*}
    por lo tanto
    \begin{align*}
        I&=\frac{m}{A}\iint_\gamma a^2\sen^2 \phi a^2\sen \phi d\phi d\theta\\
            &=\frac{m}{4\pi a^2} \int_{0}^{2\pi} \int_{0}^{\pi} a^4 \sen^3 \phi d\phi d\theta\\ 
            &=\frac{2}{3}ma^2
    \end{align*}
\end{respuesta}


%%%%%%%%%%%%%%%%%%%%%%%%%%%%%%%%%%%%%%%%%%% FC
%%%%%%%%%%%%%%%%%%%%%%%%%%%%%%%%%%%%%%%%%%%
\item Calcule la masa de una capa delgada de densidad $\rho (x,y,z)=\dfrac{1}{z^2}$ cortada del cono de ecuación $z=\sqrt{x^2+y^2}$ y por los planos de ecuación $z=1$ y $z=2$.
\begin{respuesta}
    Se define el dominio $\Omega$ como,
    \[
    \Omega=\lbrace{(r,\theta)\in \R^2:1 \leq r \leq 2, 0 \leq \theta \leq 2 \pi\rbrace},
    \]
    y la parametrización,
    \[
    \funcion{\alpha}{\Omega}{\R^3}{(r,\theta)}{(r \cos \theta, r \sen \theta, r)}
    \]
    (note que la altura esta definida por $r$ y ésta varía entre 1 y 2), su vector normal,
    \[
    \begin{aligned}
    \|N(r,\theta)\|&=T_r(r,\theta) \times T_{\theta}(r,\theta)\\
    &=(\cos \theta, \sen \theta,1) \times (-r \sen \theta, r \cos \theta, 0)\\
    &=(- r \cos \theta,r \sen \theta, r)
    \end{aligned}
    \]
    y su norma
    \[
    \|N(r,\theta)\|=\sqrt{2}r
    \]
    De donde se tiene que la masa está dada por la integral:
    \[
    \begin{aligned}
    M=\iint_\alpha \rho(r,\theta) \|N(r,\theta)\| dr d\theta&=\int_{0}^{2 \pi} \int_{1}^{2} \dfrac{1}{r^2} \sqrt{2} r dr d\theta\\
    &=2 \pi \sqrt{2} \ln 2\\
    &\approx 6.15 
    \end{aligned}
    \]
\end{respuesta}

%%%%%%%%%%%%%%%%%%%%%%%%%%%%%%%%%%%%%%%%%%%
%%%%%%%%%%%%%%%%%%%%%%%%%%%%%%%%%%%%%%%%%%%
%%%%%%%%%%%%%%%%%%%%%%%%%%%%%%%%%%%%%%%%%%%
% Tema 16: Integrales de superficie
\item
    Determinar el flujo, a través de una superficie esférica de radio $a$ centrada en el origen, debido al campo vectorial
    \[
        \funcion{F}{\R^3}{\R^3}{(x,y,z)}{(x,y,z).}
    \]
\begin{respuesta}
    La superficie es parametrizada por
    \[
        \funcion{\alpha}{[0,2\pi]\times [0,\pi]}{\R^3}{(\theta,\phi)}{(a\cos \theta \sen \phi,a\sen \theta \sen \phi,a \cos \phi)}
    \]
    de aquí se tiene que
    \[
        \frac{\partial \alpha}{\partial \theta}(\theta,\phi)=(-a\sen \theta \sen \phi,a\cos \theta \sen \phi,0)
    \]
    y
    \[
        \frac{\partial \alpha}{\partial \phi}(\theta,\phi)=(a\cos \theta \cos \phi,a\sen \theta \cos \phi,-a \sen \phi),
    \]
    entonces
    \[
        N(\theta,\phi)=\frac{\partial \alpha}{\partial \theta} \times \frac{\partial \alpha}{\partial \phi}=(-a^2\cos\theta\sen^2\phi,-a^2\sen\theta\sen^2\phi,-a^2\sen\phi\cos\phi).
    \]
    Además el flujo a través de la superficie viene dado por
    \[
        \int_\alpha FdS=\iint_\alpha F(\alpha(\theta,\phi))\cdot N(\theta,\phi)d\phi d\theta,
    \]
    nótese que 
    \[
        F(\alpha(\theta,\phi))=(a\cos \theta \sen \phi,a\sen \theta \sen \phi,a \cos \phi)
    \]
    entonces
    \[
        F(\alpha(\theta,\phi))\cdot N(\theta,\phi)=-a^3\sen\phi,
    \]
    de aquí
    \[
        \iint_\alpha F(\alpha(\theta,\phi))\cdot N(\theta,\phi)d\phi d\theta=-a^3\int_0^{2\pi}\int_0^{\pi}\sen\phi d\phi d\theta=-4\pi a^3.
    \]
\end{respuesta}

%%%%%%%%%%%%%%%%%%%%%%%%%%%%%%%%%%%%%%%%%%%
%%%%%%%%%%%%%%%%%%%%%%%%%%%%%%%%%%%%%%%%%%%
%%%%%%%%%%%%%%%%%%%%%%%%%%%%%%%%%%%%%%%%%%%
% Tema 16: Integrales de superficie
\item 
    Obtener el flujo a través del trozo de plano de ecuación $x+2y+z=1$ que se obtiene cuando $0 \leq x \leq 2$ y $0 \leq y \leq 2$ para el campo vectorial $f$ orientado según el vector normal ascendente. 
	\[
    	\funcion{f}{\R^3}{\R^3}{(x,y,z)}{\l(2e^{xy},-e^{xy},4xyz\r)}
    \]

\begin{respuesta}
    La integral de superficie de un campo vectorial está dada por:
	\[
        \iint_\alpha f dS=\iint_D f\l(\alpha (u,v) \r) \cdot N(u,v)dudv
    \]
    La parametrización sugerida es:
	\[
        \begin{aligned}
            x &= u \\
            y &= v\\
            z &= 1-u-2v \\
            \alpha (u,v)&=(u,v,1-u-2v)
        \end{aligned}
    \]
    Sabiendo que:
	\[
        \begin{aligned}
            T_u(u,v)&=\dfrac{\partial \alpha}{\partial u}(u,v)\\
            T_v(u,v)&=\dfrac{\partial \alpha}{\partial v}(u,v)
        \end{aligned}
    \]
    y que,
    \[
        N(u,v) = T_u(u,v) \times T_v(u,v)
    \]
    tenemos,
    \[
        \begin{aligned}
            N(u,v) &= (1,0,-1)\times(0,1,-2) \\
            N(u,v) &= (1,2,1)
        \end{aligned}
    \]
    La integral de superficie será,
	\[
        \begin{aligned}
            \iint_\alpha f dS&=\int_0^2 \int_0^2 \l(2e^{uv},-e^{uv},4uv\l(1-u-2v\r)\r) \cdot (1,2,1) dudv\\
                &=-48
        \end{aligned}
    \]
\end{respuesta}

%%%%%%%%%%%%%%%%%%%%%%%%%%%%%%%%%%%%%%%%%%%
%%%%%%%%%%%%%%%%%%%%%%%%%%%%%%%%%%%%%%%%%%%
%%%%%%%%%%%%%%%%%%%%%%%%%%%%%%%%%%%%%%%%%%%
% Tema 17: Teoremas integrales
\item 
    	Sea el campo vectorial
		\[
		\funcion{f}{\R^2}{\R^2}{(x,y)}{(y^2e^{x},x^2e^{y})}
		\]
		y sea $C$ la curva formada por el segmento de recta desde $(-1,1)$ hasta $(1,1)$ seguido por el arco de parábola definida por la ecuación $y=2-x^2$ desde $(1,1)$ hasta $(-1,1)$. Verifique que se cumple el Teorema de Green.

\begin{respuesta}
		Como $C$ es una curva cerrada; sea $D$ la región límitada por $C$. Primero debemos parametrizar la curva $C$ en sentido antihorario, para ello utilizamos dos funciones
		\[
		\funcion{\alpha_1}{\lcl-1,1\rcl}{\R^2}{t}{(t,1)}
		\yds
		\funcion{\alpha_2}{\lcl-1,1\rcl}{\R^2}{t}{(-t,2-t^2).}
		\]
		Como la curva $C$ es cerrada, tenemos que
			\begin{align*}
				\int_C f \cdot \ ds
				&=\int_{\alpha_1} f \cdot \ ds+
				\int_{\alpha_2} f \cdot \ ds\\
				&=\int_{-1}^1 f(\alpha_1(t)) \cdot \alpha_1'(t) dt+
				\int_{-1}^1 f(\alpha_2(t)) \cdot \alpha_2'(t) dt\\
				&=\int_{-1}^1 f(t,1) \cdot (1,0) dt+
				\int_{-1}^1 f(-t,2-t^2) \cdot (-1,-2t) dt\\
				&=\int_{-1}^1 e^t dt+
				\int_{-1}^1 -(2-t^2)^2e^{-t}-2t^3e^{2-t^2} dt\\
				&=48e^{-1}-8e\approx -4.1.
			\end{align*}
		Por otro lado, tenemos que 
		\[
		\iint_D\rot(f)\ dA
		=\int_{-1}^1\int_1^{2-x^2}2xe^y-2ye^x\ dydx=48e^{-1}-8e\approx -4.1.
		\]
		Como se puede observar ambas integrales coinciden; es decir, 
		\[
		\int_{\partial D} f \cdot \ ds
		=\iint_D\rot(f)\ dA.\qedhere
		\]
\end{respuesta}

%%%%%%%%%%%%%%%%%%%%%%%%%%%%%%%%%%%%%%%%%%%
%%%%%%%%%%%%%%%%%%%%%%%%%%%%%%%%%%%%%%%%%%%
%%%%%%%%%%%%%%%%%%%%%%%%%%%%%%%%%%%%%%%%%%%
% Tema 17: Teoremas integrales
\item Calcule la siguiente integral de línea aplicando el Teorema de Green:
	\[
     \oint_\alpha f dS
    \]
siendo $f$ el campo vectorial,
		\[
        	\funcion{f}{\R^2}{\R^2}{(x,y)}{\l(xy,x^2y^3\r)} 
        \]
y $\alpha$ el triángulo cuyos vértices son $A=(1,1)$, $B=(2,2)$ y $C=(3,0)$. 

\begin{respuesta}
La curva $\alpha$ es el triángulo formado por las rectas: $L_1$ entre $A$ y $B$, $L_2$ entre $A$ y $C$, y $L_3$ entre $B$ y $C$. Las rectas tienen las ecuaciones:
		\[
        	\funcion{L_1}{\R}{\R}{x}{x} 
        \]
        \[
        	\funcion{L_2}{\R}{\R}{x}{6-2x} 
        \]
        \[
        	\funcion{L_3}{\R}{\R}{x}{\frac{3}{2}-\frac{1}{2}x} 
        \]
\begin{center}
%		\includegraphics[scale=0.4]{images/green_1}
\end{center}

Al tratarse de una curva cerrada y tomando el sentido antihorario podemos aplicar el Teorema de Green:
	\[
    \oint_\alpha f dS=\iint_D \rot \l(f\r) dA
    \]
\\    
Tenemos que,
	\[
    \begin{aligned}
    \rot(f)(x,y)&=\dfrac{\partial f_2}{\partial x}(x,y)-\dfrac{\partial f_1}{\partial y}(x,y)\\
    	  &=2xy^3-x
    \end{aligned}
    \]
La región $D$ es de tipo 3, y se la resolverá como tipo 1. Se divide a $D$ en dos intervalos.    

	\[
    \begin{aligned}
    \iint_D \rot \l(f\r) dA &=\int_1^2 \int_{\frac{3}{2}-\frac{1}{2}x}^{x} \l(2xy^3-x\r) dydx + \int_2^3 \int_{\frac{3}{2}-\frac{1}{2}x}^{6-2x} \l(2xy^3-x\r) dydx\\
    &\approx3.747 + 1.703\\
    &\approx5.45
    \end{aligned}
    \]
\end{respuesta}

%%%%%%%%%%%%%%%%%%%%%%%%%%%%%%%%%%%%%%%%%%%
%%%%%%%%%%%%%%%%%%%%%%%%%%%%%%%%%%%%%%%%%%%
%%%%%%%%%%%%%%%%%%%%%%%%%%%%%%%%%%%%%%%%%%%
% Tema 17: Teoremas integrales
\item Calcular el trabajo realizado por la fuerza 
		\[
        	\funcion{F}{\R^2}{\R^2}{(x,y)}{\l(2xy^2-\dfrac{y^2}{2}, 2x^2y+\dfrac{x^2}{2}\r)} 
        \]
al mover una partícula a lo largo de la circunferencia con centro en $(3,5)$ y radio $2$, en sentido antihorario.

\begin{respuesta}
Al tratarse de una curva cerrada y tomando el sentido antihorario podemos aplicar el Teorema de Green.
	\[
    \oint_\alpha F dS=\iint_D \rot \l(F\r) dA
    \]
\\    
Tenemos que,
	\[
    \begin{aligned}
    \rot(F)(x,y)&=\dfrac{\partial F_2}{\partial x}(x,y)-\dfrac{\partial F_1}{\partial y}(x,y)\\
    	  &=\l(4xy+x\r)-\l(4xy-y\r)\\
          &=x+y
    \end{aligned}
    \]
    
La región $D$ es el círculo de ecuación $(x-3)^2+(y-5)^2=4$.

\begin{center}
%		\includegraphics[scale=0.5]{images/green_2}
\end{center}
Tenemos una región tipo 3 pero se la resolverá como tipo 1, para ello las semicircunferencias que limitan al círculo tendrán las ecuaciones $y=5\pm \sqrt{4-(x-3)^2}$.

La integral de línea será entonces,

	\[
    \begin{aligned}
    \iint_D \rot \l(f\r) dA &=\int_1^5 \int_{5- \sqrt{4-(x-3)^2}}^{5 + \sqrt{4-(x-3)^2}} \l(x+y\r) dydx\\ 
    &=32 \pi.
    \end{aligned}
    \]
    Se sugiere el cambio de variable
    \[
        \funcion{T}{[0,2]\times[0,2\pi]}{\R^2}{(r,t)}{(r\cos t+3,r\sen t+5 )}
    \]
    donde $J_T(r,t)=r$, entonces
    \[
        \iint_D \rot \l(f\r) dA =\int_0^{2\pi}\int_0^2 (r\cos t + r\sen t +8)rdrdt
    \]
    y verificar que el resultado es $32\pi$.
    
    
\end{respuesta}
%%%%%%%%%%%%%%%%%%%%%%%%%%%%%%%%%%%%%%%%%%%
%%%%%%%%%%%%%%%%%%%%%%%%%%%%%%%%%%%%%%%%%%%
%%%%%%%%%%%%%%%%%%%%%%%%%%%%%%%%%%%%%%%%%%%

%%%%%%%%%%%%%%%%%%%%%%%%%%%%%%%%%%%%%%%%%%%
%%%%%%%%%%%%%%%%%%%%%%%%%%%%%%%%%%%%%%%%%%%
%%%%%%%%%%%%%%%%%%%%%%%%%%%%%%%%%%%%%%%%%%%
% Tema 17: Teoremas integrales
\item
    Sea $\Omega\subset\R^{2}$, considere el campo escalar $\func{f}{\Omega}{\R}$ dos veces diferenciable en $\Omega$, tal que verifica la siguiente ecuación:
    \[
        D_{22}f(x, y) + D_{11}f(x, y) =  0
    \]
    para cada $(x, y)\in\Omega$. Demuestre que si $\mathcal{C}$ es cualquier curva cerrada a la que se aplica el teorema de Green, entonces
    \[
        \oint_{\mathcal{C}}{(D_{2}f, - D_{1}f)\cdot ds} = 0.
    \]

\begin{respuesta}\hspace{0pt}

A partir de la integral de línea definamos el campo vectorial $g$ por
\[
    \funcion{g}{\Omega}{\R^{2}}{(x, y)}{(D_{2}f(x, y), - D_{1}f(x, y)).}
\]
Notemos que $g$ es diferenciable y que su rotacional es tal que $\text{rot}(g)(x,y) = 0$ para cada $(x,y)\in\Omega$. Sea $\Omega_{1}\subset\R^{2}$ la región cuya frontera es $\mathcal{C}$, gracias a la aplicación del teorema de Green, obtenemos que
\[
    \oint_{\mathcal{C}}{(D_{2}f, - D_{1}f)\cdot ds} = \pm\iint_{\Omega_{1}}{\rot(g)(x, y)\,dA} = 0.
\]
Notemos que el signo más y menos se debe a que a priori no se conoce la orientación de $\mathcal{C}$.\qedhere    
\end{respuesta}
%%%%%%%%%%%%%%%%%%%%%%%%%%%%%%%%%%%%%
%%%%%%%%%%%%%%%%%%%%%%%%%%%%%%%%%%%%%
\item
    Utilice el teorema de Green para calcular el área entre la elipse de ecuación $\dfrac{x^{2}}{9} + \dfrac{y^{2}}{4} = 1$ y el círculo de ecuación $x^{2} + y^{2} = 25$.

\begin{respuesta}\hspace{0pt}

Notemos que tanto la elipse como el círculo se encuentran centrados en el origen
%\begin{center}
%	\includegraphics[scale=0.3]{Graficos/HE16_fig00.png}
%\end{center}

%No pude agregar la imagen, está comentado el código.

Definimos las parametrizaciones $\alpha_{1}$ y $\alpha_{2}$ para representar la elipse y el círculo, respectivamente.
\[
    \funcion{\alpha_{1}}{[0, 2\pi]}{\R^{2}}{t}{(5\cos(t), 5\sen(t))}
    \texty
    \funcion{\alpha_{2}}{[0, 2\pi]}{\R^{2}}{t}{(3\cos(t), -2\sen(t)).}
\]
Luego, el teorema de Green puede ser aplicado, y, por tanto, el área $A$ se obtiene por:
\begin{align*}
    A &= \dfrac{1}{2}\oint_{\alpha}{(-y\,dx + x\,dy)}\\
      &= \dfrac{1}{2}\oint_{\alpha_{1}}{-y\,dx + x\,dy} + \dfrac{1}{2}\oint_{\alpha_{2}}{-y\,dx + x\,dy}\\
      &= \dfrac{1}{2}\int_{0}^{2\pi}{(-5\sen(t), 5\cos(t))\cdot(-5\sen(t), 5\cos(t))\,dt} \\
      &+ \dfrac{1}{2}\int_{0}^{2\pi}{(2\sen(t), 3\cos(t))\cdot(-3\sen(t), -2\cos(t))\,dt}\\
      &= 19\pi.
\end{align*}
Por tanto, el área entre la elipse de ecuación $\dfrac{x^{2}}{9} + \dfrac{y^{2}}{4} = 1$ y el círculo de ecuación $x^{2} + y^{2} = 25$ es igual a $19\pi$ unidades cuadradas.\qedhere    
\end{respuesta}
%%%%%%%%%%%%%%%%%%%%%%%%%%%%%%%%%%%%%%%%%
%%%%%%%%%%%%%%%%%%%%%%%%%%%%%%%%%%%%%%%%%
\item Compruebe que se cumple el teorema de Stokes para la superficie $S$ parametrizada por $\alpha(u,v)=(u\cos(v),u\sen(v),v)$; $0\leq u \leq 1$; $0\leq v \leq \frac{\pi}{2}$
y el campo vectorial 
\[
    f(x,y,z)=z\veci+x\vecj+y\veck
\]

\begin{respuesta} 
    \item[i)] La superfice S es la helicoidal, comenzamos calculando
    \[
        rot(f)=\nabla \times f=\left| \begin{matrix}
        \veci &\vecj &\veck\\
        \frac{\delta}{\,dx} &\frac{\delta}{\,dy} &\frac{\delta}{\,dz} \\
        z &x &y
        \end{matrix} \right|=\veci+\vecj+\veck
            \]
    ahora calculamos el vector normal estándar, para lo que necesitamos los vectores tangentes
    \[
        T_x(u,v)=\frac{\partial\alpha}{\partial u}(u,v)=(\cos(v),\sen (v),0)       
            \]
    \[
        T_y(u,v)=\frac{\partial\alpha}{\partial v}(u,v)=(-u\sen(v),u\cos (v),1)       
            \]        
    entonces
    \[
        N(u,v)=T_x(u,v)\times   T_y(u,v)
            \]
            \[
        N(u,v)= (\sen(v), -\cos(v), u)
            \]
    por tanto
    \begin{align*}
        \iint_S\nabla \times f\cdot\,ds&=\iint_D(1,1,1)\cdot(\sen(v),-\cos(v),u)\,dv\,du\\
        &=\int_0^1\int_0^{\frac{\pi}{2}}(\sen(v)-\cos(v)+u)\,dv\,du\\
        &=\int_0^1\frac{\pi}{2}u\,du=\frac{\pi}{4}u^2|_0^1\\
        &=\frac{\pi}{4}
    \end{align*}
    \item[ii)] Por otro lado, $\delta S$ consta de cuatro partes, las cuales las parametrizamos mediante
    \begin{align*}
        \alpha_1(u)&=(u,0,0);\quad 0\leq u \leq 1\\
        \alpha_2(v)&=(\cos(v),\sen(v),v);\quad 0\leq v \leq \frac{\pi}{2}\\
        \alpha_3(u)&=(0,1-u,\frac{\pi}{2});\quad 0\leq u \leq 1\\
        \alpha_4(v)&=(0,0,\frac{\pi}{2}-v);\quad 0\leq v \leq \frac{\pi}{2}
    \end{align*}      
    entonces 
    \begin{align*}
     \oint_{\delta S}f\cdot \,dS=&\int_0^1(0,u,0)\cdot (1,0,0)\,du\\ &+\int_0^{\frac{\pi}{2}}(v,\cos(v),\sen(v))\cdot(-\sen(v),\cos(v),1)\,dv\\
     &+\int_0^1(\frac{\pi}{2},0,1-u)\cdot(0,-1,0)\,dv\\  
     &+\int_0^\frac{\pi}{2}(\frac{\pi}{2}-v,0,0)\cdot(0,0,-1)\,dv\\  
     =&\int_0^10\,du+\int_0^{\frac{\pi}{2}}(-v\sen(v)+\cos^2(v)+\sen(v)\,dv)+\int_0^10\,du+\int_0^{\frac{\pi}{2}}0\,dv\\
     =&\frac{\pi}{4}
    \end{align*}
    Las dos respuestas concuerdan 
\end{respuesta}


%%%%%%%%%%%%%%%%%%%%%%%%%%%%%%%%%%%%%%%%%%%
%%%%%%%%%%%%%%%%%%%%%%%%%%%%%%%%%%%%%%%%%%%
%%%%%%%%%%%%%%%%%%%%%%%%%%%%%%%%%%%%%%%%%%%
% Tema 17: Teoremas integrales
\item 
    	Sea el campo vectorial
		\[
		\funcion{f}{\R^3}{\R^3}{(x,y,z)}{(x^2z,xy^2,z^2)}
		\]
		y sea la curva $C$ la intersección del plano definido por $x+y+z=1$ y el cilindro definido por $x^2+y^2=9$ orientado en sentido antihorario. Use el Teorema de Stokes para calcular
		\[
		\int_C f\cdot dr.
		\]

\begin{respuesta}
		Por el Teorema de Stokes se tiene que
		\[
		\int_C f\cdot dr=\iint_S \rot(f) \cdot dS
		\]
		donde $S$ es la elipse formada por la intersección de las ecuaciones $x+y+z=1$ y $x^2+y^2=9$. Por definición se tiene que
		\[
		\rot f(x,y,z)=(0-0,x^2-0,y^2-0)=(0,x^2,y^2)
		\]
		para todo $x,y,z\in\R$. Por otra parte, una representación paramétrica de $S$ es
		\[
		\funcion{\alpha}{D}{\R^3}{(u,v)}{(u,v,1-u-v)}
		\]
		donde $D$ es el círculo centrado en el origen de radio 3, $D=\{(u,v)\in\R^2:-3\leq u \leq 3,-\sqrt{9-u^2}\leq y \leq \sqrt{9-u^2} \}$, entonces
		\[
		N(u,v)=(1,0,-1)\times(0,1,-1)=(1,1,1)
		\]
		para cada $(u,v)\in D$. Por todo esto se tiene que
			\begin{align*}
				\int_C f\cdot dr
				&=\iint_S \rot(f) \cdot dS\\
				&=\iint_D \rot(f)(\alpha(u,v)) \cdot N(u,v)\ dA\\
				&=\iint_D \rot(f)(u,v,1-u-v) \cdot (1,1,1)\ dA\\
				&=\iint_D (0,u^2,v^2) \cdot (1,1,1)\ dA\\
				&=\iint_D u^2+v^2\ dudv,\\
			\end{align*}
		realizando el cambio de variable
		\[
		    u=r\cos\theta \texty v=r\sen \theta
		\]
		donde $J(r,\theta)=r$, para $(r,\theta)\in [0,3]\times[0,2\pi]$, se obtiene
		    \begin{align*}
		        &=\int_0^{2\pi} \int_0^3 (r^2\cos^2\theta+r^2\sen^2\theta)rdrd\theta\\
		        &=\int_0^{2\pi} \int_0^3 r^3\ drd\theta\\
				&=\dfrac{81\pi}{2}.
		    \end{align*}
		Por lo tanto,
		\[
		\int_C f\cdot dr=\dfrac{81\pi}{2}.\qedhere
		\]
		
		Alternativamente, otra parametrización de $S$ es
		\[
		    \funcion{\beta}{[0,3]\times[0,2\pi]}{\R^3}{(r,\theta)}{(r\cos\theta,r\sen\theta)}
		\]
		donde $N(r,\theta)=(r,r,r)$, además
		\[
		    \rot(f_{(\beta)})=(0,r^2\cos^2\theta,r^2\sen^2\theta).
		\]
		Entonces
		\begin{align*}
		    	\int_C f\cdot dr&=\iint_{\beta}\rot(f_{(\beta)})\cdot N(r,\theta) dA\\
		    	&=\int_o^3\int_0^{2\pi}r^3d\theta dr\\
		    	&=\dfrac{81}{2}\pi.
		\end{align*}
\end{respuesta}


%%%%%%%%%%%%%%%%%%%%%%%%%%%%%%%%%%%%%%%%%%%
%%%%%%%%%%%%%%%%%%%%%%%%%%%%%%%%%%%%%%%%%%%
%%%%%%%%%%%%%%%%%%%%%%%%%%%%%%%%%%%%%%%%%%%
% Tema 17: Teoremas integrales
\item Dado el campo vectorial
    \[
    \funcion{f}{\R^3}{\R^3}
    {(x,y,z)}{f(x,y,z)=(x+y^2,y+z^2,x^2+z)}
	\]
    y la superficie $S$ con orientación positiva tal que $\delta S$ es el triangulo de vértices $(1,0,0)$, $(0,1,0)$ y $(0,0,1)$. Muestre que se cumple el teorema de Stokes.
    
\begin{respuesta}
    La curva $\delta S$ es una curva suave que esta conformada por segmentos simples y es cerrada, es decir, la superficie $S$ es la parte de un plano con ecuación $x+y+z=1$, cuya proyección respecto al plano $xy$ es un triángulo rectángulo que puede ser vista como una región $D$ dada por
    \[
    D=\{(x,y)\in\R^2\,:\,0\leq x \leq 1, \, 0 \leq y \leq 1 - x\}.
    \]
    Además, el campo vectorial posee primeras derivadas parciales continuas para cada una de sus componentes.
    
    \noindent Entonces, se verifican las hipótesis del teorema de Stokes, y por tanto
    \[
    \iint_{S}{\mbox{rot}(f)\cdot dS} = \oint_{\delta S}{f\cdot ds}.
    \]
    
    \noindent Calculemos ahora, cada uno de los términos de la ecuación anterior.     Para el termino $\iint_{S}{\mbox{rot}(f)\cdot dS}$, notemos que
    \[
    \mbox{rot}(f)(x,y,z) = (-2z,-2x,-2y),
    \]
    mientras que una parametrización de la superficie $S$ está dada por
    \[
    \funcion{\beta}{D}{\R^3}
    {(u,v)}{\beta(u,v)=(u,v,1-u-v)}
	\]
    cuyos vectores tangentes son
    \[
    \begin{array}{ccc}
    \dfrac{\partial \beta}{\partial u}(u,v) = (1,0,-1), & & \dfrac{\partial \beta}{\partial v}(u,v) = (0,1,-1),
    \end{array}
    \]
    con vector normal $N(u,v)=(1,1,1)$, para todo $(u,v)\in D$. Luego,
    \[
      \mbox{rot}(f)(\beta(u,v))\cdot N(u,v) = (-2+2u+2v,-2u,-2v)\cdot(1,1,1)=-2,
    \]
    por lo cual
    \begin{align*}
    \iint_{S}{\mbox{rot}(f)\cdot dS} &= \iint_{D}{\mbox{rot}(f)(\beta(u,v))\cdot N(u,v)dA}\\
    &= \int_{0}^{1}{\l(\int_{0}^{1-x}{-2\, dy dx}\r)}\\
    &= -1.
	\end{align*}
    
    \noindent Para el término $\oint_{\delta S}{f\cdot ds}$, considere las parametrizaciones $\func{\alpha_1, \alpha_2, \alpha_3}{[0,1]}{\R^3}$ definidas por 
    \begin{align*}
    \alpha_1(t) &= (1-t,t,0)\\
    \alpha_2(t) &= (0,1-t,t)\\
    \alpha_3(t) &= (t,0,1-t)
    \end{align*}    
    que unen los puntos (1,0,0) con (0,0,1), (0,1,0) con (0,0,1) y (0,0,1) con (1,0,0) respectivamente. A partir de lo anterior, se puede comprobar que 
    \begin{align*}
    \oint_{\delta S}{f\cdot ds} &= \int_{\alpha_1}{f\cdot ds} + \int_{\alpha_2}{f\cdot ds} + \int_{\alpha_3}{f\cdot ds}\\
    &= 3\int_{0}^{1}{(-1+2t-t^2)\,dt}\\
    &= -1.
    \end{align*}
    Entonces, $\iint_{S}{\mbox{rot}(f)\cdot dS} = \oint_{\delta S}{f\cdot ds}.$     Por lo cual, se ha mostrado que para el campo vectorial dado y la superficie $S$ se cumple el teorema de Stokes.
\end{respuesta}
%%%%%%%%%%%%%%%%%%%%%%%%%%%%%%%%%%%%%%%%%%% FC
%%%%%%%%%%%%%%%%%%%%%%%%%%%%%%%%%%%%%%%%%%%
\item 
Sean
	\[
	\funcion{F}{\R^3}{\R^3}{(x,y,z)}{\l(x-y,y-z,z-x\r)}
	\]
y
	\[
	\funcion{\Omega}{[0,5]\times[0,2\pi]}{\R^3}{(r,\theta)}{\l(r \cos \theta, r \sen \theta, 5-r\r)}
	\]
Verificar que se cumple el teorema de Stokes

\begin{respuesta}
Se puede ver que esta superficie es un cono cuya frontera es la circunferencia en el plano $xy$ de radio 5 centrada en el origen.

Para verificar que se cumple este teorema se debe comparar ambos lados de la igualdad
\[
    \oint_{\partial \Omega} F \cdot dS=\iint_{\Omega} \rot (F) \cdot dS.
\]
Se tiene que el vector normal estándar de $\Omega$ es
	\[
	    \begin{aligned}
		N(r,\theta)&=T_r(r,\theta) \times T_{\theta}(r,\theta)\\
		&=\l(\cos \theta, \sen \theta, -1 \r)\times \l(-r \sen \theta, r \cos \theta, 0\r)\\
		&= \l(r \cos \theta, r \sen \theta, r\r)
		\end{aligned}
	\]
y el rotacional de $F$ es
	\[
	    \begin{aligned}
		\rot(F)&=\nabla \times F\\
		&=\l(1,1,1\r),
		\end{aligned}
	\]
de aquí
	\begin{align*}
		\iint_{\Omega} \rot (F) \cdot dS&=\iint_{\Omega} \rot(F)_{\l(\Omega(r,\theta)\r)} \cdot N(r,\theta)\ dA\\
			&=\int_0^{2\pi} \int_0^5 \l(1,1,1\r) \cdot \l(r \cos \theta, r \sen \theta, r\r) dr d\theta\\
			&=\int_0^{2\pi} \int_0^5 \l(r \cos \theta + r \sen \theta + r\r) dr d\theta\\
			&=25\pi.
	\end{align*}
Nótese que la frontera de la superficie cónica $\Omega$ puede ser parametrizada por
    \[
        \funcion{\partial \Omega}{[0,2\pi]}{\R^3}{t}{(5\cos t,5\sen t,0)}
    \]
de aquí
\[
    \partial \Omega'(t)=(-5\sen t,5\cos t,0)
\]
y
\[
    F(\partial \Omega)=(5\cos t-5\sen t,5\sen t,-5\cos t)
\]
para $t\in [0,2\pi]$.
Finalmente,
\begin{align*}
    \oint_{\partial \Omega} F \cdot dS&=\int_0^{2\pi} F(\partial \Omega(t)) \cdot \partial \Omega'(t) dt\\
    &=\int_0^{2\pi} 25\sen^2(t)dt\\
    &=25\pi
\end{align*}
con lo cual se verifica que se cumple el teorema de Stokes.

Como trabajo adicional, dada la superficie (semiesfera) parametrizada por
\[
    \funcion{D}{[0,2\pi]\times[0,\pi/2]}{\R^3}{(u,v)}{(5\cos(u)\sen(v),5\sen(u)\sen(v),5\cos(v))}
\]
cuya frontera $\partial D$ es igual que $\partial \Omega$, entonces
\[
    \oint_{\partial \Omega} F \cdot dS=\oint_{\partial D} F \cdot dS=25\pi.
\]

Verificar que
\[
    \iint_{D} \rot (F) \cdot dS=25\pi.
\]
\end{respuesta}

%%%%%%%%%%%%%%%%%%%%%%%%%%%%%%%%%%%%%%%%%%%
%%%%%%%%%%%%%%%%%%%%%%%%%%%%%%%%%%%%%%%%%%%
%%%%%%%%%%%%%%%%%%%%%%%%%%%%%%%%%%%%%%%%%%%
% Tema 17: Teoremas integrales
\item Dado el campo vectorial
    \[
    \funcion{f}{\R^3}{\R^3}
    {(x,y,z)}{f(x,y,z)=(3xy^2,xe^z,z^3).}
	\]
    Determine el flujo de $f$ a través de $D$, en donde $D$ es la superficie del sólido acotado por el cilindro de ecuación $y^2 + z^2 = 1$ y los planos de ecuación $x=-1$ y $x=2$.

\begin{respuesta}
    Notemos que $f$ posee primeras derivadas parciales continuas en cada una de sus componentes y que $D$ es una región tal que  
\[
D = \l\{(x,y,z)\in\R^3\,:\,-1\leq x \leq 2, -1 \leq y \leq 1, -\sqrt{1-y^2}\leq z \leq \sqrt{1-y^2}\r\}.
\]

\noindent Entonces, para determinar el flujo, $\iint_{D}{f\cdot dS}$, utilizaremos el Teorema de Gauss, es decir, 
    \begin{align*}
    \iint_{D}{f\cdot dS} &= \iiint_{D}{\mbox{div}(f)(x,y,z)\,dV}\\
    &= \int_{-1}^{2}{\int_{-1}^{1}{\int_{-\sqrt{1-y^2}}^{\sqrt{1-y^2}}{3(y^2 + z^2)\,dz}dy}dx}.
    \end{align*}
    Con el cambio de variables $x=t$, $y=r\cos\theta$ y $z=r\sen\theta$, se obtiene que 
    \[
    \iint_{D}{f\cdot dS} = 3\int_{-1}^{2}{\int_{0}^{2\pi}{\int_{0}^{1}{(r^2)rdr}d\theta}dt} = \dfrac{9}{2}\pi.
    \]
    
    \noindent Por lo tanto, el flujo de $f$ a través de $D$ es igual a $\dfrac{9}{2}\pi$.
\end{respuesta}

%%%%%%%%%%%%%%%%%%%%%%%%%%%%%%%%%%%%%%%%%%%
%%%%%%%%%%%%%%%%%%%%%%%%%%%%%%%%%%%%%%%%%%%
%%%%%%%%%%%%%%%%%%%%%%%%%%%%%%%%%%%%%%%%%%%
% Tema 17: Teoremas integrales
\item
    Dado el campo vectorial
    \[
    \funcion{f}{\R^3}{\R^3}
    {(x,y,z)}{f(x,y,z)=(3x,xy,3xz).}
	\]
    Calcule $\iint_{D}{f\cdot dS}$, donde $D$ es un cubo limitado por los planos de ecuaciones $x=0$, $x=1$, $y=0$, $y=1$, $z=0$ y $z=1$.

\begin{respuesta}
Notemos que $f$ posee primeras derivadas parciales continuas en cada una de sus componentes y que $D$ es una región tal que  
\[
D = \l\{(x,y,z)\in\R^3\,:\,0\leq x \leq 1, 0 \leq y \leq 1, 0 \leq z \leq 1\r\}.
\]

\noindent Entonces, para determinar $\iint_{D}{f\cdot dS}$, utilizaremos el Teorema de Gauss, es decir, 
    \begin{align*}
    \iint_{D}{f\cdot dS} &= \iiint_{D}{\mbox{div}(f)(x,y,z)\,dV}\\
    &= \int_{0}^{1}{\int_{0}^{1}{\int_{0}^{1}{(3+3x)dz}dy}dx}\\
    \dfrac{9}{2}.
    \end{align*}
    
    \noindent Por lo tanto, $\iint_{D}{f\cdot dS}$ es igual a  $\dfrac{9}{2}$.
\end{respuesta}

%%%%%%%%%%%%%%%%%%%%%%%%%%%%%%%%%%%%%%%%%%%
%%%%%%%%%%%%%%%%%%%%%%%%%%%%%%%%%%%%%%%%%%%
%%%%%%%%%%%%%%%%%%%%%%%%%%%%%%%%%%%%%%%%%%%
% Tema 17: Teoremas integrales
\item Dado el campo vectorial
    \[
        \funcion{F}{\R^3}{\R^3}{(x,y,z)}{(x+z,y+z,x+y).}
    \]
    Verificar que se cumple el Teorema de Gauss, si se considera el sólido limitado por las superficies de ecuación $y^2+z^2=1$, $x=0$ y $X=2$.

\begin{respuesta} 
    Para verificar que se cumple el teorema de Gauss se debe comparar ambos lados de la igualdad
    \[
        \oint_{\delta R}FdS=\iiint_R div(F)dV.
    \]
	El sólido $R$ es un cilindro, tal que
	\begin{align*}
	    x&=x,\\
	    y&=r\cos t, \\
	    z&=r\sen t 
	\end{align*}
	para $(x,r,t)\in [0,2]\times[0,1]\times[0,2\pi]$ y $J(x,r,t)=r$.
    Además
    \begin{align*}
         div(f)(x,y,z)&=\frac{\partial}{\partial x}(x+z)+\frac{\partial}{\partial y}(y+z)+\frac{\partial}{\partial z}(x+y)\\
         &=1+1+0\\
         &=2
    \end{align*}
    entonces
    \begin{align*}
        \iiint_D div(f)\,dV&=\int_0^2\int_0^{2\pi}\int_0^12r\,dr\,dt\,dx\\
        &=\int_0^2\int_0^{2\pi}r^2|_0^1\,dt \,dx\\
        &=\int_0^2\int_0^{2\pi}\,dt\,dx\\
        &=\int_0^22\pi\,dx\\
        &=4\pi
    \end{align*}
    Ahora calculamos la integral de superficie de $F$ sobre la frontera del sólido $R$. La frontera de $R$ está constituida por 3 superficies cuyas representaciones paramétricas son
    \[
        \funcion{\alpha}{[0,1]\times[0,2\pi]}{\R^3}{(r,t)}{(0,-r\cos t,r\sen t),}
    \]
    
    \[
        \funcion{\beta}{[0,1]\times[0,2\pi]}{\R^3}{(r,t)}{(2,r\cos t,r\sen t)}
    \]
    y
    \[
        \funcion{\gamma}{[0,2]\times[0,2\pi]}{\R^3}{(x,t)}{(-x,\cos t,\sen t)}
    \]
    cuyos vectores normales son respectivamente $N_\alpha=(-r,0,0)$, $N_\beta=(r,0,0)$ y $N_\gamma=(0,\cos t,\sen t)$.
    
    Finalmente
   \[
       \oint_{\delta R}FdS=\int_\alpha FdS+\int_\beta FdS+\int_\gamma FdS
    \]
    donde
    \begin{align*}
        \int_\alpha FdS&=\iint_\alpha F(\alpha(r,t))\cdot N_\alpha dtdr\\
        &=\int_0^1\int_0^{2\pi} -2r\sen (t) dtdr\\
        &=0,
    \end{align*}
    
    \begin{align*}
        \int_\beta FdS&=\iint_\beta F(\beta(r,t))\cdot N_\beta dtdr\\
        &=\int_0^1\int_0^{2\pi} r(2+r\sen (t))dtdr\\
        &=2\pi
    \end{align*} 
    y
    \begin{align*}
        \int_\gamma FdS&=\iint_\gamma F(\gamma(r,t))\cdot N_\gamma dtdx\\
        &=\int_0^2\int_0^{2\pi} (\cos^2(t)+2\sen(t)\cos(t)-x\sen(t)) dtdx\\
        &=2\pi.
    \end{align*}
    Entonces 
    \[
        \oint_{\delta R}FdS=4\pi,
    \]
    con lo cual se verifica lo propuesto.

   
   
\end{respuesta}
%%%%%%%%%%%%%%%%%%%%%%%%%%%%%
%%%%%%%%%%%%%%%%%%%%%%%%%%%%%
\item 
    Sea el campo vectorial
    \[
        \funcion{F}{\R^2}{\R^2}{(x,y)}{(x^3-y,x).}
    \]
    Verifique que se cumple el Teorema de Gauss en el Plano, al considerar la circunferencia centrada en el origen de radio 1.

\begin{respuesta}
    Sea la parametrización
    \[
        \funcion{\alpha}{[0,2\pi]}{\R^2}{t}{(\cos(t),\sen(t))}
    \]
    de la circunferencia centrada en el origen y de radio 1; en sentido antihorario. El Teorema de Gauss en el plano establece que
    \[
        \int_\alpha F \cdot n ds=\iint_\Omega \dive(F) dA;
    \]
    para verificar que este teorema se cumple, vamos a calcular cada integral por separado.
    \begin{itemize}
        \item Se tiene que el vector unitario normal de $\alpha$ es
        \[
            n(t)=\dfrac{(\alpha_2'(t),-\alpha_1'(t))}{\|\alpha'(t)\|}
        \]
        para todo $t\in [0,2\pi]$. Por lo tanto, se tiene que
        \begin{align*}
            \int_\alpha F\cdot n ds
            &=\int_0^{2\pi} F(\alpha(t))\cdot n(t) \|\alpha'(t)\| dt\\
            &=\int_0^{2\pi} F(\cos(t),\sen(t))\cdot (\alpha_2'(t),-\alpha_1'(t)) dt\\
            &=\int_0^{2\pi} (\cos^3(t)-\sen(t),\cos(t))\cdot (\alpha_2'(t),-\alpha_1'(t)) dt\\
            &=\int_0^{2\pi} (\cos^3(t)-\sen(t),\cos(t))\cdot (\cos(t),\sen(t)) dt\\
            &=\int_0^{2\pi} \cos^4(t) dt\\
            &=\dfrac{3\pi}{4}.
        \end{align*}
        
        \item De igual manera, se tiene que, pasando a polares,
        \begin{align*}
            \iint_{\Omega} \dive(F) dA
            &= \iint_{\Omega} 3x^2 dA\\
            &= \int_0^{2\pi} \int_0^1 3r^3\cos^2(t) dr dt\\
            &= \dfrac{3}{4}\int_0^{2\pi} \cos^2(t) dt\\
            &= \dfrac{3}{4}\pi.
        \end{align*}        
    \end{itemize}
    Estos resultados verifican que 
    \[
        \iint_{\Omega} \dive(F) dA=\int_{\partial \Omega} F \cdot n ds. \qedhere 
    \]
\end{respuesta}


%%%%%%%%%%%%%%%%%%%%%%%%%%%%%%%%%%%%%%%%%%%
%%%%%%%%%%%%%%%%%%%%%%%%%%%%%%%%%%%%%%%%%%%
%%%%%%%%%%%%%%%%%%%%%%%%%%%%%%%%%%%%%%%%%%%
% Tema 17: Teoremas integrales
\item 
    	Sea el campo vectorial
		\[
		\funcion{f}{\R^3}{\R^3}{(x,y,z)}{(x^2z,xy^2,z^2)}
		\]
		y sea la curva $C$ la intersección del plano definido por $x+y+z=1$ y el cilindro definido por $x^2+y^2=9$ orientado en sentido antihorario. Use el Teorema de Stokes para calcular
		\[
		\int_C f\cdot dr.
		\]

\begin{respuesta}
		Por el Teorema de Stokes se tiene que
		\[
		\int_C f\cdot dr=\iint_S \rot(f) \cdot dS
		\]
		donde $S$ es la elipse formada por la intersección de las ecuaciones $x+y+z=1$ y $x^2+y^2=9$. Por definición se tiene que
		\[
		\rot f(x,y,z)=(0-0,x^2-0,y^2-0)=(0,x^2,y^2)
		\]
		para todo $x,y,z\in\R$. Por otra parte, sea
		\[
		\funcion{\alpha}{D}{\R^3}{(u,v)}{(u,v,1-u-v)}
		\]
		donde $D$ es el círculo centrado en el origen de radio 3 y además
		\[
		N(u,v)=(1,0,-1)\times(0,1,-1)=(1,1,1)
		\]
		para cada $(u,v)\in D$. Por todo esto se tiene que
			\begin{align*}
				\int_C f\cdot dr
				&=\iint_S \rot(f) \cdot dS\\
				&=\iint_D \rot(f)(\alpha(u,v)) \cdot N(u,v)\ dA\\
				&=\iint_D \rot(f)(u,v,1-u-v) \cdot (1,1,1)\ dA\\
				&=\iint_D (0,u^2,v^2) \cdot (1,1,1)\ dA\\
				&=\iint_D u^2+v^2\ dudv\\
				&=\int_0^{2\pi} \int_0^3 r^3\ drd\theta\\
				&=\dfrac{81\pi}{2}.
			\end{align*}
		Por lo tanto,
		\[
		\int_C f\cdot dr=\dfrac{81\pi}{2}.\qedhere
		\]
\end{respuesta}


%%%%%%%%%%%%%%%%%%%%%%%%%%%%%%%%%%%%%%%%%%%
%%%%%%%%%%%%%%%%%%%%%%%%%%%%%%%%%%%%%%%%%%%
%%%%%%%%%%%%%%%%%%%%%%%%%%%%%%%%%%%%%%%%%%%
% Tema 17: Teoremas integrales
\item Dado el campo vectorial
    \[
    \funcion{f}{\R^3}{\R^3}
    {(x,y,z)}{f(x,y,z)=(x+y^2,y+z^2,x^2+z)}
	\]
    y la superficie $S$ con orientación positiva tal que $\delta S$ es el triangulo de vértices $(1,0,0)$, $(0,1,0)$ y $(0,0,1)$. Muestre que se cumple el teorema de Stokes.
    
\begin{respuesta}
    La curva $\delta S$ es una curva suave que esta conformada por segmentos simples y es cerrada, es decir, la superficie $S$ es la parte de un plano con ecuación $x+y+z=1$, cuya proyección respecto al plano $xy$ es un triángulo rectángulo que puede ser vista como una región $D$ dada por
    \[
    D=\{(x,y)\in\R^2\,:\,0\leq x \leq 1, \, 0 \leq y \leq 1 - x\}.
    \]
    Además, el campo vectorial posee primeras derivadas parciales continuas para cada una de sus componentes.
    
    \noindent Entonces, se verifican las hipótesis del teorema de Stokes, y por tanto
    \[
    \iint_{S}{\mbox{rot}(f)\cdot dS} = \oint_{\delta S}{f\cdot ds}.
    \]
    
    \noindent Calculemos ahora, cada uno de los términos de la ecuación anterior.     Para el termino $\iint_{S}{\mbox{rot}(f)\cdot dS}$, notemos que
    \[
    \mbox{rot}(f)(x,y,z) = (-2z,-2x,-2y),
    \]
    mientras que una parametrización de la superficie $S$ está dada por
    \[
    \funcion{\beta}{[0,1]\times[0,1]}{\R^3}
    {(u,v)}{\beta(u,v)=(u,v,1-u-v)}
	\]
    cuyos vectores tangentes son
    \[
    \begin{array}{ccc}
    \dfrac{\partial \beta}{\partial u}(u,v) = (1,0,-1), & & \dfrac{\partial \beta}{\partial v}(u,v) = (0,1,-1),
    \end{array}
    \]
    con vector normal $N(u,v)=(1,1,1)$, para todo $(u,v)\in[0,1]\times[0,1]$. Luego,
    \[
      \mbox{rot}(f)(\beta(u,v))\cdot N(u,v) = (-2+2u+2v,-2u,-2v)\cdot(1,1,1)=-2,
    \]
    por lo cual
    \begin{align*}
    \iint_{S}{\mbox{rot}(f)\cdot dS} &= \iint_{D}{\mbox{rot}(f)(\beta(u,v))\cdot N(u,v)dA}\\
    &= \int_{0}^{1}{\l(\int_{0}^{1-x}{-2\, dy dx}\r)}\\
    &= -1.
	\end{align*}
    
    \noindent Para el término $\oint_{\delta S}{f\cdot ds}$, considere las parametrizaciones $\func{\alpha_1, \alpha_2, \alpha_3}{[0,1]}{\R^3}$ definidas por 
    \begin{align*}
    \alpha_1(t) &= (1-t,t,0)\\
    \alpha_2(t) &= (0,1-t,t)\\
    \alpha_3(t) &= (t,0,1-t)
    \end{align*}    
    que unen los puntos (1,0,0) con (0,0,1), (0,1,0) con (0,0,1) y (0,0,1) con (1,0,0) respectivamente. A partir de lo anterior, se puede comprobar que 
    \begin{align*}
    \oint_{\delta S}{f\cdot ds} &= \int_{\alpha_1}{f\cdot ds} + \int_{\alpha_2}{f\cdot ds} + \int_{\alpha_3}{f\cdot ds}\\
    &= 3\int_{0}^{1}{(-1+2t-t^2)\,dt}\\
    &= -1.
    \end{align*}
    Entonces, $\iint_{S}{\mbox{rot}(f)\cdot dS} = \oint_{\delta S}{f\cdot ds}.$     Por lo cual, se ha mostrado que para el campo vectorial dado y la superficie $S$ se cumple el teorema de Stokes.
\end{respuesta}

%%%%%%%%%%%%%%%%%%%%%%%%%%%%%%%%%%%%%%%%%%%
%%%%%%%%%%%%%%%%%%%%%%%%%%%%%%%%%%%%%%%%%%%
%%%%%%%%%%%%%%%%%%%%%%%%%%%%%%%%%%%%%%%%%%%
% Tema 17: Teoremas integrales
\item Dado el campo vectorial
    \[
    \funcion{f}{\R^3}{\R^3}
    {(x,y,z)}{f(x,y,z)=(3xy^2,xe^z,z^3).}
	\]
    Determine el flujo de $f$ a través de $D$, en donde $D$ es la superficie del sólido acotado por el cilindro de ecuación $y^2 + z^2 = 1$ y los planos de ecuación $x=-1$ y $x=2$.

\begin{respuesta}
    Notemos que $f$ posee primeras derivadas parciales continuas en cada una de sus componentes y que $D$ es una región tal que  
\[
D = \l\{(x,y,z)\in\R^3\,:\,-1\leq x \leq 2, -1 \leq y \leq 1, -\sqrt{1-y^2}\leq z \leq \sqrt{1-y^2}\r\}.
\]

\noindent Entonces, para determinar el flujo, $\iint_{D}{f\cdot dS}$, utilizaremos el Teorema de Gauss, es decir, 
    \begin{align*}
    \iint_{D}{f\cdot dS} &= \iiint_{D}{\mbox{div}(f)(x,y,z)\,dV}\\
    &= \int_{-1}^{2}{\int_{-1}^{1}{\int_{-\sqrt{1-y^2}}^{\sqrt{1-y^2}}{3(y^2 + z^2)\,dz}dy}dx}.
    \end{align*}
    Con el cambio de variables $x=t$, $y=r\cos\theta$ y $z=r\sen\theta$, se obtiene que 
    \[
    \iint_{D}{f\cdot dS} = 3\int_{-1}^{2}{\int_{0}^{2\pi}{\int_{0}^{1}{(r^2)rdr}d\theta}dt} = \dfrac{9}{2}\pi.
    \]
    
    \noindent Por lo tanto, el flujo de $f$ a través de $D$ es igual a $\dfrac{9}{2}\pi$.
\end{respuesta}

%%%%%%%%%%%%%%%%%%%%%%%%%%%%%%%%%%%%%%%%%%%
%%%%%%%%%%%%%%%%%%%%%%%%%%%%%%%%%%%%%%%%%%%
%%%%%%%%%%%%%%%%%%%%%%%%%%%%%%%%%%%%%%%%%%%
% Tema 17: Teoremas integrales
\item
    Dado el campo vectorial
    \[
    \funcion{f}{\R^3}{\R^3}
    {(x,y,z)}{f(x,y,z)=(3x,xy,2xz).}
	\]
    Calcule $\iint_{D}{f\cdot dS}$, donde $D$ es un cubo limitado por los planos de ecuaciones $x=0$, $x=1$, $y=0$, $y=1$, $z=0$ y $z=1$.

\begin{respuesta}
Notemos que $f$ posee primeras derivadas parciales continuas en cada una de sus componentes y que $D$ es una región tal que  
\[
D = \l\{(x,y,z)\in\R^3\,:\,0\leq x \leq 1, 0 \leq y \leq 1, 0 \leq z \leq 1\r\}.
\]

\noindent Entonces, para determinar $\iint_{D}{f\cdot dS}$, utilizaremos el Teorema de Gauss, es decir, 
    \begin{align*}
    \iint_{D}{f\cdot dS} &= \iiint_{D}{\mbox{div}(f)(x,y,z)\,dV}\\
    &= \int_{0}^{1}{\int_{0}^{1}{\int_{0}^{1}{(3+3x)dz}dy}dx}\\
    \dfrac{9}{2}.
    \end{align*}
    
    \noindent Por lo tanto, $\iint_{D}{f\cdot dS}$ es igual a  $\dfrac{9}{2}$.
\end{respuesta}

%%%%%%%%%%%%%%%%%%%%%%%%%%%%%%%%%%%%%%%%%%%
%%%%%%%%%%%%%%%%%%%%%%%%%%%%%%%%%%%%%%%%%%%
%%%%%%%%%%%%%%%%%%%%%%%%%%%%%%%%%%%%%%%%%%%
% Tema 17: Teoremas integrales
\item Compruebe el teorema de Gauss para el campo vectorial $f(x,y,z)=(x+z,y+z,x+y)$ sobre el sólido acotado por 
\[
    0 \leq y^2+z^2 \leq 1 \quad y \quad 0 \leq x \leq 2
        \]
	
\begin{respuesta} 
	El sólido acotado $R$ es un cilindro, tal que en coordenadas cilíndricas
    \[
        0 \leq r \leq 1 \quad ; \quad 0\leq \theta \leq 2\pi \quad ; \quad 0\leq x \leq 2        
            \]
    calculamos la divergencia del campo vectorial
    \begin{align*}
         div(f)(x,y,z)&=\sum_{i=1}^3\frac{\partial f_i (x,y,z)}{\partial x_i}\\
         &=\frac{\partial}{\partial x}(x+z)+\frac{\partial}{\partial y}(y+z)+\frac{\partial}{\partial z}(x+y)\\
         &=1+1+0\\
         &=2
    \end{align*}
    entonces
    \begin{align*}
        \iiint_D div(f)\,dV&=\int_0^2\int_0^{2\pi}\int_0^12r\,dr\,d\theta\,dx\\
        &=\int_0^2\int_0^{2\pi}r^2|_0^1\,d\theta \,dx\\
        &=\int_0^2\int_0^{2\pi}\,d\theta\,dx\\
        &=\int_0^22\pi\,dx\\
        &=4\pi
    \end{align*}
    Ahora calculamos las integrales de superfice sobre el sólido $R$, tales que \\
    $S_1: x=0$\\
   $S_2: x=2$\\
   $S_3: y^2+z^2=1$
   \item[i)] para $S_1:N_1(-1,0,0);dS_1=\,dy\,dz$
   luego
   \begin{align*}
        \iint_{S_1}f\cdot \,ds&=\iint_R(-x-z)\,dy\,dz=\iint_R-z\,dy\,dz\\
        &=\int_0^1\int_0^2-r^2\cos(\theta)\,d\theta\,dr\\
        &=\int_0^1r^2\sen(\theta)|_0^{2\pi}\,dr\\
        &=\int_0^10\,dr\\
        &=0
   \end{align*}
   \item[ii)] para  $S_2:N_2(1,0,0);dS_2=\,dy\,dz$\\
    luego
    \begin{align*}
        \iint_{S_2}f\cdot\,ds&=\int_R\int(x+z)\,dy\,dz\\
        &=\int_R\int(x+2)\,dy\,dz\\
        &=\int_0^1\int_0^{2\pi}(r^2\cos(\theta)+2r)\,d\theta\,dr\\
        &=\int_0^1(r^2\sen(\theta+2r\theta)|_0^{2\pi})\,dr\\
        &=\int_0^14\pi r\,dr=2\pi r^2|_0^1\\
        &=2\pi
    \end{align*}
\end{respuesta}











\end{preguntas}

\end{document}
